% Generated human-review packet template.  Source records, Lean statements,
% and raw-source-to-Spec screening fields are substituted by
% scripts/review_dashboard_packet.py.
% Do not hand-edit generated packets; annotate the PDF or reconcile notes into
% the paper-local human-review log.
\documentclass[10pt]{article}
\usepackage[margin=0.43in]{geometry}
\usepackage{fontspec}
\setmainfont{Latin Modern Roman}
\setmonofont{DejaVu Sans Mono}
\usepackage{hyperref}
\usepackage{amssymb}
\usepackage{fvextra}
\usepackage{enumitem}
\setlength{\parindent}{0pt}
\setlength{\parskip}{0.15em}
\linespread{0.92}
\hypersetup{colorlinks=true,linkcolor=blue,urlcolor=blue}
\DefineVerbatimEnvironment{ReviewVerbatim}{Verbatim}{breaklines=true,breakanywhere=true,fontsize=\scriptsize}
\newcommand{\reviewerbox}{%
  \par\noindent\fbox{\parbox{0.96\linewidth}{\rule{0pt}{0.38in}}}\par
}
\newcommand{\reviewmatch}[1]{%
  \par\noindent\textbf{Reviewer check:} $\square$\; #1\par
  \noindent\emph{If left unchecked, explain the discrepancy or uncertainty below.}\par
}

\begin{document}
\fontsize{9}{10}\selectfont

\begin{center}
{\LARGE Human review packet: KR21Monoculture}\\[0.35em]
\small Generated from byte-pinned source excerpts, the expanded PaperInterface
specifications, and hash-bound raw-source-to-Spec screening records.
\end{center}

\noindent\textbf{Paper:} KR21Monoculture\\
\textbf{Paper id:} \texttt{KR21Monoculture}\\
\textbf{Source version:} arXiv:2101.05853\\
\textbf{Source record:} \texttt{audit/paper\_statement\_map.json}\\
\textbf{Rows:} 57\\
\textbf{Generated:} 2026-08-18\\
\textbf{Source URL:} \href{https://arxiv.org/abs/2101.05853}{official source}

\noindent\fbox{\parbox{0.96\linewidth}{\textbf{Review status:} 0/57 source-claim screens; 0/56 library prerequisite screens. This is a review worksheet while those direct checks remain pending; it is not a final validation receipt.}}\par



\section*{How to use this packet}
For each source claim, compare the exact source excerpts with the expanded
PaperInterface specification. Mark up this PDF or fill the blank reviewer box in the TeX source;
return the annotated file for reconciliation into the paper-local human-review
log.

\section*{Open the interactive dashboard}
The interactive dashboard is an optional alternative to using this PDF; it
presents the same review surface rather than an additional review requirement.
After forking and cloning (or pulling) EconCSLib, run the following from the
repository root. The cache command is needed only the first time in a new clone.
\begin{ReviewVerbatim}
cd <your-EconCSLib-checkout>
lake exe cache get
python3 scripts/review_dashboard.py --paper KR21Monoculture --serve
\end{ReviewVerbatim}
Open \url{http://127.0.0.1:8765/} in a browser and keep that terminal running
while reviewing. The dashboard records saved annotations in this paper's
reviewer-owned review record.

\section*{Contents}
\begin{itemize}[leftmargin=1.2em]
\item \hyperlink{library-section-227ed3312974d38f}{Semantic library prerequisites (56; not paper claims)}
\begin{itemize}[leftmargin=1.2em]
\item \hyperlink{library-prerequisite-f8200dc56888f8cf}{EconCSLib.Probability.StrictlyWellOrderedNoise}
\item \hyperlink{library-prerequisite-3ca32c22c4f1786f}{EconCSLib.Probability.WeaklyWellOrderedNoise}
\item \hyperlink{library-prerequisite-a80eaaa76d4081cc}{EconCSLib.Probability.gaussianNoiseKernel}
\item \hyperlink{library-prerequisite-d1924ea126b64a99}{EconCSLib.Probability.gaussianVarianceFromStd}
\item \hyperlink{library-prerequisite-96b93cedc528228b}{EconCSLib.Probability.independentGaussianPairMeasureWithStd}
\item \hyperlink{library-prerequisite-0f700e1385fe3352}{EconCSLib.Probability.laplacianNoiseKernel}
\item \hyperlink{library-prerequisite-dfa732ba2fb3e45a}{EconCSLib.SocialChoice.Ranking.Candidate}
\item \hyperlink{library-prerequisite-8f403b33d88960b7}{EconCSLib.SocialChoice.Ranking.Ranking}
\item \hyperlink{library-prerequisite-7b06b37dd64e13d0}{EconCSLib.SocialChoice.Ranking.firstChoice}
\item \hyperlink{library-prerequisite-2f14612a40825584}{EconCSLib.SocialChoice.Ranking.secondChoice}
\item \hyperlink{library-prerequisite-90e8cd1cdd9f27a4}{EconCSLib.SocialChoice.Ranking.bestRemainingAfter}
\item \hyperlink{library-prerequisite-2a2a10773f9fb97f}{EconCSLib.SocialChoice.Ranking.secondMoverUtility}
\item \hyperlink{library-prerequisite-e5c7eedc40c11b7a}{EconCSLib.pmfExp}
\item \hyperlink{library-prerequisite-8c5104efac7b45e2}{EconCSLib.pmfPairExp}
\item \hyperlink{library-prerequisite-79ca8c9a62821d3f}{EconCSLib.SocialChoice.Ranking.expectedSecondMoverIndependent}
\item \hyperlink{library-prerequisite-67c0534cc3c0f0cd}{EconCSLib.SocialChoice.Ranking.expectedSecondMoverShared}
\item \hyperlink{library-prerequisite-a891c9f796ceb814}{EconCSLib.SocialChoice.Ranking.PrefersIndependentReranking}
\item \hyperlink{library-prerequisite-8db06726f0f2eba3}{EconCSLib.SocialChoice.Ranking.PrefersWeakerCompetition}
\item \hyperlink{library-prerequisite-9e722a7603183d67}{EconCSLib.SocialChoice.Ranking.RankingPair}
\item \hyperlink{library-prerequisite-ffa701659b5a9425}{EconCSLib.SocialChoice.Ranking.rankOf}
\item \hyperlink{library-prerequisite-150c4e692433d71a}{EconCSLib.SocialChoice.Ranking.StrictlyOrderedBy}
\item \hyperlink{library-prerequisite-8b25302a88abab0b}{EconCSLib.SocialChoice.Ranking.WeaklyOrderedBy}
\item \hyperlink{library-prerequisite-73e694b78a9c1a15}{EconCSLib.SocialChoice.Ranking.disagreementEvent}
\item \hyperlink{library-prerequisite-ee43c221688d5474}{EconCSLib.SocialChoice.Ranking.decidableDisagreementEvent}
\item \hyperlink{library-prerequisite-5875c858a5dee439}{EconCSLib.pmfPairProb}
\item \hyperlink{library-prerequisite-f084d92aa50ed2fc}{EconCSLib.SocialChoice.Ranking.disagreementProb}
\item \hyperlink{library-prerequisite-550af61d293e17c6}{EconCSLib.SocialChoice.Ranking.expectedFirstMoverUtility}
\item \hyperlink{library-prerequisite-7fec9dedef722fc8}{EconCSLib.pmfProb}
\item \hyperlink{library-prerequisite-60202fed8e52e293}{EconCSLib.SocialChoice.Ranking.firstChoiceProb}
\item \hyperlink{library-prerequisite-ce32083db4914ab9}{EconCSLib.SocialChoice.Ranking.firstChoiceMissProb}
\item \hyperlink{library-prerequisite-dab3c208272a262f}{EconCSLib.SocialChoice.Ranking.instMeasurableSpaceRanking}
\item \hyperlink{library-prerequisite-3b92294c8958cd6a}{EconCSLib.SocialChoice.Ranking.invertedPair}
\item \hyperlink{library-prerequisite-0bcb643e4e8f308c}{EconCSLib.SocialChoice.Ranking.inversionFinset}
\item \hyperlink{library-prerequisite-bba935eca4d406d8}{EconCSLib.SocialChoice.Ranking.kendallTau}
\item \hyperlink{library-prerequisite-02920becedba7938}{EconCSLib.SocialChoice.Ranking.rerankingGainOnPair}
\item \hyperlink{library-prerequisite-d3b381d476267785}{EconCSLib.SocialChoice.Ranking.pairRerankingGain}
\item \hyperlink{library-prerequisite-4d0316233c2913da}{EconCSLib.SocialChoice.Ranking.rankByScore}
\item \hyperlink{library-prerequisite-3d857727212e6236}{EconCSLib.SocialChoice.Ranking.rankingPMFOfMeasure}
\item \hyperlink{library-prerequisite-9282609ac289eae1}{EconCSLib.SocialChoice.Ranking.rum3Ranking012}
\item \hyperlink{library-prerequisite-3e9037b4cf76ac78}{EconCSLib.SocialChoice.Ranking.rum3Ranking021}
\item \hyperlink{library-prerequisite-d6840f3f435bf0a3}{EconCSLib.SocialChoice.Ranking.rum3Ranking102}
\item \hyperlink{library-prerequisite-17cdcd166ffcfa57}{EconCSLib.SocialChoice.Ranking.rum3Ranking120}
\item \hyperlink{library-prerequisite-88da1ce8c8de0862}{EconCSLib.SocialChoice.Ranking.rum3Ranking201}
\item \hyperlink{library-prerequisite-11023713299f4bb8}{EconCSLib.SocialChoice.Ranking.rum3Ranking210}
\item \hyperlink{library-prerequisite-2b57daef13f9e5a2}{EconCSLib.SocialChoice.Ranking.rum3RankByScores}
\item \hyperlink{library-prerequisite-d3a1fe6e17d1ebe8}{EconCSLib.finCons}
\item \hyperlink{library-prerequisite-4aa064b5cbd66ed3}{EconCSLib.finiteAvailableWeight}
\item \hyperlink{library-prerequisite-637b06ac0adfcac1}{EconCSLib.finiteFreshList}
\item \hyperlink{library-prerequisite-517099bc7eacc408}{EconCSLib.finiteFreshListCons}
\item \hyperlink{library-prerequisite-010f5218e34f7dbc}{EconCSLib.finiteFreshListNil}
\item \hyperlink{library-prerequisite-8598d8001b9c1939}{EconCSLib.finiteWeightedPMF}
\item \hyperlink{library-prerequisite-478178a796f1a098}{EconCSLib.finiteWeightedPMFAvailable}
\item \hyperlink{library-prerequisite-17cb3d8406f594aa}{EconCSLib.finiteWithoutReplacementPMF}
\item \hyperlink{library-prerequisite-7ae22d5411bd7792}{EconCSLib.measureProb}
\item \hyperlink{library-prerequisite-197097ef6d41e4d7}{EconCSLib.pmfPairIndicatorExp}
\item \hyperlink{library-prerequisite-f185c8e6394c9317}{EconCSLib.pmfProd}
\end{itemize}
\item \hyperlink{source-section-f10b5f68e4acf3dc}{Source claims (57)}
\begin{itemize}[leftmargin=1.2em]
\item \hyperlink{source-claim-70a4d7b84f526fcd}{source\_model\_candidates\_values\_and\_ranking\_lawSpec}
\item \hyperlink{source-claim-fbe3ef1f5678507e}{source\_definition1\_iffSpec}
\item \hyperlink{source-claim-413fc0a7e71398bc}{source\_equation1\_removal\_monotonicitySpec}
\item \hyperlink{source-claim-f4aa43d846eb3ae7}{source\_two\_firm\_selection\_game\_semanticsSpec}
\item \hyperlink{source-claim-fcc1dae9c26a2b31}{source\_model\_outer\_distribution\_probability\_and\_point\_massSpec}
\item \hyperlink{source-claim-a71947e08e3fd64a}{source\_definition2\_literal\_outer\_joint\_semantic\_completeSpec}
\item \hyperlink{source-claim-faa08c37fc761736}{source\_definition3\_literal\_outer\_semantic\_completeSpec}
\item \hyperlink{source-claim-f88d6727dfcb8617}{theorem1\_raw\_outer\_literal\_payoff\_terminal\_conclusionSpec}
\item \hyperlink{source-claim-0a74bbe90425faa9}{equation2\_outer\_joint\_payoff\_equivalence\_semantic\_completeSpec}
\item \hyperlink{source-claim-36b115610854f337}{equation3\_outer\_joint\_payoff\_identity\_semantic\_completeSpec}
\item \hyperlink{source-claim-ea5f4fa858fd5854}{equation4\_algorithm\_best\_response\_against\_algorithm\_iffSpec}
\item \hyperlink{source-claim-9d6309ccb42ce9df}{equation5\_from\_literal\_definition1\_removalSpec}
\item \hyperlink{source-claim-fa14b138c169a341}{equation6\_outer\_payoff\_indifference\_from\_source\_conditionsSpec}
\item \hyperlink{source-claim-bd2d52e2663b165c}{section31\_iidFinitePositiveVariance\_normalizationSpec}
\item \hyperlink{source-claim-cc8111af16950cc8}{theorem2\_gaussian\_laplace\_source\_semantic\_completeSpec}
\item \hyperlink{source-claim-8940f397fa4678ac}{equation7\_plackettLuce\_choice\_probabilitySpec}
\item \hyperlink{source-claim-0156e1eedd228cc4}{section31\_literalUnitVarianceGumbel\_outer\_zero\_effect\_semantic\_completeSpec}
\item \hyperlink{source-claim-8dd2c310e7ffc8a6}{equation8\_source\_concrete\_mallows\_probabilitySpec}
\item \hyperlink{source-claim-fd702590ac30f669}{theorem3\_source\_complete\_semantic\_completeSpec}
\item \hyperlink{source-claim-e42fc5562761542b}{theorem4\_source\_mallows\_outer\_horizon\_lt\_literal\_semantic\_completeSpec}
\item \hyperlink{source-claim-54f72d882fb14d8c}{appendixA\_theorem5\_corrected\_source\_completeSpec}
\item \hyperlink{source-claim-9835a680c6142d09}{equationA1\_source\_w11\_iid\_literal\_event\_conditionalTail\_integralSpec}
\item \hyperlink{source-claim-aa9b8f734e8009d3}{appendixB1\_source\_discrete\_definition2\_counterexample\_semantic\_completeSpec}
\item \hyperlink{source-claim-8c4914b3b2e7a592}{appendixB2\_source\_discrete\_definition3\_counterexample\_semantic\_completeSpec}
\item \hyperlink{source-claim-36e8cf848060682f}{equationB1\_counterexample\_first\_choice\_x1\_semantic\_completeSpec}
\item \hyperlink{source-claim-a3b878a82854afb8}{equationB2\_counterexample\_first\_choice\_x2\_semantic\_completeSpec}
\item \hyperlink{source-claim-34bc63da82092219}{appendixC1\_source\_pairwise\_full\_eventsSpec}
\item \hyperlink{source-claim-6a2a92d325b21a46}{appendixC2\_source\_pairwise\_full\_eventsSpec}
\item \hyperlink{source-claim-d11ceae944c696a3}{definition4\_strictlyWellOrderedNoise\_iffSpec}
\item \hyperlink{source-claim-231c9ba63e6b1424}{theorem6\_source\_raw\_iid\_density\_literal\_semantic\_completeSpec}
\item \hyperlink{source-claim-d7ed4659ca9b332e}{lemma1\_corrected\_gaussian\_laplace\_kernel\_targetSpec}
\item \hyperlink{source-claim-e077005933af7923}{lemma2\_source\_arbitraryFinite\_iid\_bottom\_first\_probability\_semantic\_completeSpec}
\item \hyperlink{source-claim-a48d69bf76046da3}{lemma3\_source\_arbitraryFinite\_iid\_delta\_le\_top\_delta\_semantic\_completeSpec}
\item \hyperlink{source-claim-aceb5e75cb75d8b8}{theorem7\_source\_semantic\_completeSpec}
\item \hyperlink{source-claim-fb611bb534f07b11}{equationC3\_laplace\_strict\_conditional\_ratio\_eq\_density\_cdf\_semantic\_completeSpec}
\item \hyperlink{source-claim-29d2b7baa5996f3d}{equationC4\_laplace\_case3\_expression\_posSpec}
\item \hyperlink{source-claim-33bdbb546218e99a}{theorem8\_source\_semantic\_completeSpec}
\item \hyperlink{source-claim-43364c7dea10aac0}{equationC5\_gaussian\_source\_semantic\_completeSpec}
\item \hyperlink{source-claim-c25f59d1a0e8be35}{equationC6\_gaussian\_reduced\_expression\_posSpec}
\item \hyperlink{source-claim-4d39c9a7a45156fa}{equationC7\_gaussian\_reduced\_expression\_tendsto\_atBot\_zeroSpec}
\item \hyperlink{source-claim-2c9d403cd319b538}{equationC8\_gaussian\_reduced\_expression\_hasDerivAtSpec}
\item \hyperlink{source-claim-4139a609895b0de6}{equationC8\_C9\_corrected\_gaussian\_targetSpec}
\item \hyperlink{source-claim-7a440699a84fb997}{equationC10\_gaussian\_g\_inv\_derivative\_lower\_boundSpec}
\item \hyperlink{source-claim-d732fbc16f2ecdeb}{theorem9\_source\_mallows\_definition1\_literal\_semantic\_completeSpec}
\item \hyperlink{source-claim-7751231944a77b74}{equationD1\_source\_phi\_likelihood\_ratio\_completeSpec}
\item \hyperlink{source-claim-c7560a9682e8906d}{equationE1\_source\_phi\_literal\_conditional\_semantic\_completeSpec}
\item \hyperlink{source-claim-73dd5d1917ec1a7c}{equationE2\_source\_phi\_literal\_conditional\_semantic\_completeSpec}
\item \hyperlink{source-claim-027748e89e133285}{equationE3\_source\_phi\_completeSpec}
\item \hyperlink{source-claim-f5654e8ee63c4d5d}{lemma4\_weighted\_average\_lt\_of\_cross\_ratioSpec}
\item \hyperlink{source-claim-7e7ac58d3e77e0cc}{equationF1\_source\_phi\_completeSpec}
\item \hyperlink{source-claim-7ef5e5254a6c73ab}{lemma6\_source\_phi\_rank\_power\_completeSpec}
\item \hyperlink{source-claim-4e3a949fd1ab09f3}{equationF2\_source\_phi\_closed\_form\_completeSpec}
\item \hyperlink{source-claim-1e4808a386b72297}{lemma7\_source\_phi\_completeSpec}
\item \hyperlink{source-claim-961880e0b1d1c772}{lemma8\_source\_mallows\_phi\_pairwise\_correct\_probability\_ltSpec}
\item \hyperlink{source-claim-2d4c1840ef47dce5}{section31\_corrected\_gumbel\_plackett\_luce\_targetSpec}
\item \hyperlink{source-claim-e12ec3045f35c0cc}{section31\_plackettLuce\_outer\_best\_available\_weakly\_dominatesSpec}
\item \hyperlink{source-claim-4cfc28215a7bcccf}{appendixB\_smoothing\_source\_coreSpec}
\end{itemize}
\end{itemize}

\clearpage
\hypertarget{library-section-227ed3312974d38f}{}
\section*{Material library prerequisites}
\hypertarget{library-prerequisite-f8200dc56888f8cf}{}
\subsection*{\small\ttfamily\raggedright EconCSLib.\allowbreak{}Probability.\allowbreak{}StrictlyWellOrderedNoise}
\textbf{Source connection:} no explicit byte-pinned paper source connection is registered.
\paragraph{Lean-expanded library semantic target}
\begin{ReviewVerbatim}
∀ (f : ℝ → ℝ) {{a b c d : ℝ}}, b < a → d < c → f (a - c) * f (b - d) > f (a - d) * f (b - c)
\end{ReviewVerbatim}

\paragraph{Exact Lean library declaration}
\begin{ReviewVerbatim}
def StrictlyWellOrderedNoise (f : ℝ → ℝ) : Prop :=
  ∀ {{a b c d : ℝ}}, b < a → d < c →
    f (a - c) * f (b - d) > f (a - d) * f (b - c)
\end{ReviewVerbatim}

\paragraph{Recorded source-to-library screening}
\textbf{Verdict:} \texttt{not recorded}
\reviewmatch{Matches source input}
\noindent\textbf{Reviewer annotation}\par
\reviewerbox
\clearpage
\hypertarget{library-prerequisite-3ca32c22c4f1786f}{}
\subsection*{\small\ttfamily\raggedright EconCSLib.\allowbreak{}Probability.\allowbreak{}WeaklyWellOrderedNoise}
\textbf{Source connection:} no explicit byte-pinned paper source connection is registered.
\paragraph{Lean-expanded library semantic target}
\begin{ReviewVerbatim}
∀ (f : ℝ → ℝ) {{a b c d : ℝ}}, b < a → d < c → f (a - d) * f (b - c) ≤ f (a - c) * f (b - d)
\end{ReviewVerbatim}

\paragraph{Exact Lean library declaration}
\begin{ReviewVerbatim}
def WeaklyWellOrderedNoise (f : ℝ → ℝ) : Prop :=
  ∀ {{a b c d : ℝ}}, b < a → d < c →
    f (a - d) * f (b - c) ≤ f (a - c) * f (b - d)
\end{ReviewVerbatim}

\paragraph{Recorded source-to-library screening}
\textbf{Verdict:} \texttt{not recorded}
\reviewmatch{Matches source input}
\noindent\textbf{Reviewer annotation}\par
\reviewerbox
\clearpage
\hypertarget{library-prerequisite-a80eaaa76d4081cc}{}
\subsection*{\small\ttfamily\raggedright EconCSLib.\allowbreak{}Probability.\allowbreak{}gaussianNoiseKernel}
\textbf{Source connection:} no explicit byte-pinned paper source connection is registered.
\paragraph{Lean-expanded library semantic target}
\begin{ReviewVerbatim}
(κ x : ℝ) → Real.exp (-κ * x ^ 2)
\end{ReviewVerbatim}

\paragraph{Exact Lean library declaration}
\begin{ReviewVerbatim}
def gaussianNoiseKernel (κ : ℝ) (x : ℝ) : ℝ :=
  Real.exp (-κ * x ^ 2)
\end{ReviewVerbatim}

\paragraph{Recorded source-to-library screening}
\textbf{Verdict:} \texttt{not recorded}
\reviewmatch{Matches source input}
\noindent\textbf{Reviewer annotation}\par
\reviewerbox
\clearpage
\hypertarget{library-prerequisite-d1924ea126b64a99}{}
\subsection*{\small\ttfamily\raggedright EconCSLib.\allowbreak{}Probability.\allowbreak{}gaussianVarianceFromStd}
\textbf{Source connection:} no explicit byte-pinned paper source connection is registered.
\paragraph{Lean-expanded library semantic target}
\begin{ReviewVerbatim}
(σ : ℝ) → NNReal.mk (σ ^ 2) ⋯
\end{ReviewVerbatim}

\paragraph{Exact Lean library declaration}
\begin{ReviewVerbatim}
def gaussianVarianceFromStd (σ : ℝ) : ℝ≥0 :=
  NNReal.mk (σ ^ 2) (sq_nonneg σ)
\end{ReviewVerbatim}

\paragraph{Recorded source-to-library screening}
\textbf{Verdict:} \texttt{not recorded}
\reviewmatch{Matches source input}
\noindent\textbf{Reviewer annotation}\par
\reviewerbox
\clearpage
\hypertarget{library-prerequisite-96b93cedc528228b}{}
\subsection*{\small\ttfamily\raggedright EconCSLib.\allowbreak{}Probability.\allowbreak{}independentGaussianPairMeasureWithStd}
\textbf{Source connection:} no explicit byte-pinned paper source connection is registered.
\paragraph{Lean-expanded library semantic target}
\begin{ReviewVerbatim}
(σ xi xj : ℝ) →
  (ProbabilityTheory.gaussianReal xi (EconCSLib.Probability.gaussianVarianceFromStd σ)).prod
    (ProbabilityTheory.gaussianReal xj (EconCSLib.Probability.gaussianVarianceFromStd σ))
\end{ReviewVerbatim}

\paragraph{Exact Lean library declaration}
\begin{ReviewVerbatim}
noncomputable def independentGaussianPairMeasureWithStd
    (σ xi xj : ℝ) : Measure (ℝ × ℝ) :=
  (ProbabilityTheory.gaussianReal xi (gaussianVarianceFromStd σ)).prod
    (ProbabilityTheory.gaussianReal xj (gaussianVarianceFromStd σ))
\end{ReviewVerbatim}

\paragraph{Recorded source-to-library screening}
\textbf{Verdict:} \texttt{not recorded}
\reviewmatch{Matches source input}
\noindent\textbf{Reviewer annotation}\par
\reviewerbox
\clearpage
\hypertarget{library-prerequisite-0f700e1385fe3352}{}
\subsection*{\small\ttfamily\raggedright EconCSLib.\allowbreak{}Probability.\allowbreak{}laplacianNoiseKernel}
\textbf{Source connection:} no explicit byte-pinned paper source connection is registered.
\paragraph{Lean-expanded library semantic target}
\begin{ReviewVerbatim}
(lam x : ℝ) → Real.exp (-lam * |x|)
\end{ReviewVerbatim}

\paragraph{Exact Lean library declaration}
\begin{ReviewVerbatim}
def laplacianNoiseKernel (lam : ℝ) (x : ℝ) : ℝ :=
  Real.exp (-lam * |x|)
\end{ReviewVerbatim}

\paragraph{Recorded source-to-library screening}
\textbf{Verdict:} \texttt{not recorded}
\reviewmatch{Matches source input}
\noindent\textbf{Reviewer annotation}\par
\reviewerbox
\clearpage
\hypertarget{library-prerequisite-dfa732ba2fb3e45a}{}
\subsection*{\small\ttfamily\raggedright EconCSLib.\allowbreak{}SocialChoice.\allowbreak{}Ranking.\allowbreak{}Candidate}
\textbf{Source connection:} no explicit byte-pinned paper source connection is registered.
\paragraph{Lean-expanded library semantic target}
\begin{ReviewVerbatim}
(n : ℕ) → Fin (n + 2)
\end{ReviewVerbatim}

\paragraph{Exact Lean library declaration}
\begin{ReviewVerbatim}
abbrev Candidate (n : ℕ) := Fin (n + 2)
\end{ReviewVerbatim}

\paragraph{Recorded source-to-library screening}
\textbf{Verdict:} \texttt{not recorded}
\reviewmatch{Matches source input}
\noindent\textbf{Reviewer annotation}\par
\reviewerbox
\clearpage
\hypertarget{library-prerequisite-8f403b33d88960b7}{}
\subsection*{\small\ttfamily\raggedright EconCSLib.\allowbreak{}SocialChoice.\allowbreak{}Ranking.\allowbreak{}Ranking}
\textbf{Source connection:} no explicit byte-pinned paper source connection is registered.
\paragraph{Lean-expanded library semantic target}
\begin{ReviewVerbatim}
(n : ℕ) → Equiv.Perm (EconCSLib.SocialChoice.Ranking.Candidate n)
\end{ReviewVerbatim}

\paragraph{Exact Lean library declaration}
\begin{ReviewVerbatim}
abbrev Ranking (n : ℕ) := Equiv.Perm (Candidate n)
\end{ReviewVerbatim}

\paragraph{Recorded source-to-library screening}
\textbf{Verdict:} \texttt{not recorded}
\reviewmatch{Matches source input}
\noindent\textbf{Reviewer annotation}\par
\reviewerbox
\clearpage
\hypertarget{library-prerequisite-7b06b37dd64e13d0}{}
\subsection*{\small\ttfamily\raggedright EconCSLib.\allowbreak{}SocialChoice.\allowbreak{}Ranking.\allowbreak{}firstChoice}
\textbf{Source connection:} no explicit byte-pinned paper source connection is registered.
\paragraph{Lean-expanded library semantic target}
\begin{ReviewVerbatim}
{n : ℕ} → (π : EconCSLib.SocialChoice.Ranking.Ranking n) → π 0
\end{ReviewVerbatim}

\paragraph{Exact Lean library declaration}
\begin{ReviewVerbatim}
def firstChoice {n : ℕ} (π : Ranking n) : Candidate n :=
  π 0
\end{ReviewVerbatim}

\paragraph{Recorded source-to-library screening}
\textbf{Verdict:} \texttt{not recorded}
\reviewmatch{Matches source input}
\noindent\textbf{Reviewer annotation}\par
\reviewerbox
\clearpage
\hypertarget{library-prerequisite-2f14612a40825584}{}
\subsection*{\small\ttfamily\raggedright EconCSLib.\allowbreak{}SocialChoice.\allowbreak{}Ranking.\allowbreak{}secondChoice}
\textbf{Source connection:} no explicit byte-pinned paper source connection is registered.
\paragraph{Lean-expanded library semantic target}
\begin{ReviewVerbatim}
{n : ℕ} → (π : EconCSLib.SocialChoice.Ranking.Ranking n) → π 1
\end{ReviewVerbatim}

\paragraph{Exact Lean library declaration}
\begin{ReviewVerbatim}
def secondChoice {n : ℕ} (π : Ranking n) : Candidate n :=
  π 1
\end{ReviewVerbatim}

\paragraph{Recorded source-to-library screening}
\textbf{Verdict:} \texttt{not recorded}
\reviewmatch{Matches source input}
\noindent\textbf{Reviewer annotation}\par
\reviewerbox
\clearpage
\hypertarget{library-prerequisite-90e8cd1cdd9f27a4}{}
\subsection*{\small\ttfamily\raggedright EconCSLib.\allowbreak{}SocialChoice.\allowbreak{}Ranking.\allowbreak{}bestRemainingAfter}
\textbf{Source connection:} no explicit byte-pinned paper source connection is registered.
\paragraph{Lean-expanded library semantic target}
\begin{ReviewVerbatim}
{n : ℕ} →
  (π : EconCSLib.SocialChoice.Ranking.Ranking n) →
    (c : EconCSLib.SocialChoice.Ranking.Candidate n) →
      if EconCSLib.SocialChoice.Ranking.firstChoice π = c then EconCSLib.SocialChoice.Ranking.secondChoice π
      else EconCSLib.SocialChoice.Ranking.firstChoice π
\end{ReviewVerbatim}

\paragraph{Exact Lean library declaration}
\begin{ReviewVerbatim}
def bestRemainingAfter {n : ℕ} (π : Ranking n) (c : Candidate n) : Candidate n :=
  if firstChoice π = c then secondChoice π else firstChoice π
\end{ReviewVerbatim}

\paragraph{Recorded source-to-library screening}
\textbf{Verdict:} \texttt{not recorded}
\reviewmatch{Matches source input}
\noindent\textbf{Reviewer annotation}\par
\reviewerbox
\clearpage
\hypertarget{library-prerequisite-2a2a10773f9fb97f}{}
\subsection*{\small\ttfamily\raggedright EconCSLib.\allowbreak{}SocialChoice.\allowbreak{}Ranking.\allowbreak{}secondMoverUtility}
\textbf{Source connection:} no explicit byte-pinned paper source connection is registered.
\paragraph{Lean-expanded library semantic target}
\begin{ReviewVerbatim}
{n : ℕ} →
  (value : EconCSLib.SocialChoice.Ranking.Candidate n → ℝ) →
    (π σ : EconCSLib.SocialChoice.Ranking.Ranking n) →
      value (EconCSLib.SocialChoice.Ranking.bestRemainingAfter π (EconCSLib.SocialChoice.Ranking.firstChoice σ))
\end{ReviewVerbatim}

\paragraph{Exact Lean library declaration}
\begin{ReviewVerbatim}
def secondMoverUtility {n : ℕ} (value : Candidate n → ℝ)
    (π σ : Ranking n) : ℝ :=
  value (bestRemainingAfter π (firstChoice σ))
\end{ReviewVerbatim}

\paragraph{Recorded source-to-library screening}
\textbf{Verdict:} \texttt{not recorded}
\reviewmatch{Matches source input}
\noindent\textbf{Reviewer annotation}\par
\reviewerbox
\clearpage
\hypertarget{library-prerequisite-e5c7eedc40c11b7a}{}
\subsection*{\small\ttfamily\raggedright EconCSLib.\allowbreak{}pmfExp}
\textbf{Source connection:} no explicit byte-pinned paper source connection is registered.
\paragraph{Lean-expanded library semantic target}
\begin{ReviewVerbatim}
{α : Type u_1} → [inst : Fintype α] → [DecidableEq α] → (μ : PMF α) → (f : α → ℝ) → ∑ a, (μ a).toReal * f a
\end{ReviewVerbatim}

\paragraph{Exact Lean library declaration}
\begin{ReviewVerbatim}
noncomputable def pmfExp {α : Type*} [Fintype α] [DecidableEq α]
    (μ : PMF α) (f : α → ℝ) : ℝ :=
  ∑ a, (μ a).toReal * f a
\end{ReviewVerbatim}

\paragraph{Recorded source-to-library screening}
\textbf{Verdict:} \texttt{not recorded}
\reviewmatch{Matches source input}
\noindent\textbf{Reviewer annotation}\par
\reviewerbox
\clearpage
\hypertarget{library-prerequisite-8c5104efac7b45e2}{}
\subsection*{\small\ttfamily\raggedright EconCSLib.\allowbreak{}pmfPairExp}
\textbf{Source connection:} no explicit byte-pinned paper source connection is registered.
\paragraph{Lean-expanded library semantic target}
\begin{ReviewVerbatim}
{α : Type u_1} →
  {β : Type u_2} →
    [inst : Fintype α] →
      [inst_1 : DecidableEq α] →
        [inst_2 : Fintype β] →
          [inst_3 : DecidableEq β] →
            (μ : PMF α) → (ν : PMF β) → (f : α → β → ℝ) → EconCSLib.pmfExp μ fun a => EconCSLib.pmfExp ν fun b => f a b
\end{ReviewVerbatim}

\paragraph{Exact Lean library declaration}
\begin{ReviewVerbatim}
noncomputable def pmfPairExp {α β : Type*}
    [Fintype α] [DecidableEq α] [Fintype β] [DecidableEq β]
    (μ : PMF α) (ν : PMF β) (f : α → β → ℝ) : ℝ :=
  pmfExp μ (fun a => pmfExp ν (fun b => f a b))
\end{ReviewVerbatim}

\paragraph{Recorded source-to-library screening}
\textbf{Verdict:} \texttt{not recorded}
\reviewmatch{Matches source input}
\noindent\textbf{Reviewer annotation}\par
\reviewerbox
\clearpage
\hypertarget{library-prerequisite-79ca8c9a62821d3f}{}
\subsection*{\small\ttfamily\raggedright EconCSLib.\allowbreak{}SocialChoice.\allowbreak{}Ranking.\allowbreak{}expectedSecondMoverIndependent}
\textbf{Source connection:} no explicit byte-pinned paper source connection is registered.
\paragraph{Lean-expanded library semantic target}
\begin{ReviewVerbatim}
{n : ℕ} →
  (μ₂ μ₁ : PMF (EconCSLib.SocialChoice.Ranking.Ranking n)) →
    (value : EconCSLib.SocialChoice.Ranking.Candidate n → ℝ) →
      EconCSLib.pmfPairExp μ₂ μ₁ fun π σ => EconCSLib.SocialChoice.Ranking.secondMoverUtility value π σ
\end{ReviewVerbatim}

\paragraph{Exact Lean library declaration}
\begin{ReviewVerbatim}
def expectedSecondMoverIndependent {n : ℕ}
    (μ₂ μ₁ : PMF (Ranking n)) (value : Candidate n → ℝ) : ℝ :=
  pmfPairExp μ₂ μ₁ (fun π σ => secondMoverUtility value π σ)
\end{ReviewVerbatim}

\paragraph{Recorded source-to-library screening}
\textbf{Verdict:} \texttt{not recorded}
\reviewmatch{Matches source input}
\noindent\textbf{Reviewer annotation}\par
\reviewerbox
\clearpage
\hypertarget{library-prerequisite-67c0534cc3c0f0cd}{}
\subsection*{\small\ttfamily\raggedright EconCSLib.\allowbreak{}SocialChoice.\allowbreak{}Ranking.\allowbreak{}expectedSecondMoverShared}
\textbf{Source connection:} no explicit byte-pinned paper source connection is registered.
\paragraph{Lean-expanded library semantic target}
\begin{ReviewVerbatim}
{n : ℕ} →
  (μ : PMF (EconCSLib.SocialChoice.Ranking.Ranking n)) →
    (value : EconCSLib.SocialChoice.Ranking.Candidate n → ℝ) →
      EconCSLib.pmfExp μ fun π => value (EconCSLib.SocialChoice.Ranking.secondChoice π)
\end{ReviewVerbatim}

\paragraph{Exact Lean library declaration}
\begin{ReviewVerbatim}
def expectedSecondMoverShared {n : ℕ}
    (μ : PMF (Ranking n)) (value : Candidate n → ℝ) : ℝ :=
  pmfExp μ (fun π => value (secondChoice π))
\end{ReviewVerbatim}

\paragraph{Recorded source-to-library screening}
\textbf{Verdict:} \texttt{not recorded}
\reviewmatch{Matches source input}
\noindent\textbf{Reviewer annotation}\par
\reviewerbox
\clearpage
\hypertarget{library-prerequisite-a891c9f796ceb814}{}
\subsection*{\small\ttfamily\raggedright EconCSLib.\allowbreak{}SocialChoice.\allowbreak{}Ranking.\allowbreak{}PrefersIndependentReranking}
\textbf{Source connection:} no explicit byte-pinned paper source connection is registered.
\paragraph{Lean-expanded library semantic target}
\begin{ReviewVerbatim}
∀ {n : ℕ} (μ : PMF (EconCSLib.SocialChoice.Ranking.Ranking n)) (value : EconCSLib.SocialChoice.Ranking.Candidate n → ℝ),
  EconCSLib.SocialChoice.Ranking.expectedSecondMoverShared μ value <
    EconCSLib.SocialChoice.Ranking.expectedSecondMoverIndependent μ μ value
\end{ReviewVerbatim}

\paragraph{Exact Lean library declaration}
\begin{ReviewVerbatim}
def PrefersIndependentReranking {n : ℕ}
    (μ : PMF (Ranking n)) (value : Candidate n → ℝ) : Prop :=
  expectedSecondMoverShared μ value < expectedSecondMoverIndependent μ μ value
\end{ReviewVerbatim}

\paragraph{Recorded source-to-library screening}
\textbf{Verdict:} \texttt{not recorded}
\reviewmatch{Matches source input}
\noindent\textbf{Reviewer annotation}\par
\reviewerbox
\clearpage
\hypertarget{library-prerequisite-8db06726f0f2eba3}{}
\subsection*{\small\ttfamily\raggedright EconCSLib.\allowbreak{}SocialChoice.\allowbreak{}Ranking.\allowbreak{}PrefersWeakerCompetition}
\textbf{Source connection:} no explicit byte-pinned paper source connection is registered.
\paragraph{Lean-expanded library semantic target}
\begin{ReviewVerbatim}
∀ {n : ℕ} (μBetter μWorse : PMF (EconCSLib.SocialChoice.Ranking.Ranking n))
  (value : EconCSLib.SocialChoice.Ranking.Candidate n → ℝ),
  EconCSLib.SocialChoice.Ranking.expectedSecondMoverIndependent μWorse μBetter value <
    EconCSLib.SocialChoice.Ranking.expectedSecondMoverIndependent μWorse μWorse value
\end{ReviewVerbatim}

\paragraph{Exact Lean library declaration}
\begin{ReviewVerbatim}
def PrefersWeakerCompetition {n : ℕ}
    (μBetter μWorse : PMF (Ranking n)) (value : Candidate n → ℝ) : Prop :=
  expectedSecondMoverIndependent μWorse μBetter value <
    expectedSecondMoverIndependent μWorse μWorse value
\end{ReviewVerbatim}

\paragraph{Recorded source-to-library screening}
\textbf{Verdict:} \texttt{not recorded}
\reviewmatch{Matches source input}
\noindent\textbf{Reviewer annotation}\par
\reviewerbox
\clearpage
\hypertarget{library-prerequisite-9e722a7603183d67}{}
\subsection*{\small\ttfamily\raggedright EconCSLib.\allowbreak{}SocialChoice.\allowbreak{}Ranking.\allowbreak{}RankingPair}
\textbf{Source connection:} no explicit byte-pinned paper source connection is registered.
\paragraph{Lean-expanded library semantic target}
\begin{ReviewVerbatim}
(n : ℕ) → EconCSLib.SocialChoice.Ranking.Ranking n × EconCSLib.SocialChoice.Ranking.Ranking n
\end{ReviewVerbatim}

\paragraph{Exact Lean library declaration}
\begin{ReviewVerbatim}
abbrev RankingPair (n : ℕ) := Ranking n × Ranking n
\end{ReviewVerbatim}

\paragraph{Recorded source-to-library screening}
\textbf{Verdict:} \texttt{not recorded}
\reviewmatch{Matches source input}
\noindent\textbf{Reviewer annotation}\par
\reviewerbox
\clearpage
\hypertarget{library-prerequisite-ffa701659b5a9425}{}
\subsection*{\small\ttfamily\raggedright EconCSLib.\allowbreak{}SocialChoice.\allowbreak{}Ranking.\allowbreak{}rankOf}
\textbf{Source connection:} no explicit byte-pinned paper source connection is registered.
\paragraph{Lean-expanded library semantic target}
\begin{ReviewVerbatim}
{n : ℕ} →
  (π : EconCSLib.SocialChoice.Ranking.Ranking n) → (c : EconCSLib.SocialChoice.Ranking.Candidate n) → (Equiv.symm π) c
\end{ReviewVerbatim}

\paragraph{Exact Lean library declaration}
\begin{ReviewVerbatim}
def rankOf {n : ℕ} (π : Ranking n) (c : Candidate n) : Candidate n :=
  π.symm c
\end{ReviewVerbatim}

\paragraph{Recorded source-to-library screening}
\textbf{Verdict:} \texttt{not recorded}
\reviewmatch{Matches source input}
\noindent\textbf{Reviewer annotation}\par
\reviewerbox
\clearpage
\hypertarget{library-prerequisite-150c4e692433d71a}{}
\subsection*{\small\ttfamily\raggedright EconCSLib.\allowbreak{}SocialChoice.\allowbreak{}Ranking.\allowbreak{}StrictlyOrderedBy}
\textbf{Source connection:} no explicit byte-pinned paper source connection is registered.
\paragraph{Lean-expanded library semantic target}
\begin{ReviewVerbatim}
∀ {n : ℕ} (ρ : EconCSLib.SocialChoice.Ranking.Ranking n) (value : EconCSLib.SocialChoice.Ranking.Candidate n → ℝ)
  {a b : EconCSLib.SocialChoice.Ranking.Candidate n},
  EconCSLib.SocialChoice.Ranking.rankOf ρ a < EconCSLib.SocialChoice.Ranking.rankOf ρ b → value b < value a
\end{ReviewVerbatim}

\paragraph{Exact Lean library declaration}
\begin{ReviewVerbatim}
def StrictlyOrderedBy {n : ℕ} (ρ : Ranking n) (value : Candidate n → ℝ) : Prop :=
  ∀ {a b : Candidate n}, rankOf ρ a < rankOf ρ b → value b < value a
\end{ReviewVerbatim}

\paragraph{Recorded source-to-library screening}
\textbf{Verdict:} \texttt{not recorded}
\reviewmatch{Matches source input}
\noindent\textbf{Reviewer annotation}\par
\reviewerbox
\clearpage
\hypertarget{library-prerequisite-8b25302a88abab0b}{}
\subsection*{\small\ttfamily\raggedright EconCSLib.\allowbreak{}SocialChoice.\allowbreak{}Ranking.\allowbreak{}WeaklyOrderedBy}
\textbf{Source connection:} no explicit byte-pinned paper source connection is registered.
\paragraph{Lean-expanded library semantic target}
\begin{ReviewVerbatim}
∀ {n : ℕ} (ρ : EconCSLib.SocialChoice.Ranking.Ranking n) (value : EconCSLib.SocialChoice.Ranking.Candidate n → ℝ)
  {a b : EconCSLib.SocialChoice.Ranking.Candidate n},
  EconCSLib.SocialChoice.Ranking.rankOf ρ a < EconCSLib.SocialChoice.Ranking.rankOf ρ b → value b ≤ value a
\end{ReviewVerbatim}

\paragraph{Exact Lean library declaration}
\begin{ReviewVerbatim}
def WeaklyOrderedBy {n : ℕ} (ρ : Ranking n) (value : Candidate n → ℝ) : Prop :=
  ∀ {a b : Candidate n}, rankOf ρ a < rankOf ρ b → value b ≤ value a
\end{ReviewVerbatim}

\paragraph{Recorded source-to-library screening}
\textbf{Verdict:} \texttt{not recorded}
\reviewmatch{Matches source input}
\noindent\textbf{Reviewer annotation}\par
\reviewerbox
\clearpage
\hypertarget{library-prerequisite-73e694b78a9c1a15}{}
\subsection*{\small\ttfamily\raggedright EconCSLib.\allowbreak{}SocialChoice.\allowbreak{}Ranking.\allowbreak{}disagreementEvent}
\textbf{Source connection:} no explicit byte-pinned paper source connection is registered.
\paragraph{Lean-expanded library semantic target}
\begin{ReviewVerbatim}
∀ {n : ℕ} (a : EconCSLib.SocialChoice.Ranking.RankingPair n),
  EconCSLib.SocialChoice.Ranking.firstChoice a.1 ≠ EconCSLib.SocialChoice.Ranking.firstChoice a.2
\end{ReviewVerbatim}

\paragraph{Exact Lean library declaration}
\begin{ReviewVerbatim}
def disagreementEvent {n : ℕ} : RankingPair n → Prop :=
  fun p => firstChoice p.1 ≠ firstChoice p.2
\end{ReviewVerbatim}

\paragraph{Recorded source-to-library screening}
\textbf{Verdict:} \texttt{not recorded}
\reviewmatch{Matches source input}
\noindent\textbf{Reviewer annotation}\par
\reviewerbox
\clearpage
\hypertarget{library-prerequisite-ee43c221688d5474}{}
\subsection*{\small\ttfamily\raggedright EconCSLib.\allowbreak{}SocialChoice.\allowbreak{}Ranking.\allowbreak{}decidableDisagreementEvent}
\textbf{Source connection:} no explicit byte-pinned paper source connection is registered.
\paragraph{Lean-expanded library semantic target}
\begin{ReviewVerbatim}
{n : ℕ} → (a : EconCSLib.SocialChoice.Ranking.RankingPair n) → id inferInstance
\end{ReviewVerbatim}

\paragraph{Exact Lean library declaration}
\begin{ReviewVerbatim}
library declaration is not registered or discoverable
\end{ReviewVerbatim}

\paragraph{Recorded source-to-library screening}
\textbf{Verdict:} \texttt{not recorded}
\reviewmatch{Matches source input}
\noindent\textbf{Reviewer annotation}\par
\reviewerbox
\clearpage
\hypertarget{library-prerequisite-5875c858a5dee439}{}
\subsection*{\small\ttfamily\raggedright EconCSLib.\allowbreak{}pmfPairProb}
\textbf{Source connection:} no explicit byte-pinned paper source connection is registered.
\paragraph{Lean-expanded library semantic target}
\begin{ReviewVerbatim}
{α : Type u_1} →
  {β : Type u_2} →
    [inst : Fintype α] →
      [inst_1 : DecidableEq α] →
        [inst_2 : Fintype β] →
          [inst_3 : DecidableEq β] →
            (μ : PMF α) →
              (ν : PMF β) →
                (p : α × β → Prop) →
                  [inst_4 : DecidablePred p] → EconCSLib.pmfPairExp μ ν fun a b => if p (a, b) then 1 else 0
\end{ReviewVerbatim}

\paragraph{Exact Lean library declaration}
\begin{ReviewVerbatim}
noncomputable def pmfPairProb {α β : Type*}
    [Fintype α] [DecidableEq α] [Fintype β] [DecidableEq β]
    (μ : PMF α) (ν : PMF β) (p : α × β → Prop) [DecidablePred p] : ℝ :=
  pmfPairExp μ ν (fun a b => if p (a, b) then 1 else 0)
\end{ReviewVerbatim}

\paragraph{Recorded source-to-library screening}
\textbf{Verdict:} \texttt{not recorded}
\reviewmatch{Matches source input}
\noindent\textbf{Reviewer annotation}\par
\reviewerbox
\clearpage
\hypertarget{library-prerequisite-f084d92aa50ed2fc}{}
\subsection*{\small\ttfamily\raggedright EconCSLib.\allowbreak{}SocialChoice.\allowbreak{}Ranking.\allowbreak{}disagreementProb}
\textbf{Source connection:} no explicit byte-pinned paper source connection is registered.
\paragraph{Lean-expanded library semantic target}
\begin{ReviewVerbatim}
{n : ℕ} →
  (μ : PMF (EconCSLib.SocialChoice.Ranking.Ranking n)) →
    EconCSLib.pmfPairProb μ μ EconCSLib.SocialChoice.Ranking.disagreementEvent
\end{ReviewVerbatim}

\paragraph{Exact Lean library declaration}
\begin{ReviewVerbatim}
def disagreementProb {n : ℕ} (μ : PMF (Ranking n)) : ℝ :=
  pmfPairProb μ μ disagreementEvent
\end{ReviewVerbatim}

\paragraph{Recorded source-to-library screening}
\textbf{Verdict:} \texttt{not recorded}
\reviewmatch{Matches source input}
\noindent\textbf{Reviewer annotation}\par
\reviewerbox
\clearpage
\hypertarget{library-prerequisite-550af61d293e17c6}{}
\subsection*{\small\ttfamily\raggedright EconCSLib.\allowbreak{}SocialChoice.\allowbreak{}Ranking.\allowbreak{}expectedFirstMoverUtility}
\textbf{Source connection:} no explicit byte-pinned paper source connection is registered.
\paragraph{Lean-expanded library semantic target}
\begin{ReviewVerbatim}
{n : ℕ} →
  (μ : PMF (EconCSLib.SocialChoice.Ranking.Ranking n)) →
    (value : EconCSLib.SocialChoice.Ranking.Candidate n → ℝ) →
      EconCSLib.pmfExp μ fun π => value (EconCSLib.SocialChoice.Ranking.firstChoice π)
\end{ReviewVerbatim}

\paragraph{Exact Lean library declaration}
\begin{ReviewVerbatim}
def expectedFirstMoverUtility {n : ℕ}
    (μ : PMF (Ranking n)) (value : Candidate n → ℝ) : ℝ :=
  pmfExp μ (fun π => value (firstChoice π))
\end{ReviewVerbatim}

\paragraph{Recorded source-to-library screening}
\textbf{Verdict:} \texttt{not recorded}
\reviewmatch{Matches source input}
\noindent\textbf{Reviewer annotation}\par
\reviewerbox
\clearpage
\hypertarget{library-prerequisite-7fec9dedef722fc8}{}
\subsection*{\small\ttfamily\raggedright EconCSLib.\allowbreak{}pmfProb}
\textbf{Source connection:} no explicit byte-pinned paper source connection is registered.
\paragraph{Lean-expanded library semantic target}
\begin{ReviewVerbatim}
{α : Type u_1} →
  [inst : Fintype α] →
    [inst_1 : DecidableEq α] →
      (μ : PMF α) → (p : α → Prop) → [inst_2 : DecidablePred p] → EconCSLib.pmfExp μ fun a => if p a then 1 else 0
\end{ReviewVerbatim}

\paragraph{Exact Lean library declaration}
\begin{ReviewVerbatim}
noncomputable def pmfProb {α : Type*} [Fintype α] [DecidableEq α]
    (μ : PMF α) (p : α → Prop) [DecidablePred p] : ℝ :=
  pmfExp μ (fun a => if p a then 1 else 0)
\end{ReviewVerbatim}

\paragraph{Recorded source-to-library screening}
\textbf{Verdict:} \texttt{not recorded}
\reviewmatch{Matches source input}
\noindent\textbf{Reviewer annotation}\par
\reviewerbox
\clearpage
\hypertarget{library-prerequisite-60202fed8e52e293}{}
\subsection*{\small\ttfamily\raggedright EconCSLib.\allowbreak{}SocialChoice.\allowbreak{}Ranking.\allowbreak{}firstChoiceProb}
\textbf{Source connection:} no explicit byte-pinned paper source connection is registered.
\paragraph{Lean-expanded library semantic target}
\begin{ReviewVerbatim}
{n : ℕ} →
  (μ : PMF (EconCSLib.SocialChoice.Ranking.Ranking n)) →
    (c : EconCSLib.SocialChoice.Ranking.Candidate n) →
      EconCSLib.pmfProb μ fun π => c = EconCSLib.SocialChoice.Ranking.firstChoice π
\end{ReviewVerbatim}

\paragraph{Exact Lean library declaration}
\begin{ReviewVerbatim}
def firstChoiceProb {n : ℕ}
    (μ : PMF (Ranking n)) (c : Candidate n) : ℝ :=
  EconCSLib.pmfProb μ (fun π => c = firstChoice π)
\end{ReviewVerbatim}

\paragraph{Recorded source-to-library screening}
\textbf{Verdict:} \texttt{not recorded}
\reviewmatch{Matches source input}
\noindent\textbf{Reviewer annotation}\par
\reviewerbox
\clearpage
\hypertarget{library-prerequisite-ce32083db4914ab9}{}
\subsection*{\small\ttfamily\raggedright EconCSLib.\allowbreak{}SocialChoice.\allowbreak{}Ranking.\allowbreak{}firstChoiceMissProb}
\textbf{Source connection:} no explicit byte-pinned paper source connection is registered.
\paragraph{Lean-expanded library semantic target}
\begin{ReviewVerbatim}
{n : ℕ} →
  (μ : PMF (EconCSLib.SocialChoice.Ranking.Ranking n)) →
    (c : EconCSLib.SocialChoice.Ranking.Candidate n) → 1 - EconCSLib.SocialChoice.Ranking.firstChoiceProb μ c
\end{ReviewVerbatim}

\paragraph{Exact Lean library declaration}
\begin{ReviewVerbatim}
def firstChoiceMissProb {n : ℕ}
    (μ : PMF (Ranking n)) (c : Candidate n) : ℝ :=
  1 - firstChoiceProb μ c
\end{ReviewVerbatim}

\paragraph{Recorded source-to-library screening}
\textbf{Verdict:} \texttt{not recorded}
\reviewmatch{Matches source input}
\noindent\textbf{Reviewer annotation}\par
\reviewerbox
\clearpage
\hypertarget{library-prerequisite-dab3c208272a262f}{}
\subsection*{\small\ttfamily\raggedright EconCSLib.\allowbreak{}SocialChoice.\allowbreak{}Ranking.\allowbreak{}instMeasurableSpaceRanking}
\textbf{Source connection:} no explicit byte-pinned paper source connection is registered.
\paragraph{Lean-expanded library semantic target}
\begin{ReviewVerbatim}
(n : ℕ) → ⊤
\end{ReviewVerbatim}

\paragraph{Exact Lean library declaration}
\begin{ReviewVerbatim}
library declaration is not registered or discoverable
\end{ReviewVerbatim}

\paragraph{Recorded source-to-library screening}
\textbf{Verdict:} \texttt{not recorded}
\reviewmatch{Matches source input}
\noindent\textbf{Reviewer annotation}\par
\reviewerbox
\clearpage
\hypertarget{library-prerequisite-3b92294c8958cd6a}{}
\subsection*{\small\ttfamily\raggedright EconCSLib.\allowbreak{}SocialChoice.\allowbreak{}Ranking.\allowbreak{}invertedPair}
\textbf{Source connection:} no explicit byte-pinned paper source connection is registered.
\paragraph{Lean-expanded library semantic target}
\begin{ReviewVerbatim}
∀ {n : ℕ} (ρ π : EconCSLib.SocialChoice.Ranking.Ranking n)
  (ab : EconCSLib.SocialChoice.Ranking.Candidate n × EconCSLib.SocialChoice.Ranking.Candidate n),
  EconCSLib.SocialChoice.Ranking.rankOf ρ ab.1 < EconCSLib.SocialChoice.Ranking.rankOf ρ ab.2 ∧
    EconCSLib.SocialChoice.Ranking.rankOf π ab.2 < EconCSLib.SocialChoice.Ranking.rankOf π ab.1
\end{ReviewVerbatim}

\paragraph{Exact Lean library declaration}
\begin{ReviewVerbatim}
def invertedPair {n : ℕ} (ρ π : Ranking n) (ab : Candidate n × Candidate n) : Prop :=
  rankOf ρ ab.1 < rankOf ρ ab.2 ∧ rankOf π ab.2 < rankOf π ab.1
\end{ReviewVerbatim}

\paragraph{Recorded source-to-library screening}
\textbf{Verdict:} \texttt{not recorded}
\reviewmatch{Matches source input}
\noindent\textbf{Reviewer annotation}\par
\reviewerbox
\clearpage
\hypertarget{library-prerequisite-0bcb643e4e8f308c}{}
\subsection*{\small\ttfamily\raggedright EconCSLib.\allowbreak{}SocialChoice.\allowbreak{}Ranking.\allowbreak{}inversionFinset}
\textbf{Source connection:} no explicit byte-pinned paper source connection is registered.
\paragraph{Lean-expanded library semantic target}
\begin{ReviewVerbatim}
{n : ℕ} → (ρ π : EconCSLib.SocialChoice.Ranking.Ranking n) → {ab | EconCSLib.SocialChoice.Ranking.invertedPair ρ π ab}
\end{ReviewVerbatim}

\paragraph{Exact Lean library declaration}
\begin{ReviewVerbatim}
noncomputable def inversionFinset {n : ℕ} (ρ π : Ranking n) :
    Finset (Candidate n × Candidate n) := by
  classical
  exact Finset.univ.filter (fun ab => invertedPair ρ π ab)
\end{ReviewVerbatim}

\paragraph{Recorded source-to-library screening}
\textbf{Verdict:} \texttt{not recorded}
\reviewmatch{Matches source input}
\noindent\textbf{Reviewer annotation}\par
\reviewerbox
\clearpage
\hypertarget{library-prerequisite-bba935eca4d406d8}{}
\subsection*{\small\ttfamily\raggedright EconCSLib.\allowbreak{}SocialChoice.\allowbreak{}Ranking.\allowbreak{}kendallTau}
\textbf{Source connection:} no explicit byte-pinned paper source connection is registered.
\paragraph{Lean-expanded library semantic target}
\begin{ReviewVerbatim}
{n : ℕ} → (ρ π : EconCSLib.SocialChoice.Ranking.Ranking n) → (EconCSLib.SocialChoice.Ranking.inversionFinset ρ π).card
\end{ReviewVerbatim}

\paragraph{Exact Lean library declaration}
\begin{ReviewVerbatim}
noncomputable def kendallTau {n : ℕ} (ρ π : Ranking n) : ℕ :=
  (inversionFinset ρ π).card
\end{ReviewVerbatim}

\paragraph{Recorded source-to-library screening}
\textbf{Verdict:} \texttt{not recorded}
\reviewmatch{Matches source input}
\noindent\textbf{Reviewer annotation}\par
\reviewerbox
\clearpage
\hypertarget{library-prerequisite-02920becedba7938}{}
\subsection*{\small\ttfamily\raggedright EconCSLib.\allowbreak{}SocialChoice.\allowbreak{}Ranking.\allowbreak{}rerankingGainOnPair}
\textbf{Source connection:} no explicit byte-pinned paper source connection is registered.
\paragraph{Lean-expanded library semantic target}
\begin{ReviewVerbatim}
{n : ℕ} →
  (value : EconCSLib.SocialChoice.Ranking.Candidate n → ℝ) →
    (π σ : EconCSLib.SocialChoice.Ranking.Ranking n) →
      if EconCSLib.SocialChoice.Ranking.firstChoice π = EconCSLib.SocialChoice.Ranking.firstChoice σ then 0
      else value (EconCSLib.SocialChoice.Ranking.firstChoice π) - value (EconCSLib.SocialChoice.Ranking.secondChoice π)
\end{ReviewVerbatim}

\paragraph{Exact Lean library declaration}
\begin{ReviewVerbatim}
def rerankingGainOnPair {n : ℕ} (value : Candidate n → ℝ)
    (π σ : Ranking n) : ℝ :=
  if firstChoice π = firstChoice σ then 0
  else value (firstChoice π) - value (secondChoice π)
\end{ReviewVerbatim}

\paragraph{Recorded source-to-library screening}
\textbf{Verdict:} \texttt{not recorded}
\reviewmatch{Matches source input}
\noindent\textbf{Reviewer annotation}\par
\reviewerbox
\clearpage
\hypertarget{library-prerequisite-d3b381d476267785}{}
\subsection*{\small\ttfamily\raggedright EconCSLib.\allowbreak{}SocialChoice.\allowbreak{}Ranking.\allowbreak{}pairRerankingGain}
\textbf{Source connection:} no explicit byte-pinned paper source connection is registered.
\paragraph{Lean-expanded library semantic target}
\begin{ReviewVerbatim}
{n : ℕ} →
  (value : EconCSLib.SocialChoice.Ranking.Candidate n → ℝ) →
    (a : EconCSLib.SocialChoice.Ranking.RankingPair n) →
      EconCSLib.SocialChoice.Ranking.rerankingGainOnPair value a.1 a.2
\end{ReviewVerbatim}

\paragraph{Exact Lean library declaration}
\begin{ReviewVerbatim}
def pairRerankingGain {n : ℕ} (value : Candidate n → ℝ) :
    RankingPair n → ℝ :=
  fun p => rerankingGainOnPair value p.1 p.2
\end{ReviewVerbatim}

\paragraph{Recorded source-to-library screening}
\textbf{Verdict:} \texttt{not recorded}
\reviewmatch{Matches source input}
\noindent\textbf{Reviewer annotation}\par
\reviewerbox
\clearpage
\hypertarget{library-prerequisite-4d0316233c2913da}{}
\subsection*{\small\ttfamily\raggedright EconCSLib.\allowbreak{}SocialChoice.\allowbreak{}Ranking.\allowbreak{}rankByScore}
\textbf{Source connection:} no explicit byte-pinned paper source connection is registered.
\paragraph{Lean-expanded library semantic target}
\begin{ReviewVerbatim}
{n : ℕ} → (score : EconCSLib.SocialChoice.Ranking.Candidate n → ℝ) → Tuple.sort fun i => -score i
\end{ReviewVerbatim}

\paragraph{Exact Lean library declaration}
\begin{ReviewVerbatim}
noncomputable def rankByScore {n : ℕ} (score : Candidate n → ℝ) : Ranking n :=
  Tuple.sort (fun i : Candidate n => -score i)
\end{ReviewVerbatim}

\paragraph{Recorded source-to-library screening}
\textbf{Verdict:} \texttt{not recorded}
\reviewmatch{Matches source input}
\noindent\textbf{Reviewer annotation}\par
\reviewerbox
\clearpage
\hypertarget{library-prerequisite-3d857727212e6236}{}
\subsection*{\small\ttfamily\raggedright EconCSLib.\allowbreak{}SocialChoice.\allowbreak{}Ranking.\allowbreak{}rankingPMFOfMeasure}
\textbf{Source connection:} no explicit byte-pinned paper source connection is registered.
\paragraph{Lean-expanded library semantic target}
\begin{ReviewVerbatim}
{n : ℕ} →
  {Ω : Type u_1} →
    [inst : MeasurableSpace Ω] →
      (μ : MeasureTheory.Measure Ω) →
        [inst_1 : MeasureTheory.IsProbabilityMeasure μ] →
          (rank : Ω → EconCSLib.SocialChoice.Ranking.Ranking n) →
            (hrank : Measurable rank) → (MeasureTheory.Measure.map rank μ).toPMF
\end{ReviewVerbatim}

\paragraph{Exact Lean library declaration}
\begin{ReviewVerbatim}
def rankingPMFOfMeasure
    {n : ℕ} {Ω : Type*} [MeasurableSpace Ω]
    (μ : MeasureTheory.Measure Ω) [MeasureTheory.IsProbabilityMeasure μ]
    (rank : Ω → Ranking n) (hrank : Measurable rank) : PMF (Ranking n) :=
  @MeasureTheory.Measure.toPMF (Ranking n) _ _ _ (μ.map rank)
    (MeasureTheory.Measure.isProbabilityMeasure_map hrank.aemeasurable)
\end{ReviewVerbatim}

\paragraph{Recorded source-to-library screening}
\textbf{Verdict:} \texttt{not recorded}
\reviewmatch{Matches source input}
\noindent\textbf{Reviewer annotation}\par
\reviewerbox
\clearpage
\hypertarget{library-prerequisite-9282609ac289eae1}{}
\subsection*{\small\ttfamily\raggedright EconCSLib.\allowbreak{}SocialChoice.\allowbreak{}Ranking.\allowbreak{}rum3Ranking012}
\textbf{Source connection:} no explicit byte-pinned paper source connection is registered.
\paragraph{Lean-expanded library semantic target}
\begin{ReviewVerbatim}
Equiv.refl (EconCSLib.SocialChoice.Ranking.Candidate 1)
\end{ReviewVerbatim}

\paragraph{Exact Lean library declaration}
\begin{ReviewVerbatim}
def rum3Ranking012 : Ranking 1 :=
  Equiv.refl (Candidate 1)
\end{ReviewVerbatim}

\paragraph{Recorded source-to-library screening}
\textbf{Verdict:} \texttt{not recorded}
\reviewmatch{Matches source input}
\noindent\textbf{Reviewer annotation}\par
\reviewerbox
\clearpage
\hypertarget{library-prerequisite-3e9037b4cf76ac78}{}
\subsection*{\small\ttfamily\raggedright EconCSLib.\allowbreak{}SocialChoice.\allowbreak{}Ranking.\allowbreak{}rum3Ranking021}
\textbf{Source connection:} no explicit byte-pinned paper source connection is registered.
\paragraph{Lean-expanded library semantic target}
\begin{ReviewVerbatim}
Equiv.swap 1 2
\end{ReviewVerbatim}

\paragraph{Exact Lean library declaration}
\begin{ReviewVerbatim}
def rum3Ranking021 : Ranking 1 :=
  Equiv.swap (1 : Candidate 1) (2 : Candidate 1)
\end{ReviewVerbatim}

\paragraph{Recorded source-to-library screening}
\textbf{Verdict:} \texttt{not recorded}
\reviewmatch{Matches source input}
\noindent\textbf{Reviewer annotation}\par
\reviewerbox
\clearpage
\hypertarget{library-prerequisite-d6840f3f435bf0a3}{}
\subsection*{\small\ttfamily\raggedright EconCSLib.\allowbreak{}SocialChoice.\allowbreak{}Ranking.\allowbreak{}rum3Ranking102}
\textbf{Source connection:} no explicit byte-pinned paper source connection is registered.
\paragraph{Lean-expanded library semantic target}
\begin{ReviewVerbatim}
Equiv.swap 0 1
\end{ReviewVerbatim}

\paragraph{Exact Lean library declaration}
\begin{ReviewVerbatim}
def rum3Ranking102 : Ranking 1 :=
  Equiv.swap (0 : Candidate 1) (1 : Candidate 1)
\end{ReviewVerbatim}

\paragraph{Recorded source-to-library screening}
\textbf{Verdict:} \texttt{not recorded}
\reviewmatch{Matches source input}
\noindent\textbf{Reviewer annotation}\par
\reviewerbox
\clearpage
\hypertarget{library-prerequisite-17cdcd166ffcfa57}{}
\subsection*{\small\ttfamily\raggedright EconCSLib.\allowbreak{}SocialChoice.\allowbreak{}Ranking.\allowbreak{}rum3Ranking120}
\textbf{Source connection:} no explicit byte-pinned paper source connection is registered.
\paragraph{Lean-expanded library semantic target}
\begin{ReviewVerbatim}
Equiv.trans (Equiv.swap 1 2) (Equiv.swap 0 1)
\end{ReviewVerbatim}

\paragraph{Exact Lean library declaration}
\begin{ReviewVerbatim}
def rum3Ranking120 : Ranking 1 :=
  (Equiv.swap (1 : Candidate 1) (2 : Candidate 1)).trans
    (Equiv.swap (0 : Candidate 1) (1 : Candidate 1))
\end{ReviewVerbatim}

\paragraph{Recorded source-to-library screening}
\textbf{Verdict:} \texttt{not recorded}
\reviewmatch{Matches source input}
\noindent\textbf{Reviewer annotation}\par
\reviewerbox
\clearpage
\hypertarget{library-prerequisite-88da1ce8c8de0862}{}
\subsection*{\small\ttfamily\raggedright EconCSLib.\allowbreak{}SocialChoice.\allowbreak{}Ranking.\allowbreak{}rum3Ranking201}
\textbf{Source connection:} no explicit byte-pinned paper source connection is registered.
\paragraph{Lean-expanded library semantic target}
\begin{ReviewVerbatim}
Equiv.trans (Equiv.swap 0 2) (Equiv.swap 0 1)
\end{ReviewVerbatim}

\paragraph{Exact Lean library declaration}
\begin{ReviewVerbatim}
def rum3Ranking201 : Ranking 1 :=
  (Equiv.swap (0 : Candidate 1) (2 : Candidate 1)).trans
    (Equiv.swap (0 : Candidate 1) (1 : Candidate 1))
\end{ReviewVerbatim}

\paragraph{Recorded source-to-library screening}
\textbf{Verdict:} \texttt{not recorded}
\reviewmatch{Matches source input}
\noindent\textbf{Reviewer annotation}\par
\reviewerbox
\clearpage
\hypertarget{library-prerequisite-11023713299f4bb8}{}
\subsection*{\small\ttfamily\raggedright EconCSLib.\allowbreak{}SocialChoice.\allowbreak{}Ranking.\allowbreak{}rum3Ranking210}
\textbf{Source connection:} no explicit byte-pinned paper source connection is registered.
\paragraph{Lean-expanded library semantic target}
\begin{ReviewVerbatim}
Equiv.swap 0 2
\end{ReviewVerbatim}

\paragraph{Exact Lean library declaration}
\begin{ReviewVerbatim}
def rum3Ranking210 : Ranking 1 :=
  Equiv.swap (0 : Candidate 1) (2 : Candidate 1)
\end{ReviewVerbatim}

\paragraph{Recorded source-to-library screening}
\textbf{Verdict:} \texttt{not recorded}
\reviewmatch{Matches source input}
\noindent\textbf{Reviewer annotation}\par
\reviewerbox
\clearpage
\hypertarget{library-prerequisite-2b57daef13f9e5a2}{}
\subsection*{\small\ttfamily\raggedright EconCSLib.\allowbreak{}SocialChoice.\allowbreak{}Ranking.\allowbreak{}rum3RankByScores}
\textbf{Source connection:} no explicit byte-pinned paper source connection is registered.
\paragraph{Lean-expanded library semantic target}
\begin{ReviewVerbatim}
(s1 s2 s3 : ℝ) →
  if h0 : s2 ≤ s1 ∧ s3 ≤ s1 then
    if s3 ≤ s2 then EconCSLib.SocialChoice.Ranking.rum3Ranking012 else EconCSLib.SocialChoice.Ranking.rum3Ranking021
  else
    if h1 : s1 < s2 ∧ s3 ≤ s2 then
      if s3 ≤ s1 then EconCSLib.SocialChoice.Ranking.rum3Ranking102 else EconCSLib.SocialChoice.Ranking.rum3Ranking120
    else
      if s2 ≤ s1 then EconCSLib.SocialChoice.Ranking.rum3Ranking201 else EconCSLib.SocialChoice.Ranking.rum3Ranking210
\end{ReviewVerbatim}

\paragraph{Exact Lean library declaration}
\begin{ReviewVerbatim}
noncomputable def rum3RankByScores (s1 s2 s3 : ℝ) : Ranking 1 :=
  if h0 : s2 ≤ s1 ∧ s3 ≤ s1 then
    if s3 ≤ s2 then rum3Ranking012 else rum3Ranking021
  else if h1 : s1 < s2 ∧ s3 ≤ s2 then
    if s3 ≤ s1 then rum3Ranking102 else rum3Ranking120
  else
    if s2 ≤ s1 then rum3Ranking201 else rum3Ranking210
\end{ReviewVerbatim}

\paragraph{Recorded source-to-library screening}
\textbf{Verdict:} \texttt{not recorded}
\reviewmatch{Matches source input}
\noindent\textbf{Reviewer annotation}\par
\reviewerbox
\clearpage
\hypertarget{library-prerequisite-d3a1fe6e17d1ebe8}{}
\subsection*{\small\ttfamily\raggedright EconCSLib.\allowbreak{}finCons}
\textbf{Source connection:} no explicit byte-pinned paper source connection is registered.
\paragraph{Lean-expanded library semantic target}
\begin{ReviewVerbatim}
{α : Type u_1} →
  {k : ℕ} →
    (head : α) →
      (tail : Fin k → α) →
        (a : Fin (k + 1)) →
          match a with
          | ⟨0, _h⟩ => head
          | ⟨n.succ, h⟩ => tail ⟨n, ⋯⟩
\end{ReviewVerbatim}

\paragraph{Exact Lean library declaration}
\begin{ReviewVerbatim}
def finCons {α : Type*} {k : ℕ} (head : α) (tail : Fin k → α) :
    Fin (k + 1) → α
  | ⟨0, _h⟩ => head
  | ⟨n + 1, h⟩ => tail ⟨n, Nat.succ_lt_succ_iff.mp h⟩

@[simp]
theorem finCons_zero {α : Type*} {k : ℕ} (head : α) (tail : Fin k → α) :
    finCons head tail ⟨0, Nat.succ_pos k⟩ = head := rfl
\end{ReviewVerbatim}

\paragraph{Recorded source-to-library screening}
\textbf{Verdict:} \texttt{not recorded}
\reviewmatch{Matches source input}
\noindent\textbf{Reviewer annotation}\par
\reviewerbox
\clearpage
\hypertarget{library-prerequisite-4aa064b5cbd66ed3}{}
\subsection*{\small\ttfamily\raggedright EconCSLib.\allowbreak{}finiteAvailableWeight}
\textbf{Source connection:} no explicit byte-pinned paper source connection is registered.
\paragraph{Lean-expanded library semantic target}
\begin{ReviewVerbatim}
{α : Type u_1} →
  [inst : Fintype α] →
    [inst_1 : DecidableEq α] →
      (baseWeight : α → ℝ) → (forbidden : Finset α) → ∑ a, if a ∈ forbidden then 0 else baseWeight a
\end{ReviewVerbatim}

\paragraph{Exact Lean library declaration}
\begin{ReviewVerbatim}
noncomputable def finiteAvailableWeight {α : Type*} [Fintype α] [DecidableEq α]
    (baseWeight : α → ℝ) (forbidden : Finset α) : ℝ :=
  ∑ a : α, if a ∈ forbidden then 0 else baseWeight a
\end{ReviewVerbatim}

\paragraph{Recorded source-to-library screening}
\textbf{Verdict:} \texttt{not recorded}
\reviewmatch{Matches source input}
\noindent\textbf{Reviewer annotation}\par
\reviewerbox
\clearpage
\hypertarget{library-prerequisite-637b06ac0adfcac1}{}
\subsection*{\small\ttfamily\raggedright EconCSLib.\allowbreak{}finiteFreshList}
\textbf{Source connection:} no explicit byte-pinned paper source connection is registered.
\paragraph{Lean-expanded library semantic target}
\begin{ReviewVerbatim}
(α : Type u_1) →
  (k : ℕ) → (forbidden : Finset α) → { list // Function.Injective list ∧ ∀ (slot : Fin k), list slot ∉ forbidden }
\end{ReviewVerbatim}

\paragraph{Exact Lean library declaration}
\begin{ReviewVerbatim}
abbrev finiteFreshList (α : Type*) (k : ℕ) (forbidden : Finset α) :=
  {list : Fin k → α // Function.Injective list ∧ ∀ slot, list slot ∉ forbidden}
\end{ReviewVerbatim}

\paragraph{Recorded source-to-library screening}
\textbf{Verdict:} \texttt{not recorded}
\reviewmatch{Matches source input}
\noindent\textbf{Reviewer annotation}\par
\reviewerbox
\clearpage
\hypertarget{library-prerequisite-517099bc7eacc408}{}
\subsection*{\small\ttfamily\raggedright EconCSLib.\allowbreak{}finiteFreshListCons}
\textbf{Source connection:} no explicit byte-pinned paper source connection is registered.
\paragraph{Lean-expanded library semantic target}
\begin{ReviewVerbatim}
{α : Type u_1} →
  [inst : DecidableEq α] →
    {k : ℕ} →
      {forbidden : Finset α} →
        (head : { a // a ∉ forbidden }) →
          (tail : EconCSLib.finiteFreshList α k (insert (↑head) forbidden)) → ⟨EconCSLib.finCons ↑head ↑tail, ⋯⟩
\end{ReviewVerbatim}

\paragraph{Exact Lean library declaration}
\begin{ReviewVerbatim}
def finiteFreshListCons {α : Type*} [DecidableEq α] {k : ℕ}
    {forbidden : Finset α}
    (head : {a // a ∉ forbidden})
    (tail : finiteFreshList α k (insert head.1 forbidden)) :
    finiteFreshList α (k + 1) forbidden :=
  ⟨finCons head.1 tail.1,
    by
      refine finCons_injective ?_ tail.2.1
      intro i htail_eq_head
      exact tail.2.2 i (by
        rw [htail_eq_head]
        simp)
    ,
    by
      intro slot
      cases slot using Fin.cases with
      | zero =>
          exact head.2
      | succ i =>
          intro hforbidden
          exact tail.2.2 i
            (Finset.mem_insert.mpr (Or.inr hforbidden))⟩
\end{ReviewVerbatim}

\paragraph{Recorded source-to-library screening}
\textbf{Verdict:} \texttt{not recorded}
\reviewmatch{Matches source input}
\noindent\textbf{Reviewer annotation}\par
\reviewerbox
\clearpage
\hypertarget{library-prerequisite-010f5218e34f7dbc}{}
\subsection*{\small\ttfamily\raggedright EconCSLib.\allowbreak{}finiteFreshListNil}
\textbf{Source connection:} no explicit byte-pinned paper source connection is registered.
\paragraph{Lean-expanded library semantic target}
\begin{ReviewVerbatim}
(α : Type u_1) → (forbidden : Finset α) → ⟨fun i => i.elim0, ⋯⟩
\end{ReviewVerbatim}

\paragraph{Exact Lean library declaration}
\begin{ReviewVerbatim}
def finiteFreshListNil (α : Type*) (forbidden : Finset α) :
    finiteFreshList α 0 forbidden :=
  ⟨fun i => Fin.elim0 i,
    by
      intro i _j _h
      exact Fin.elim0 i,
    by
      intro i
      exact Fin.elim0 i⟩
\end{ReviewVerbatim}

\paragraph{Recorded source-to-library screening}
\textbf{Verdict:} \texttt{not recorded}
\reviewmatch{Matches source input}
\noindent\textbf{Reviewer annotation}\par
\reviewerbox
\clearpage
\hypertarget{library-prerequisite-8598d8001b9c1939}{}
\subsection*{\small\ttfamily\raggedright EconCSLib.\allowbreak{}finiteWeightedPMF}
\textbf{Source connection:} no explicit byte-pinned paper source connection is registered.
\paragraph{Lean-expanded library semantic target}
\begin{ReviewVerbatim}
{α : Type u_1} →
  [inst : Fintype α] →
    (weight : α → ℝ) →
      (hweight_nonneg : ∀ (a : α), 0 ≤ weight a) →
        (htotal_pos : 0 < ∑ a, weight a) → PMF.ofFintype (fun a => ENNReal.ofReal (weight a / ∑ b, weight b)) ⋯
\end{ReviewVerbatim}

\paragraph{Exact Lean library declaration}
\begin{ReviewVerbatim}
noncomputable def finiteWeightedPMF {α : Type*} [Fintype α]
    (weight : α → ℝ)
    (hweight_nonneg : ∀ a, 0 ≤ weight a)
    (htotal_pos : 0 < ∑ a, weight a) : PMF α :=
  PMF.ofFintype
    (fun a => ENNReal.ofReal (weight a / ∑ b, weight b))
    (by
      have hterm_nonneg : ∀ a : α, 0 ≤ weight a / ∑ b, weight b := by
        intro a
        exact div_nonneg (hweight_nonneg a) (le_of_lt htotal_pos)
      have hsum_real :
          (∑ a : α, weight a / ∑ b, weight b) = 1 := by
        rw [← Finset.sum_div]
        field_simp [ne_of_gt htotal_pos]
      calc
        ∑ a : α, ENNReal.ofReal (weight a / ∑ b, weight b)
            = ENNReal.ofReal (∑ a : α, weight a / ∑ b, weight b) := by
              symm
              exact ENNReal.ofReal_sum_of_nonneg
                (s := (Finset.univ : Finset α))
                (f := fun a => weight a / ∑ b, weight b)
                (by intro a _; exact hterm_nonneg a)
        _ = ENNReal.ofReal 1 := by rw [hsum_real]
        _ = 1 := by norm_num)
\end{ReviewVerbatim}

\paragraph{Recorded source-to-library screening}
\textbf{Verdict:} \texttt{not recorded}
\reviewmatch{Matches source input}
\noindent\textbf{Reviewer annotation}\par
\reviewerbox
\clearpage
\hypertarget{library-prerequisite-478178a796f1a098}{}
\subsection*{\small\ttfamily\raggedright EconCSLib.\allowbreak{}finiteWeightedPMFAvailable}
\textbf{Source connection:} no explicit byte-pinned paper source connection is registered.
\paragraph{Lean-expanded library semantic target}
\begin{ReviewVerbatim}
{α : Type u_1} →
  [inst : Fintype α] →
    [inst_1 : DecidableEq α] →
      (baseWeight : α → ℝ) →
        (forbidden : Finset α) →
          (hbase_nonneg : ∀ (a : α), 0 ≤ baseWeight a) →
            (havailable_pos : 0 < EconCSLib.finiteAvailableWeight baseWeight forbidden) →
              EconCSLib.finiteWeightedPMF (fun a => baseWeight ↑a) ⋯ ⋯
\end{ReviewVerbatim}

\paragraph{Exact Lean library declaration}
\begin{ReviewVerbatim}
noncomputable def finiteWeightedPMFAvailable
    {α : Type*} [Fintype α] [DecidableEq α]
    (baseWeight : α → ℝ) (forbidden : Finset α)
    (hbase_nonneg : ∀ a, 0 ≤ baseWeight a)
    (havailable_pos : 0 < finiteAvailableWeight baseWeight forbidden) :
    PMF {a // a ∉ forbidden} :=
  finiteWeightedPMF
    (fun a : {a // a ∉ forbidden} => baseWeight a.1)
    (by intro a; exact hbase_nonneg a.1)
    (by
      rw [← finiteAvailableWeight_eq_subtype_sum baseWeight forbidden]
      exact havailable_pos)
\end{ReviewVerbatim}

\paragraph{Recorded source-to-library screening}
\textbf{Verdict:} \texttt{not recorded}
\reviewmatch{Matches source input}
\noindent\textbf{Reviewer annotation}\par
\reviewerbox
\clearpage
\hypertarget{library-prerequisite-17cb3d8406f594aa}{}
\subsection*{\small\ttfamily\raggedright EconCSLib.\allowbreak{}finiteWithoutReplacementPMF}
\textbf{Source connection:} no explicit byte-pinned paper source connection is registered.
\paragraph{Lean-expanded library semantic target}
\begin{ReviewVerbatim}
{α : Type u_1} →
  [inst : Fintype α] →
    [inst_1 : DecidableEq α] →
      (baseWeight : α → ℝ) →
        (hbase_nonneg : ∀ (a : α), 0 ≤ baseWeight a) →
          (havailable :
              ∀ (forbidden : Finset α),
                forbidden.card < Fintype.card α → 0 < EconCSLib.finiteAvailableWeight baseWeight forbidden) →
            (k : ℕ) →
              (forbidden : Finset α) →
                (a : forbidden.card + k ≤ Fintype.card α) →
                  WellFounded.Nat.fix
                    (fun x => PSigma.casesOn x fun k forbidden => PSigma.casesOn forbidden fun forbidden a => k)
                    (fun _x a =>
                      PSigma.casesOn (motive := fun _x =>
                        ((y : (k : ℕ) ×' (forbidden : Finset α) ×' forbidden.card + k ≤ Fintype.card α) →
                            InvImage (fun x1 x2 => x1 < x2)
                                (fun x =>
                                  PSigma.casesOn x fun k forbidden => PSigma.casesOn forbidden fun forbidden a => k)
                                y _x →
                              PMF (EconCSLib.finiteFreshList α y.1 y.2.1)) →
                          PMF (EconCSLib.finiteFreshList α _x.1 _x.2.1))
                        _x
                        (fun k forbidden a =>
                          PSigma.casesOn (motive := fun forbidden =>
                            ((y : (k : ℕ) ×' (forbidden : Finset α) ×' forbidden.card + k ≤ Fintype.card α) →
                                InvImage (fun x1 x2 => x1 < x2)
                                    (fun x =>
                                      PSigma.casesOn x fun k forbidden => PSigma.casesOn forbidden fun forbidden a => k)
                                    y ⟨k, forbidden⟩ →
                                  PMF (EconCSLib.finiteFreshList α y.1 y.2.1)) →
                              PMF (EconCSLib.finiteFreshList α ⟨k, forbidden⟩.1 ⟨k, forbidden⟩.2.1))
                            forbidden
                            (fun forbidden a a_1 =>
                              (match (motive :=
                                  (x : ℕ) →
                                    (x_1 : Finset α) →
                                      (x_2 : x_1.card + x ≤ Fintype.card α) →
                                        ((y :
                                              (k : ℕ) ×'
                                                (forbidden : Finset α) ×' forbidden.card + k ≤ Fintype.card α) →
                                            InvImage (fun x1 x2 => x1 < x2)
                                                (fun x =>
                                                  PSigma.casesOn x fun k forbidden =>
                                                    PSigma.casesOn forbidden fun forbidden a => k)
                                                y ⟨x, ⟨x_1, x_2⟩⟩ →
                                              PMF (EconCSLib.finiteFreshList α y.1 y.2.1)) →
                                          PMF (EconCSLib.finiteFreshList α x x_1))
                                  k, forbidden, a with
                                | 0, forbidden, _hbudget => fun x => PMF.pure (EconCSLib.finiteFreshListNil α forbidden)
                                | k.succ, forbidden, hbudget => fun x =>
                                  (EconCSLib.finiteWeightedPMFAvailable baseWeight forbidden hbase_nonneg ⋯).bind
                                    fun head =>
                                    PMF.map (EconCSLib.finiteFreshListCons head)
                                      (x ⟨k, ⟨insert (↑head) forbidden, ⋯⟩⟩ ⋯))
                                a_1)
                            a)
                        a)
                    ⟨k, ⟨forbidden, a⟩⟩
\end{ReviewVerbatim}

\paragraph{Exact Lean library declaration}
\begin{ReviewVerbatim}
noncomputable def finiteWithoutReplacementPMF
    {α : Type*} [Fintype α] [DecidableEq α]
    (baseWeight : α → ℝ)
    (hbase_nonneg : ∀ a, 0 ≤ baseWeight a)
    (havailable : ∀ forbidden : Finset α,
      forbidden.card < Fintype.card α →
        0 < finiteAvailableWeight baseWeight forbidden) :
    (k : ℕ) → (forbidden : Finset α) →
      forbidden.card + k ≤ Fintype.card α →
        PMF (finiteFreshList α k forbidden)
  | 0, forbidden, _hbudget => PMF.pure (finiteFreshListNil α forbidden)
  | k + 1, forbidden, hbudget =>
      (finiteWeightedPMFAvailable
        baseWeight forbidden hbase_nonneg
        (havailable forbidden (by omega))).bind fun head =>
        (finiteWithoutReplacementPMF
          baseWeight hbase_nonneg havailable
          k (insert head.1 forbidden)
          (by
            rw [Finset.card_insert_of_notMem head.2]
            omega)).map (finiteFreshListCons head)
termination_by k forbidden hbudget => k

/--
The first draw of the recursive without-replacement sampler has exactly the
available weighted law from the current forbidden set.
-/
theorem finiteWithoutReplacementPMF_head_prob
    {α : Type*} [Fintype α] [DecidableEq α]
    (baseWeight : α → ℝ)
    (hbase_nonneg : ∀ a, 0 ≤ baseWeight a)
    (havailable : ∀ forbidden : Finset α,
      forbidden.card < Fintype.card α →
        0 < finiteAvailableWeight baseWeight forbidden)
    {k : ℕ} (forbidden : Finset α)
    (hbudget : forbidden.card + (k + 1) ≤ Fintype.card α)
    (head₀ : {a // a ∉ forbidden}) :
    pmfProb
        (finiteWithoutReplacementPMF
          baseWeight hbase_nonneg havailable
          (k + 1) forbidden hbudget)
        (fun sample =>
          sample.1 ⟨0, Nat.succ_pos k⟩ = head₀.1) =
      (finiteWeightedPMFAvailable
        baseWeight forbidden hbase_nonneg
        (havailable forbidden (by omega)) head₀).toReal := by
  classical
  rw [finiteWithoutReplacementPMF]
  rw [pmfProb_bind]
  let headLaw :=
    finiteWeightedPMFAvailable
      baseWeight forbidden hbase_nonneg
      (havailable forbidden (by omega))
  have hinner : ∀ head : {a // a ∉ forbidden},
      pmfProb
          ((finiteWithoutReplacementPMF
            baseWeight hbase_nonneg havailable
            k (insert head.1 forbidden)
            (by
              rw [Finset.card_insert_of_notMem head.2]
              omega)).map (finiteFreshListCons head))
          (fun sample =>
            sample.1 ⟨0, Nat.succ_pos k⟩ = head₀.1) =
        if head = head₀ then (1 : ℝ) else 0 := by
    intro head
    rw [pmfProb_map]
    by_cases hhead : head = head₀
    · subst head
      simpa using
        (pmfProb_eq_one_of_forall _ _ (by
          intro tail
          rfl))
    · have hval_ne : head.1 ≠ head₀.1 := by
        intro hval
        exact hhead (Subtype.ext hval)
      simpa [hhead] using
        (pmfProb_eq_zero_of_no_mass _ _ (by
        intro tail htail
        exact False.elim (hval_ne (by simpa using htail)))
        )
  calc
    pmfExp headLaw
        (fun head =>
          pmfProb
            ((finiteWithoutReplacementPMF
              baseWeight hbase_nonneg havailable
              k (insert head.1 forbidden)
              (by
                rw [Finset.card_insert_of_notMem head.2]
                omega)).map (finiteFreshListCons head))
            (fun sample =>
              sample.1 ⟨0, Nat.succ_pos k⟩ = head₀.1))
        = pmfExp headLaw (fun head => if head = head₀ then (1 : ℝ) else 0) := by
            refine pmfExp_congr headLaw ?_
            intro head
            exact hinner head
    _ = pmfProb headLaw (fun head => head = head₀) := rfl
    _ = (headLaw head₀).toReal := pmfProb_singleton headLaw head₀
    _ = (finiteWeightedPMFAvailable
        baseWeight forbidden hbase_nonneg
        (havailable forbidden (by omega)) head₀).toReal := rfl
\end{ReviewVerbatim}

\paragraph{Recorded source-to-library screening}
\textbf{Verdict:} \texttt{not recorded}
\reviewmatch{Matches source input}
\noindent\textbf{Reviewer annotation}\par
\reviewerbox
\clearpage
\hypertarget{library-prerequisite-7ae22d5411bd7792}{}
\subsection*{\small\ttfamily\raggedright EconCSLib.\allowbreak{}measureProb}
\textbf{Source connection:} no explicit byte-pinned paper source connection is registered.
\paragraph{Lean-expanded library semantic target}
\begin{ReviewVerbatim}
{α : Type u_1} → [inst : MeasurableSpace α] → (μ : MeasureTheory.Measure α) → (p : α → Prop) → (μ {a | p a}).toReal
\end{ReviewVerbatim}

\paragraph{Exact Lean library declaration}
\begin{ReviewVerbatim}
noncomputable def measureProb {α : Type*} [MeasurableSpace α]
    (μ : Measure α) (p : α → Prop) : ℝ :=
  (μ {a | p a}).toReal
\end{ReviewVerbatim}

\paragraph{Recorded source-to-library screening}
\textbf{Verdict:} \texttt{not recorded}
\reviewmatch{Matches source input}
\noindent\textbf{Reviewer annotation}\par
\reviewerbox
\clearpage
\hypertarget{library-prerequisite-197097ef6d41e4d7}{}
\subsection*{\small\ttfamily\raggedright EconCSLib.\allowbreak{}pmfPairIndicatorExp}
\textbf{Source connection:} no explicit byte-pinned paper source connection is registered.
\paragraph{Lean-expanded library semantic target}
\begin{ReviewVerbatim}
{α : Type u_1} →
  {β : Type u_2} →
    [inst : Fintype α] →
      [inst_1 : DecidableEq α] →
        [inst_2 : Fintype β] →
          [inst_3 : DecidableEq β] →
            (μ : PMF α) →
              (ν : PMF β) →
                (p : α × β → Prop) →
                  [inst_4 : DecidablePred p] →
                    (f : α × β → ℝ) → EconCSLib.pmfPairExp μ ν fun a b => if p (a, b) then f (a, b) else 0
\end{ReviewVerbatim}

\paragraph{Exact Lean library declaration}
\begin{ReviewVerbatim}
noncomputable def pmfPairIndicatorExp {α β : Type*}
    [Fintype α] [DecidableEq α] [Fintype β] [DecidableEq β]
    (μ : PMF α) (ν : PMF β) (p : α × β → Prop) [DecidablePred p]
    (f : α × β → ℝ) : ℝ :=
  pmfPairExp μ ν (fun a b => if p (a, b) then f (a, b) else 0)
\end{ReviewVerbatim}

\paragraph{Recorded source-to-library screening}
\textbf{Verdict:} \texttt{not recorded}
\reviewmatch{Matches source input}
\noindent\textbf{Reviewer annotation}\par
\reviewerbox
\clearpage
\hypertarget{library-prerequisite-f185c8e6394c9317}{}
\subsection*{\small\ttfamily\raggedright EconCSLib.\allowbreak{}pmfProd}
\textbf{Source connection:} no explicit byte-pinned paper source connection is registered.
\paragraph{Lean-expanded library semantic target}
\begin{ReviewVerbatim}
{α : Type u_1} →
  {β : Type u_2} →
    [inst : Fintype α] → [inst_1 : Fintype β] → (μ : PMF α) → (ν : PMF β) → PMF.ofFintype (fun p => μ p.1 * ν p.2) ⋯
\end{ReviewVerbatim}

\paragraph{Exact Lean library declaration}
\begin{ReviewVerbatim}
noncomputable def pmfProd {α β : Type*} [Fintype α] [Fintype β]
    (μ : PMF α) (ν : PMF β) : PMF (α × β) :=
  PMF.ofFintype (fun p : α × β => μ p.1 * ν p.2) (by
    classical
    have hμ : ∑ a : α, μ a = 1 := by
      rw [← PMF.tsum_coe μ, tsum_fintype]
    have hν : ∑ b : β, ν b = 1 := by
      rw [← PMF.tsum_coe ν, tsum_fintype]
    calc
      ∑ p : α × β, μ p.1 * ν p.2
          = ∑ a : α, ∑ b : β, μ a * ν b := by
            rw [Fintype.sum_prod_type]
      _ = ∑ a : α, μ a * (∑ b : β, ν b) := by
            simp [Finset.mul_sum]
      _ = 1 := by
            simp [hμ, hν])
\end{ReviewVerbatim}

\paragraph{Recorded source-to-library screening}
\textbf{Verdict:} \texttt{not recorded}
\reviewmatch{Matches source input}
\noindent\textbf{Reviewer annotation}\par
\reviewerbox

\clearpage
\hypertarget{source-claim-70a4d7b84f526fcd}{}
\hypertarget{source-section-f10b5f68e4acf3dc}{}
\section*{Source claims}
\subsection*{Review row 1}
\textbf{Source locator:} {\footnotesize\raggedright cited publication, lines 133-\allowbreak{}138\par}
\paragraph{Verbatim source input}
\begin{ReviewVerbatim}
[No record available.]
\end{ReviewVerbatim}


\paragraph{Expanded PaperInterface specification}
\begin{ReviewVerbatim}
∀ {n : ℕ} (trueRanking : EconCSLib.SocialChoice.Ranking.Ranking n) (F : KR21Monoculture.AccuracyFamily n),
  EconCSLib.SocialChoice.Ranking.StrictlyOrderedBy trueRanking F.value →
    Fintype.card (EconCSLib.SocialChoice.Ranking.Candidate n) = n + 2 ∧
      2 ≤ Fintype.card (EconCSLib.SocialChoice.Ranking.Candidate n) ∧
        (∀ {a b : EconCSLib.SocialChoice.Ranking.Candidate n},
            EconCSLib.SocialChoice.Ranking.rankOf trueRanking a < EconCSLib.SocialChoice.Ranking.rankOf trueRanking b →
              F.value b < F.value a) ∧
          ∀ (theta : ℝ), 0 < theta → (F.dist theta).toMeasure Set.univ = 1
\end{ReviewVerbatim}

\paragraph{Lean proof endpoint (built separately)}\begin{ReviewVerbatim}
KR21Monoculture.PaperInterface.source_model_candidates_values_and_ranking_law
\end{ReviewVerbatim}

\paragraph{Recorded source-to-Spec screening}
\textbf{Verdict:} \texttt{not recorded}
\\
\textbf{Reason:} not recorded
\reviewmatch{Matches source input}
\noindent\textbf{Reviewer annotation}\par
\reviewerbox
\clearpage
\hypertarget{source-claim-fbe3ef1f5678507e}{}
\section*{Review row 2}
\textbf{Source locator:} {\footnotesize\raggedright cited publication, lines 139-\allowbreak{}165\par}
\paragraph{Verbatim source input}
\begin{ReviewVerbatim}
Definition 1 (Noisy permutation family). A noisy permutation family Fθ is a family of distributions over
permutations that satisfies the following conditions for any θ > 0 and set of candidates x:
1. (Differentiability) For any permutation π, PrFθ [π] is continuous and differentiable in θ.
2. (Asymptotic optimality) For the true ranking π ∗ , limθ→∞ PrFθ [π ∗ ] = 1.
3. (Monotonicity) For any (possibly empty) S ⊂ x, let π (−S) be the partial ranking produced by removing
(−S)
the items in S from π. Let π1
denote the value of the top-ranked candidate according to π (−S) . For
0
any θ > θ,
h
i
h
i
(−S)
(−S)
EFθ0 π1
≥ EFθ π1
.
(1)
Moreover, for S = ∅, (1) holds with strict inequality.
θ serves as an “accuracy parameter”: for large θ, the noisy ranking converges to the true ranking over
candidates. The monotonicity condition states that a higher value of θ leads to a better first choice, even
if some of the candidates are removed after ranking. Removal after ranking (as opposed to before) is
important because some of the ranking models we will consider later do not satisfy Independence of Irrelevant
Alternatives. Examples of noisy permutation families include Random Utility Models [23] and the Mallows
Model [14], both of which we will discuss in detail later.
\end{ReviewVerbatim}


\paragraph{Expanded PaperInterface specification}
\begin{ReviewVerbatim}
∀ {n : ℕ} (F : KR21Monoculture.AccuracyFamily n) (trueRanking : EconCSLib.SocialChoice.Ranking.Ranking n),
  EconCSLib.SocialChoice.Ranking.StrictlyOrderedBy trueRanking F.value →
    (KR21Monoculture.SourceDefinition1 F trueRanking ↔
      (∀ (theta : ℝ),
          0 < theta →
            ∀ (pi : EconCSLib.SocialChoice.Ranking.Ranking n),
              ContinuousAt (fun theta' => ((F.dist theta') pi).toReal) theta ∧
                DifferentiableAt ℝ (fun theta' => ((F.dist theta') pi).toReal) theta) ∧
        Filter.Tendsto (fun theta => ((F.dist theta) trueRanking).toReal) Filter.atTop (nhds 1) ∧
          ∀ (thetaA thetaH : ℝ),
            0 < thetaH →
              thetaH < thetaA →
                (∀ (remaining : Finset (EconCSLib.SocialChoice.Ranking.Candidate n)),
                    remaining.Nonempty →
                      KR21Monoculture.expectedBestInSet (F.dist thetaH) F.value remaining ≤
                        KR21Monoculture.expectedBestInSet (F.dist thetaA) F.value remaining) ∧
                  KR21Monoculture.expectedBestInSet (F.dist thetaH) F.value Finset.univ <
                    KR21Monoculture.expectedBestInSet (F.dist thetaA) F.value Finset.univ)
\end{ReviewVerbatim}

\paragraph{Lean proof endpoint (built separately)}\begin{ReviewVerbatim}
KR21Monoculture.PaperInterface.source_definition1_iff
\end{ReviewVerbatim}

\paragraph{Recorded source-to-Spec screening}
\textbf{Verdict:} \texttt{not recorded}
\\
\textbf{Reason:} not recorded
\reviewmatch{Matches source input}
\noindent\textbf{Reviewer annotation}\par
\reviewerbox
\clearpage
\hypertarget{source-claim-413fc0a7e71398bc}{}
\section*{Review row 3}
\textbf{Source locator:} {\footnotesize\raggedright cited publication, lines 143-\allowbreak{}159\par}
\paragraph{Verbatim source input}
\begin{ReviewVerbatim}
3. (Monotonicity) For any (possibly empty) S ⊂ x, let π (−S) be the partial ranking produced by removing
(−S)
the items in S from π. Let π1
denote the value of the top-ranked candidate according to π (−S) . For
0
any θ > θ,
h
i
h
i
(−S)
(−S)
EFθ0 π1
≥ EFθ π1
.
(1)
Moreover, for S = ∅, (1) holds with strict inequality.
\end{ReviewVerbatim}


\paragraph{Expanded PaperInterface specification}
\begin{ReviewVerbatim}
∀ {n : ℕ} {F : KR21Monoculture.AccuracyFamily n} {trueRanking : EconCSLib.SocialChoice.Ranking.Ranking n},
  EconCSLib.SocialChoice.Ranking.StrictlyOrderedBy trueRanking F.value →
    KR21Monoculture.SourceDefinition1 F trueRanking →
      ∀ (thetaA thetaH : ℝ),
        0 < thetaH →
          thetaH < thetaA →
            (∀ (remaining : Finset (EconCSLib.SocialChoice.Ranking.Candidate n)),
                remaining.Nonempty →
                  KR21Monoculture.expectedBestInSet (F.dist thetaH) F.value remaining ≤
                    KR21Monoculture.expectedBestInSet (F.dist thetaA) F.value remaining) ∧
              KR21Monoculture.expectedBestInSet (F.dist thetaH) F.value Finset.univ <
                KR21Monoculture.expectedBestInSet (F.dist thetaA) F.value Finset.univ
\end{ReviewVerbatim}

\paragraph{Lean proof endpoint (built separately)}\begin{ReviewVerbatim}
KR21Monoculture.PaperInterface.source_equation1_removal_monotonicity
\end{ReviewVerbatim}

\paragraph{Recorded source-to-Spec screening}
\textbf{Verdict:} \texttt{not recorded}
\\
\textbf{Reason:} not recorded
\reviewmatch{Matches source input}
\noindent\textbf{Reviewer annotation}\par
\reviewerbox
\clearpage
\hypertarget{source-claim-f4aa43d846eb3ae7}{}
\section*{Review row 4}
\textbf{Source locator:} {\footnotesize\raggedright cited publication, lines 176-\allowbreak{}192\par}
\paragraph{Verbatim source input}
\begin{ReviewVerbatim}
[No record available.]
\end{ReviewVerbatim}


\paragraph{Expanded PaperInterface specification}
\begin{ReviewVerbatim}
∀ {n : ℕ} (M : KR21Monoculture.Model n),
  M.firstMoverEU KR21Monoculture.Strategy.algorithm =
      EconCSLib.SocialChoice.Ranking.expectedFirstMoverUtility M.algorithmRanking M.value ∧
    M.firstMoverEU KR21Monoculture.Strategy.human =
        EconCSLib.SocialChoice.Ranking.expectedFirstMoverUtility M.humanRanking M.value ∧
      M.secondMoverEU KR21Monoculture.Strategy.algorithm KR21Monoculture.Strategy.algorithm =
          EconCSLib.SocialChoice.Ranking.expectedSecondMoverShared M.algorithmRanking M.value ∧
        M.secondMoverEU KR21Monoculture.Strategy.algorithm KR21Monoculture.Strategy.human =
            EconCSLib.SocialChoice.Ranking.expectedSecondMoverIndependent M.humanRanking M.algorithmRanking M.value ∧
          M.secondMoverEU KR21Monoculture.Strategy.human KR21Monoculture.Strategy.algorithm =
              EconCSLib.SocialChoice.Ranking.expectedSecondMoverIndependent M.algorithmRanking M.humanRanking M.value ∧
            M.secondMoverEU KR21Monoculture.Strategy.human KR21Monoculture.Strategy.human =
                EconCSLib.SocialChoice.Ranking.expectedSecondMoverIndependent M.humanRanking M.humanRanking M.value ∧
              ∀ (self other : KR21Monoculture.Strategy),
                M.payoffAgainst self other = (M.firstMoverEU self + M.secondMoverEU other self) / 2
\end{ReviewVerbatim}

\paragraph{Lean proof endpoint (built separately)}\begin{ReviewVerbatim}
KR21Monoculture.PaperInterface.source_two_firm_selection_game_semantics
\end{ReviewVerbatim}

\paragraph{Recorded source-to-Spec screening}
\textbf{Verdict:} \texttt{not recorded}
\\
\textbf{Reason:} not recorded
\reviewmatch{Matches source input}
\noindent\textbf{Reviewer annotation}\par
\reviewerbox
\clearpage
\hypertarget{source-claim-fcc1dae9c26a2b31}{}
\section*{Review row 5}
\textbf{Source locator:} {\footnotesize\raggedright cited publication, lines 190-\allowbreak{}198\par}
\paragraph{Verbatim source input}
\begin{ReviewVerbatim}
[No record available.]
\end{ReviewVerbatim}


\paragraph{Expanded PaperInterface specification}
\begin{ReviewVerbatim}
∀ {n : ℕ} (trueRanking : EconCSLib.SocialChoice.Ranking.Ranking n)
  (D : MeasureTheory.Measure (KR21Monoculture.ValueProfile n)) [MeasureTheory.IsProbabilityMeasure D],
  (∀ᵐ (realizedValue : EconCSLib.SocialChoice.Ranking.Candidate n → ℝ) ∂D,
      EconCSLib.SocialChoice.Ranking.StrictlyOrderedBy trueRanking realizedValue) →
    ∀ (value : KR21Monoculture.ValueProfile n),
      EconCSLib.SocialChoice.Ranking.StrictlyOrderedBy trueRanking value →
        D Set.univ = 1 ∧
          (MeasureTheory.Measure.dirac value) Set.univ = 1 ∧
            MeasureTheory.IsProbabilityMeasure (MeasureTheory.Measure.dirac value)
\end{ReviewVerbatim}

\paragraph{Lean proof endpoint (built separately)}\begin{ReviewVerbatim}
KR21Monoculture.PaperInterface.source_model_outer_distribution_probability_and_point_mass
\end{ReviewVerbatim}

\paragraph{Recorded source-to-Spec screening}
\textbf{Verdict:} \texttt{not recorded}
\\
\textbf{Reason:} not recorded
\reviewmatch{Matches source input}
\noindent\textbf{Reviewer annotation}\par
\reviewerbox
\clearpage
\hypertarget{source-claim-a71947e08e3fd64a}{}
\section*{Review row 6}
\textbf{Source locator:} {\footnotesize\raggedright cited publication, lines 210-\allowbreak{}216\par}
\paragraph{Verbatim source input}
\begin{ReviewVerbatim}
[No record available.]
\end{ReviewVerbatim}


\paragraph{Expanded PaperInterface specification}
\begin{ReviewVerbatim}
∀ {n : ℕ} (F : KR21Monoculture.DistributionalAccuracyFamily n)
  (D : MeasureTheory.Measure (KR21Monoculture.ValueProfile n)) [MeasureTheory.IsProbabilityMeasure D] (theta : ℝ),
  0 < theta →
    ∀
      (hatom :
        ∀ (ranking : EconCSLib.SocialChoice.Ranking.Ranking n), Measurable fun value => (F.dist theta value) ranking),
      MeasureTheory.Integrable (fun x => EconCSLib.SocialChoice.Ranking.secondMoverUtility x.1 x.2.1 x.2.1)
          (F.outerIndependentPairJointLaw D theta hatom) →
        MeasureTheory.Integrable (fun x => EconCSLib.SocialChoice.Ranking.secondMoverUtility x.1 x.2.1 x.2.2)
            (F.outerIndependentPairJointLaw D theta hatom) →
          0 <
              ∫ (x : KR21Monoculture.ValueProfile n × KR21Monoculture.RankingPair n),
                if
                    EconCSLib.SocialChoice.Ranking.firstChoice x.2.1 ≠
                      EconCSLib.SocialChoice.Ranking.firstChoice x.2.2 then
                  1
                else 0 ∂F.outerIndependentPairJointLaw D theta hatom →
            have J := F.outerIndependentPairJointLaw D theta hatom;
            have d2Numerator :=
              ∫ (x : KR21Monoculture.ValueProfile n × KR21Monoculture.RankingPair n),
                if
                    EconCSLib.SocialChoice.Ranking.firstChoice x.2.1 ≠
                      EconCSLib.SocialChoice.Ranking.firstChoice x.2.2 then
                  x.1 (EconCSLib.SocialChoice.Ranking.firstChoice x.2.1) -
                    x.1 (EconCSLib.SocialChoice.Ranking.secondChoice x.2.1)
                else 0 ∂J;
            have d2Denominator :=
              ∫ (x : KR21Monoculture.ValueProfile n × KR21Monoculture.RankingPair n),
                if
                    EconCSLib.SocialChoice.Ranking.firstChoice x.2.1 ≠
                      EconCSLib.SocialChoice.Ranking.firstChoice x.2.2 then
                  1
                else 0 ∂J;
            have d2Shared :=
              ∫ (x : KR21Monoculture.ValueProfile n × KR21Monoculture.RankingPair n),
                EconCSLib.SocialChoice.Ranking.secondMoverUtility x.1 x.2.1 x.2.1 ∂J;
            have d2Independent :=
              ∫ (x : KR21Monoculture.ValueProfile n × KR21Monoculture.RankingPair n),
                EconCSLib.SocialChoice.Ranking.secondMoverUtility x.1 x.2.1 x.2.2 ∂J;
            J = D.compProd (F.independentPairKernel theta hatom) ∧
              (∀ (value : KR21Monoculture.ValueProfile n),
                  (F.independentPairKernel theta hatom) value =
                    (EconCSLib.pmfProd (F.dist theta value) (F.dist theta value)).toMeasure) ∧
                (KR21Monoculture.SourceDefinition2ConditionalAt F D theta hatom ↔ 0 < d2Numerator / d2Denominator) ∧
                  (0 < d2Numerator / d2Denominator ↔ d2Shared < d2Independent)
\end{ReviewVerbatim}

\paragraph{Lean proof endpoint (built separately)}\begin{ReviewVerbatim}
KR21Monoculture.PaperInterface.source_definition2_literal_outer_joint_semantic_complete
\end{ReviewVerbatim}

\paragraph{Recorded source-to-Spec screening}
\textbf{Verdict:} \texttt{not recorded}
\\
\textbf{Reason:} not recorded
\reviewmatch{Matches source input}
\noindent\textbf{Reviewer annotation}\par
\reviewerbox
\clearpage
\hypertarget{source-claim-faa08c37fc761736}{}
\section*{Review row 7}
\textbf{Source locator:} {\footnotesize\raggedright cited publication, lines 217-\allowbreak{}242\par}
\paragraph{Verbatim source input}
\begin{ReviewVerbatim}
Definition 3 (Preference for weaker competition.). A candidate distribution D and noisy permutation family
Fθ , exhibits a preference for weaker competition if the following holds: for all θ1 > θ2 , σ ∼ Fθ1 and
π, τ ∼ Fθ2 ,
h
i
h
i
(−{σ1 })

E π1

(−{τ1 })

< E π1

.

Intuitively, suppose we have a higher accuracy parameter θ1 and a lower accuracy parameter θ2 < θ1 ; we
draw a permutation π from Fθ2 ; and we then derive two permutations from π: π (−{σ1 }) obtained by deleting
the first-ranked element of a permutation σ drawn from the more accurate distribution Fθ1 , and π (−{τ1 })
obtained by deleting the first-ranked element of a permutation τ drawn from the less accurate distribution
Fθ2 .
Then the expected value of the first-ranked candidate in π (−{τ1 }) is strictly greater than the expected
value of the first-ranked candidate in π (−{σ1 }) — that is, when a random candidate is removed from π, the
best remaining candidate is better in expectation when the randomly removed candidate is chosen based on
a noisier ranking.
\end{ReviewVerbatim}


\paragraph{Expanded PaperInterface specification}
\begin{ReviewVerbatim}
∀ {n : ℕ} (F : KR21Monoculture.DistributionalAccuracyFamily n)
  (D : MeasureTheory.Measure (KR21Monoculture.ValueProfile n)) [MeasureTheory.IsProbabilityMeasure D]
  (thetaA thetaH : ℝ),
  0 < thetaH →
    thetaH < thetaA →
      (∀ (c : EconCSLib.SocialChoice.Ranking.Candidate n), MeasureTheory.Integrable (fun value => value c) D) →
        (∀ (pi : EconCSLib.SocialChoice.Ranking.Ranking n),
            MeasureTheory.AEStronglyMeasurable (fun value => ((F.dist thetaH value) pi).toReal) D) →
          (∀ (pi : EconCSLib.SocialChoice.Ranking.Ranking n),
              MeasureTheory.AEStronglyMeasurable (fun value => ((F.dist thetaA value) pi).toReal) D) →
            have d3Algorithm :=
              ∫ (value : KR21Monoculture.ValueProfile n),
                EconCSLib.pmfPairExp (F.dist thetaH value) (F.dist thetaA value) fun pi sigma =>
                  EconCSLib.SocialChoice.Ranking.secondMoverUtility value pi sigma ∂D;
            have d3Human :=
              ∫ (value : KR21Monoculture.ValueProfile n),
                EconCSLib.pmfPairExp (F.dist thetaH value) (F.dist thetaH value) fun pi sigma =>
                  EconCSLib.SocialChoice.Ranking.secondMoverUtility value pi sigma ∂D;
            MeasureTheory.Integrable
                (fun value =>
                  EconCSLib.pmfPairExp (F.dist thetaH value) (F.dist thetaA value) fun pi sigma =>
                    EconCSLib.SocialChoice.Ranking.secondMoverUtility value pi sigma)
                D ∧
              MeasureTheory.Integrable
                  (fun value =>
                    EconCSLib.pmfPairExp (F.dist thetaH value) (F.dist thetaH value) fun pi sigma =>
                      EconCSLib.SocialChoice.Ranking.secondMoverUtility value pi sigma)
                  D ∧
                (KR21Monoculture.SourceDefinition3At F D thetaA thetaH ↔ d3Algorithm < d3Human)
\end{ReviewVerbatim}

\paragraph{Lean proof endpoint (built separately)}\begin{ReviewVerbatim}
KR21Monoculture.PaperInterface.source_definition3_literal_outer_semantic_complete
\end{ReviewVerbatim}

\paragraph{Recorded source-to-Spec screening}
\textbf{Verdict:} \texttt{not recorded}
\\
\textbf{Reason:} not recorded
\reviewmatch{Matches source input}
\noindent\textbf{Reviewer annotation}\par
\reviewerbox
\clearpage
\hypertarget{source-claim-f88d6727dfcb8617}{}
\section*{Review row 8}
\textbf{Source locator:} {\footnotesize\raggedright cited publication, lines 176-\allowbreak{}182\par}
\paragraph{Verbatim source input}
\begin{ReviewVerbatim}
Each firm in our model has access to the same underlying pool of n candidates, which they rank using
a randomized mechanism R to get a permutation π as described above. Then, in a random order, each
firm hires the highest-ranked remaining candidate according to their ranking. Thus, if two firms both rank
candidate i first, only one of them can hire i; the other must hire the next available candidate according to
their ranking. In our model, candidates automatically accept the offer they get from a firm. For the sake
of simplicity, throughout this paper, we restrict ourselves to the case where there are two firms hiring one
candidate each, although our model readily generalizes to more complex cases.

[Next verbatim source excerpt]

Using these two conditions, we can state our theorem.
Theorem 1. Suppose that a given candidate distribution D and noisy permutation family Fθ satisfy Definition 2 (preference for the first position) and Definition 3 (preference for weaker competition).
Then, for any θH , there exists θA > θH such that using the algorithmic ranking is a strictly dominant
strategy for both firms, but social welfare would be higher if both firms used human evaluators.

[Next verbatim source excerpt]

Proving Theorem 1

With this intuition, we are ready to prove Theorem 1.
Proof of Theorem 1. For given values of θA and θH , using the algorithmic ranking is a strictly dominant
strategy as long as
UA (θA , θH ) + UAA (θA , θH ) > UH (θA , θH ) + UAH (θA , θH )

(4)

UA (θA , θH ) + UHA (θA , θH ) > UH (θA , θH ) + UHH (θA , θH )

(5)

Note that (5) is always true for θA > θH by the monotonicity assumption on Fθ : UA (θA , θH ) ≥ UH (θA , θH )
because a more accurate ranking produces a top-ranked candidate with higher expected value, and UHA (θA , θH ) ≥
UHH (θA , θH ) because this holds even conditioned on removing any candidate from the pool (in this case, the
5


[form-feed in source extraction]
candidate randomly selected by the firm that hires first). Crucially, in (5), the first firm’s random selection
is independent from the second firm’s selection; the same logic could not be used to argue that (4) always
holds for θA ≥ θH . Moreover, when θA > θH , UA (θA , θH ) > UH (θA , θH ) by the monotonicity assumption,
meaning (5) holds.
Let Ws1 s2 (θA , θH ) denote social welfare when the two firms employ strategies s1 , s2 ∈ {A, H}. Then,
when both firms use the algorithmic ranking, social welfare is
WAA (θA , θH ) = UA (θA , θH ) + UAA (θA , θH ).
By (2), Definition 2 implies that for any θ, UAA (θ, θ) < UAH (θ, θ), implying
UA (θH , θH ) + UAA (θH , θH ) < UH (θH , θH ) + UAH (θH , θH ).
However, by the optimality assumption on Fθ in Definition 1, for sufficiently large θ̂A ,
UA (θ̂A , θH ) + UAA (θ̂A , θH ) > UH (θ̂A , θH ) + UAH (θ̂A , θH ).
Note that Us1 (θA , θH ) and Us1 s2 (θA , θH ) are continuous with respect to θA for any s1 , s2 ∈ {A, H} since
they are expectations over discrete distributions with probabilities that are by assumption differentiable with
∗
> θH
respect to θA . Therefore, by the Differentiability assumption on Fθ from Definition 1, there is some θA
such that
∗
∗
∗
∗
UA (θA
, θH ) + UAA (θA
, θH ) = UH (θA
, θH ) + UAH (θA
, θH ),
(6)
i.e., given that its competitor uses the algorithmic ranking, a firm is indifferent between the two strategies.
∗
, using the algorithmic ranking is still a weakly dominant strategy. By definition of WAA ,
For such θA
∗
∗
∗
WAA (θA
, θH ) = UH (θA
, θH ) + UAH (θA
, θH ).

If both firms had instead used human evaluators, social welfare would be
∗
∗
∗
WHH (θA
, θH ) = UH (θA
, θH ) + UHH (θA
, θH ).

By Definition 3, for σ ∼ FθA∗ and π, τ ∼ FθH ,
h
i
h
i
(−{σ })
(−{τ })
E π1 1
< E π1 1 .
Note that
h
i
(−{σ })
∗
UAH (θA
, θH ) = E π1 1
h
i
(−{τ })
∗
UHH (θA
, θH ) = E π1 1
∗
∗
, θH ). As a result for θA∗ > θH , using
Thus, Definition 3 implies that for θA∗ > θH , UHH (θA
, θH ) > UAH (θA
the algorithmic ranking is a weakly dominant strategy, but
∗
∗
∗
WHH (θA
, θH ) = UH (θA
, θH ) + UHH (θA
, θH )
∗
∗
> UH (θA
, θH ) + UAH (θA
, θH )
∗
∗
= UA (θA
, θH ) + UAA (θA
, θH )
∗
= WAA (θA
, θH ),

meaning social welfare would have been higher had both firms used human evaluators.
0
We can show that this effect persists for a value θA
such that using the algorithmic ranking is a strictly
∗
dominant strategy. Intuitively, this is simply by slightly increasing θA
so the algorithmic ranking is strictly
dominant. For fixed θH , define
f (θA ) = UA (θA , θH ) + UAA (θA , θH )
g(θA ) = UH (θA , θH ) + UAH (θA , θH )
h(θA ) = UH (θA , θH ) + UHH (θA , θH )
6


[form-feed in source extraction]
0
0
0
Because (5) always holds for θA > θH , it suffices to show that there exists θA
such that g(θA
) < f (θA
)<
0
0
0
0
0
0
h(θA ). This is because g(θA ) < f (θA ) is equivalent to (4) and f (θA ) < h(θA ) is equivalent to WAA (θA , θH ) <
0
WHH (θA
, θH ).
First, note that h(θA ) is a constant, and by Definition 3, g(θA ) < h(θA ) for all θA > θH . By the
optimality assumption of Definition 1, there exists sufficiently large θ̂A such that f (θ̂A ) > g(θ̂A ). Recall
∗
∗
∗
that by definition of θA
, f (θA
) = g(θA
). Both f and g are continuous by the Differentiability assumption in
0
∗
0
0
0
Definition 1. Thus, there must exist some θA
> θA
such that g(θA
) < f (θA
) < h(θA
). This means that for
0
θA , using the algorithmic ranking is a strictly dominant strategy, but social welfare would still be larger if
both firms used human evaluators.
\end{ReviewVerbatim}


\paragraph{Expanded PaperInterface specification}
\begin{ReviewVerbatim}
∀ {n : ℕ} (F : KR21Monoculture.DistributionalAccuracyFamily n)
  (D : MeasureTheory.Measure (KR21Monoculture.ValueProfile n)),
  MeasureTheory.IsProbabilityMeasure D →
    ∀ (center : EconCSLib.SocialChoice.Ranking.Ranking n) (thetaH : ℝ),
      0 < thetaH →
        (∀ (c : EconCSLib.SocialChoice.Ranking.Candidate n), MeasureTheory.Integrable (fun value => value c) D) →
          (∀ (theta : ℝ) (pi : EconCSLib.SocialChoice.Ranking.Ranking n),
              MeasureTheory.AEStronglyMeasurable (fun value => ((F.dist theta value) pi).toReal) D) →
            (∀ (value : KR21Monoculture.ValueProfile n) (theta : ℝ),
                0 < theta →
                  ∀ (pi : EconCSLib.SocialChoice.Ranking.Ranking n),
                    ContinuousAt (fun theta' => ((F.dist theta' value) pi).toReal) theta) →
              (∀ (value : KR21Monoculture.ValueProfile n) (theta : ℝ),
                  0 < theta →
                    ∀ (pi : EconCSLib.SocialChoice.Ranking.Ranking n),
                      DifferentiableAt ℝ (fun theta' => ((F.dist theta' value) pi).toReal) theta) →
                (∀ (value : KR21Monoculture.ValueProfile n),
                    Filter.Tendsto (fun theta => ((F.dist theta value) center).toReal) Filter.atTop (nhds 1)) →
                  (∀ᵐ (value : EconCSLib.SocialChoice.Ranking.Candidate n → ℝ) ∂D,
                      EconCSLib.SocialChoice.Ranking.StrictlyOrderedBy center value) →
                    ∀
                      (hatom_measurable :
                        ∀ (theta : ℝ) (pi : EconCSLib.SocialChoice.Ranking.Ranking n),
                          Measurable fun value => (F.dist theta value) pi),
                      (∀ (theta : ℝ),
                          0 < theta →
                            MeasureTheory.Integrable
                              (fun value =>
                                EconCSLib.pmfPairExp (F.dist theta value) (F.dist theta value) fun pi sigma =>
                                  if
                                      EconCSLib.SocialChoice.Ranking.firstChoice pi ≠
                                        EconCSLib.SocialChoice.Ranking.firstChoice sigma then
                                    1
                                  else 0)
                              D) →
                        (∀ (theta : ℝ),
                            0 < theta →
                              MeasureTheory.Integrable
                                (fun value =>
                                  EconCSLib.pmfExp (F.dist theta value) fun pi =>
                                    value (EconCSLib.SocialChoice.Ranking.secondChoice pi))
                                D) →
                          (∀ (theta : ℝ),
                              0 < theta →
                                MeasureTheory.Integrable
                                  (fun value =>
                                    EconCSLib.pmfPairExp (F.dist theta value) (F.dist theta value) fun pi sigma =>
                                      EconCSLib.SocialChoice.Ranking.secondMoverUtility value pi sigma)
                                  D) →
                            (∀ (theta : ℝ),
                                0 < theta →
                                  MeasureTheory.Integrable
                                    (fun x => EconCSLib.SocialChoice.Ranking.secondMoverUtility x.1 x.2.1 x.2.1)
                                    (F.outerIndependentPairJointLaw D theta ⋯)) →
                              (∀ (theta : ℝ),
                                  0 < theta →
                                    MeasureTheory.Integrable
                                      (fun x => EconCSLib.SocialChoice.Ranking.secondMoverUtility x.1 x.2.1 x.2.2)
                                      (F.outerIndependentPairJointLaw D theta ⋯)) →
                                (∀ (theta : ℝ),
                                    0 < theta →
                                      0 <
                                        ∫ (x : KR21Monoculture.ValueProfile n × KR21Monoculture.RankingPair n),
                                          if
                                              EconCSLib.SocialChoice.Ranking.firstChoice x.2.1 ≠
                                                EconCSLib.SocialChoice.Ranking.firstChoice x.2.2 then
                                            1
                                          else 0 ∂F.outerIndependentPairJointLaw D theta ⋯) →
                                  (∀ (theta : ℝ),
                                      0 < theta →
                                        0 <
                                          (∫ (x : KR21Monoculture.ValueProfile n × KR21Monoculture.RankingPair n),
                                              if
                                                  EconCSLib.SocialChoice.Ranking.firstChoice x.2.1 ≠
                                                    EconCSLib.SocialChoice.Ranking.firstChoice x.2.2 then
                                                x.1 (EconCSLib.SocialChoice.Ranking.firstChoice x.2.1) -
                                                  x.1 (EconCSLib.SocialChoice.Ranking.secondChoice x.2.1)
                                              else 0 ∂F.outerIndependentPairJointLaw D theta ⋯) /
                                            ∫ (x : KR21Monoculture.ValueProfile n × KR21Monoculture.RankingPair n),
                                              if
                                                  EconCSLib.SocialChoice.Ranking.firstChoice x.2.1 ≠
                                                    EconCSLib.SocialChoice.Ranking.firstChoice x.2.2 then
                                                1
                                              else 0 ∂F.outerIndependentPairJointLaw D theta ⋯) →
                                    (∀ (thetaA thetaH : ℝ),
                                        0 < thetaH →
                                          thetaH < thetaA →
                                            MeasureTheory.Integrable
                                              (fun value =>
                                                EconCSLib.pmfPairExp (F.dist thetaH value) (F.dist thetaA value)
                                                  fun pi sigma =>
                                                  EconCSLib.SocialChoice.Ranking.secondMoverUtility value pi sigma)
                                              D) →
                                      (∀ (thetaA thetaH : ℝ),
                                          0 < thetaH →
                                            thetaH < thetaA →
                                              MeasureTheory.Integrable
                                                (fun value =>
                                                  EconCSLib.pmfPairExp (F.dist thetaH value) (F.dist thetaH value)
                                                    fun pi sigma =>
                                                    EconCSLib.SocialChoice.Ranking.secondMoverUtility value pi sigma)
                                                D) →
                                        (∀ (thetaA thetaH : ℝ),
                                            0 < thetaH →
                                              thetaH < thetaA →
                                                ∫ (value : KR21Monoculture.ValueProfile n),
                                                    EconCSLib.pmfPairExp (F.dist thetaH value) (F.dist thetaA value)
                                                      fun pi sigma =>
                                                      EconCSLib.SocialChoice.Ranking.secondMoverUtility value pi
                                                        sigma ∂D <
                                                  ∫ (value : KR21Monoculture.ValueProfile n),
                                                    EconCSLib.pmfPairExp (F.dist thetaH value) (F.dist thetaH value)
                                                      fun pi sigma =>
                                                      EconCSLib.SocialChoice.Ranking.secondMoverUtility value pi
                                                        sigma ∂D) →
                                          (∀ (thetaA thetaH : ℝ),
                                              0 < thetaH →
                                                thetaH < thetaA →
                                                  ∀ (value : KR21Monoculture.ValueProfile n)
                                                    (remaining : Finset (EconCSLib.SocialChoice.Ranking.Candidate n)),
                                                    remaining.Nonempty →
                                                      KR21Monoculture.expectedBestInSet (F.dist thetaH value) value
                                                          remaining ≤
                                                        KR21Monoculture.expectedBestInSet (F.dist thetaA value) value
                                                          remaining) →
                                            (∀ (thetaA thetaH : ℝ),
                                                0 < thetaH →
                                                  thetaH < thetaA →
                                                    ∀ (value : KR21Monoculture.ValueProfile n),
                                                      KR21Monoculture.expectedBestInSet (F.dist thetaH value) value
                                                          Finset.univ <
                                                        KR21Monoculture.expectedBestInSet (F.dist thetaA value) value
                                                          Finset.univ) →
                                              ∃ thetaA,
                                                thetaH < thetaA ∧
                                                  ∫ (value : KR21Monoculture.ValueProfile n),
                                                          ((F.pointFamily value).modelAt thetaA thetaH).firstMoverEU
                                                            KR21Monoculture.Strategy.human ∂D +
                                                        ∫ (value : KR21Monoculture.ValueProfile n),
                                                          ((F.pointFamily value).modelAt thetaA thetaH).secondMoverEU
                                                            KR21Monoculture.Strategy.algorithm
                                                            KR21Monoculture.Strategy.human ∂D <
                                                      ∫ (value : KR21Monoculture.ValueProfile n),
                                                          ((F.pointFamily value).modelAt thetaA thetaH).firstMoverEU
                                                            KR21Monoculture.Strategy.algorithm ∂D +
                                                        ∫ (value : KR21Monoculture.ValueProfile n),
                                                          ((F.pointFamily value).modelAt thetaA thetaH).secondMoverEU
                                                            KR21Monoculture.Strategy.algorithm
                                                            KR21Monoculture.Strategy.algorithm ∂D ∧
                                                    ∫ (value : KR21Monoculture.ValueProfile n),
                                                            ((F.pointFamily value).modelAt thetaA thetaH).firstMoverEU
                                                              KR21Monoculture.Strategy.human ∂D +
                                                          ∫ (value : KR21Monoculture.ValueProfile n),
                                                            ((F.pointFamily value).modelAt thetaA thetaH).secondMoverEU
                                                              KR21Monoculture.Strategy.human
                                                              KR21Monoculture.Strategy.human ∂D <
                                                        ∫ (value : KR21Monoculture.ValueProfile n),
                                                            ((F.pointFamily value).modelAt thetaA thetaH).firstMoverEU
                                                              KR21Monoculture.Strategy.algorithm ∂D +
                                                          ∫ (value : KR21Monoculture.ValueProfile n),
                                                            ((F.pointFamily value).modelAt thetaA thetaH).secondMoverEU
                                                              KR21Monoculture.Strategy.human
                                                              KR21Monoculture.Strategy.algorithm ∂D ∧
                                                      ∫ (value : KR21Monoculture.ValueProfile n),
                                                            ((F.pointFamily value).modelAt thetaA thetaH).firstMoverEU
                                                              KR21Monoculture.Strategy.algorithm ∂D +
                                                          ∫ (value : KR21Monoculture.ValueProfile n),
                                                            ((F.pointFamily value).modelAt thetaA thetaH).secondMoverEU
                                                              KR21Monoculture.Strategy.algorithm
                                                              KR21Monoculture.Strategy.algorithm ∂D <
                                                        ∫ (value : KR21Monoculture.ValueProfile n),
                                                            ((F.pointFamily value).modelAt thetaA thetaH).firstMoverEU
                                                              KR21Monoculture.Strategy.human ∂D +
                                                          ∫ (value : KR21Monoculture.ValueProfile n),
                                                            ((F.pointFamily value).modelAt thetaA thetaH).secondMoverEU
                                                              KR21Monoculture.Strategy.human
                                                              KR21Monoculture.Strategy.human ∂D
\end{ReviewVerbatim}

\paragraph{Lean proof endpoint (built separately)}\begin{ReviewVerbatim}
KR21Monoculture.PaperInterface.theorem1_raw_outer_literal_payoff_terminal_conclusion
\end{ReviewVerbatim}

\paragraph{Recorded source-to-Spec screening}
\textbf{Verdict:} \texttt{not recorded}
\\
\textbf{Reason:} not recorded
\reviewmatch{Matches source input}
\noindent\textbf{Reviewer annotation}\par
\reviewerbox
\clearpage
\hypertarget{source-claim-0a74bbe90425faa9}{}
\section*{Review row 9}
\textbf{Source locator:} {\footnotesize\raggedright cited publication, lines 263-\allowbreak{}276\par}
\paragraph{Verbatim source input}
\begin{ReviewVerbatim}
To do so, we need to introduce some notation. Recall that the two firms hire in a random order. Given
a candidate distribution D, let Us (θA , θH ) denote the expected utility of the first firm to hire a candidate
when using ranking s, where s ∈ {A, H} is either the algorithmic ranking or the ranking generated by a
human evaluator respectively. Similarly, let Us1 s2 (θA , θH ) be the expected utility of the second firm to hire
given that the first firm used strategy s1 and the second firm uses strategy s2 , where again s1 , s2 ∈ {A, H}.
Finally, let π, σ ∼ Fθ .
In what follows, we will show that for any θ,
E [π1 − π2 | π1 6= σ1 ] > 0 ⇐⇒ UAH (θ, θ) > UAA (θ, θ).

(2)

In other words, whenever a ranking model meets Definition 2, the firm chosen to select second will prefer to
use an independent ranking mechanism from it’s competitor, given that the ranking mechanisms are equally
accurate.
\end{ReviewVerbatim}


\paragraph{Expanded PaperInterface specification}
\begin{ReviewVerbatim}
∀ {n : ℕ} (F : KR21Monoculture.DistributionalAccuracyFamily n)
  (D : MeasureTheory.Measure (KR21Monoculture.ValueProfile n)) [MeasureTheory.IsProbabilityMeasure D] (theta : ℝ),
  0 < theta →
    ∀
      (hatom :
        ∀ (ranking : EconCSLib.SocialChoice.Ranking.Ranking n), Measurable fun value => (F.dist theta value) ranking),
      MeasureTheory.Integrable (fun x => EconCSLib.SocialChoice.Ranking.secondMoverUtility x.1 x.2.1 x.2.1)
          (F.outerIndependentPairJointLaw D theta hatom) →
        MeasureTheory.Integrable (fun x => EconCSLib.SocialChoice.Ranking.secondMoverUtility x.1 x.2.1 x.2.2)
            (F.outerIndependentPairJointLaw D theta hatom) →
          0 <
              ∫ (x : KR21Monoculture.ValueProfile n × KR21Monoculture.RankingPair n),
                if
                    EconCSLib.SocialChoice.Ranking.firstChoice x.2.1 ≠
                      EconCSLib.SocialChoice.Ranking.firstChoice x.2.2 then
                  1
                else 0 ∂F.outerIndependentPairJointLaw D theta hatom →
            have J := F.outerIndependentPairJointLaw D theta hatom;
            have numerator :=
              ∫ (x : KR21Monoculture.ValueProfile n × KR21Monoculture.RankingPair n),
                if
                    EconCSLib.SocialChoice.Ranking.firstChoice x.2.1 ≠
                      EconCSLib.SocialChoice.Ranking.firstChoice x.2.2 then
                  x.1 (EconCSLib.SocialChoice.Ranking.firstChoice x.2.1) -
                    x.1 (EconCSLib.SocialChoice.Ranking.secondChoice x.2.1)
                else 0 ∂J;
            have denominator :=
              ∫ (x : KR21Monoculture.ValueProfile n × KR21Monoculture.RankingPair n),
                if
                    EconCSLib.SocialChoice.Ranking.firstChoice x.2.1 ≠
                      EconCSLib.SocialChoice.Ranking.firstChoice x.2.2 then
                  1
                else 0 ∂J;
            have UAA :=
              ∫ (x : KR21Monoculture.ValueProfile n × KR21Monoculture.RankingPair n),
                EconCSLib.SocialChoice.Ranking.secondMoverUtility x.1 x.2.1 x.2.1 ∂J;
            have UAH :=
              ∫ (x : KR21Monoculture.ValueProfile n × KR21Monoculture.RankingPair n),
                EconCSLib.SocialChoice.Ranking.secondMoverUtility x.1 x.2.1 x.2.2 ∂J;
            J = D.compProd (F.independentPairKernel theta hatom) ∧
              (∀ (value : KR21Monoculture.ValueProfile n),
                  (F.independentPairKernel theta hatom) value =
                    (EconCSLib.pmfProd (F.dist theta value) (F.dist theta value)).toMeasure) ∧
                (0 < numerator / denominator ↔ UAA < UAH)
\end{ReviewVerbatim}

\paragraph{Lean proof endpoint (built separately)}\begin{ReviewVerbatim}
KR21Monoculture.PaperInterface.equation2_outer_joint_payoff_equivalence_semantic_complete
\end{ReviewVerbatim}

\paragraph{Recorded source-to-Spec screening}
\textbf{Verdict:} \texttt{not recorded}
\\
\textbf{Reason:} not recorded
\reviewmatch{Matches source input}
\noindent\textbf{Reviewer annotation}\par
\reviewerbox
\clearpage
\hypertarget{source-claim-36b115610854f337}{}
\section*{Review row 10}
\textbf{Source locator:} {\footnotesize\raggedright cited publication, lines 277-\allowbreak{}290\par}
\paragraph{Verbatim source input}
\begin{ReviewVerbatim}
First, we can write
UAH (θA , θH ) = E [π1 · 1π1 6=σ1 + π2 · 1π1 =σ1 ]
UAA (θA , θH ) = E [σ2 ]
= E [σ2 · 1π1 6=σ1 + σ2 · 1π1 =σ1 ]
Thus,
UAH (θA , θH ) − UAA (θA , θH )
= E [(π1 − σ2 ) · 1π1 6=σ1 + (π2 − σ2 ) · 1π1 =σ1 ] .
Conditioned on either π1 = σ1 or π1 6= σ1 , π2 and σ2 are identically distributed and therefore have equal
expectations. As a result,
UAH (θA , θH ) − UAA (θA , θH ) = E [(π1 − π2 ) · 1π1 6=σ1 ] ,

(3)

which implies (2). Thus, whenever a ranking model meets Definition 2, firms rationally prefer independent
\end{ReviewVerbatim}


\paragraph{Expanded PaperInterface specification}
\begin{ReviewVerbatim}
∀ {n : ℕ} (F : KR21Monoculture.DistributionalAccuracyFamily n)
  (D : MeasureTheory.Measure (KR21Monoculture.ValueProfile n)) [MeasureTheory.IsProbabilityMeasure D] (theta : ℝ)
  (hatom :
    ∀ (ranking : EconCSLib.SocialChoice.Ranking.Ranking n), Measurable fun value => (F.dist theta value) ranking),
  MeasureTheory.Integrable (fun x => EconCSLib.SocialChoice.Ranking.secondMoverUtility x.1 x.2.1 x.2.1)
      (F.outerIndependentPairJointLaw D theta hatom) →
    MeasureTheory.Integrable (fun x => EconCSLib.SocialChoice.Ranking.secondMoverUtility x.1 x.2.1 x.2.2)
        (F.outerIndependentPairJointLaw D theta hatom) →
      have J := F.outerIndependentPairJointLaw D theta hatom;
      have numerator :=
        ∫ (x : KR21Monoculture.ValueProfile n × KR21Monoculture.RankingPair n),
          if EconCSLib.SocialChoice.Ranking.firstChoice x.2.1 ≠ EconCSLib.SocialChoice.Ranking.firstChoice x.2.2 then
            x.1 (EconCSLib.SocialChoice.Ranking.firstChoice x.2.1) -
              x.1 (EconCSLib.SocialChoice.Ranking.secondChoice x.2.1)
          else 0 ∂J;
      have UAA :=
        ∫ (x : KR21Monoculture.ValueProfile n × KR21Monoculture.RankingPair n),
          EconCSLib.SocialChoice.Ranking.secondMoverUtility x.1 x.2.1 x.2.1 ∂J;
      have UAH :=
        ∫ (x : KR21Monoculture.ValueProfile n × KR21Monoculture.RankingPair n),
          EconCSLib.SocialChoice.Ranking.secondMoverUtility x.1 x.2.1 x.2.2 ∂J;
      J = D.compProd (F.independentPairKernel theta hatom) ∧
        (∀ (value : KR21Monoculture.ValueProfile n),
            (F.independentPairKernel theta hatom) value =
              (EconCSLib.pmfProd (F.dist theta value) (F.dist theta value)).toMeasure) ∧
          UAH - UAA = numerator
\end{ReviewVerbatim}

\paragraph{Lean proof endpoint (built separately)}\begin{ReviewVerbatim}
KR21Monoculture.PaperInterface.equation3_outer_joint_payoff_identity_semantic_complete
\end{ReviewVerbatim}

\paragraph{Recorded source-to-Spec screening}
\textbf{Verdict:} \texttt{not recorded}
\\
\textbf{Reason:} not recorded
\reviewmatch{Matches source input}
\noindent\textbf{Reviewer annotation}\par
\reviewerbox
\clearpage
\hypertarget{source-claim-ea5f4fa858fd5854}{}
\section*{Review row 11}
\textbf{Source locator:} {\footnotesize\raggedright cited publication, lines 176-\allowbreak{}182\par}
\paragraph{Verbatim source input}
\begin{ReviewVerbatim}
Each firm in our model has access to the same underlying pool of n candidates, which they rank using
a randomized mechanism R to get a permutation π as described above. Then, in a random order, each
firm hires the highest-ranked remaining candidate according to their ranking. Thus, if two firms both rank
candidate i first, only one of them can hire i; the other must hire the next available candidate according to
their ranking. In our model, candidates automatically accept the offer they get from a firm. For the sake
of simplicity, throughout this paper, we restrict ourselves to the case where there are two firms hiring one
candidate each, although our model readily generalizes to more complex cases.

[Next verbatim source excerpt]

With this intuition, we are ready to prove Theorem 1.
Proof of Theorem 1. For given values of θA and θH , using the algorithmic ranking is a strictly dominant
strategy as long as
UA (θA , θH ) + UAA (θA , θH ) > UH (θA , θH ) + UAH (θA , θH )

(4)
\end{ReviewVerbatim}


\paragraph{Expanded PaperInterface specification}
\begin{ReviewVerbatim}
∀ {n : ℕ} (M : KR21Monoculture.Model n),
  M.firstMoverEU KR21Monoculture.Strategy.algorithm +
        M.secondMoverEU KR21Monoculture.Strategy.algorithm KR21Monoculture.Strategy.algorithm >
      M.firstMoverEU KR21Monoculture.Strategy.human +
        M.secondMoverEU KR21Monoculture.Strategy.algorithm KR21Monoculture.Strategy.human ↔
    M.payoffAgainst KR21Monoculture.Strategy.algorithm KR21Monoculture.Strategy.algorithm >
      M.payoffAgainst KR21Monoculture.Strategy.human KR21Monoculture.Strategy.algorithm
\end{ReviewVerbatim}

\paragraph{Lean proof endpoint (built separately)}\begin{ReviewVerbatim}
KR21Monoculture.PaperInterface.equation4_algorithm_best_response_against_algorithm_iff
\end{ReviewVerbatim}

\paragraph{Recorded source-to-Spec screening}
\textbf{Verdict:} \texttt{not recorded}
\\
\textbf{Reason:} not recorded
\reviewmatch{Matches source input}
\noindent\textbf{Reviewer annotation}\par
\reviewerbox
\clearpage
\hypertarget{source-claim-9d6309ccb42ce9df}{}
\section*{Review row 12}
\textbf{Source locator:} {\footnotesize\raggedright cited publication, lines 176-\allowbreak{}182\par}
\paragraph{Verbatim source input}
\begin{ReviewVerbatim}
Each firm in our model has access to the same underlying pool of n candidates, which they rank using
a randomized mechanism R to get a permutation π as described above. Then, in a random order, each
firm hires the highest-ranked remaining candidate according to their ranking. Thus, if two firms both rank
candidate i first, only one of them can hire i; the other must hire the next available candidate according to
their ranking. In our model, candidates automatically accept the offer they get from a firm. For the sake
of simplicity, throughout this paper, we restrict ourselves to the case where there are two firms hiring one
candidate each, although our model readily generalizes to more complex cases.

[Next verbatim source excerpt]


UA (θA , θH ) + UHA (θA , θH ) > UH (θA , θH ) + UHH (θA , θH )

(5)

Note that (5) is always true for θA > θH by the monotonicity assumption on Fθ : UA (θA , θH ) ≥ UH (θA , θH )
because a more accurate ranking produces a top-ranked candidate with higher expected value, and UHA (θA , θH ) ≥
UHH (θA , θH ) because this holds even conditioned on removing any candidate from the pool (in this case, the
5


[form-feed in source extraction]
candidate randomly selected by the firm that hires first). Crucially, in (5), the first firm’s random selection
is independent from the second firm’s selection; the same logic could not be used to argue that (4) always
holds for θA ≥ θH . Moreover, when θA > θH , UA (θA , θH ) > UH (θA , θH ) by the monotonicity assumption,
meaning (5) holds.
\end{ReviewVerbatim}


\paragraph{Expanded PaperInterface specification}
\begin{ReviewVerbatim}
∀ {n : ℕ} (F : KR21Monoculture.AccuracyFamily n) (thetaA thetaH : ℝ),
  0 < thetaH →
    thetaH < thetaA →
      (∀ (remaining : Finset (EconCSLib.SocialChoice.Ranking.Candidate n)),
          remaining.Nonempty →
            KR21Monoculture.expectedBestInSet (F.dist thetaH) F.value remaining ≤
              KR21Monoculture.expectedBestInSet (F.dist thetaA) F.value remaining) →
        KR21Monoculture.expectedBestInSet (F.dist thetaH) F.value Finset.univ <
            KR21Monoculture.expectedBestInSet (F.dist thetaA) F.value Finset.univ →
          (F.modelAt thetaA thetaH).firstMoverEU KR21Monoculture.Strategy.algorithm +
              (F.modelAt thetaA thetaH).secondMoverEU KR21Monoculture.Strategy.human
                KR21Monoculture.Strategy.algorithm >
            (F.modelAt thetaA thetaH).firstMoverEU KR21Monoculture.Strategy.human +
              (F.modelAt thetaA thetaH).secondMoverEU KR21Monoculture.Strategy.human KR21Monoculture.Strategy.human
\end{ReviewVerbatim}

\paragraph{Lean proof endpoint (built separately)}\begin{ReviewVerbatim}
KR21Monoculture.PaperInterface.equation5_from_literal_definition1_removal
\end{ReviewVerbatim}

\paragraph{Recorded source-to-Spec screening}
\textbf{Verdict:} \texttt{not recorded}
\\
\textbf{Reason:} not recorded
\reviewmatch{Matches source input}
\noindent\textbf{Reviewer annotation}\par
\reviewerbox
\clearpage
\hypertarget{source-claim-fa14b138c169a341}{}
\section*{Review row 13}
\textbf{Source locator:} {\footnotesize\raggedright cited publication, lines 176-\allowbreak{}182\par}
\paragraph{Verbatim source input}
\begin{ReviewVerbatim}
Each firm in our model has access to the same underlying pool of n candidates, which they rank using
a randomized mechanism R to get a permutation π as described above. Then, in a random order, each
firm hires the highest-ranked remaining candidate according to their ranking. Thus, if two firms both rank
candidate i first, only one of them can hire i; the other must hire the next available candidate according to
their ranking. In our model, candidates automatically accept the offer they get from a firm. For the sake
of simplicity, throughout this paper, we restrict ourselves to the case where there are two firms hiring one
candidate each, although our model readily generalizes to more complex cases.

[Next verbatim source excerpt]

The choice of which ranking mechanism to use leads to a game-theoretic setting: both firms know the
accuracy parameters of the human evaluators (θH ) and the algorithm (θA ), and they must decide whether to
3


[form-feed in source extraction]
use a human evaluator or the algorithm. This choice introduces a subtlety: for many ranking models, a firm’s
rational behavior depends not only on the accuracy of the ranking mechanism, but also on the underlying
candidate values x1 , . . . , xn . Thus, to fully specify a firm’s behavior, we assume that x1 , . . . , xn are drawn
from a known joint distribution D. Our main results will hold for any D, meaning they apply even when the
candidate values (but not their identities) are deterministically known.

[Next verbatim source excerpt]

Definition 2 (Preference for the first position.). A candidate distribution D and noisy permutation family
Fθ exhibits a preference for the first position if for all θ > 0, if π, σ ∼ Fθ ,
E [π1 − π2 | π1 6= σ1 ] > 0.

In other words, for any θ > 0, suppose we draw two permutations π and σ independently from Fθ , and
suppose that the first-ranked candidates differ in π and σ. Then the expected value of the first-ranked
candidate in π is strictly greater than the expected value of the second-ranked candidate in π.

[Next verbatim source excerpt]

To do so, we need to introduce some notation. Recall that the two firms hire in a random order. Given
a candidate distribution D, let Us (θA , θH ) denote the expected utility of the first firm to hire a candidate
when using ranking s, where s ∈ {A, H} is either the algorithmic ranking or the ranking generated by a
human evaluator respectively. Similarly, let Us1 s2 (θA , θH ) be the expected utility of the second firm to hire
given that the first firm used strategy s1 and the second firm uses strategy s2 , where again s1 , s2 ∈ {A, H}.
Finally, let π, σ ∼ Fθ .
In what follows, we will show that for any θ,
E [π1 − π2 | π1 6= σ1 ] > 0 ⇐⇒ UAH (θ, θ) > UAA (θ, θ).

(2)

In other words, whenever a ranking model meets Definition 2, the firm chosen to select second will prefer to

[Next verbatim source excerpt]

Let Ws1 s2 (θA , θH ) denote social welfare when the two firms employ strategies s1 , s2 ∈ {A, H}. Then,
when both firms use the algorithmic ranking, social welfare is
WAA (θA , θH ) = UA (θA , θH ) + UAA (θA , θH ).
By (2), Definition 2 implies that for any θ, UAA (θ, θ) < UAH (θ, θ), implying
UA (θH , θH ) + UAA (θH , θH ) < UH (θH , θH ) + UAH (θH , θH ).
However, by the optimality assumption on Fθ in Definition 1, for sufficiently large θ̂A ,
UA (θ̂A , θH ) + UAA (θ̂A , θH ) > UH (θ̂A , θH ) + UAH (θ̂A , θH ).
Note that Us1 (θA , θH ) and Us1 s2 (θA , θH ) are continuous with respect to θA for any s1 , s2 ∈ {A, H} since
they are expectations over discrete distributions with probabilities that are by assumption differentiable with
∗
> θH
respect to θA . Therefore, by the Differentiability assumption on Fθ from Definition 1, there is some θA
such that
∗
∗
∗
∗
UA (θA
, θH ) + UAA (θA
, θH ) = UH (θA
, θH ) + UAH (θA
, θH ),
(6)
i.e., given that its competitor uses the algorithmic ranking, a firm is indifferent between the two strategies.
∗
, using the algorithmic ranking is still a weakly dominant strategy. By definition of WAA ,
For such θA
∗
∗
∗
WAA (θA
, θH ) = UH (θA
, θH ) + UAH (θA
, θH ).
\end{ReviewVerbatim}


\paragraph{Expanded PaperInterface specification}
\begin{ReviewVerbatim}
∀ {n : ℕ} (F : KR21Monoculture.DistributionalAccuracyFamily n)
  (D : MeasureTheory.Measure (KR21Monoculture.ValueProfile n)),
  MeasureTheory.IsProbabilityMeasure D →
    ∀ (center : EconCSLib.SocialChoice.Ranking.Ranking n) (thetaH : ℝ),
      0 < thetaH →
        (∀ (c : EconCSLib.SocialChoice.Ranking.Candidate n), MeasureTheory.Integrable (fun value => value c) D) →
          (∀ (theta : ℝ) (pi : EconCSLib.SocialChoice.Ranking.Ranking n),
              MeasureTheory.AEStronglyMeasurable (fun value => ((F.dist theta value) pi).toReal) D) →
            (∀ (value : KR21Monoculture.ValueProfile n) (theta : ℝ),
                0 < theta →
                  ∀ (pi : EconCSLib.SocialChoice.Ranking.Ranking n),
                    ContinuousAt (fun theta' => ((F.dist theta' value) pi).toReal) theta) →
              (∀ (value : KR21Monoculture.ValueProfile n) (theta : ℝ),
                  0 < theta →
                    ∀ (pi : EconCSLib.SocialChoice.Ranking.Ranking n),
                      DifferentiableAt ℝ (fun theta' => ((F.dist theta' value) pi).toReal) theta) →
                (∀ (value : KR21Monoculture.ValueProfile n),
                    Filter.Tendsto (fun theta => ((F.dist theta value) center).toReal) Filter.atTop (nhds 1)) →
                  (∀ᵐ (value : EconCSLib.SocialChoice.Ranking.Candidate n → ℝ) ∂D,
                      EconCSLib.SocialChoice.Ranking.StrictlyOrderedBy center value) →
                    ∀
                      (hatom_measurable :
                        ∀ (theta : ℝ) (pi : EconCSLib.SocialChoice.Ranking.Ranking n),
                          Measurable fun value => (F.dist theta value) pi),
                      (∀ (theta : ℝ),
                          0 < theta →
                            MeasureTheory.Integrable
                              (fun value => KR21Monoculture.disagreementProb (F.dist theta value)) D) →
                        (∀ (theta : ℝ),
                            0 < theta →
                              MeasureTheory.Integrable
                                (fun value =>
                                  EconCSLib.SocialChoice.Ranking.expectedSecondMoverShared (F.dist theta value) value)
                                D) →
                          (∀ (theta : ℝ),
                              0 < theta →
                                MeasureTheory.Integrable
                                  (fun value =>
                                    EconCSLib.SocialChoice.Ranking.expectedSecondMoverIndependent (F.dist theta value)
                                      (F.dist theta value) value)
                                  D) →
                            (∀ (theta : ℝ),
                                0 < theta →
                                  MeasureTheory.Integrable
                                    KR21Monoculture.DistributionalAccuracyFamily.jointSharedSecondMoverPayoff
                                    (F.outerIndependentPairJointLaw D theta ⋯)) →
                              (∀ (theta : ℝ),
                                  0 < theta →
                                    MeasureTheory.Integrable
                                      KR21Monoculture.DistributionalAccuracyFamily.jointIndependentSecondMoverPayoff
                                      (F.outerIndependentPairJointLaw D theta ⋯)) →
                                (∀ (theta : ℝ),
                                    0 < theta →
                                      0 <
                                        ∫ (x : KR21Monoculture.ValueProfile n × KR21Monoculture.RankingPair n),
                                          if KR21Monoculture.disagreementEvent x.2 then 1
                                          else 0 ∂F.outerIndependentPairJointLaw D theta ⋯) →
                                  (∀ (theta : ℝ), 0 < theta → 0 < F.jointLawDisagreementConditionalGain D theta ⋯) →
                                    ∃ thetaA,
                                      thetaH < thetaA ∧
                                        ∫ (value : KR21Monoculture.ValueProfile n),
                                              ((F.pointFamily value).modelAt thetaA thetaH).firstMoverEU
                                                KR21Monoculture.Strategy.algorithm ∂D +
                                            ∫ (value : KR21Monoculture.ValueProfile n),
                                              ((F.pointFamily value).modelAt thetaA thetaH).secondMoverEU
                                                KR21Monoculture.Strategy.algorithm
                                                KR21Monoculture.Strategy.algorithm ∂D =
                                          ∫ (value : KR21Monoculture.ValueProfile n),
                                              ((F.pointFamily value).modelAt thetaA thetaH).firstMoverEU
                                                KR21Monoculture.Strategy.human ∂D +
                                            ∫ (value : KR21Monoculture.ValueProfile n),
                                              ((F.pointFamily value).modelAt thetaA thetaH).secondMoverEU
                                                KR21Monoculture.Strategy.algorithm KR21Monoculture.Strategy.human ∂D
\end{ReviewVerbatim}

\paragraph{Lean proof endpoint (built separately)}\begin{ReviewVerbatim}
KR21Monoculture.PaperInterface.equation6_outer_payoff_indifference_from_source_conditions
\end{ReviewVerbatim}

\paragraph{Recorded source-to-Spec screening}
\textbf{Verdict:} \texttt{not recorded}
\\
\textbf{Reason:} not recorded
\reviewmatch{Matches source input}
\noindent\textbf{Reviewer annotation}\par
\reviewerbox
\clearpage
\hypertarget{source-claim-bd2d52e2663b165c}{}
\section*{Review row 14}
\textbf{Source locator:} {\footnotesize\raggedright cited publication, lines 493-\allowbreak{}504\par}
\paragraph{Verbatim source input}
\begin{ReviewVerbatim}
[No record available.]
\end{ReviewVerbatim}


\paragraph{Expanded PaperInterface specification}
\begin{ReviewVerbatim}
∀ {n : ℕ} (E : MeasureTheory.Measure ℝ) [inst : MeasureTheory.IsProbabilityMeasure E],
  MeasureTheory.MemLp id 2 E →
    0 < ProbabilityTheory.variance id E →
      ∀ (value : EconCSLib.SocialChoice.Ranking.Candidate n → ℝ) (theta : ℝ),
        0 < theta →
          ProbabilityTheory.iIndepFun (fun i epsilon => epsilon i)
              (KR21Monoculture.sourceRUMNormalizedIIDNoiseLaw E √(ProbabilityTheory.variance id E)) ∧
            MeasureTheory.MemLp id 2
                (KR21Monoculture.sourceRUMNormalizedScalarNoiseLaw E √(ProbabilityTheory.variance id E)) ∧
              (∀ (i : EconCSLib.SocialChoice.Ranking.Candidate n),
                  ProbabilityTheory.variance (fun epsilon => epsilon i)
                      (KR21Monoculture.sourceRUMNormalizedIIDNoiseLaw E √(ProbabilityTheory.variance id E)) =
                    1) ∧
                KR21Monoculture.paper_appendixA_scaledNoiseRankingPMF (MeasureTheory.Measure.pi fun x => E) value
                    theta =
                  KR21Monoculture.paper_appendixA_scaledNoiseRankingPMF
                    (KR21Monoculture.sourceRUMNormalizedIIDNoiseLaw E √(ProbabilityTheory.variance id E)) value
                    (theta / √(ProbabilityTheory.variance id E))
\end{ReviewVerbatim}

\paragraph{Lean proof endpoint (built separately)}\begin{ReviewVerbatim}
KR21Monoculture.PaperInterface.section31_iidFinitePositiveVariance_normalization
\end{ReviewVerbatim}

\paragraph{Recorded source-to-Spec screening}
\textbf{Verdict:} \texttt{not recorded}
\\
\textbf{Reason:} not recorded
\reviewmatch{Matches source input}
\noindent\textbf{Reviewer annotation}\par
\reviewerbox
\clearpage
\hypertarget{source-claim-cc8111af16950cc8}{}
\section*{Review row 15}
\textbf{Source locator:} {\footnotesize\raggedright cited publication, lines 176-\allowbreak{}182\par}
\paragraph{Verbatim source input}
\begin{ReviewVerbatim}
Each firm in our model has access to the same underlying pool of n candidates, which they rank using
a randomized mechanism R to get a permutation π as described above. Then, in a random order, each
firm hires the highest-ranked remaining candidate according to their ranking. Thus, if two firms both rank
candidate i first, only one of them can hire i; the other must hire the next available candidate according to
their ranking. In our model, candidates automatically accept the offer they get from a firm. For the sake
of simplicity, throughout this paper, we restrict ourselves to the case where there are two firms hiring one
candidate each, although our model readily generalizes to more complex cases.

[Next verbatim source excerpt]

Because RUMs are notoriously difficult to work with analytically, we restrict ourselves to the case where
n = 3, i.e., there are 3 candidates. Under this restriction, we can show that for Gaussian and Laplacian noise
7


[form-feed in source extraction]
Figure 3: Regions for different equilibria. When human evaluators are more accurate than the algorithm,
both firms decide to employ humans (HH). When the algorithm is significantly more accurate, both firms
use the algorithm (AA). When the algorithm is slightly more accurate than human evaluators, two possible
equilibria exist: (1) one firm uses the algorithm and the other employs a human (AH), or (2) both decide
whether to use the algorithm with some probability p. The shaded portion of the green AA region depicts
where social welfare is smaller at the AA equilibrium than it would be if both firms used human evaluators.
distributions, Definition 2 and 3 — the two conditions of Theorem 1 — are met, regardless of the candidate
distribution D. We defer the proof to Appendix C.
Theorem 2. Let Fθ be the family of RUMs with either Gaussian or Laplacian noise with standard deviation
1/θ. Then, for any candidate distribution D over 3 candidates, the conditions of Theorem 1 are satisfied.
It might be tempting to generalize Theorem 2 to other distributions and more candidates; however,
certain noise and candidate distributions violate the conditions of Theorem 1. Even for 3-candidate RUMs,
there exist distributions for which each of the conditions is violated; we provide such examples in Appendix B.
\end{ReviewVerbatim}


\paragraph{Expanded PaperInterface specification}
\begin{ReviewVerbatim}
(∀ {theta x1 x2 x3 : ℝ},
    0 < theta →
      MeasureTheory.MeasurePreserving (⇑(KR21Monoculture.rightTripleToCandidateMeasurableEquiv ℝ))
          ((ProbabilityTheory.gaussianReal 0 1).prod
            ((ProbabilityTheory.gaussianReal 0 1).prod (ProbabilityTheory.gaussianReal 0 1)))
          (MeasureTheory.Measure.pi fun x => ProbabilityTheory.gaussianReal 0 1) ∧
        ((KR21Monoculture.theorem2GaussianScaledNoiseFamily x1 x2 x3).dist theta).toMeasure =
            MeasureTheory.Measure.map
              (fun epsilon =>
                EconCSLib.SocialChoice.Ranking.rankByScore fun i =>
                  KR21Monoculture.threeCandidateValueProfile x1 x2 x3 i +
                    KR21Monoculture.rightTripleToCandidateFunction epsilon i / theta)
              ((ProbabilityTheory.gaussianReal 0 1).prod
                ((ProbabilityTheory.gaussianReal 0 1).prod (ProbabilityTheory.gaussianReal 0 1))) ∧
          (KR21Monoculture.theorem2GaussianScaledNoiseFamily x1 x2 x3).dist theta =
            KR21Monoculture.gaussianThreeCandidateRankingLaw theta x1 x2 x3) ∧
  ProbabilityTheory.variance id (ProbabilityTheory.gaussianReal 0 1) = 1 ∧
    ProbabilityTheory.variance id
          (KR21Monoculture.w11BaseNoiseLaw KR21Monoculture.sourceUnitVarianceLaplaceBaseDensity) =
        1 ∧
      (∀ {theta x1 x2 x3 : ℝ},
          0 < theta →
            MeasureTheory.MeasurePreserving (⇑(KR21Monoculture.rightTripleToCandidateMeasurableEquiv ℝ))
                ((KR21Monoculture.theorem7LaplaceMeasure (√2) 0).prod
                  ((KR21Monoculture.theorem7LaplaceMeasure (√2) 0).prod
                    (KR21Monoculture.theorem7LaplaceMeasure (√2) 0)))
                (MeasureTheory.Measure.pi fun x => KR21Monoculture.theorem7LaplaceMeasure (√2) 0) ∧
              ((KR21Monoculture.sourceUnitVarianceLaplaceScaledNoiseFamily x1 x2 x3).dist theta).toMeasure =
                  MeasureTheory.Measure.map
                    (fun epsilon =>
                      EconCSLib.SocialChoice.Ranking.rankByScore fun i =>
                        KR21Monoculture.threeCandidateValueProfile x1 x2 x3 i +
                          KR21Monoculture.rightTripleToCandidateFunction epsilon i / theta)
                    ((KR21Monoculture.theorem7LaplaceMeasure (√2) 0).prod
                      ((KR21Monoculture.theorem7LaplaceMeasure (√2) 0).prod
                        (KR21Monoculture.theorem7LaplaceMeasure (√2) 0))) ∧
                (KR21Monoculture.sourceUnitVarianceLaplaceScaledNoiseFamily x1 x2 x3).dist theta =
                  KR21Monoculture.sourceUnitVarianceLaplaceThreeCandidateRankingLaw theta x1 x2 x3) ∧
        (∀ {x1 x2 x3 : ℝ},
            x2 < x1 →
              x3 < x2 →
                (∀ (theta : ℝ),
                    0 < theta →
                      ∀ (pi : EconCSLib.SocialChoice.Ranking.Ranking 1),
                        ContinuousAt
                            (fun theta' =>
                              (((KR21Monoculture.theorem2GaussianScaledNoiseFamily x1 x2 x3).dist theta') pi).toReal)
                            theta ∧
                          DifferentiableAt ℝ
                            (fun theta' =>
                              (((KR21Monoculture.theorem2GaussianScaledNoiseFamily x1 x2 x3).dist theta') pi).toReal)
                            theta) ∧
                  Filter.Tendsto
                      (fun theta =>
                        (((KR21Monoculture.theorem2GaussianScaledNoiseFamily x1 x2 x3).dist theta)
                            KR21Monoculture.rum3Ranking012).toReal)
                      Filter.atTop (nhds 1) ∧
                    ∀ (thetaA thetaH : ℝ),
                      0 < thetaH →
                        thetaH < thetaA →
                          (∀ (remaining : Finset (EconCSLib.SocialChoice.Ranking.Candidate 1)),
                              remaining.Nonempty →
                                KR21Monoculture.expectedBestInSet
                                    ((KR21Monoculture.theorem2GaussianScaledNoiseFamily x1 x2 x3).dist thetaH)
                                    (KR21Monoculture.threeCandidateValueProfile x1 x2 x3) remaining ≤
                                  KR21Monoculture.expectedBestInSet
                                    ((KR21Monoculture.theorem2GaussianScaledNoiseFamily x1 x2 x3).dist thetaA)
                                    (KR21Monoculture.threeCandidateValueProfile x1 x2 x3) remaining) ∧
                            KR21Monoculture.expectedBestInSet
                                ((KR21Monoculture.theorem2GaussianScaledNoiseFamily x1 x2 x3).dist thetaH)
                                (KR21Monoculture.threeCandidateValueProfile x1 x2 x3) Finset.univ <
                              KR21Monoculture.expectedBestInSet
                                ((KR21Monoculture.theorem2GaussianScaledNoiseFamily x1 x2 x3).dist thetaA)
                                (KR21Monoculture.threeCandidateValueProfile x1 x2 x3) Finset.univ) ∧
          (∀ {x1 x2 x3 : ℝ},
              x2 < x1 →
                x3 < x2 →
                  (∀ (theta : ℝ),
                      0 < theta →
                        ∀ (pi : EconCSLib.SocialChoice.Ranking.Ranking 1),
                          ContinuousAt
                              (fun theta' =>
                                (((KR21Monoculture.sourceUnitVarianceLaplaceScaledNoiseFamily x1 x2 x3).dist theta')
                                    pi).toReal)
                              theta ∧
                            DifferentiableAt ℝ
                              (fun theta' =>
                                (((KR21Monoculture.sourceUnitVarianceLaplaceScaledNoiseFamily x1 x2 x3).dist theta')
                                    pi).toReal)
                              theta) ∧
                    Filter.Tendsto
                        (fun theta =>
                          (((KR21Monoculture.sourceUnitVarianceLaplaceScaledNoiseFamily x1 x2 x3).dist theta)
                              KR21Monoculture.rum3Ranking012).toReal)
                        Filter.atTop (nhds 1) ∧
                      ∀ (thetaA thetaH : ℝ),
                        0 < thetaH →
                          thetaH < thetaA →
                            (∀ (remaining : Finset (EconCSLib.SocialChoice.Ranking.Candidate 1)),
                                remaining.Nonempty →
                                  KR21Monoculture.expectedBestInSet
                                      ((KR21Monoculture.sourceUnitVarianceLaplaceScaledNoiseFamily x1 x2 x3).dist
                                        thetaH)
                                      (KR21Monoculture.threeCandidateValueProfile x1 x2 x3) remaining ≤
                                    KR21Monoculture.expectedBestInSet
                                      ((KR21Monoculture.sourceUnitVarianceLaplaceScaledNoiseFamily x1 x2 x3).dist
                                        thetaA)
                                      (KR21Monoculture.threeCandidateValueProfile x1 x2 x3) remaining) ∧
                              KR21Monoculture.expectedBestInSet
                                  ((KR21Monoculture.sourceUnitVarianceLaplaceScaledNoiseFamily x1 x2 x3).dist thetaH)
                                  (KR21Monoculture.threeCandidateValueProfile x1 x2 x3) Finset.univ <
                                KR21Monoculture.expectedBestInSet
                                  ((KR21Monoculture.sourceUnitVarianceLaplaceScaledNoiseFamily x1 x2 x3).dist thetaA)
                                  (KR21Monoculture.threeCandidateValueProfile x1 x2 x3) Finset.univ) ∧
            (∀ (D : MeasureTheory.Measure (KR21Monoculture.ValueProfile 1)) [MeasureTheory.IsProbabilityMeasure D],
                (∀ (c : EconCSLib.SocialChoice.Ranking.Candidate 1),
                    MeasureTheory.Integrable (fun value => value c) D) →
                  (∀ᵐ (value : KR21Monoculture.ValueProfile 1) ∂D, value 1 < value 0 ∧ value 2 < value 1) →
                    (∀ (theta : ℝ),
                        0 < theta →
                          ∀
                            (regularity :
                              KR21Monoculture.gaussianThreeCandidateDistributionalFamily.OuterIndependentRerankingJointRegularity
                                D theta),
                            0 <
                                ∫ (x : KR21Monoculture.ValueProfile 1 × KR21Monoculture.RankingPair 1),
                                  if KR21Monoculture.disagreementEvent x.2 then 1
                                  else
                                    0 ∂KR21Monoculture.gaussianThreeCandidateDistributionalFamily.outerIndependentPairJointLaw
                                    D theta ⋯ ∧
                              0 <
                                KR21Monoculture.gaussianThreeCandidateDistributionalFamily.jointLawDisagreementConditionalGain
                                  D theta ⋯) ∧
                      ∀ (thetaA thetaH : ℝ),
                        0 < thetaH →
                          thetaH < thetaA →
                            ∫ (value : KR21Monoculture.ValueProfile 1),
                                EconCSLib.SocialChoice.Ranking.expectedSecondMoverIndependent
                                  (KR21Monoculture.gaussianThreeCandidateDistributionalFamily.dist thetaH value)
                                  (KR21Monoculture.gaussianThreeCandidateDistributionalFamily.dist thetaA value)
                                  value ∂D <
                              ∫ (value : KR21Monoculture.ValueProfile 1),
                                EconCSLib.SocialChoice.Ranking.expectedSecondMoverIndependent
                                  (KR21Monoculture.gaussianThreeCandidateDistributionalFamily.dist thetaH value)
                                  (KR21Monoculture.gaussianThreeCandidateDistributionalFamily.dist thetaH value)
                                  value ∂D) ∧
              (∀ (D : MeasureTheory.Measure (KR21Monoculture.ValueProfile 1)) [MeasureTheory.IsProbabilityMeasure D],
                  (∀ (c : EconCSLib.SocialChoice.Ranking.Candidate 1),
                      MeasureTheory.Integrable (fun value => value c) D) →
                    (∀ᵐ (value : KR21Monoculture.ValueProfile 1) ∂D, value 1 < value 0 ∧ value 2 < value 1) →
                      (∀ (theta : ℝ),
                          0 < theta →
                            ∀
                              (regularity :
                                KR21Monoculture.sourceUnitVarianceLaplaceThreeCandidateDistributionalFamily.OuterIndependentRerankingJointRegularity
                                  D theta),
                              0 <
                                  ∫ (x : KR21Monoculture.ValueProfile 1 × KR21Monoculture.RankingPair 1),
                                    if KR21Monoculture.disagreementEvent x.2 then 1
                                    else
                                      0 ∂KR21Monoculture.sourceUnitVarianceLaplaceThreeCandidateDistributionalFamily.outerIndependentPairJointLaw
                                      D theta ⋯ ∧
                                0 <
                                  KR21Monoculture.sourceUnitVarianceLaplaceThreeCandidateDistributionalFamily.jointLawDisagreementConditionalGain
                                    D theta ⋯) ∧
                        ∀ (thetaA thetaH : ℝ),
                          0 < thetaH →
                            thetaH < thetaA →
                              ∫ (value : KR21Monoculture.ValueProfile 1),
                                  EconCSLib.SocialChoice.Ranking.expectedSecondMoverIndependent
                                    (KR21Monoculture.sourceUnitVarianceLaplaceThreeCandidateDistributionalFamily.dist
                                      thetaH value)
                                    (KR21Monoculture.sourceUnitVarianceLaplaceThreeCandidateDistributionalFamily.dist
                                      thetaA value)
                                    value ∂D <
                                ∫ (value : KR21Monoculture.ValueProfile 1),
                                  EconCSLib.SocialChoice.Ranking.expectedSecondMoverIndependent
                                    (KR21Monoculture.sourceUnitVarianceLaplaceThreeCandidateDistributionalFamily.dist
                                      thetaH value)
                                    (KR21Monoculture.sourceUnitVarianceLaplaceThreeCandidateDistributionalFamily.dist
                                      thetaH value)
                                    value ∂D) ∧
                (∀ (D : MeasureTheory.Measure (KR21Monoculture.ValueProfile 1)) [MeasureTheory.IsProbabilityMeasure D]
                    (thetaH : ℝ),
                    0 < thetaH →
                      (∀ (c : EconCSLib.SocialChoice.Ranking.Candidate 1),
                          MeasureTheory.Integrable (fun value => value c) D) →
                        (∀ᵐ (value : KR21Monoculture.ValueProfile 1) ∂D, value 1 < value 0 ∧ value 2 < value 1) →
                          KR21Monoculture.PaperInterface.literal_outer_payoff_paradox
                            KR21Monoculture.gaussianThreeCandidateDistributionalFamily D thetaH) ∧
                  ∀ (D : MeasureTheory.Measure (KR21Monoculture.ValueProfile 1)) [MeasureTheory.IsProbabilityMeasure D]
                    (thetaH : ℝ),
                    0 < thetaH →
                      (∀ (c : EconCSLib.SocialChoice.Ranking.Candidate 1),
                          MeasureTheory.Integrable (fun value => value c) D) →
                        (∀ᵐ (value : KR21Monoculture.ValueProfile 1) ∂D, value 1 < value 0 ∧ value 2 < value 1) →
                          KR21Monoculture.PaperInterface.literal_outer_payoff_paradox
                            KR21Monoculture.sourceUnitVarianceLaplaceThreeCandidateDistributionalFamily D thetaH
\end{ReviewVerbatim}

\paragraph{Lean proof endpoint (built separately)}\begin{ReviewVerbatim}
KR21Monoculture.PaperInterface.theorem2_gaussian_laplace_source_semantic_complete_literal_payoff
\end{ReviewVerbatim}

\paragraph{Recorded source-to-Spec screening}
\textbf{Verdict:} \texttt{not recorded}
\\
\textbf{Reason:} not recorded
\reviewmatch{Matches source input}
\noindent\textbf{Reviewer annotation}\par
\reviewerbox
\clearpage
\hypertarget{source-claim-8940f397fa4678ac}{}
\section*{Review row 16}
\textbf{Source locator:} {\footnotesize\raggedright cited publication, lines 530-\allowbreak{}540\par}
\paragraph{Verbatim source input}
\begin{ReviewVerbatim}
Finally, there exist noise distributions that violate Definition 2 for any candidate distribution D. In
particular, the RUM family defined by the Gumbel distribution is well-known to be equivalent to the PlackettLuce model of ranking, which is generated by sequentially selecting candidate i with probability
exp(θxi )
,
j∈S exp(θxj )

P

(7)

where S is the set of remaining candidates [12, 4]. Under the Plackett-Luce model, for any θ, UAH (θ, θ) =
\end{ReviewVerbatim}


\paragraph{Expanded PaperInterface specification}
\begin{ReviewVerbatim}
∀ {n : ℕ} (theta : ℝ) (value : EconCSLib.SocialChoice.Ranking.Candidate n → ℝ),
  0 < theta →
    ∀ (remaining : Finset (EconCSLib.SocialChoice.Ranking.Candidate n))
      {i : EconCSLib.SocialChoice.Ranking.Candidate n},
      i ∈ remaining →
        KR21Monoculture.plackettLuceChoiceProb theta value remaining i =
          Real.exp (theta * value i) / ∑ j ∈ remaining, Real.exp (theta * value j)
\end{ReviewVerbatim}

\paragraph{Lean proof endpoint (built separately)}\begin{ReviewVerbatim}
KR21Monoculture.PaperInterface.equation7_plackettLuce_choice_probability
\end{ReviewVerbatim}

\paragraph{Recorded source-to-Spec screening}
\textbf{Verdict:} \texttt{not recorded}
\\
\textbf{Reason:} not recorded
\reviewmatch{Matches source input}
\noindent\textbf{Reviewer annotation}\par
\reviewerbox
\clearpage
\hypertarget{source-claim-0156e1eedd228cc4}{}
\section*{Review row 17}
\textbf{Source locator:} {\footnotesize\raggedright cited publication, lines 530-\allowbreak{}548\par}
\paragraph{Verbatim source input}
\begin{ReviewVerbatim}
Finally, there exist noise distributions that violate Definition 2 for any candidate distribution D. In
particular, the RUM family defined by the Gumbel distribution is well-known to be equivalent to the PlackettLuce model of ranking, which is generated by sequentially selecting candidate i with probability
exp(θxi )
,
j∈S exp(θxj )

P

(7)

where S is the set of remaining candidates [12, 4]. Under the Plackett-Luce model, for any θ, UAH (θ, θ) =
UAA (θ, θ). To see this, suppose the firm that hires first selects candidate i∗ . Then, the firm that hires second
8


[form-feed in source extraction]
gets each candidate i with probability given by (7) with S = {1, . . . , n}\i∗ . As a result, by (3), if π, σ ∼ Fθ ,
E [π1 − π2 | π1 6= σ1 ] = 0
for any candidate distribution D, meaning the Plackett-Luce model never meets Definition 2. Thus, under
the Plackett-Luce model, monoculture has no effect—the optimal strategy is always to use the best available
ranking, regardless of competitors’ strategies.
\end{ReviewVerbatim}


\paragraph{Expanded PaperInterface specification}
\begin{ReviewVerbatim}
∀ {n : ℕ} (D : MeasureTheory.Measure (KR21Monoculture.ValueProfile n)) [MeasureTheory.IsProbabilityMeasure D]
  {location theta : ℝ},
  0 < theta →
    (∀ (c : EconCSLib.SocialChoice.Ranking.Candidate n), MeasureTheory.Integrable (fun value => value c) D) →
      (KR21Monoculture.sourceUnitVarianceGumbelNoiseLaw location =
          MeasureTheory.Measure.pi fun x =>
            MeasureTheory.Measure.map (fun arrival => location + -(√6 / Real.pi) * Real.log arrival)
              (ProbabilityTheory.expMeasure 1)) ∧
        (∀ (profile epsilon : EconCSLib.SocialChoice.Ranking.Candidate n → ℝ)
            (i : EconCSLib.SocialChoice.Ranking.Candidate n),
            KR21Monoculture.sourceUnitVarianceGumbelScores theta profile epsilon i = profile i + epsilon i / theta) ∧
          (∀ (profile : KR21Monoculture.ValueProfile n),
              (KR21Monoculture.sourceUnitVarianceGumbelRUMRankingPMF location theta profile).toMeasure =
                MeasureTheory.Measure.map
                  (fun epsilon => EconCSLib.SocialChoice.Ranking.rankByScore fun i => profile i + epsilon i / theta)
                  (MeasureTheory.Measure.pi fun x =>
                    MeasureTheory.Measure.map (fun arrival => location + -(√6 / Real.pi) * Real.log arrival)
                      (ProbabilityTheory.expMeasure 1))) ∧
            (∀ (profile : KR21Monoculture.ValueProfile n),
                KR21Monoculture.sourceUnitVarianceGumbelRUMRankingPMF location theta profile =
                  KR21Monoculture.plackettLuceRankingPMF (theta / KR21Monoculture.unitVarianceGumbelScale) profile) ∧
              have F :=
                KR21Monoculture.DistributionalAccuracyFamily.sourceUnitVarianceGumbelRUMDistributionalFamily location;
              ∃ (hatom :
                ∀ (ranking : EconCSLib.SocialChoice.Ranking.Ranking n),
                  Measurable fun value => (F.dist theta value) ranking),
                have J := F.outerIndependentPairJointLaw D theta hatom;
                have numerator :=
                  ∫ (x : KR21Monoculture.ValueProfile n × KR21Monoculture.RankingPair n),
                    if
                        EconCSLib.SocialChoice.Ranking.firstChoice x.2.1 ≠
                          EconCSLib.SocialChoice.Ranking.firstChoice x.2.2 then
                      x.1 (EconCSLib.SocialChoice.Ranking.firstChoice x.2.1) -
                        x.1 (EconCSLib.SocialChoice.Ranking.secondChoice x.2.1)
                    else 0 ∂J;
                have denominator :=
                  ∫ (x : KR21Monoculture.ValueProfile n × KR21Monoculture.RankingPair n),
                    if
                        EconCSLib.SocialChoice.Ranking.firstChoice x.2.1 ≠
                          EconCSLib.SocialChoice.Ranking.firstChoice x.2.2 then
                      1
                    else 0 ∂J;
                J = D.compProd (F.independentPairKernel theta hatom) ∧
                  (∀ (profile : KR21Monoculture.ValueProfile n),
                      (F.independentPairKernel theta hatom) profile =
                        (EconCSLib.pmfProd (F.dist theta profile) (F.dist theta profile)).toMeasure) ∧
                    MeasureTheory.Integrable
                        (fun value =>
                          EconCSLib.SocialChoice.Ranking.expectedSecondMoverIndependent (F.dist theta value)
                            (F.dist theta value) value)
                        D ∧
                      MeasureTheory.Integrable
                          (fun value =>
                            EconCSLib.SocialChoice.Ranking.expectedSecondMoverShared (F.dist theta value) value)
                          D ∧
                        0 < denominator ∧
                          numerator / denominator = 0 ∧
                            ∫ (value : KR21Monoculture.ValueProfile n),
                                EconCSLib.SocialChoice.Ranking.expectedSecondMoverIndependent (F.dist theta value)
                                  (F.dist theta value) value ∂D =
                              ∫ (value : KR21Monoculture.ValueProfile n),
                                EconCSLib.SocialChoice.Ranking.expectedSecondMoverShared (F.dist theta value) value ∂D
\end{ReviewVerbatim}

\paragraph{Lean proof endpoint (built separately)}\begin{ReviewVerbatim}
KR21Monoculture.PaperInterface.section31_literalUnitVarianceGumbel_outer_zero_effect_semantic_complete
\end{ReviewVerbatim}

\paragraph{Recorded source-to-Spec screening}
\textbf{Verdict:} \texttt{not recorded}
\\
\textbf{Reason:} not recorded
\reviewmatch{Matches source input}
\noindent\textbf{Reviewer annotation}\par
\reviewerbox
\clearpage
\hypertarget{source-claim-8dd2c310e7ffc8a6}{}
\section*{Review row 18}
\textbf{Source locator:} {\footnotesize\raggedright cited publication, lines 574-\allowbreak{}587\par}
\paragraph{Verbatim source input}
\begin{ReviewVerbatim}
The Mallows Model also appears frequently in the ranking literature [8, 11], and is much more analytically
tractable than RUMs. Under the Mallows Model, the likelihood of a permutation is related to its distance
from the true ranking π ∗ :
∗
1
(8)
Pr[π] = φ−d(π,π ) ,
Z
where Z is a normalizing constant. In this model, φ > 1 is the accuracy parameter: the larger φ is, the
more likely the ranking procedure is to output a ranking π that is close to the true ranking r. To instantiate
this model, we need a notion of distance d(·, ·) over permutations. For this, we’ll use Kendall tau distance,
another standard notion in the literature, which is simply the number of pairs of elements in π that are
incorrectly ordered [10]. In Appendix D, we verify that the family of distributions Fθ given by the Mallows
Model satisfies Definition 1, defining θ = φ − 1 (for consistency, so θ is well-defined on (0, ∞)).
\end{ReviewVerbatim}


\paragraph{Expanded PaperInterface specification}
\begin{ReviewVerbatim}
∀ {n : ℕ} (center : EconCSLib.SocialChoice.Ranking.Ranking n) (phi theta : ℝ),
  1 < phi →
    theta = phi - 1 →
      ∀ (pi : EconCSLib.SocialChoice.Ranking.Ranking n),
        have M := KR21Monoculture.concreteMallowsSpec center theta;
        M.q = phi⁻¹ ∧
          (M.law pi).toReal =
            phi⁻¹ ^ EconCSLib.SocialChoice.Ranking.kendallTau center pi /
              ∑ tau, phi⁻¹ ^ EconCSLib.SocialChoice.Ranking.kendallTau center tau
\end{ReviewVerbatim}

\paragraph{Lean proof endpoint (built separately)}\begin{ReviewVerbatim}
KR21Monoculture.PaperInterface.equation8_source_concrete_mallows_probability
\end{ReviewVerbatim}

\paragraph{Recorded source-to-Spec screening}
\textbf{Verdict:} \texttt{not recorded}
\\
\textbf{Reason:} not recorded
\reviewmatch{Matches source input}
\noindent\textbf{Reviewer annotation}\par
\reviewerbox
\clearpage
\hypertarget{source-claim-fd702590ac30f669}{}
\section*{Review row 19}
\textbf{Source locator:} {\footnotesize\raggedright cited publication, lines 133-\allowbreak{}138\par}
\paragraph{Verbatim source input}
\begin{ReviewVerbatim}
More formally, we model the n candidates as having intrinsic values x1 , . . . , xn , where any employer would
derive utility xi from hiring candidate i. Throughout the paper, we assume without loss of generality that
x1 > x2 > · · · > xn . These values, however, are unknown to the employer; instead, they must use some
noisy procedure to rank the candidates. We model such a procedure as a randomized mechanism R that
takes in the true candidate values and draws a permutation π over those candidates from some distribution.
Our main results hold for families of distributions over permutations as defined below:

[Next verbatim source excerpt]

Each firm in our model has access to the same underlying pool of n candidates, which they rank using
a randomized mechanism R to get a permutation π as described above. Then, in a random order, each
firm hires the highest-ranked remaining candidate according to their ranking. Thus, if two firms both rank
candidate i first, only one of them can hire i; the other must hire the next available candidate according to
their ranking. In our model, candidates automatically accept the offer they get from a firm. For the sake
of simplicity, throughout this paper, we restrict ourselves to the case where there are two firms hiring one
candidate each, although our model readily generalizes to more complex cases.

[Next verbatim source excerpt]

The Mallows Model also appears frequently in the ranking literature [8, 11], and is much more analytically
tractable than RUMs. Under the Mallows Model, the likelihood of a permutation is related to its distance
from the true ranking π ∗ :
∗
1
(8)
Pr[π] = φ−d(π,π ) ,
Z
where Z is a normalizing constant. In this model, φ > 1 is the accuracy parameter: the larger φ is, the
more likely the ranking procedure is to output a ranking π that is close to the true ranking r. To instantiate
this model, we need a notion of distance d(·, ·) over permutations. For this, we’ll use Kendall tau distance,
another standard notion in the literature, which is simply the number of pairs of elements in π that are
incorrectly ordered [10]. In Appendix D, we verify that the family of distributions Fθ given by the Mallows
Model satisfies Definition 1, defining θ = φ − 1 (for consistency, so θ is well-defined on (0, ∞)).

[Next verbatim source excerpt]

Theorem 3. Let Fθ be the family of Mallows Model distributions with parameter θ = φ − 1. Then, for any
candidate distribution D, the conditions of Theorem 1 are satisfied.

[Next verbatim source excerpt]

E

Proof of Theorem 3

E.1

Verifying Definition 2

We must show that when π, τ ∼ Fθ ,
E [π1 − π2 | π1 6= τ1 ] > 0.

(E.1)

We begin by expanding:
E [π1 − π2 | π1 6= τ1 ] =

n
n X
X

(xi − xj ) Pr[π1 = xi ∩ π2 = xj | π1 6= τ1 ]

i=1 j=1

=

n−1
XX

(xi − xj ) (Pr[π1 = xi ∩ π2 = xj | π1 6= τ1 ] − Pr[π1 = xj ∩ π2 = xi | π1 6= τ1 ])

i=1 j>i

Since xi > xj for i < j, it suffices to show that for all i < j,
Pr[π1 = xi ∩ π2 = xj | π1 6= τ1 ] ≥ Pr[π1 = xj ∩ π2 = xi | π1 6= τ1 ],

(E.2)

and that this holds strictly for some i < j. We simplify (E.2) as follows:
Pr[π1 = xi ∩ π2 = xj | π1 6= τ1 ] > Pr[π1 = xj ∩ π2 = xi | π1 6= τ1 ]
Pr[π1 = xi ∩ π2 = xj ∩ π1 6= τ1 ]
Pr[π1 = xj ∩ π2 = xi ∩ π1 6= τ1 ]
>
Pr[π1 6= τ1 ]
Pr[π1 6= τ1 ]
⇐⇒ Pr[π1 = xi ∩ π2 = xj ∩ π1 6= τ1 ] > Pr[π1 = xj ∩ π2 = xi ∩ π1 6= τ1 ]
⇐⇒

⇐⇒ Pr[π1 = xi ∩ π2 = xj ∩ τ1 6= xi ] > Pr[π1 = xj ∩ π2 = xi ∩ τ1 6= xj ]
⇐⇒ Pr[π1 = xi ∩ π2 = xj ] Pr[τ1 6= xi ] > Pr[π1 = xj ∩ π2 = xi ] Pr[τ1 6= xj ]

(E.3)

We can simplify (E.3) using Lemmas 5 and 6. Let |i − j| denote the difference in rank between xi and xj .
Pr[π1 = xi ∩ π2 = xj ] Pr[τ1 6= xi ] − Pr[π1 = xj ∩ π2 = xi ] Pr[τ1 6= xj ]
= Pr[π1 = xi ∩ π2 = xj ](1 − Pr[τ1 = xi ]) − φ−1 Pr[π1 = xi ∩ π2 = xj ](1 − Pr[τ1 = xj ])
= Pr[π1 = xi ∩ π2 = xj ](1 − Pr[τ1 = xi ]) − φ−1 Pr[π1 = xi ∩ π2 = xj ](1 − φ−|i−j| Pr[τ1 = xi ])
= Pr[π1 = xi ∩ π2 = xj ](1 − Pr[τ1 = xi ] − φ−1 − φ−|i−j|−1 Pr[τ1 = xi ]))
29


[form-feed in source extraction]
This is positive if and only if
1 − Pr[τ1 = xi ] − φ−1 − φ−|i−j|−1 Pr[τ1 = xi ] > 0
[U+0011 control character]
[U+0010 control character]
⇐⇒ Pr[τ1 = xi ] 1 − φ−|i−j|−1 < 1 − φ−1
1 − φ−1
1 − φ−|i−j|−1
−1
1−φ
1 − φ−1
⇐⇒ i−1
<
−n
φ (1 − φ )
1 − φ−|i−j|−1
⇐⇒ Pr[τ1 = xi ] <

⇐⇒ φi−1 (1 − φ−n ) > 1 − φ−|i−j|−1
This is weakly true for any i < j because φi−1 ≥ 1 and |i − j| + 1 ≤ n, and it is strictly true for any i, j
other than 1 and n. Thus, E [π1 − π2 | π1 6= τ1 ] > 0.

E.2

Verifying Definition 3

Recall that Definition 3 is equivalent to UAH (θA , θH ) < UHH (θA , θH ) for θA > θH . Let τ be the algorithmic
ranking, and let π be a ranking from a human evaluator. Recall that UH (θA , θH ) = E [π1 ]. Throughout this
proof, we will drop the (θA , θH ) notation and simply write UH , UAH , and UHH .
UAH =

n
X
(Pr[π1 = xi ∩ τ1 6= xi ] + Pr[π2 = xi ∩ π1 = τ1 ])xi
i=1

=
=

n
X
i=1
n
X

Pr[π1 = xi ∩ τ1 6= xi ]xi +

n
X

Pr[π2 = xi ∩ π1 = τ1 ]xi

i=1

(P r[π1 = xi ] − Pr[π1 = xi ∩ τ1 = xi ]) xi +

i=1

= UH −
= UH +

n X
X

Pr[π1 = xj ∩ τ1 = xj ∩ π2 = xi ]xi

i=1 j6=i
n
X
i=1
n
X

Pr[π1 = xi ∩ τ1 = xi ]xi +

n
X

Pr[π1 = xi ∩ τ1 = xi ]E [π2 | π1 = xi ∩ τ1 = xi ]

i=1

Pr[π1 = xi ] Pr[τ1 = xi ] (E [π2 | π1 = xi ] − xi )

i=1

Similarly, because two human evaluators are independent,
UHH = UH +

n
X

Pr[π1 = xi ]2 (E [π2 | π1 = xi ] − xi ) .

i=1

Let V−i = E [π2 | π1 = xi ]. Note that conditioned on π1 = xi , the remaining elements of π1 follow a Mallows
model distribution over n − 1 candidates. Because the Mallows model is value-independent, increasing any
item value increases the expected value of the top-ranked item (and in fact, the item ranked at any position).
Thus, V−i increases as i increases (since xi , the value of the unavailable candidate, decreases). Moreover, xi
is strictly decreasing in i, so V−i − xi is strictly increasing in i. With this, we have
UAH − UH =
UHH − UH =

n
X
i=1
n
X

Pr[π1 = xi ] Pr[τ1 = xi ] (V−i − xi )
Pr[π1 = xi ]2 (V−i − xi )

i=1

Pn

Pn
Let CA = Pr[π1 = τ1 ] = i=1 Pr[π1 = xi ] Pr[τ1 = xi ], and similarly let CH = i=1 Pr[π1 = xi ]2 . CA > CH
by Lemma 4 with yi0 = Pr[π1 = xn−i+1 ], p0i = Pr[π1 = xn−i+1 ] and qi0 = Pr[τ1 = xn−i+1 ].

30


[form-feed in source extraction]
Let pi = Pr[π10 = i] Pr[π1 = i]/CA , qi = Pr[π10 = i]2 /CH , and yi = V−i − xi . Then, we have
n
X
UAH − UH
=
pi yi
CA
i=1
n
X
UHH − UH
=
qi yi
CH
i=1

With φA = θA + 1 and φH = θH + 1,
CH
pi
=
·
qi
CA

1−φ−1
A
i−1
φA (1−φ−n
A )
1−φ−1
H
−n
φi−1
H (1−φH )

Pn

which is decreasing in i since φH < φA . By Lemma 4,
0 by Lemma 7, so
n
X

pi yi <

n
X

i=1

∝

φi−1
H
,
φi−1
A

i=1 pi yi <

Pn

i=1 qi yi . Finally, note that UHH −UH <

qi yi

i=1

UHH − UH
UAH − UH
<
CA
CH
CH (UAH − UH )
< UHH − UH
CA
UAH − UH < UHH − UH

(CA > CH , and UHH − UH < 0)

UAH < UHH
\end{ReviewVerbatim}

\paragraph{Approved corrected review target}
The archival source above is not asserted equivalent to this replacement.
\begin{ReviewVerbatim}
For Candidate n = Fin(n+2) with n > 0, a fixed rank-labelled Mallows center, and a probability law D on value profiles with coordinatewise finite first moments and D-a.e. strict center order: for every phi > 1 and theta = phi - 1, the fixed-center source Mallows law has q = phi^(-1) and, for every value profile v and ranking pi, mass phi^(-KendallTau(center, pi)) / sum_tau phi^(-KendallTau(center, tau)). For every theta > 0 the actual outer-law conditionally-iid-ranking top-disagreement gain is positive, and for every thetaA > thetaH > 0 the Definition 3 weaker-competition conclusion holds.
\end{ReviewVerbatim}

\textbf{Archival source anchor:} {\footnotesize\raggedright cited publication, lines 590-\allowbreak{}591\par}
\paragraph{Recorded basis}
User instruction 2026-\allowbreak{}07-\allowbreak{}25:\allowbreak{} assume whatever convention makes the result true, provided the convention is explicit and the archival claim is not represented as proved.\allowbreak{}
\paragraph{Expanded PaperInterface specification}
\begin{ReviewVerbatim}
∀ {n : ℕ} (D : MeasureTheory.Measure (KR21Monoculture.ValueProfile n)) [MeasureTheory.IsProbabilityMeasure D]
  (center : EconCSLib.SocialChoice.Ranking.Ranking n),
  0 < n →
    (∀ (c : EconCSLib.SocialChoice.Ranking.Candidate n), MeasureTheory.Integrable (fun value => value c) D) →
      (∀ᵐ (value : EconCSLib.SocialChoice.Ranking.Candidate n → ℝ) ∂D,
          EconCSLib.SocialChoice.Ranking.StrictlyOrderedBy center value) →
        (∀ (phi theta : ℝ),
            1 < phi →
              theta = phi - 1 →
                (KR21Monoculture.concreteMallowsSpec center theta).q = phi⁻¹ ∧
                  ∀ (value : KR21Monoculture.ValueProfile n) (pi : EconCSLib.SocialChoice.Ranking.Ranking n),
                    (((KR21Monoculture.fixedCenterMallowsDistributionalFamily center).dist theta value) pi).toReal =
                      phi⁻¹ ^ EconCSLib.SocialChoice.Ranking.kendallTau center pi /
                        ∑ tau, phi⁻¹ ^ EconCSLib.SocialChoice.Ranking.kendallTau center tau) ∧
          (∀ (value : KR21Monoculture.ValueProfile n),
              EconCSLib.SocialChoice.Ranking.StrictlyOrderedBy center value →
                (∀ (theta : ℝ),
                    0 < theta →
                      ∀ (pi : EconCSLib.SocialChoice.Ranking.Ranking n),
                        ContinuousAt
                            (fun theta' =>
                              (((KR21Monoculture.fixedCenterMallowsPointFamily center value).dist theta') pi).toReal)
                            theta ∧
                          DifferentiableAt ℝ
                            (fun theta' =>
                              (((KR21Monoculture.fixedCenterMallowsPointFamily center value).dist theta') pi).toReal)
                            theta) ∧
                  Filter.Tendsto
                      (fun theta =>
                        (((KR21Monoculture.fixedCenterMallowsPointFamily center value).dist theta) center).toReal)
                      Filter.atTop (nhds 1) ∧
                    ∀ (thetaA thetaH : ℝ),
                      0 < thetaH →
                        thetaH < thetaA →
                          (∀ (remaining : Finset (EconCSLib.SocialChoice.Ranking.Candidate n)),
                              remaining.Nonempty →
                                KR21Monoculture.expectedBestInSet
                                    ((KR21Monoculture.fixedCenterMallowsPointFamily center value).dist thetaH) value
                                    remaining ≤
                                  KR21Monoculture.expectedBestInSet
                                    ((KR21Monoculture.fixedCenterMallowsPointFamily center value).dist thetaA) value
                                    remaining) ∧
                            KR21Monoculture.expectedBestInSet
                                ((KR21Monoculture.fixedCenterMallowsPointFamily center value).dist thetaH) value
                                Finset.univ <
                              KR21Monoculture.expectedBestInSet
                                ((KR21Monoculture.fixedCenterMallowsPointFamily center value).dist thetaA) value
                                Finset.univ) ∧
            (∀ (theta : ℝ),
                0 < theta →
                  have J :=
                    (KR21Monoculture.fixedCenterMallowsDistributionalFamily center).outerIndependentPairJointLaw D theta
                      ⋯;
                  0 <
                      ∫ (x : KR21Monoculture.ValueProfile n × KR21Monoculture.RankingPair n),
                        if
                            EconCSLib.SocialChoice.Ranking.firstChoice x.2.1 ≠
                              EconCSLib.SocialChoice.Ranking.firstChoice x.2.2 then
                          1
                        else 0 ∂J ∧
                    0 <
                      (∫ (x : KR21Monoculture.ValueProfile n × KR21Monoculture.RankingPair n),
                          if
                              EconCSLib.SocialChoice.Ranking.firstChoice x.2.1 ≠
                                EconCSLib.SocialChoice.Ranking.firstChoice x.2.2 then
                            x.1 (EconCSLib.SocialChoice.Ranking.firstChoice x.2.1) -
                              x.1 (EconCSLib.SocialChoice.Ranking.secondChoice x.2.1)
                          else 0 ∂J) /
                        ∫ (x : KR21Monoculture.ValueProfile n × KR21Monoculture.RankingPair n),
                          if
                              EconCSLib.SocialChoice.Ranking.firstChoice x.2.1 ≠
                                EconCSLib.SocialChoice.Ranking.firstChoice x.2.2 then
                            1
                          else 0 ∂J) ∧
              (∀ (thetaA thetaH : ℝ),
                  0 < thetaH →
                    thetaH < thetaA →
                      ∫ (value : KR21Monoculture.ValueProfile n),
                          EconCSLib.pmfPairExp
                            ((KR21Monoculture.fixedCenterMallowsDistributionalFamily center).dist thetaH value)
                            ((KR21Monoculture.fixedCenterMallowsDistributionalFamily center).dist thetaA value)
                            fun pi sigma => EconCSLib.SocialChoice.Ranking.secondMoverUtility value pi sigma ∂D <
                        ∫ (value : KR21Monoculture.ValueProfile n),
                          EconCSLib.pmfPairExp
                            ((KR21Monoculture.fixedCenterMallowsDistributionalFamily center).dist thetaH value)
                            ((KR21Monoculture.fixedCenterMallowsDistributionalFamily center).dist thetaH value)
                            fun pi sigma => EconCSLib.SocialChoice.Ranking.secondMoverUtility value pi sigma ∂D) ∧
                ∀ (thetaH : ℝ),
                  0 < thetaH →
                    have F := KR21Monoculture.fixedCenterMallowsDistributionalFamily center;
                    ∃ thetaA,
                      thetaH < thetaA ∧
                        ∫ (value : KR21Monoculture.ValueProfile n),
                                ((F.pointFamily value).modelAt thetaA thetaH).firstMoverEU
                                  KR21Monoculture.Strategy.human ∂D +
                              ∫ (value : KR21Monoculture.ValueProfile n),
                                ((F.pointFamily value).modelAt thetaA thetaH).secondMoverEU
                                  KR21Monoculture.Strategy.algorithm KR21Monoculture.Strategy.human ∂D <
                            ∫ (value : KR21Monoculture.ValueProfile n),
                                ((F.pointFamily value).modelAt thetaA thetaH).firstMoverEU
                                  KR21Monoculture.Strategy.algorithm ∂D +
                              ∫ (value : KR21Monoculture.ValueProfile n),
                                ((F.pointFamily value).modelAt thetaA thetaH).secondMoverEU
                                  KR21Monoculture.Strategy.algorithm KR21Monoculture.Strategy.algorithm ∂D ∧
                          ∫ (value : KR21Monoculture.ValueProfile n),
                                  ((F.pointFamily value).modelAt thetaA thetaH).firstMoverEU
                                    KR21Monoculture.Strategy.human ∂D +
                                ∫ (value : KR21Monoculture.ValueProfile n),
                                  ((F.pointFamily value).modelAt thetaA thetaH).secondMoverEU
                                    KR21Monoculture.Strategy.human KR21Monoculture.Strategy.human ∂D <
                              ∫ (value : KR21Monoculture.ValueProfile n),
                                  ((F.pointFamily value).modelAt thetaA thetaH).firstMoverEU
                                    KR21Monoculture.Strategy.algorithm ∂D +
                                ∫ (value : KR21Monoculture.ValueProfile n),
                                  ((F.pointFamily value).modelAt thetaA thetaH).secondMoverEU
                                    KR21Monoculture.Strategy.human KR21Monoculture.Strategy.algorithm ∂D ∧
                            ∫ (value : KR21Monoculture.ValueProfile n),
                                  ((F.pointFamily value).modelAt thetaA thetaH).firstMoverEU
                                    KR21Monoculture.Strategy.algorithm ∂D +
                                ∫ (value : KR21Monoculture.ValueProfile n),
                                  ((F.pointFamily value).modelAt thetaA thetaH).secondMoverEU
                                    KR21Monoculture.Strategy.algorithm KR21Monoculture.Strategy.algorithm ∂D <
                              ∫ (value : KR21Monoculture.ValueProfile n),
                                  ((F.pointFamily value).modelAt thetaA thetaH).firstMoverEU
                                    KR21Monoculture.Strategy.human ∂D +
                                ∫ (value : KR21Monoculture.ValueProfile n),
                                  ((F.pointFamily value).modelAt thetaA thetaH).secondMoverEU
                                    KR21Monoculture.Strategy.human KR21Monoculture.Strategy.human ∂D
\end{ReviewVerbatim}

\paragraph{Lean proof endpoint (built separately)}\begin{ReviewVerbatim}
KR21Monoculture.PaperInterface.theorem3_source_complete_semantic_complete_literal_payoff
\end{ReviewVerbatim}

\paragraph{Recorded source-to-Spec screening}
\textbf{Verdict:} \texttt{not recorded}
\\
\textbf{Reason:} not recorded
\reviewmatch{Matches approved corrected target}
\noindent\textbf{Reviewer annotation}\par
\reviewerbox
\clearpage
\hypertarget{source-claim-e42fc5562761542b}{}
\section*{Review row 20}
\textbf{Source locator:} {\footnotesize\raggedright cited publication, lines 133-\allowbreak{}138\par}
\paragraph{Verbatim source input}
\begin{ReviewVerbatim}
[No record available.]
\end{ReviewVerbatim}


\paragraph{Expanded PaperInterface specification}
\begin{ReviewVerbatim}
∀ {n k : ℕ} (D : MeasureTheory.Measure (KR21Monoculture.ValueProfile n)) [MeasureTheory.IsProbabilityMeasure D]
  (center : EconCSLib.SocialChoice.Ranking.Ranking n) (phiA thetaA phiH thetaH : ℝ),
  1 < phiA →
    1 < phiH →
      thetaA = phiA - 1 →
        thetaH = phiH - 1 →
          k < Fintype.card (EconCSLib.SocialChoice.Ranking.Candidate n) →
            (∀ (c : EconCSLib.SocialChoice.Ranking.Candidate n), MeasureTheory.Integrable (fun value => value c) D) →
              have algorithm := KR21Monoculture.concreteMallowsSpec center thetaA;
              have human := KR21Monoculture.concreteMallowsSpec center thetaH;
              algorithm.q = phiA⁻¹ ∧
                human.q = phiH⁻¹ ∧
                  (∀ (pi : EconCSLib.SocialChoice.Ranking.Ranking n),
                      (algorithm.law pi).toReal =
                        phiA⁻¹ ^ EconCSLib.SocialChoice.Ranking.kendallTau center pi /
                          ∑ tau, phiA⁻¹ ^ EconCSLib.SocialChoice.Ranking.kendallTau center tau) ∧
                    (∀ (pi : EconCSLib.SocialChoice.Ranking.Ranking n),
                        (human.law pi).toReal =
                          phiH⁻¹ ^ EconCSLib.SocialChoice.Ranking.kendallTau center pi /
                            ∑ tau, phiH⁻¹ ^ EconCSLib.SocialChoice.Ranking.kendallTau center tau) ∧
                      (phiA ≤ phiH →
                          (∀ᵐ (value : EconCSLib.SocialChoice.Ranking.Candidate n → ℝ) ∂D,
                              EconCSLib.SocialChoice.Ranking.WeaklyOrderedBy center value) →
                            ∀ (i : Fin k) (hired : Finset (EconCSLib.SocialChoice.Ranking.Candidate n)),
                              hired.card = ↑i →
                                ∀ (strategy : KR21Monoculture.Strategy),
                                  ∫ (value : KR21Monoculture.ValueProfile n),
                                      ∑ pi,
                                        ((match strategy with
                                              | KR21Monoculture.Strategy.algorithm => algorithm.law
                                              | KR21Monoculture.Strategy.human => human.law)
                                              pi).toReal *
                                          value (KR21Monoculture.bestInSet pi (Finset.univ \ hired)) ∂D ≤
                                    ∫ (value : KR21Monoculture.ValueProfile n),
                                      ∑ pi,
                                        (human.law pi).toReal *
                                          value (KR21Monoculture.bestInSet pi (Finset.univ \ hired)) ∂D) ∧
                        (phiA < phiH →
                          (∀ᵐ (value : EconCSLib.SocialChoice.Ranking.Candidate n → ℝ) ∂D,
                              EconCSLib.SocialChoice.Ranking.StrictlyOrderedBy center value) →
                            ∀ (i : Fin k) (hired : Finset (EconCSLib.SocialChoice.Ranking.Candidate n)),
                              hired.card = ↑i →
                                ∀ (strategy : KR21Monoculture.Strategy),
                                  (∀ (alternative : KR21Monoculture.Strategy),
                                      ∫ (value : KR21Monoculture.ValueProfile n),
                                          ∑ pi,
                                            ((match alternative with
                                                  | KR21Monoculture.Strategy.algorithm => algorithm.law
                                                  | KR21Monoculture.Strategy.human => human.law)
                                                  pi).toReal *
                                              value (KR21Monoculture.bestInSet pi (Finset.univ \ hired)) ∂D ≤
                                        ∫ (value : KR21Monoculture.ValueProfile n),
                                          ∑ pi,
                                            ((match strategy with
                                                  | KR21Monoculture.Strategy.algorithm => algorithm.law
                                                  | KR21Monoculture.Strategy.human => human.law)
                                                  pi).toReal *
                                              value (KR21Monoculture.bestInSet pi (Finset.univ \ hired)) ∂D) →
                                    strategy = KR21Monoculture.Strategy.human)
\end{ReviewVerbatim}

\paragraph{Lean proof endpoint (built separately)}\begin{ReviewVerbatim}
KR21Monoculture.PaperInterface.theorem4_source_mallows_outer_horizon_lt_literal_semantic_complete
\end{ReviewVerbatim}

\paragraph{Recorded source-to-Spec screening}
\textbf{Verdict:} \texttt{not recorded}
\\
\textbf{Reason:} not recorded
\reviewmatch{Matches source input}
\noindent\textbf{Reviewer annotation}\par
\reviewerbox
\clearpage
\hypertarget{source-claim-54f72d882fb14d8c}{}
\section*{Review row 21}
\textbf{Source locator:} {\footnotesize\raggedright cited publication, lines 782-\allowbreak{}951\par}
\paragraph{Verbatim source input}
\begin{ReviewVerbatim}
Random Utility Models satisfying Definition 1

Theorem 5. Let f be the pdf of E. The family of RUMs Fθ given by ranking xi + εθi with εi ∼ E satisfies
the conditions of Definition 1 if:
• f is differentiable
• f has positive support on (−∞, ∞)
Proof. We need to show that Fθ satisfies the differentiability, asymptotic optimality, and monotonicity
conditions in Definition 1.
Qn
Differentiability: The probability density of any realization of the n noise samples εi /θ is i=1 f (εi /θ).
Let ε = [ε1 /θ, . . . , εn /θ] be the vector of noise values and let M (π) ⊆ Rn be the region such that any ε ∈ M (π)
will produces the ranking π. The probability of any permutation π is
n
[U+0010 control character]ε [U+0011 control character]
Y
i
dn z.
Pr[π] =
f
θ
θ
M (π) i=1

Z

Because f is differentiable,
[U+0010 control character]x[U+0011 control character] [U+0010 control character] x [U+0011 control character]
d [U+0010 control character]x[U+0011 control character]
f
= f0
· − 2
dθ
θ
θ
θ
Because Pr θ(π) is an integral of the product of differentiable functions over a fixed region, it is differentiable.
Asymptotic optimality: We will show that for any pair of elements and any δ > 0, there exists
sufficiently large θ such that the probability that they incorrectly ranked is at most δ. We will conclude with
a union bound over the n − 1 pairs of adjacent candidates that there exists sufficiently large θ such that the
probability of outputting the correct ranking must be at least 1 − (n − 1)δ.
Consider two candidates xi > xi+1 . Let ν be the difference xi − xi+1 . Then, they will be correctly ranked
if
εi
ν
>−
θ
2
ν
εi+1
<
θ
2
Let q and q be the 1 − 2δ and 2δ quantiles of E respectively, and let q = max(|q|, |q|). For θ > 2q
ν ,
[U+0014 control character]
[U+0015 control character]
νi
νθ
<−
= Pr εi < −
Pr
θ
2
2
hε

i

< Pr [εi < −q]
[U+0002 control character]
[U+0003 control character]
≤ Pr εi < q
δ
=
2 [U+0014 control character]
[U+0015 control character]
hε
νi
νθ
i+1
Pr
>
= Pr εi+1 >
θ
2
2
< Pr [εi+1 > q]
≤ Pr [εi+1 > q]
δ
=
2
Thus, for sufficiently large θ, the probability that xi and xi+1 are incorrectly ordered is at most δ.
Repeating this analysis for all n − 1 pairs of adjacent elements, taking the maximum of all the θ’s, and
taking a union bound yields that the probability of incorrectly ordering any pair of elements is at most
(n − 1)δ, meaning the probability of outputting the correct ranking is at least 1 − (n − 1)δ. Since δ is
arbitrary, this probability can be made arbitrarily close to 1, satisfying the asymptotic optimality condition.
14


[form-feed in source extraction]
Monotonicity: The removal of any elements does not alter the distribution of the remaining elements,
meaning that the distribution of π (−S) is equivalent to a RUM with n − |S| elements. Thus, it suffices to
show that for a RUM with positive support on (−∞, ∞), the probability of ranking the best candidate first
strictly increases with θ.
Recall that by definition, the candidates are ranked according to xi + εθi . The probability that x1 is
ranked first is
[U+0014 control character]
[U+0015 control character]
[U+0014 control character]
[U+0015 control character]
εi
εi
ε1
ε1
Pr x1 +
> max xi +
= Pr
> max xi − x1 +
2≤i≤n
2≤i≤n
θ
θ
θ
θ
[U+0014 control character]
[U+0015 control character]
= Pr ε1 > max θ(xi − x1 ) + εi
2≤i≤n
[U+0014 control character]
[U+0015 control character]
= Eε2 ,...,εn Pr ε1 > max θ(xi − x1 ) + εi | ε2 , . . . , εn
(A.1)
2≤i≤n

We want to show that (A.1) is increasing in θ. Intuitively, this is because as θ increases, the right hand
side of the inequality inside the probability decreases. To prove this formally, it suffices to show that the
subderivative of (A.1) with respect to θ only includes strictly positive numbers. First, we have
[U+0014 control character]
[U+0015 control character]
∂
Eε2 ,...,εn Pr ε1 > max θ(xi − x1 ) + εi | ε2 , . . . , εn ⊂ R>0
2≤i≤n
∂θ
[U+0014 control character]
[U+0015 control character]
∂
Pr ε1 > max θ(xi − x1 ) + εi | ε2 , . . . , εn ⊂ R>0
⇐⇒
2≤i≤n
∂θ
Let F and f be the cumulative density function and probability density function of E respectively. Then,
[U+0014 control character]
[U+0015 control character]
[U+0012 control character]
[U+0013 control character]
Pr ε1 > max θ(xi − x1 ) + εi | ε2 , . . . , εn = 1 − F max θ(xi − x1 ) + εi
2≤i≤n

2≤i≤n

Note that F (·) is strictly increasing (since f is assumed to have positive support on (−∞, ∞)), so it suffices
to show that
∂
max θ(xi − x1 ) + εi ⊂ R<0
∂θ 2≤i≤n
For any i,

d
θ(xi − x1 ) + εi = xi − x1 < 0.
dθ
Thus, the subderivative of the max of such functions includes only strictly negative numbers, which completes
the proof.
\end{ReviewVerbatim}

\paragraph{Approved corrected review target}
The archival source above is not asserted equivalent to this replacement.
\begin{ReviewVerbatim}
Let Candidate n = Fin(n+2). Let f : R -> R be measurable, integrable, everywhere strictly positive, normalized by integral ENNReal.ofReal(f) = 1, and absolutely continuous on every real interval; let f_prime be integrable and agree almost everywhere with deriv f. For a value profile strictly descending along a center ranking, the iid product density law is a probability law, and there is a ranking law F_theta such that for every theta > 0 it is exactly the pushforward of that product law under epsilon |-> rank(value_i + epsilon_i/theta). Every ranking atom of F_theta is continuous and differentiable at every positive theta, every atom tends to the pure center-ranking mass as theta tends to infinity, expected value of the best candidate in every nonempty remaining set is weakly increasing with theta, and full-set expected value is strictly increasing when theta_A > theta_H > 0.
\end{ReviewVerbatim}

\textbf{Archival source anchor:} {\footnotesize\raggedright cited publication, lines 784-\allowbreak{}787\par}
\paragraph{Recorded basis}
User instruction 2026-\allowbreak{}07-\allowbreak{}25:\allowbreak{} assume whatever convention makes the result true, provided the convention is explicit and the archival claim is not represented as proved.\allowbreak{}
\paragraph{Expanded PaperInterface specification}
\begin{ReviewVerbatim}
∀ {n : ℕ} (f derivative : ℝ → ℝ),
  MeasureTheory.Integrable f MeasureTheory.volume →
    MeasureTheory.Integrable derivative MeasureTheory.volume →
      Measurable f →
        (∀ (x : ℝ), 0 < f x) →
          (∀ (a b : ℝ), AbsolutelyContinuousOnInterval f a b) →
            derivative =ᵐ[MeasureTheory.volume] deriv f →
              ∫⁻ (x : ℝ), ENNReal.ofReal (f x) = 1 →
                ∀ (value : EconCSLib.SocialChoice.Ranking.Candidate n → ℝ)
                  (center : EconCSLib.SocialChoice.Ranking.Ranking n),
                  (∀ (i : Fin (n + 1)), value (center i.succ) < value (center i.castSucc)) →
                    Fintype.card (EconCSLib.SocialChoice.Ranking.Candidate n) = n + 2 ∧
                      MeasureTheory.IsProbabilityMeasure
                          (MeasureTheory.Measure.pi fun x =>
                            MeasureTheory.volume.withDensity fun x => ENNReal.ofReal (f x)) ∧
                        ∃ rankingLaw,
                          (∀ (theta : ℝ),
                              0 < theta →
                                (rankingLaw theta).toMeasure =
                                  MeasureTheory.Measure.map
                                    (fun epsilon =>
                                      EconCSLib.SocialChoice.Ranking.rankByScore fun i => value i + epsilon i / theta)
                                    (MeasureTheory.Measure.pi fun x =>
                                      MeasureTheory.volume.withDensity fun x => ENNReal.ofReal (f x))) ∧
                            (∀ (theta : ℝ),
                                0 < theta →
                                  ∀ (pi : EconCSLib.SocialChoice.Ranking.Ranking n),
                                    ContinuousAt (fun theta' => ((rankingLaw theta') pi).toReal) theta ∧
                                      DifferentiableAt ℝ (fun theta' => ((rankingLaw theta') pi).toReal) theta) ∧
                              (∀ (pi : EconCSLib.SocialChoice.Ranking.Ranking n),
                                  Filter.Tendsto (fun theta => ((rankingLaw theta) pi).toReal) Filter.atTop
                                    (nhds ((PMF.pure center) pi).toReal)) ∧
                                (∀ (thetaA thetaH : ℝ),
                                    0 < thetaH →
                                      thetaH < thetaA →
                                        ∀ (remaining : Finset (EconCSLib.SocialChoice.Ranking.Candidate n)),
                                          remaining.Nonempty →
                                            KR21Monoculture.expectedBestInSet (rankingLaw thetaH) value remaining ≤
                                              KR21Monoculture.expectedBestInSet (rankingLaw thetaA) value remaining) ∧
                                  ∀ (thetaA thetaH : ℝ),
                                    0 < thetaH →
                                      thetaH < thetaA →
                                        KR21Monoculture.expectedBestInSet (rankingLaw thetaH) value Finset.univ <
                                          KR21Monoculture.expectedBestInSet (rankingLaw thetaA) value Finset.univ
\end{ReviewVerbatim}

\paragraph{Lean proof endpoint (built separately)}\begin{ReviewVerbatim}
KR21Monoculture.PaperInterface.appendixA_theorem5_corrected_source_complete
\end{ReviewVerbatim}

\paragraph{Recorded source-to-Spec screening}
\textbf{Verdict:} \texttt{not recorded}
\\
\textbf{Reason:} not recorded
\reviewmatch{Matches approved corrected target}
\noindent\textbf{Reviewer annotation}\par
\reviewerbox
\clearpage
\hypertarget{source-claim-9835a680c6142d09}{}
\section*{Review row 22}
\textbf{Source locator:} {\footnotesize\raggedright cited publication, lines 493-\allowbreak{}504\par}
\paragraph{Verbatim source input}
\begin{ReviewVerbatim}
In Random Utility Models, the underlying candidate values xi are perturbed by independent and identically
distributed noise εi ∼ E, and the perturbed values are ranked to produce π. Originally conceived in the
psychology literature [23], this model has been well-studied over nearly a century, [7, 4, 9, 24, 16, 22],
including more recently in the computer science and machine learning literature [2, 1, 19, 25, 13].
First, we must define a family of RUMs that satisfies the conditions of Definition 1. Assume without loss of
generality that the noise distribution E has unit variance. Then, consider the family of RUMs parameterized
by θ in which candidates are ranked according to xi + εθi . By this definition, the standard deviation of the
noise for a particular value of θ is simply 1/θ. Intuitively, larger values of θ reduce the effect of the noise,
making the ranking more accurate. In Theorem 5 in Appendix A, we show as long as the noise distribution
E has positive support on (−∞, ∞), this definition of Fθ meets the differentiability, asymptotic optimality,
and monotonicity conditions in Definition 1. For distributions with finite support, many of our results can
be generalized by relaxing strict inequalities in Definition 1 and Theorem 1 to weak inequalities.

[Next verbatim source excerpt]


[form-feed in source extraction]
Monotonicity: The removal of any elements does not alter the distribution of the remaining elements,
meaning that the distribution of π (−S) is equivalent to a RUM with n − |S| elements. Thus, it suffices to
show that for a RUM with positive support on (−∞, ∞), the probability of ranking the best candidate first
strictly increases with θ.
Recall that by definition, the candidates are ranked according to xi + εθi . The probability that x1 is
ranked first is
[U+0014 control character]
[U+0015 control character]
[U+0014 control character]
[U+0015 control character]
εi
εi
ε1
ε1
Pr x1 +
> max xi +
= Pr
> max xi − x1 +
2≤i≤n
2≤i≤n
θ
θ
θ
θ
[U+0014 control character]
[U+0015 control character]
= Pr ε1 > max θ(xi − x1 ) + εi
2≤i≤n
[U+0014 control character]
[U+0015 control character]
= Eε2 ,...,εn Pr ε1 > max θ(xi − x1 ) + εi | ε2 , . . . , εn
(A.1)
2≤i≤n

We want to show that (A.1) is increasing in θ. Intuitively, this is because as θ increases, the right hand
side of the inequality inside the probability decreases. To prove this formally, it suffices to show that the
subderivative of (A.1) with respect to θ only includes strictly positive numbers. First, we have
[U+0014 control character]
[U+0015 control character]
∂
Eε2 ,...,εn Pr ε1 > max θ(xi − x1 ) + εi | ε2 , . . . , εn ⊂ R>0
2≤i≤n
∂θ
[U+0014 control character]
[U+0015 control character]
∂
Pr ε1 > max θ(xi − x1 ) + εi | ε2 , . . . , εn ⊂ R>0
⇐⇒
2≤i≤n
∂θ
Let F and f be the cumulative density function and probability density function of E respectively. Then,
[U+0014 control character]
[U+0015 control character]
[U+0012 control character]
[U+0013 control character]
Pr ε1 > max θ(xi − x1 ) + εi | ε2 , . . . , εn = 1 − F max θ(xi − x1 ) + εi
2≤i≤n

2≤i≤n

Note that F (·) is strictly increasing (since f is assumed to have positive support on (−∞, ∞)), so it suffices
to show that
∂
max θ(xi − x1 ) + εi ⊂ R<0
∂θ 2≤i≤n
For any i,

d
θ(xi − x1 ) + εi = xi − x1 < 0.
dθ
Thus, the subderivative of the max of such functions includes only strictly negative numbers, which completes
\end{ReviewVerbatim}


\paragraph{Expanded PaperInterface specification}
\begin{ReviewVerbatim}
∀ {n : ℕ} (f : ℝ → ℝ),
  Measurable f →
    (∀ (x : ℝ), 0 ≤ f x) →
      ∫⁻ (x : ℝ), ENNReal.ofReal (f x) = 1 →
        ∀ (value : EconCSLib.SocialChoice.Ranking.Candidate n → ℝ),
          (∀ (d : Fin (n + 1)), value d.succ < value 0) →
            ∀ {theta : ℝ},
              0 < theta →
                (EconCSLib.measureProb
                    ((KR21Monoculture.sourceAppendixARestNoiseLaw n (KR21Monoculture.w11BaseNoiseLaw f)).prod
                      (KR21Monoculture.w11BaseNoiseLaw f))
                    fun z =>
                    EconCSLib.SocialChoice.Ranking.firstChoice
                        (EconCSLib.SocialChoice.Ranking.rankByScore fun i =>
                          value i + KR21Monoculture.sourceAppendixAProductNoise z i / theta) =
                      0) =
                  ∫ (rest : Fin (n + 1) → ℝ),
                    EconCSLib.measureProb (KR21Monoculture.w11BaseNoiseLaw f) fun epsilon =>
                      ∀ (d : Fin (n + 1)),
                        theta * (value d.succ - value 0) + rest d <
                          epsilon ∂KR21Monoculture.sourceAppendixARestNoiseLaw n (KR21Monoculture.w11BaseNoiseLaw f)
\end{ReviewVerbatim}

\paragraph{Lean proof endpoint (built separately)}\begin{ReviewVerbatim}
KR21Monoculture.PaperInterface.equationA1_source_w11_iid_literal_event_conditionalTail_integral
\end{ReviewVerbatim}

\paragraph{Recorded source-to-Spec screening}
\textbf{Verdict:} \texttt{not recorded}
\\
\textbf{Reason:} not recorded
\reviewmatch{Matches source input}
\noindent\textbf{Reviewer annotation}\par
\reviewerbox
\clearpage
\hypertarget{source-claim-aa9b8f734e8009d3}{}
\section*{Review row 23}
\textbf{Source locator:} {\footnotesize\raggedright cited publication, lines 190-\allowbreak{}198\par}
\paragraph{Verbatim source input}
\begin{ReviewVerbatim}
The choice of which ranking mechanism to use leads to a game-theoretic setting: both firms know the
accuracy parameters of the human evaluators (θH ) and the algorithm (θA ), and they must decide whether to
3


[form-feed in source extraction]
use a human evaluator or the algorithm. This choice introduces a subtlety: for many ranking models, a firm’s
rational behavior depends not only on the accuracy of the ranking mechanism, but also on the underlying
candidate values x1 , . . . , xn . Thus, to fully specify a firm’s behavior, we assume that x1 , . . . , xn are drawn
from a known joint distribution D. Our main results will hold for any D, meaning they apply even when the
candidate values (but not their identities) are deterministically known.

[Next verbatim source excerpt]

B

3-candidate RUM Counterexamples

B.1

Violating Definition 2

Here, we provide a noise mode E, accuracy parameter θ, and candidate distribution D such that UAH < UAA .
Choose the noise distribution E and accuracy parameter θ such that
[U+F8F1 delimiter glyph]
[U+F8F4 delimiter glyph]1
w.p. 2δ
ε [U+F8F2 delimiter glyph]
= 0
w.p.1 − δ
θ [U+F8F4 delimiter glyph]
[U+F8F3 delimiter glyph]
−1 w.p. 2δ
Note that this distribution does not satisfy Definition 1 because it is neither differentiable nor supported on
(−∞, ∞); however, we can provide a “smooth” approximation to this distribution by expressing it as the
sum of arbitrarily tightly concentrated Gaussians with the same results.

15


[form-feed in source extraction]
We choose the candidate distribution D such that x1 − 1 > x2 > x3 > x1 − 2. For example,
7
4
1
x2 =
2
x3 = 0
x1 =

Under this condition, assuming x3 = 0 without loss of generality,
UAH (θ, θ) − UAA (θ, θ) =

[U+0001 control character]
δ2 3
δ x1 − 4δ 2 x1 + 4δx1 + 2δ 3 x2 − 14δ 2 x2 + 20δx2 − 8x2
32
2

Notice that the lowest-power δ term is − δ 4x2 . Therefore, for sufficiently small δ, this is negative. For
example, plugging in the values given above with δ = .1, UAH (θ, θ) − UAA (θ, θ) ≈ −0.00076.

B.2
\end{ReviewVerbatim}


\paragraph{Expanded PaperInterface specification}
\begin{ReviewVerbatim}
have componentLaw := KR21Monoculture.appendixB1NoisePMF;
have rawNoiseLaw := EconCSLib.pmfProd (EconCSLib.pmfProd componentLaw componentLaw) componentLaw;
have rawRank := fun noise =>
  EconCSLib.SocialChoice.Ranking.rankByScore fun c =>
    KR21Monoculture.appendixB1Value c +
      KR21Monoculture.appendixB1NoiseValue (KR21Monoculture.appendixB1NoiseTripleFunction noise c);
have rawLaw := PMF.map rawRank rawNoiseLaw;
(KR21Monoculture.appendixB1Value 0 = 7 / 4 ∧
    KR21Monoculture.appendixB1Value 1 = 1 / 2 ∧ KR21Monoculture.appendixB1Value 2 = 0) ∧
  (KR21Monoculture.appendixB1NoiseValue KR21Monoculture.AppendixB1NoiseAtom.plusOne = 1 ∧
      KR21Monoculture.appendixB1NoiseValue KR21Monoculture.AppendixB1NoiseAtom.zero = 0 ∧
        KR21Monoculture.appendixB1NoiseValue KR21Monoculture.AppendixB1NoiseAtom.minusOne = -1) ∧
    (componentLaw KR21Monoculture.AppendixB1NoiseAtom.plusOne).toReal = 1 / 20 ∧
      (componentLaw KR21Monoculture.AppendixB1NoiseAtom.zero).toReal = 9 / 10 ∧
        (componentLaw KR21Monoculture.AppendixB1NoiseAtom.minusOne).toReal = 1 / 20 ∧
          (∀ (noise : KR21Monoculture.AppendixB1NoiseTriple),
              rawRank noise = KR21Monoculture.appendixB1DiscreteTripleRank noise) ∧
            rawLaw = KR21Monoculture.appendixB1RankingPMF ∧
              EconCSLib.SocialChoice.Ranking.expectedSecondMoverIndependent rawLaw rawLaw
                      KR21Monoculture.appendixB1Value -
                    EconCSLib.SocialChoice.Ranking.expectedSecondMoverShared rawLaw KR21Monoculture.appendixB1Value =
                  -(9749 / 12800000) ∧
                ¬KR21Monoculture.Model.PrefersIndependentReranking rawLaw KR21Monoculture.appendixB1Value
\end{ReviewVerbatim}

\paragraph{Lean proof endpoint (built separately)}\begin{ReviewVerbatim}
KR21Monoculture.PaperInterface.appendixB1_source_discrete_definition2_counterexample_semantic_complete
\end{ReviewVerbatim}

\paragraph{Recorded source-to-Spec screening}
\textbf{Verdict:} \texttt{not recorded}
\\
\textbf{Reason:} not recorded
\reviewmatch{Matches source input}
\noindent\textbf{Reviewer annotation}\par
\reviewerbox
\clearpage
\hypertarget{source-claim-8c4914b3b2e7a592}{}
\section*{Review row 24}
\textbf{Source locator:} {\footnotesize\raggedright cited publication, lines 190-\allowbreak{}198\par}
\paragraph{Verbatim source input}
\begin{ReviewVerbatim}
The choice of which ranking mechanism to use leads to a game-theoretic setting: both firms know the
accuracy parameters of the human evaluators (θH ) and the algorithm (θA ), and they must decide whether to
3


[form-feed in source extraction]
use a human evaluator or the algorithm. This choice introduces a subtlety: for many ranking models, a firm’s
rational behavior depends not only on the accuracy of the ranking mechanism, but also on the underlying
candidate values x1 , . . . , xn . Thus, to fully specify a firm’s behavior, we assume that x1 , . . . , xn are drawn
from a known joint distribution D. Our main results will hold for any D, meaning they apply even when the
candidate values (but not their identities) are deterministically known.

[Next verbatim source excerpt]

Violating Definition 3

Next, we’ll give a 3-candidate RUM for which UAH < UHH does not hold in general. Consider the following
3-candidate example.
x1 = 3
x2 = 2
x3 = 0
Choose E and θ such that

[U+F8F1 delimiter glyph]
1
[U+F8F4 delimiter glyph]
[U+F8F4 delimiter glyph]
[U+F8F4 delimiter glyph]
[U+F8F2 delimiter glyph]
−1
ε
=
θ [U+F8F4 delimiter glyph]
10
[U+F8F4 delimiter glyph]
[U+F8F4 delimiter glyph]
[U+F8F3 delimiter glyph]
−10

w.p. 1−δ
2
w.p. 1−δ
2
w.p. 2δ
w.p. 2δ

Again, while this noise model doesn’t satisfy Definition 1, we can approximate it arbitrarily closely with the
sum of tightly concentrated Gaussians. Let the θA = 1.1θ and θH = 0.9θ.
We will show that for these parameters, UAH (θA , θH ) > UHH (θA , θH ), i.e., it is somehow better to choose
after a better opponent than after a worse opponent. At a high level, the reasoning for this is as follows:
1. When choosing first, the only difference between the algorithm and the human evaluator is that the
algorithm is more likely to choose x2 than x3 . Both strategies have identical probabilities of selecting
x1 .
2. When choosing second, the human evaluator’s utility is higher when x2 has already been chosen than
when x3 has already been chosen. This is because when x2 is unavailable, the human evaluator
is almost guaranteed to get x1 ; when x3 is unavailable, the human evaluator will choose x2 with
probability ≈ 1/4.
Let τ and π be rankings generated by the algorithm and human evaluator respectively. First, we will
show that
Pr[τ1 = x1 ] = Pr[π1 = x1 ]

(B.1)

Pr[τ1 = x2 ] > Pr[π1 = x2 ]

(B.2)

To do so, consider the realizations of ε1 , ε2 , ε3 that result in different rankings under θA and θH . In fact,
the only set of realizations that result in different rankings are when ε2 /θ = −1 and ε3 /θ = 1. Thus,
the algorithm and human evaluator always rank x1 in the same position, conditioned on a realization,
which proves (B.1); the only difference is that the algorithm sometimes ranks x2 above x3 when the human
16


[form-feed in source extraction]
evaluator does not. Moreover, whenever ε1 /θ = −10, x2 is more strictly more likely to be ranked first under
the algorithm than the human evaluator, which proves (B.2).
Next, we must show that when choosing second, the human evaluator is better off when x2 is unavailable
than when x3 is unavailable. This is clearly true because for the human evaluator,
[U+0014 control character]
[U+0015 control character]
ε1 θH
ε3 θH
Pr x1 +
> x3 +
≈ 1 − O(δ)
θ
θ
[U+0014 control character]
[U+0015 control character]
ε1 θH
ε3 θH
3
Pr x1 +
> x2 +
≈
θ
θ
4
Thus, conditioned on x2 being unavailable, the human evaluator gets utility ≈ 3, whereas when x3 is
unavailable, the human evaluator gets utility ≈ 2.75. Let u−i be the expected utility for the human evaluator
when xi is unavailable. Putting this together, we get
UAH (θA , θH ) − UHH (θA , θH )
=

3
X

(Pr[τ1 = xi ] − Pr[π1 = xi ])u−i

i=1

= (Pr[τ1 = x1 ] − Pr[π1 = x1 ])u−1 + (Pr[τ1 = x2 ] − Pr[π1 = x2 ])u−2 + (Pr[τ1 = x3 ] − Pr[π1 = x3 ])u−3
= (Pr[τ1 = x2 ] − Pr[π1 = x2 ])u−2 + (Pr[τ1 = x3 ] − Pr[π1 = x3 ])u−3
(Pr[τ1 = x1 = Pr[π1 = x1 ])
P3
P3
= (Pr[τ1 = x2 ] − Pr[π1 = x2 ])(u−2 − u−3 )
( i=1 Pr[τ1 = xi ] = i=1 Pr[π1 = xi ])
>0
The last step follows from (B.2) and because u−2 > u−3 .
\end{ReviewVerbatim}


\paragraph{Expanded PaperInterface specification}
\begin{ReviewVerbatim}
have componentLaw := KR21Monoculture.appendixB2NoisePMF;
have rawNoiseLaw := EconCSLib.pmfProd (EconCSLib.pmfProd componentLaw componentLaw) componentLaw;
have algorithmRank := fun noise =>
  EconCSLib.SocialChoice.Ranking.rankByScore fun c =>
    KR21Monoculture.appendixB2Value c +
      10 / 11 * KR21Monoculture.appendixB2NoiseValue (KR21Monoculture.appendixB2NoiseTripleFunction noise c);
have humanRank := fun noise =>
  EconCSLib.SocialChoice.Ranking.rankByScore fun c =>
    KR21Monoculture.appendixB2Value c +
      10 / 9 * KR21Monoculture.appendixB2NoiseValue (KR21Monoculture.appendixB2NoiseTripleFunction noise c);
have algorithmLaw := PMF.map algorithmRank rawNoiseLaw;
have humanLaw := PMF.map humanRank rawNoiseLaw;
(KR21Monoculture.appendixB2Value 0 = 3 ∧
    KR21Monoculture.appendixB2Value 1 = 2 ∧ KR21Monoculture.appendixB2Value 2 = 0) ∧
  (KR21Monoculture.appendixB2NoiseValue KR21Monoculture.AppendixB2NoiseAtom.plusOne = 1 ∧
      KR21Monoculture.appendixB2NoiseValue KR21Monoculture.AppendixB2NoiseAtom.minusOne = -1 ∧
        KR21Monoculture.appendixB2NoiseValue KR21Monoculture.AppendixB2NoiseAtom.plusTen = 10 ∧
          KR21Monoculture.appendixB2NoiseValue KR21Monoculture.AppendixB2NoiseAtom.minusTen = -10) ∧
    (componentLaw KR21Monoculture.AppendixB2NoiseAtom.plusOne).toReal = 9 / 20 ∧
      (componentLaw KR21Monoculture.AppendixB2NoiseAtom.minusOne).toReal = 9 / 20 ∧
        (componentLaw KR21Monoculture.AppendixB2NoiseAtom.plusTen).toReal = 1 / 20 ∧
          (componentLaw KR21Monoculture.AppendixB2NoiseAtom.minusTen).toReal = 1 / 20 ∧
            (11 / 10 > 9 / 10 ∧ 10 / 11 * (11 / 10) = 1 ∧ 10 / 9 * (9 / 10) = 1) ∧
              (∀ (noise : KR21Monoculture.AppendixB2NoiseTriple),
                  algorithmRank noise = KR21Monoculture.appendixB2AlgorithmDiscreteTripleRank noise) ∧
                (∀ (noise : KR21Monoculture.AppendixB2NoiseTriple),
                    humanRank noise = KR21Monoculture.appendixB2HumanDiscreteTripleRank noise) ∧
                  algorithmLaw = KR21Monoculture.appendixB2AlgorithmRankingPMF ∧
                    humanLaw = KR21Monoculture.appendixB2HumanRankingPMF ∧
                      EconCSLib.SocialChoice.Ranking.expectedSecondMoverIndependent humanLaw algorithmLaw
                              KR21Monoculture.appendixB2Value -
                            EconCSLib.SocialChoice.Ranking.expectedSecondMoverIndependent humanLaw humanLaw
                              KR21Monoculture.appendixB2Value =
                          567 / 3200000 ∧
                        ¬KR21Monoculture.Model.PrefersWeakerCompetition algorithmLaw humanLaw
                            KR21Monoculture.appendixB2Value
\end{ReviewVerbatim}

\paragraph{Lean proof endpoint (built separately)}\begin{ReviewVerbatim}
KR21Monoculture.PaperInterface.appendixB2_source_discrete_definition3_counterexample_semantic_complete
\end{ReviewVerbatim}

\paragraph{Recorded source-to-Spec screening}
\textbf{Verdict:} \texttt{not recorded}
\\
\textbf{Reason:} not recorded
\reviewmatch{Matches source input}
\noindent\textbf{Reviewer annotation}\par
\reviewerbox
\clearpage
\hypertarget{source-claim-36e8cf848060682f}{}
\section*{Review row 25}
\textbf{Source locator:} {\footnotesize\raggedright cited publication, lines 1040-\allowbreak{}1057\par}
\paragraph{Verbatim source input}
\begin{ReviewVerbatim}
2. When choosing second, the human evaluator’s utility is higher when x2 has already been chosen than
when x3 has already been chosen. This is because when x2 is unavailable, the human evaluator
is almost guaranteed to get x1 ; when x3 is unavailable, the human evaluator will choose x2 with
probability ≈ 1/4.
Let τ and π be rankings generated by the algorithm and human evaluator respectively. First, we will
show that
Pr[τ1 = x1 ] = Pr[π1 = x1 ]

(B.1)

Pr[τ1 = x2 ] > Pr[π1 = x2 ]

(B.2)

To do so, consider the realizations of ε1 , ε2 , ε3 that result in different rankings under θA and θH . In fact,
the only set of realizations that result in different rankings are when ε2 /θ = −1 and ε3 /θ = 1. Thus,
the algorithm and human evaluator always rank x1 in the same position, conditioned on a realization,
which proves (B.1); the only difference is that the algorithm sometimes ranks x2 above x3 when the human
\end{ReviewVerbatim}


\paragraph{Expanded PaperInterface specification}
\begin{ReviewVerbatim}
have componentLaw := KR21Monoculture.appendixB2NoisePMF;
have rawNoiseLaw := EconCSLib.pmfProd (EconCSLib.pmfProd componentLaw componentLaw) componentLaw;
have algorithmRank := fun noise =>
  EconCSLib.SocialChoice.Ranking.rankByScore fun c =>
    KR21Monoculture.appendixB2Value c +
      10 / 11 * KR21Monoculture.appendixB2NoiseValue (KR21Monoculture.appendixB2NoiseTripleFunction noise c);
have humanRank := fun noise =>
  EconCSLib.SocialChoice.Ranking.rankByScore fun c =>
    KR21Monoculture.appendixB2Value c +
      10 / 9 * KR21Monoculture.appendixB2NoiseValue (KR21Monoculture.appendixB2NoiseTripleFunction noise c);
have algorithmLaw := PMF.map algorithmRank rawNoiseLaw;
have humanLaw := PMF.map humanRank rawNoiseLaw;
(KR21Monoculture.appendixB2Value 0 = 3 ∧
    KR21Monoculture.appendixB2Value 1 = 2 ∧ KR21Monoculture.appendixB2Value 2 = 0) ∧
  (KR21Monoculture.appendixB2NoiseValue KR21Monoculture.AppendixB2NoiseAtom.plusOne = 1 ∧
      KR21Monoculture.appendixB2NoiseValue KR21Monoculture.AppendixB2NoiseAtom.minusOne = -1 ∧
        KR21Monoculture.appendixB2NoiseValue KR21Monoculture.AppendixB2NoiseAtom.plusTen = 10 ∧
          KR21Monoculture.appendixB2NoiseValue KR21Monoculture.AppendixB2NoiseAtom.minusTen = -10) ∧
    (componentLaw KR21Monoculture.AppendixB2NoiseAtom.plusOne).toReal = 9 / 20 ∧
      (componentLaw KR21Monoculture.AppendixB2NoiseAtom.minusOne).toReal = 9 / 20 ∧
        (componentLaw KR21Monoculture.AppendixB2NoiseAtom.plusTen).toReal = 1 / 20 ∧
          (componentLaw KR21Monoculture.AppendixB2NoiseAtom.minusTen).toReal = 1 / 20 ∧
            algorithmLaw = KR21Monoculture.appendixB2AlgorithmRankingPMF ∧
              humanLaw = KR21Monoculture.appendixB2HumanRankingPMF ∧
                KR21Monoculture.firstChoiceProb algorithmLaw 0 = KR21Monoculture.firstChoiceProb humanLaw 0
\end{ReviewVerbatim}

\paragraph{Lean proof endpoint (built separately)}\begin{ReviewVerbatim}
KR21Monoculture.PaperInterface.equationB1_counterexample_first_choice_x1_semantic_complete
\end{ReviewVerbatim}

\paragraph{Recorded source-to-Spec screening}
\textbf{Verdict:} \texttt{not recorded}
\\
\textbf{Reason:} not recorded
\reviewmatch{Matches source input}
\noindent\textbf{Reviewer annotation}\par
\reviewerbox
\clearpage
\hypertarget{source-claim-a3b878a82854afb8}{}
\section*{Review row 26}
\textbf{Source locator:} {\footnotesize\raggedright cited publication, lines 1044-\allowbreak{}1105\par}
\paragraph{Verbatim source input}
\begin{ReviewVerbatim}
Let τ and π be rankings generated by the algorithm and human evaluator respectively. First, we will
show that
Pr[τ1 = x1 ] = Pr[π1 = x1 ]

(B.1)

Pr[τ1 = x2 ] > Pr[π1 = x2 ]

(B.2)

To do so, consider the realizations of ε1 , ε2 , ε3 that result in different rankings under θA and θH . In fact,
the only set of realizations that result in different rankings are when ε2 /θ = −1 and ε3 /θ = 1. Thus,
the algorithm and human evaluator always rank x1 in the same position, conditioned on a realization,
which proves (B.1); the only difference is that the algorithm sometimes ranks x2 above x3 when the human
16


[form-feed in source extraction]
evaluator does not. Moreover, whenever ε1 /θ = −10, x2 is more strictly more likely to be ranked first under
the algorithm than the human evaluator, which proves (B.2).
Next, we must show that when choosing second, the human evaluator is better off when x2 is unavailable
than when x3 is unavailable. This is clearly true because for the human evaluator,
[U+0014 control character]
[U+0015 control character]
ε1 θH
ε3 θH
Pr x1 +
> x3 +
≈ 1 − O(δ)
θ
θ
[U+0014 control character]
[U+0015 control character]
ε1 θH
ε3 θH
3
Pr x1 +
> x2 +
≈
θ
θ
4
Thus, conditioned on x2 being unavailable, the human evaluator gets utility ≈ 3, whereas when x3 is
unavailable, the human evaluator gets utility ≈ 2.75. Let u−i be the expected utility for the human evaluator
when xi is unavailable. Putting this together, we get
UAH (θA , θH ) − UHH (θA , θH )
=

3
X

(Pr[τ1 = xi ] − Pr[π1 = xi ])u−i

i=1

= (Pr[τ1 = x1 ] − Pr[π1 = x1 ])u−1 + (Pr[τ1 = x2 ] − Pr[π1 = x2 ])u−2 + (Pr[τ1 = x3 ] − Pr[π1 = x3 ])u−3
= (Pr[τ1 = x2 ] − Pr[π1 = x2 ])u−2 + (Pr[τ1 = x3 ] − Pr[π1 = x3 ])u−3
(Pr[τ1 = x1 = Pr[π1 = x1 ])
P3
P3
= (Pr[τ1 = x2 ] − Pr[π1 = x2 ])(u−2 − u−3 )
( i=1 Pr[τ1 = xi ] = i=1 Pr[π1 = xi ])
>0
The last step follows from (B.2) and because u−2 > u−3 .
\end{ReviewVerbatim}


\paragraph{Expanded PaperInterface specification}
\begin{ReviewVerbatim}
have componentLaw := KR21Monoculture.appendixB2NoisePMF;
have rawNoiseLaw := EconCSLib.pmfProd (EconCSLib.pmfProd componentLaw componentLaw) componentLaw;
have algorithmRank := fun noise =>
  EconCSLib.SocialChoice.Ranking.rankByScore fun c =>
    KR21Monoculture.appendixB2Value c +
      10 / 11 * KR21Monoculture.appendixB2NoiseValue (KR21Monoculture.appendixB2NoiseTripleFunction noise c);
have humanRank := fun noise =>
  EconCSLib.SocialChoice.Ranking.rankByScore fun c =>
    KR21Monoculture.appendixB2Value c +
      10 / 9 * KR21Monoculture.appendixB2NoiseValue (KR21Monoculture.appendixB2NoiseTripleFunction noise c);
have algorithmLaw := PMF.map algorithmRank rawNoiseLaw;
have humanLaw := PMF.map humanRank rawNoiseLaw;
(KR21Monoculture.appendixB2Value 0 = 3 ∧
    KR21Monoculture.appendixB2Value 1 = 2 ∧ KR21Monoculture.appendixB2Value 2 = 0) ∧
  (KR21Monoculture.appendixB2NoiseValue KR21Monoculture.AppendixB2NoiseAtom.plusOne = 1 ∧
      KR21Monoculture.appendixB2NoiseValue KR21Monoculture.AppendixB2NoiseAtom.minusOne = -1 ∧
        KR21Monoculture.appendixB2NoiseValue KR21Monoculture.AppendixB2NoiseAtom.plusTen = 10 ∧
          KR21Monoculture.appendixB2NoiseValue KR21Monoculture.AppendixB2NoiseAtom.minusTen = -10) ∧
    (componentLaw KR21Monoculture.AppendixB2NoiseAtom.plusOne).toReal = 9 / 20 ∧
      (componentLaw KR21Monoculture.AppendixB2NoiseAtom.minusOne).toReal = 9 / 20 ∧
        (componentLaw KR21Monoculture.AppendixB2NoiseAtom.plusTen).toReal = 1 / 20 ∧
          (componentLaw KR21Monoculture.AppendixB2NoiseAtom.minusTen).toReal = 1 / 20 ∧
            algorithmLaw = KR21Monoculture.appendixB2AlgorithmRankingPMF ∧
              humanLaw = KR21Monoculture.appendixB2HumanRankingPMF ∧
                KR21Monoculture.firstChoiceProb humanLaw 1 < KR21Monoculture.firstChoiceProb algorithmLaw 1
\end{ReviewVerbatim}

\paragraph{Lean proof endpoint (built separately)}\begin{ReviewVerbatim}
KR21Monoculture.PaperInterface.equationB2_counterexample_first_choice_x2_semantic_complete
\end{ReviewVerbatim}

\paragraph{Recorded source-to-Spec screening}
\textbf{Verdict:} \texttt{not recorded}
\\
\textbf{Reason:} not recorded
\reviewmatch{Matches source input}
\noindent\textbf{Reviewer annotation}\par
\reviewerbox
\clearpage
\hypertarget{source-claim-34bc63da82092219}{}
\section*{Review row 27}
\textbf{Source locator:} {\footnotesize\raggedright cited publication, lines 1111-\allowbreak{}1147\par}
\paragraph{Verbatim source input}
\begin{ReviewVerbatim}
C.1

Verifying Definition 2

By (2), we can equivalently show that for any θ, UAH (θ, θ) > UAA (θ, θ). Let τ and π be the algorithmic and
human-generated rankings respectively. Note that they’re identically distributed because θA = θH . Define
(
π1 π1 6= τ1
Y ,
π2 otherwise
Note that UAH (θ, θ) = E [Y ] and UAA (θ, θ) = E [τ2 ]. We want to show that UAH (θ, θ) − UAA (θ, θ) =
E [xY − xτ2 ] > 0. It is sufficient to show that for any k, E [Y − τ2 | τ1 = xk ] > 0. Let Xi = xi + εi /θ. Note
that for distinct i, j, k and xi > xj ,
E [Y − τ2 | τ1 = xk ] > 0 ⇐=

Pr[τ2 = xi | τ1 = xk ]
Pr[Y = xi | τ1 = xk ]
>
Pr[Y = xj | τ1 = xk ]
Pr[τ2 = xj | τ1 = xk ]

⇐⇒ Pr[Y = xi | τ1 = xk ] > Pr[τ2 = xi | τ1 = xk ]
(numerator and denominator sum to 1)
⇐⇒ Pr[Xi > Xj ] > Pr[Xi > Xj | Xk > Xi ∩ Xk > Xj ]
⇐⇒ Pr[Xi > Xj ] > EXk [Pr[Xi > Xj | Xk = a, Xi < a, Xj < a]].
Thus, it suffices to show that for any a,
Pr[Xi > Xj ] > Pr[Xi > Xj | Xi < a, Xj < a].

(C.1)

Since Pr[Xi > Xj ] = lima→∞ Pr[Xi > Xj | Xi < a, Xj < a], it suffices to show that for all a,
d
Pr[Xi > Xj | Xi < a, Xj < a] ≥ 0,
(C.2)
da
and that it is strictly positive for some a. In other words, the higher a is, the more likely i and j are to be
correctly ordered. In Theorems 7 and 8, we show that (C.2) holds for both Laplacian and Gaussian noise

[Next verbatim source excerpt]

Finally, we show that (C.2) holds for both Laplacian and Gaussian noise.
Theorem 7. For any a ∈ R and Xi = xi + σεi where εi is Laplacian with unit variance,
d
Pr[Xi > Xj | Xi < a, Xj < a] ≥ 0.
da
Moreover, it is strictly positive for some a.
Proof. First, we must derive an expression for Pr[Xi > Xj | Xi < a, Xj < a]. Recall that the Laplace
distribution parameterized by µ and λ has pdf
f (x; µ, λ) =
and cdf

λ
exp(−λ|x − µ|)
2

(
F (x; µ, λ) =

1
2 exp (−λ(µ − x))
1 − 12 exp (−λ(x − µ))

21

x<µ
x≥µ


[form-feed in source extraction]
Note that xi and xj be the respective means of Xi and Xj , with xi > xj . Because the Laplace distribution
is piecewise defined, we must consider 3 cases and show that in all 3 cases, (C.2) holds. Note that
Ra
f (x; xi , λ)F (x; xj , λ) dx
(C.3)
Pr[Xi > Xj | Xi < a, Xj < a] = −∞
F (a; xi , λ)F (a; xj , λ)
Case 1: a ≤ xj .
Then, the numerator of (C.3) is
Z a
Z
λ
1
λ a
exp(−λ(xi + xj − 2x)) dx
exp (−λ(xi − x)) · exp(−λ(xj − x)) dx =
2
4 −∞
−∞ 2
Z
λ exp(−λ(xi + xj )) a
exp(2λx) dx
=
4
−∞
λ exp(−λ(xi + xj )) 1
=
exp(2λa)
4
2λ
exp(−λ(xi + xj − 2a))
=
8
The denominator is
1
1
1
exp(−λ(xi − a)) · exp(−λ(xj − a)) = exp(−λ(xi + xj − 2a)).
2
2
4
Thus,
Pr[Xi > Xj | Xi < a, Xj < a] =

1
,
2

so its derivative is trivially nonnegative.
Case 2: xj < a ≤ xi .
Then, the numerator of (C.3) is
[U+0012 control character]
[U+0013 control character]
Z xj
Z a
λ
1
λ
1
exp (−λ(xi − x)) · exp(−λ(xj − x)) dx +
exp (−λ(xi − x)) 1 − exp(−λ(x − xj )) dx
2
2
−∞ 2
xj 2
Z a
Z a
exp(−λ(xi − xj )) λ
λ
+
exp(−λ(xi − x)) dx −
exp(−λ(xi − xj )) dx
=
8
2 xj
4 xj
exp(−λ(xi − xj )) λ 1
λ
+
(exp(−λ(xi − a)) − exp(−λ(xi − xj ))) − (a − xj ) exp(−λ(xi − xj ))
8
2[U+0012 control character]λ
4
[U+0013 control character]
1
3 λ
= exp(−λ(xi − a)) −
+ (a − xj ) exp(−λ(xi − xj ))
2
8
4

=

The denominator is
[U+0012 control character]
[U+0013 control character]
1
1
1
1
1 − exp(−λ(a − xj )) · exp(λ(xj − a)) = exp(−λ(xi − a)) − exp(−λ(xi − xj ))
2
2
2
4
We can factor out 41 exp(−λ(xi − xj )) from both, so
2 exp(λ(a − xj )) − 23 + λ(a − xj )
Pr[Xi > Xj | Xi < a, Xj < a] =
2 exp(λ(a − xj )) − 1
=

[U+0001 control character]

[U+0001 control character]
2 exp(λ(a − xj )) − 1 − 21 + λ(a − xj )
2 exp(λ(a − xj )) − 1

=1−

1
2 + λ(a − xj )

2 exp(λ(a − xj )) − 1

22


[form-feed in source extraction]
Thus,
d
Pr[Xi > Xj | Xi < a, Xj < a] > 0
da
1
d
2 + λ(a − xj )
⇐⇒
<0
da 2 exp(λ(a − xj )) − 1
[U+0012 control character]
[U+0013 control character]
1
⇐⇒(2 exp(λ(a − xj )) − 1)λ <
+ λ(a − xj ) 2λ exp(λ(a − xj ))
2
[U+0012 control character]
[U+0013 control character]
1
⇐⇒2 − exp(−λ(a − xj )) < 2
+ λ(a − xj )
2
⇐⇒1 − exp(−λ(a − xj )) < 2λ(a − xj )
⇐⇒ exp(−λ(a − xj )) > 1 − 2λ(a − xj )
This is true because λ(a − xj ) > 0, and for z > 0,
exp(−z) > 1 − z > 1 − 2z.
Case 3: a > xi .
Then, the numerator of (C.3) is
[U+0012 control character]
[U+0013 control character]
Z xj
Z xi
λ
1
λ
1
exp (−λ(xi − x)) · exp(−λ(xj − x)) dx +
exp (−λ(xi − x)) 1 − exp(−λ(x − xj )) dx
2
2
−∞ 2
xj 2
[U+0012 control character]
[U+0013 control character]
Z a
1
λ
exp(−λ(x − xi )) 1 − exp(−λ(x − xj )) dx
+
2
2
xi
[U+0012 control character]
[U+0013 control character]
Z
1
3 λ
1
λ a
= −
+ (xi − xj ) exp(−λ(xi − xj )) + (1 − exp(−λ(a − xi ))) −
exp(−λ(2x − xi − xj )) dx
2
8
4
2
4 xi
[U+0012 control character]
[U+0013 control character]
3 λ
1
=1−
+ (xi − xj ) exp(−λ(xi − xj )) − exp(−λ(a − xi ))
8
4
2
1
+ exp(λ(xi + xj ))(exp(−2λa) − exp(−2λxi ))
8 [U+0012 control character]
[U+0013 control character]
1
1
1 λ
+ (xi − xj ) exp(−λ(xi − xj )) − exp(−λ(a − xi )) + exp(−λ(2a − xi − xj ))
=1−
2
4
2
8
The denominator is
[U+0012 control character]
[U+0013 control character][U+0012 control character]
[U+0013 control character]
1
1
1 − exp(−λ(t − xi ))
1 − exp(−λ(t − xj ))
2
2
1
1
1
= 1 − exp(−λ(a − xi )) − exp(−λ(a − xj )) + exp(−λ(2a − xi − xj ))
2
2
4
Thus,
Pr[Xi > Xj | Xi < a, Xj < a]
[U+0001 control character]
1 − 21 + λ4 (xi − xj ) exp(−λ(xi − xj )) − 12 exp(−λ(a − xi )) + 81 exp(−λ(2a − xi − xj ))
=
1 − 12 exp(−λ(a − xi )) − 21 exp(−λ(a − xj )) + 14 exp(−λ(2a − xi − xj ))
8 − (4 + 2λ(xi − xj )) exp(−λ(xi − xj )) − 4 exp(−λ(a − xi )) + exp(−λ(2a − xi − xj ))
∝
4 − 2 exp(−λ(a − xi )) − 2 exp(−λ(a − xj )) + exp(−λ(2a − xi − xj ))

23


[form-feed in source extraction]
We’re interested in
d
Pr[Xi > Xj | Xi < a, Xj < a] > 0
da
⇐⇒ (4 − 2 exp(−λ(a − xi )) − 2 exp(−λ(a − xj )) + exp(−λ(2a − xi − xj )
· (4λ exp(−λ(a − xi )) − 2λ exp(−λ(2a − xi − xj )))
> (8 − 4 exp(−λ(a − xi )) + exp(−λ(2a − xi − xj )) − (4 + 2λ(xi − xj )) exp(−λ(xi − xj )))
· (2λ exp(−λ(a − xi )) + 2λ exp(−λ(a − xj )) − 2λ exp(−λ(2a − xi − xj )))
⇐⇒16 exp(−λ(a − xi )) − 8 exp(−λ(2a − xi − xj )) − 8 exp(−2λ(a − xi )) + 4 exp(−λ(3a − 2xi − xj ))
− 8 exp(−λ(2a − xi − xj )) + 4 exp(−λ(3a − xi − 2xj )) + 4 exp(−λ(3a − 2xi − xj ))
− 2 exp(−2λ(2a − xi − xj ))
> 16 exp(−λ(a − xi )) + 16 exp(−λ(a − xj )) − 16 exp(−λ(2a − xi − xj ))
− 8 exp(−2λ(a − xi )) − 8 exp(−λ(2a − xi − xj )) + 8 exp(−λ(3a − 2xi − xj ))
+ 2 exp(−λ(3a − 2xi − xj )) + 2 exp(−λ(3a − xi − 2xj )) − 2 exp(−2λ(2a − xi − xj ))
− 2(4 + 2λ(xi − xj )) exp(−λ(a − xj )) − 2(4 + 2λ(xi − xj )) exp(−λ(a + xi − 2xj ))
+ 2(4 + 2λ(xi − xj )) exp(−2λ(a − xj ))
⇐⇒ exp(−λ(3a − xi − 2xj ))
> 8 exp(−λ(a − xj )) − 4 exp(−λ(2a − xi − xj )) + exp(−λ(3a − 2xi − xj ))
− (4 + 2λ(xi − xj )) exp(−λ(a − xj )) − (4 + 2λ(xi − xj )) exp(−λ(a + xi − 2xj ))
+ (4 + 2λ(xi − xj )) exp(−2λ(a − xj ))
⇐⇒ exp(−λ(2a − xi − xj ))
> 8 − 4 exp(−λ(a − xi )) + exp(−λ(2a − 2xi ))
− (4 + 2λ(xi − xj )) − (4 + 2λ(xi − xj )) exp(−λ(xi − xj )) + (4 + 2λ(xi − xj )) exp(−λ(a − xj ))
⇐⇒ exp(−λ(2a − xi − xj )) − 8 + 4 exp(−λ(a − xi )) − exp(−2λ(a − xi ))
+ (4 + 2λ(xi − xj ))(1 + exp(−λ(xi − xj ))) − (4 + 2λ(xi − xj )) exp(−λ(a − xj ))
>0

(C.4)

Note that for any z ≥ 0, we have
(4 + 2z)(1 + e−z ) − 8 ≥ 0 ⇐⇒ (2 + z)(1 + e−z ) ≥ 4
⇐⇒ z + 2e−z + ze−z ≥ 2
For z = 0, this holds with equality, and the left hand side is increasing since
d
z + 2e−z + ze−z ≥ 0 ⇐⇒ 1 − 2e−z + e−z − ze−z ≥ 0
dx
⇐⇒ 1 ≥ e−z + ze−z
1
⇐⇒
≥ e−z
1+z
⇐⇒ 1 + z ≤ ez
Therefore, choosing z = λ(xi − xj ) and plugging back to (C.4), we have
exp(−λ(2a − xi − xj )) − 8 + 4 exp(−λ(a − xi )) − exp(−2λ(a − xi ))
+ (4 + 2λ(xi − xj ))(1 + exp(−λ(xi − xj ))) − (4 + 2λ(xi − xj )) exp(−λ(a − xj )) > 0
⇐= exp(−λ(2a − xi − xj )) + 4 exp(−λ(a − xi )) − exp(−2λ(a − xi )) − (4 + 2λ(xi − xj )) exp(−λ(a − xj )) > 0
⇐⇒ exp(−λ(a − xj )) + 4 − exp(−λ(a − xi )) − (4 + 2λ(xi − xj )) exp(−λ(xi − xj )) > 0
⇐⇒4(1 − exp(−λ(xi − xj ))) + exp(−λ(a − xi ))(exp(−λ(xi − xj )) − 1) − 2λ(xi − xj ) exp(−λ(xi − xj )) > 0
⇐⇒(4 − exp(−λ(a − xi )))(1 − exp(−λ(xi − xj ))) − 2λ(xi − xj ) exp(−λ(xi − xj )) > 0
⇐=3(1 − exp(−λ(xi − xj ))) − 2λ(xi − xj ) exp(−λ(xi − xj )) > 0
24

(exp(−λ(a − xi )) < 1)


[form-feed in source extraction]
Again letting z = λ(xi − xj ), this is true if and only if
3(1 − e−z ) > 2ze−z ⇐⇒ 3(ez − 1) > 2z
⇐⇒ 3ez > 3 + 2z
which is true because ez > 1 + z for z > 0. This completes the proof for Case 3.
As a result, we have that
d
Pr[Xi > Xj | Xi < a, Xj < a] ≥ 0
da
for all a, with strict inequality for some a, which proves the theorem.

[Next verbatim source excerpt]

Theorem 8. For any a ∈ R and Xi = xi + σεi where εi ∼ N (0, 1),
d
Pr[Xi > Xj | Xi < a, Xj < a] > 0.
da
√
0
Proof. Assume σ = 1/ 2. This
√ is without loss of generality because for any instance with arbitrary σ , there
is an instance with σ = 1/ 2 that yields the same distribution over rankings (simply by scaling all item
values by σ/σ 0 ). First, we have
Ra
Pr[Xi = x] Pr[Xj < x] dx
Pr[Xi > Xj | Xi < a, Xj < a] = −∞
P r[Xi < a] Pr[Xj < a]
Ra
√
exp(−(x − xi )2 )/ π · (1 + erf(x − xj ))/2 dx
−∞
=
(1 + erf(a − xi ))/2 · (1 + erf(a − xj ))/2
Ra
exp(−(x − xi )2 )(1 + erf(x − xj )) dx
2
= √ −∞
(1 + erf(a − xi )) · (1 + erf(a − xj ))
π
The derivative with respect to a is positive if and only if
(1 + erf(a − xi ))(1 + erf(a − xj )) exp(−(a − xi )2 )(1 + erf(a − xj ))
Z a
>
exp(−(x − xi )2 )(1 + erf(x − xj )) dx
−∞

[U+0001 control character]
2
· √ (1 + erf(a − xi )) exp(−(a − xj )2 ) + (1 + erf(a − xj )) exp(−(a − xi )2 )
π

(C.5)

Let t = a − xi and δ = xi − xj . Then, using the fact that
Z a

Z a−xi

2

exp(−(x − xi ) )(1 + erf(x − xj )) dx =
−∞

2

Z a−xi

exp(−x ) dx +
−∞

√
=

π
(1 + erf(a − xi )) +
2

exp(−x2 ) erf(x + δ) dx

−∞
Z a−xi

exp(−x2 ) erf(x + δ) dx,

−∞

(C.5) becomes
√
√
Z t
π
(1 + erf(t))(1 + erf(t + δ))2 exp(−t2 )
π
·
>
(1
+
erf(t))
+
exp(−x2 ) erf(x + δ) dx
2 (1 + erf(t)) exp(−(t + δ)2 ) + (1 + erf(t + δ)) exp(−t2 )
2
−∞
Z t
(1 + erf(t))(1 + erf(t + δ))2 exp(−t2 )
2
√
⇐⇒
−
(1
+
erf(t))
−
exp(−x2 ) erf(x + δ) dx > 0
(1 + erf(t)) exp(−(t + δ)2 ) + (1 + erf(t + δ)) exp(−t2 )
π −∞
(C.6)
To show that this is true, we will use the fact that f (t) > 0 whenever the following conditions are met:
1. f (t) is continuous and differentiable everywhere
25


[form-feed in source extraction]
2. limt→−∞ f (t) = 0
d
3. dt
f (t) > 0

We’ll show that these conditions hold for the LHS of (C.6).
(1 + erf(t))(1 + erf(t + δ))2 exp(−t2 )
2
− (1 + erf(t)) − √
t→−∞ (1 + erf(t)) exp(−(t + δ)2 ) + (1 + erf(t + δ)) exp(−t2 )
π

Z t

lim

(1 + erf(t))(1 + erf(t + δ))2 exp(−t2 )
t→−∞ (1 + erf(t)) exp(−(t + δ)2 ) + (1 + erf(t + δ)) exp(−t2 )

= lim

exp(−x2 ) erf(x + δ) dx

−∞

(C.7)

Observe that both the numerator and denominator of (C.7) are positive, so this limit must be at least 0.
We can upper bound it by
(1 + erf(t))(1 + erf(t + δ))2 exp(−t2 )
(1 + erf(t))(1 + erf(t + δ))2 exp(−t2 )
≤
lim
t→−∞
t→−∞ (1 + erf(t)) exp(−(t + δ)2 ) + (1 + erf(t + δ)) exp(−t2 )
(1 + erf(t + δ)) exp(−t2 )
= lim (1 + erf(t))(1 + erf(t + δ))
lim

t→−∞

=0
Thus, the limit is 0. Now, we must show that the derivative is positive. The derivative is
[U+0015 control character]
[U+0014 control character]
Z t
2
(1 + erf(t))(1 + erf(t + δ))2 exp(−t2 )
d
2
− (1 + erf(t)) − √
exp(−x ) erf(x + δ) dx
dt (1 + erf(t)) exp(−(t + δ)2 ) + (1 + erf(t + δ)) exp(−t2 )
π −∞
[U+0014 control character]
[U+0015 control character]
d
(1 + erf(t))(1 + erf(t + δ))2 exp(−t2 )
2
2
=
− √ exp(−t2 ) − √ exp(−t2 ) erf(t + δ)
dt (1 + erf(t)) exp(−(t + δ)2 ) + (1 + erf(t + δ)) exp(−t2 )
π
π
(C.8)
Taking this derivative and factoring out
√

2(1 + erf(t))(1 + erf(t + δ)) exp(4t2 )
2

π (erf (t) + 1) et2 + (erf (δ + t) + 1) e(δ+t)

[U+0001 control character]2 ,

we get that (C.8) is positive if and only if
√
δ π exp((t + δ)2 )(1 + erf(t))(1 + erf(t + δ)) − exp(2δt + t2 )(1 + erf(t + δ)) + (1 + erf(t)) > 0
√
1 + erf(t)
⇐⇒δ π exp((t + δ)2 )(1 + erf(t)) +
− exp(2δt + t2 ) > 0
1 + erf(t + δ)
√
1 + erf(t)
⇐⇒δ π exp(t2 )(1 + erf(t)) + exp(−2δt − t2 )
−1>0
1 + erf(t + δ)
[U+0014 control character]
[U+0015 control character]
√
exp(−2δt − δ 2 )
⇐⇒(1 + erf(t)) δ π exp(t2 ) +
−1>0
1 + erf(t + δ)
[U+0015 control character]
[U+0014 control character]
1 + erf(t) √
exp(−(t + δ)2 )
⇐⇒
δ
π
+
−1>0
(C.9)
exp(−t2 )
1 + erf(t + δ)
Define
g(t) ,

1 + erf(t)
.
exp(−t2 )

Then, (C.9) is
[U+0014 control character]
√
g(t) δ π +

[U+0015 control character]
1
−1>0
g(t + δ)
√
1
1
⇐⇒
−
<δ π
g(t) g(t + δ)
26


[form-feed in source extraction]
By the Mean Value Theorem,
1
1
d 1
−
= −δ
g(t) g(t + δ)
dt g(t) t=t∗
for some t ≤ t∗ ≤ t + δ. Thus, it suffices to show that
√
d 1
>− π
dt g(t)

(C.10)

for all t. To do this, consider Mills Ratio [17]
Z ∞

2

R(t) , exp(t /2)

exp(−x2 /2) dx.

t

Note that this is quite similar in functional form to g(t), and with some manipulation, we can relate the two:
Z ∞
2
R(t) = exp(t /2)
exp(−x2 /2) dx
t
Z ∞
√
R( 2t) = exp(t2 ) √ exp(−x2 /2) dx
2t
Z ∞
√
2
= 2 exp(t )
exp(−x2 ) dx
t
Z −t
√
2
= 2 exp(t )
exp(−x2 ) dx
(exp(−x2 ) is symmetric)
−∞

√

R(− 2t) =

√

2 exp(t2 )

Z t

exp(−x2 ) dx

−∞

√

√

π
2 exp(t ) ·
(1 + erf(t))
2
r [U+0012 control character]
[U+0013 control character]
π 1 + erf(t)
=
2 exp(−t2 )
r
√
π
R(− 2t) =
g(t)
2
=

2

d 1
Sampford [21, Eq. (3)] proved that dt
R(t) < 1 for any t. Thus,

d 1
d
1
q
=
=
dt g(t)
dt 2 R(−√2t)
π

r

π d
1
√
>
2 dt R(− 2t)

r

√
√
π
· 1 · − 2 = − π,
2

which proves (C.10) and completes the proof.
\end{ReviewVerbatim}


\paragraph{Expanded PaperInterface specification}
\begin{ReviewVerbatim}
∀ {theta xi xj a : ℝ},
  0 < theta →
    xj < xi →
      have laplace := MeasureTheory.volume.withDensity fun z => ENNReal.ofReal (√2 / 2 * Real.exp (-√2 * |z|));
      have gaussian := ProbabilityTheory.gaussianReal 0 1;
      have laplaceDenominator :=
        ((laplace.prod laplace) {epsilon | xi + epsilon.1 / theta < a ∧ xj + epsilon.2 / theta < a}).toReal;
      have gaussianDenominator :=
        ((gaussian.prod gaussian) {epsilon | xi + epsilon.1 / theta < a ∧ xj + epsilon.2 / theta < a}).toReal;
      (0 < laplaceDenominator ∧
          ((laplace.prod laplace)
                  {epsilon |
                    xi + epsilon.1 / theta < a ∧
                      xj + epsilon.2 / theta < a ∧ xj + epsilon.2 / theta < xi + epsilon.1 / theta}).toReal /
              laplaceDenominator <
            ((laplace.prod laplace) {epsilon | xj + epsilon.2 / theta < xi + epsilon.1 / theta}).toReal) ∧
        0 < gaussianDenominator ∧
          ((gaussian.prod gaussian)
                  {epsilon |
                    xi + epsilon.1 / theta < a ∧
                      xj + epsilon.2 / theta < a ∧ xj + epsilon.2 / theta < xi + epsilon.1 / theta}).toReal /
              gaussianDenominator <
            ((gaussian.prod gaussian) {epsilon | xj + epsilon.2 / theta < xi + epsilon.1 / theta}).toReal
\end{ReviewVerbatim}

\paragraph{Lean proof endpoint (built separately)}\begin{ReviewVerbatim}
KR21Monoculture.PaperInterface.appendixC1_source_pairwise_full_events
\end{ReviewVerbatim}

\paragraph{Recorded source-to-Spec screening}
\textbf{Verdict:} \texttt{not recorded}
\\
\textbf{Reason:} not recorded
\reviewmatch{Matches source input}
\noindent\textbf{Reviewer annotation}\par
\reviewerbox
\clearpage
\hypertarget{source-claim-6a2a92d325b21a46}{}
\section*{Review row 28}
\textbf{Source locator:} {\footnotesize\raggedright cited publication, lines 1139-\allowbreak{}1147\par}
\paragraph{Verbatim source input}
\begin{ReviewVerbatim}
(C.1)

Since Pr[Xi > Xj ] = lima→∞ Pr[Xi > Xj | Xi < a, Xj < a], it suffices to show that for all a,
d
Pr[Xi > Xj | Xi < a, Xj < a] ≥ 0,
(C.2)
da
and that it is strictly positive for some a. In other words, the higher a is, the more likely i and j are to be
correctly ordered. In Theorems 7 and 8, we show that (C.2) holds for both Laplacian and Gaussian noise

[Next verbatim source excerpt]

Finally, we show that (C.2) holds for both Laplacian and Gaussian noise.
Theorem 7. For any a ∈ R and Xi = xi + σεi where εi is Laplacian with unit variance,
d
Pr[Xi > Xj | Xi < a, Xj < a] ≥ 0.
da
Moreover, it is strictly positive for some a.
Proof. First, we must derive an expression for Pr[Xi > Xj | Xi < a, Xj < a]. Recall that the Laplace
distribution parameterized by µ and λ has pdf
f (x; µ, λ) =
and cdf

λ
exp(−λ|x − µ|)
2

(
F (x; µ, λ) =

1
2 exp (−λ(µ − x))
1 − 12 exp (−λ(x − µ))

21

x<µ
x≥µ


[form-feed in source extraction]
Note that xi and xj be the respective means of Xi and Xj , with xi > xj . Because the Laplace distribution
is piecewise defined, we must consider 3 cases and show that in all 3 cases, (C.2) holds. Note that
Ra
f (x; xi , λ)F (x; xj , λ) dx
(C.3)
Pr[Xi > Xj | Xi < a, Xj < a] = −∞
F (a; xi , λ)F (a; xj , λ)
Case 1: a ≤ xj .
Then, the numerator of (C.3) is
Z a
Z
λ
1
λ a
exp(−λ(xi + xj − 2x)) dx
exp (−λ(xi − x)) · exp(−λ(xj − x)) dx =
2
4 −∞
−∞ 2
Z
λ exp(−λ(xi + xj )) a
exp(2λx) dx
=
4
−∞
λ exp(−λ(xi + xj )) 1
=
exp(2λa)
4
2λ
exp(−λ(xi + xj − 2a))
=
8
The denominator is
1
1
1
exp(−λ(xi − a)) · exp(−λ(xj − a)) = exp(−λ(xi + xj − 2a)).
2
2
4
Thus,
Pr[Xi > Xj | Xi < a, Xj < a] =

1
,
2

so its derivative is trivially nonnegative.
Case 2: xj < a ≤ xi .
Then, the numerator of (C.3) is
[U+0012 control character]
[U+0013 control character]
Z xj
Z a
λ
1
λ
1
exp (−λ(xi − x)) · exp(−λ(xj − x)) dx +
exp (−λ(xi − x)) 1 − exp(−λ(x − xj )) dx
2
2
−∞ 2
xj 2
Z a
Z a
exp(−λ(xi − xj )) λ
λ
+
exp(−λ(xi − x)) dx −
exp(−λ(xi − xj )) dx
=
8
2 xj
4 xj
exp(−λ(xi − xj )) λ 1
λ
+
(exp(−λ(xi − a)) − exp(−λ(xi − xj ))) − (a − xj ) exp(−λ(xi − xj ))
8
2[U+0012 control character]λ
4
[U+0013 control character]
1
3 λ
= exp(−λ(xi − a)) −
+ (a − xj ) exp(−λ(xi − xj ))
2
8
4

=

The denominator is
[U+0012 control character]
[U+0013 control character]
1
1
1
1
1 − exp(−λ(a − xj )) · exp(λ(xj − a)) = exp(−λ(xi − a)) − exp(−λ(xi − xj ))
2
2
2
4
We can factor out 41 exp(−λ(xi − xj )) from both, so
2 exp(λ(a − xj )) − 23 + λ(a − xj )
Pr[Xi > Xj | Xi < a, Xj < a] =
2 exp(λ(a − xj )) − 1
=

[U+0001 control character]

[U+0001 control character]
2 exp(λ(a − xj )) − 1 − 21 + λ(a − xj )
2 exp(λ(a − xj )) − 1

=1−

1
2 + λ(a − xj )

2 exp(λ(a − xj )) − 1

22


[form-feed in source extraction]
Thus,
d
Pr[Xi > Xj | Xi < a, Xj < a] > 0
da
1
d
2 + λ(a − xj )
⇐⇒
<0
da 2 exp(λ(a − xj )) − 1
[U+0012 control character]
[U+0013 control character]
1
⇐⇒(2 exp(λ(a − xj )) − 1)λ <
+ λ(a − xj ) 2λ exp(λ(a − xj ))
2
[U+0012 control character]
[U+0013 control character]
1
⇐⇒2 − exp(−λ(a − xj )) < 2
+ λ(a − xj )
2
⇐⇒1 − exp(−λ(a − xj )) < 2λ(a − xj )
⇐⇒ exp(−λ(a − xj )) > 1 − 2λ(a − xj )
This is true because λ(a − xj ) > 0, and for z > 0,
exp(−z) > 1 − z > 1 − 2z.
Case 3: a > xi .
Then, the numerator of (C.3) is
[U+0012 control character]
[U+0013 control character]
Z xj
Z xi
λ
1
λ
1
exp (−λ(xi − x)) · exp(−λ(xj − x)) dx +
exp (−λ(xi − x)) 1 − exp(−λ(x − xj )) dx
2
2
−∞ 2
xj 2
[U+0012 control character]
[U+0013 control character]
Z a
1
λ
exp(−λ(x − xi )) 1 − exp(−λ(x − xj )) dx
+
2
2
xi
[U+0012 control character]
[U+0013 control character]
Z
1
3 λ
1
λ a
= −
+ (xi − xj ) exp(−λ(xi − xj )) + (1 − exp(−λ(a − xi ))) −
exp(−λ(2x − xi − xj )) dx
2
8
4
2
4 xi
[U+0012 control character]
[U+0013 control character]
3 λ
1
=1−
+ (xi − xj ) exp(−λ(xi − xj )) − exp(−λ(a − xi ))
8
4
2
1
+ exp(λ(xi + xj ))(exp(−2λa) − exp(−2λxi ))
8 [U+0012 control character]
[U+0013 control character]
1
1
1 λ
+ (xi − xj ) exp(−λ(xi − xj )) − exp(−λ(a − xi )) + exp(−λ(2a − xi − xj ))
=1−
2
4
2
8
The denominator is
[U+0012 control character]
[U+0013 control character][U+0012 control character]
[U+0013 control character]
1
1
1 − exp(−λ(t − xi ))
1 − exp(−λ(t − xj ))
2
2
1
1
1
= 1 − exp(−λ(a − xi )) − exp(−λ(a − xj )) + exp(−λ(2a − xi − xj ))
2
2
4
Thus,
Pr[Xi > Xj | Xi < a, Xj < a]
[U+0001 control character]
1 − 21 + λ4 (xi − xj ) exp(−λ(xi − xj )) − 12 exp(−λ(a − xi )) + 81 exp(−λ(2a − xi − xj ))
=
1 − 12 exp(−λ(a − xi )) − 21 exp(−λ(a − xj )) + 14 exp(−λ(2a − xi − xj ))
8 − (4 + 2λ(xi − xj )) exp(−λ(xi − xj )) − 4 exp(−λ(a − xi )) + exp(−λ(2a − xi − xj ))
∝
4 − 2 exp(−λ(a − xi )) − 2 exp(−λ(a − xj )) + exp(−λ(2a − xi − xj ))

23


[form-feed in source extraction]
We’re interested in
d
Pr[Xi > Xj | Xi < a, Xj < a] > 0
da
⇐⇒ (4 − 2 exp(−λ(a − xi )) − 2 exp(−λ(a − xj )) + exp(−λ(2a − xi − xj )
· (4λ exp(−λ(a − xi )) − 2λ exp(−λ(2a − xi − xj )))
> (8 − 4 exp(−λ(a − xi )) + exp(−λ(2a − xi − xj )) − (4 + 2λ(xi − xj )) exp(−λ(xi − xj )))
· (2λ exp(−λ(a − xi )) + 2λ exp(−λ(a − xj )) − 2λ exp(−λ(2a − xi − xj )))
⇐⇒16 exp(−λ(a − xi )) − 8 exp(−λ(2a − xi − xj )) − 8 exp(−2λ(a − xi )) + 4 exp(−λ(3a − 2xi − xj ))
− 8 exp(−λ(2a − xi − xj )) + 4 exp(−λ(3a − xi − 2xj )) + 4 exp(−λ(3a − 2xi − xj ))
− 2 exp(−2λ(2a − xi − xj ))
> 16 exp(−λ(a − xi )) + 16 exp(−λ(a − xj )) − 16 exp(−λ(2a − xi − xj ))
− 8 exp(−2λ(a − xi )) − 8 exp(−λ(2a − xi − xj )) + 8 exp(−λ(3a − 2xi − xj ))
+ 2 exp(−λ(3a − 2xi − xj )) + 2 exp(−λ(3a − xi − 2xj )) − 2 exp(−2λ(2a − xi − xj ))
− 2(4 + 2λ(xi − xj )) exp(−λ(a − xj )) − 2(4 + 2λ(xi − xj )) exp(−λ(a + xi − 2xj ))
+ 2(4 + 2λ(xi − xj )) exp(−2λ(a − xj ))
⇐⇒ exp(−λ(3a − xi − 2xj ))
> 8 exp(−λ(a − xj )) − 4 exp(−λ(2a − xi − xj )) + exp(−λ(3a − 2xi − xj ))
− (4 + 2λ(xi − xj )) exp(−λ(a − xj )) − (4 + 2λ(xi − xj )) exp(−λ(a + xi − 2xj ))
+ (4 + 2λ(xi − xj )) exp(−2λ(a − xj ))
⇐⇒ exp(−λ(2a − xi − xj ))
> 8 − 4 exp(−λ(a − xi )) + exp(−λ(2a − 2xi ))
− (4 + 2λ(xi − xj )) − (4 + 2λ(xi − xj )) exp(−λ(xi − xj )) + (4 + 2λ(xi − xj )) exp(−λ(a − xj ))
⇐⇒ exp(−λ(2a − xi − xj )) − 8 + 4 exp(−λ(a − xi )) − exp(−2λ(a − xi ))
+ (4 + 2λ(xi − xj ))(1 + exp(−λ(xi − xj ))) − (4 + 2λ(xi − xj )) exp(−λ(a − xj ))
>0

(C.4)

Note that for any z ≥ 0, we have
(4 + 2z)(1 + e−z ) − 8 ≥ 0 ⇐⇒ (2 + z)(1 + e−z ) ≥ 4
⇐⇒ z + 2e−z + ze−z ≥ 2
For z = 0, this holds with equality, and the left hand side is increasing since
d
z + 2e−z + ze−z ≥ 0 ⇐⇒ 1 − 2e−z + e−z − ze−z ≥ 0
dx
⇐⇒ 1 ≥ e−z + ze−z
1
⇐⇒
≥ e−z
1+z
⇐⇒ 1 + z ≤ ez
Therefore, choosing z = λ(xi − xj ) and plugging back to (C.4), we have
exp(−λ(2a − xi − xj )) − 8 + 4 exp(−λ(a − xi )) − exp(−2λ(a − xi ))
+ (4 + 2λ(xi − xj ))(1 + exp(−λ(xi − xj ))) − (4 + 2λ(xi − xj )) exp(−λ(a − xj )) > 0
⇐= exp(−λ(2a − xi − xj )) + 4 exp(−λ(a − xi )) − exp(−2λ(a − xi )) − (4 + 2λ(xi − xj )) exp(−λ(a − xj )) > 0
⇐⇒ exp(−λ(a − xj )) + 4 − exp(−λ(a − xi )) − (4 + 2λ(xi − xj )) exp(−λ(xi − xj )) > 0
⇐⇒4(1 − exp(−λ(xi − xj ))) + exp(−λ(a − xi ))(exp(−λ(xi − xj )) − 1) − 2λ(xi − xj ) exp(−λ(xi − xj )) > 0
⇐⇒(4 − exp(−λ(a − xi )))(1 − exp(−λ(xi − xj ))) − 2λ(xi − xj ) exp(−λ(xi − xj )) > 0
⇐=3(1 − exp(−λ(xi − xj ))) − 2λ(xi − xj ) exp(−λ(xi − xj )) > 0
24

(exp(−λ(a − xi )) < 1)


[form-feed in source extraction]
Again letting z = λ(xi − xj ), this is true if and only if
3(1 − e−z ) > 2ze−z ⇐⇒ 3(ez − 1) > 2z
⇐⇒ 3ez > 3 + 2z
which is true because ez > 1 + z for z > 0. This completes the proof for Case 3.
As a result, we have that
d
Pr[Xi > Xj | Xi < a, Xj < a] ≥ 0
da
for all a, with strict inequality for some a, which proves the theorem.

[Next verbatim source excerpt]

Theorem 8. For any a ∈ R and Xi = xi + σεi where εi ∼ N (0, 1),
d
Pr[Xi > Xj | Xi < a, Xj < a] > 0.
da
√
0
Proof. Assume σ = 1/ 2. This
√ is without loss of generality because for any instance with arbitrary σ , there
is an instance with σ = 1/ 2 that yields the same distribution over rankings (simply by scaling all item
values by σ/σ 0 ). First, we have
Ra
Pr[Xi = x] Pr[Xj < x] dx
Pr[Xi > Xj | Xi < a, Xj < a] = −∞
P r[Xi < a] Pr[Xj < a]
Ra
√
exp(−(x − xi )2 )/ π · (1 + erf(x − xj ))/2 dx
−∞
=
(1 + erf(a − xi ))/2 · (1 + erf(a − xj ))/2
Ra
exp(−(x − xi )2 )(1 + erf(x − xj )) dx
2
= √ −∞
(1 + erf(a − xi )) · (1 + erf(a − xj ))
π
The derivative with respect to a is positive if and only if
(1 + erf(a − xi ))(1 + erf(a − xj )) exp(−(a − xi )2 )(1 + erf(a − xj ))
Z a
>
exp(−(x − xi )2 )(1 + erf(x − xj )) dx
−∞

[U+0001 control character]
2
· √ (1 + erf(a − xi )) exp(−(a − xj )2 ) + (1 + erf(a − xj )) exp(−(a − xi )2 )
π

(C.5)

Let t = a − xi and δ = xi − xj . Then, using the fact that
Z a

Z a−xi

2

exp(−(x − xi ) )(1 + erf(x − xj )) dx =
−∞

2

Z a−xi

exp(−x ) dx +
−∞

√
=

π
(1 + erf(a − xi )) +
2

exp(−x2 ) erf(x + δ) dx

−∞
Z a−xi

exp(−x2 ) erf(x + δ) dx,

−∞

(C.5) becomes
√
√
Z t
π
(1 + erf(t))(1 + erf(t + δ))2 exp(−t2 )
π
·
>
(1
+
erf(t))
+
exp(−x2 ) erf(x + δ) dx
2 (1 + erf(t)) exp(−(t + δ)2 ) + (1 + erf(t + δ)) exp(−t2 )
2
−∞
Z t
(1 + erf(t))(1 + erf(t + δ))2 exp(−t2 )
2
√
⇐⇒
−
(1
+
erf(t))
−
exp(−x2 ) erf(x + δ) dx > 0
(1 + erf(t)) exp(−(t + δ)2 ) + (1 + erf(t + δ)) exp(−t2 )
π −∞
(C.6)
To show that this is true, we will use the fact that f (t) > 0 whenever the following conditions are met:
1. f (t) is continuous and differentiable everywhere
25


[form-feed in source extraction]
2. limt→−∞ f (t) = 0
d
3. dt
f (t) > 0

We’ll show that these conditions hold for the LHS of (C.6).
(1 + erf(t))(1 + erf(t + δ))2 exp(−t2 )
2
− (1 + erf(t)) − √
t→−∞ (1 + erf(t)) exp(−(t + δ)2 ) + (1 + erf(t + δ)) exp(−t2 )
π

Z t

lim

(1 + erf(t))(1 + erf(t + δ))2 exp(−t2 )
t→−∞ (1 + erf(t)) exp(−(t + δ)2 ) + (1 + erf(t + δ)) exp(−t2 )

= lim

exp(−x2 ) erf(x + δ) dx

−∞

(C.7)

Observe that both the numerator and denominator of (C.7) are positive, so this limit must be at least 0.
We can upper bound it by
(1 + erf(t))(1 + erf(t + δ))2 exp(−t2 )
(1 + erf(t))(1 + erf(t + δ))2 exp(−t2 )
≤
lim
t→−∞
t→−∞ (1 + erf(t)) exp(−(t + δ)2 ) + (1 + erf(t + δ)) exp(−t2 )
(1 + erf(t + δ)) exp(−t2 )
= lim (1 + erf(t))(1 + erf(t + δ))
lim

t→−∞

=0
Thus, the limit is 0. Now, we must show that the derivative is positive. The derivative is
[U+0015 control character]
[U+0014 control character]
Z t
2
(1 + erf(t))(1 + erf(t + δ))2 exp(−t2 )
d
2
− (1 + erf(t)) − √
exp(−x ) erf(x + δ) dx
dt (1 + erf(t)) exp(−(t + δ)2 ) + (1 + erf(t + δ)) exp(−t2 )
π −∞
[U+0014 control character]
[U+0015 control character]
d
(1 + erf(t))(1 + erf(t + δ))2 exp(−t2 )
2
2
=
− √ exp(−t2 ) − √ exp(−t2 ) erf(t + δ)
dt (1 + erf(t)) exp(−(t + δ)2 ) + (1 + erf(t + δ)) exp(−t2 )
π
π
(C.8)
Taking this derivative and factoring out
√

2(1 + erf(t))(1 + erf(t + δ)) exp(4t2 )
2

π (erf (t) + 1) et2 + (erf (δ + t) + 1) e(δ+t)

[U+0001 control character]2 ,

we get that (C.8) is positive if and only if
√
δ π exp((t + δ)2 )(1 + erf(t))(1 + erf(t + δ)) − exp(2δt + t2 )(1 + erf(t + δ)) + (1 + erf(t)) > 0
√
1 + erf(t)
⇐⇒δ π exp((t + δ)2 )(1 + erf(t)) +
− exp(2δt + t2 ) > 0
1 + erf(t + δ)
√
1 + erf(t)
⇐⇒δ π exp(t2 )(1 + erf(t)) + exp(−2δt − t2 )
−1>0
1 + erf(t + δ)
[U+0014 control character]
[U+0015 control character]
√
exp(−2δt − δ 2 )
⇐⇒(1 + erf(t)) δ π exp(t2 ) +
−1>0
1 + erf(t + δ)
[U+0015 control character]
[U+0014 control character]
1 + erf(t) √
exp(−(t + δ)2 )
⇐⇒
δ
π
+
−1>0
(C.9)
exp(−t2 )
1 + erf(t + δ)
Define
g(t) ,

1 + erf(t)
.
exp(−t2 )

Then, (C.9) is
[U+0014 control character]
√
g(t) δ π +

[U+0015 control character]
1
−1>0
g(t + δ)
√
1
1
⇐⇒
−
<δ π
g(t) g(t + δ)
26


[form-feed in source extraction]
By the Mean Value Theorem,
1
1
d 1
−
= −δ
g(t) g(t + δ)
dt g(t) t=t∗
for some t ≤ t∗ ≤ t + δ. Thus, it suffices to show that
√
d 1
>− π
dt g(t)

(C.10)

for all t. To do this, consider Mills Ratio [17]
Z ∞

2

R(t) , exp(t /2)

exp(−x2 /2) dx.

t

Note that this is quite similar in functional form to g(t), and with some manipulation, we can relate the two:
Z ∞
2
R(t) = exp(t /2)
exp(−x2 /2) dx
t
Z ∞
√
R( 2t) = exp(t2 ) √ exp(−x2 /2) dx
2t
Z ∞
√
2
= 2 exp(t )
exp(−x2 ) dx
t
Z −t
√
2
= 2 exp(t )
exp(−x2 ) dx
(exp(−x2 ) is symmetric)
−∞

√

R(− 2t) =

√

2 exp(t2 )

Z t

exp(−x2 ) dx

−∞

√

√

π
2 exp(t ) ·
(1 + erf(t))
2
r [U+0012 control character]
[U+0013 control character]
π 1 + erf(t)
=
2 exp(−t2 )
r
√
π
R(− 2t) =
g(t)
2
=

2

d 1
Sampford [21, Eq. (3)] proved that dt
R(t) < 1 for any t. Thus,

d 1
d
1
q
=
=
dt g(t)
dt 2 R(−√2t)
π

r

π d
1
√
>
2 dt R(− 2t)

r

√
√
π
· 1 · − 2 = − π,
2

which proves (C.10) and completes the proof.
\end{ReviewVerbatim}


\paragraph{Expanded PaperInterface specification}
\begin{ReviewVerbatim}
∀ {theta xi xj : ℝ},
  0 < theta →
    xj < xi →
      have laplace := MeasureTheory.volume.withDensity fun z => ENNReal.ofReal (√2 / 2 * Real.exp (-√2 * |z|));
      have gaussian := ProbabilityTheory.gaussianReal 0 1;
      have laplaceRatio := fun u =>
        ((laplace.prod laplace)
              {epsilon |
                xi + epsilon.1 / theta < u ∧
                  xj + epsilon.2 / theta < u ∧ xj + epsilon.2 / theta < xi + epsilon.1 / theta}).toReal /
          ((laplace.prod laplace) {epsilon | xi + epsilon.1 / theta < u ∧ xj + epsilon.2 / theta < u}).toReal;
      have gaussianRatio := fun u =>
        ((gaussian.prod gaussian)
              {epsilon |
                xi + epsilon.1 / theta < u ∧
                  xj + epsilon.2 / theta < u ∧ xj + epsilon.2 / theta < xi + epsilon.1 / theta}).toReal /
          ((gaussian.prod gaussian) {epsilon | xi + epsilon.1 / theta < u ∧ xj + epsilon.2 / theta < u}).toReal;
      ((∀ (a : ℝ),
            0 < ((laplace.prod laplace) {epsilon | xi + epsilon.1 / theta < a ∧ xj + epsilon.2 / theta < a}).toReal ∧
              ∃ d, HasDerivAt laplaceRatio d a ∧ 0 ≤ d) ∧
          ∃ a d, HasDerivAt laplaceRatio d a ∧ 0 < d) ∧
        ∀ (a : ℝ),
          0 < ((gaussian.prod gaussian) {epsilon | xi + epsilon.1 / theta < a ∧ xj + epsilon.2 / theta < a}).toReal ∧
            ∃ d, HasDerivAt gaussianRatio d a ∧ 0 < d
\end{ReviewVerbatim}

\paragraph{Lean proof endpoint (built separately)}\begin{ReviewVerbatim}
KR21Monoculture.PaperInterface.appendixC2_source_pairwise_full_events
\end{ReviewVerbatim}

\paragraph{Recorded source-to-Spec screening}
\textbf{Verdict:} \texttt{not recorded}
\\
\textbf{Reason:} not recorded
\reviewmatch{Matches source input}
\noindent\textbf{Reviewer annotation}\par
\reviewerbox
\clearpage
\hypertarget{source-claim-d11ceae944c696a3}{}
\section*{Review row 29}
\textbf{Source locator:} {\footnotesize\raggedright cited publication, lines 1151-\allowbreak{}1162\par}
\paragraph{Verbatim source input}
\begin{ReviewVerbatim}
C.2

Verifying Definition 3

Next, we show that for both Laplacian and Gaussian distributions, UAH (θA , θH ) < UHH (θA , θH ) for all
θA > θH . In fact, for 3-candidate RUM families, we will show that this is always true for any well-ordered
distribution, defined as follows.
Definition 4. A noise model with density f (·) is well-ordered if for any a > b and c > d,
f (a − c)f (b − d) > f (a − d)f (b − c).
In other words, for a well-ordered noise model, given two numbers, two candidates are more likely to be
correctly ordered than inverted conditioned on realizing those two numbers in some order. Lemma 1 shows
that both Gaussian and Laplacian distributions are well-ordered.
\end{ReviewVerbatim}


\paragraph{Expanded PaperInterface specification}
\begin{ReviewVerbatim}
∀ (f : ℝ → ℝ),
  KR21Monoculture.PaperInterface.strictlyWellOrderedNoise f ↔
    ∀ {{a b c d : ℝ}}, b < a → d < c → f (a - c) * f (b - d) > f (a - d) * f (b - c)
\end{ReviewVerbatim}

\paragraph{Lean proof endpoint (built separately)}\begin{ReviewVerbatim}
KR21Monoculture.PaperInterface.definition4_strictlyWellOrderedNoise_iff
\end{ReviewVerbatim}

\paragraph{Recorded source-to-Spec screening}
\textbf{Verdict:} \texttt{not recorded}
\\
\textbf{Reason:} not recorded
\reviewmatch{Matches source input}
\noindent\textbf{Reviewer annotation}\par
\reviewerbox
\clearpage
\hypertarget{source-claim-231c9ba63e6b1424}{}
\section*{Review row 30}
\textbf{Source locator:} {\footnotesize\raggedright cited publication, lines 217-\allowbreak{}242\par}
\paragraph{Verbatim source input}
\begin{ReviewVerbatim}
Definition 3 (Preference for weaker competition.). A candidate distribution D and noisy permutation family
Fθ , exhibits a preference for weaker competition if the following holds: for all θ1 > θ2 , σ ∼ Fθ1 and
π, τ ∼ Fθ2 ,
h
i
h
i
(−{σ1 })

E π1

(−{τ1 })

< E π1

.

Intuitively, suppose we have a higher accuracy parameter θ1 and a lower accuracy parameter θ2 < θ1 ; we
draw a permutation π from Fθ2 ; and we then derive two permutations from π: π (−{σ1 }) obtained by deleting
the first-ranked element of a permutation σ drawn from the more accurate distribution Fθ1 , and π (−{τ1 })
obtained by deleting the first-ranked element of a permutation τ drawn from the less accurate distribution
Fθ2 .
Then the expected value of the first-ranked candidate in π (−{τ1 }) is strictly greater than the expected
value of the first-ranked candidate in π (−{σ1 }) — that is, when a random candidate is removed from π, the
best remaining candidate is better in expectation when the randomly removed candidate is chosen based on
a noisier ranking.

[Next verbatim source excerpt]

In Random Utility Models, the underlying candidate values xi are perturbed by independent and identically
distributed noise εi ∼ E, and the perturbed values are ranked to produce π. Originally conceived in the
psychology literature [23], this model has been well-studied over nearly a century, [7, 4, 9, 24, 16, 22],
including more recently in the computer science and machine learning literature [2, 1, 19, 25, 13].
First, we must define a family of RUMs that satisfies the conditions of Definition 1. Assume without loss of
generality that the noise distribution E has unit variance. Then, consider the family of RUMs parameterized
by θ in which candidates are ranked according to xi + εθi . By this definition, the standard deviation of the
noise for a particular value of θ is simply 1/θ. Intuitively, larger values of θ reduce the effect of the noise,
making the ranking more accurate. In Theorem 5 in Appendix A, we show as long as the noise distribution
E has positive support on (−∞, ∞), this definition of Fθ meets the differentiability, asymptotic optimality,
and monotonicity conditions in Definition 1. For distributions with finite support, many of our results can

[Next verbatim source excerpt]

Definition 4. A noise model with density f (·) is well-ordered if for any a > b and c > d,
f (a − c)f (b − d) > f (a − d)f (b − c).
In other words, for a well-ordered noise model, given two numbers, two candidates are more likely to be
correctly ordered than inverted conditioned on realizing those two numbers in some order. Lemma 1 shows
that both Gaussian and Laplacian distributions are well-ordered.
Thus, it suffices to show that for any 3-candidate RUM with a well-ordered noise model, UAH (θA , θH ) <
UHH (θA , θH ) when θA > θH .
Theorem 6. For 3 candidates with unique values x1 > x2 > x3 and well-ordered i.i.d. noise with support
(−∞, ∞), if θA > θH , then UAH (θA , θH ) < UHH (θA , θH ).
Proof. Define u−i to be the expected utility of the maximum element of the human-generated ranking when
i is not available. Because we’re in the 3-candidate setting, we have
u−1 = λ1 x2 + (1 − λ1 )x3
u−2 = λ2 x1 + (1 − λ2 )x3
u−3 = λ3 x1 + (1 − λ3 )x2
where 1/2 < λi < 1. This is because the noise has support everywhere, so it is impossible to correctly rank
any two candidates with probability 1, and any two candidates are more likely than not to be correctly
ordered:
[U+0014 control character]
[U+0015 control character]
i
hε
εj
εi − εj
1
i
−
> −δ = Pr[εi − εj ≥ 0] + Pr 0 >
> −δ >
Pr
θ
θ
θ
2
Note that λ2 > λ1 and λ2 > λ3 , since
λ2 = Pr[ε1 − ε3 > −θ(x1 − x3 )] > max {Pr[ε1 − ε3 > −θ(x2 − x3 )], Pr[ε1 − ε3 > −θ(x1 − x2 )]} = max{λ1 , λ3 }.
Let τ ∼ FθA and π ∼ FθH . With this, we can write
UAH (θA , θH ) =

3
X

Pr[τ1 = i]u−i

i=1

UHH (θA , θH ) =

3
X

Pr[π1 = i]u−i

i=1

Define
∆pi = Pr[τ1 = i] − Pr[π1 = i]
Using Lemmas 2, and 3, we have
∆p1 > 0

(By monotonicity of RUM families, see Appendix A)

∆p1 ≥ ∆p2
∆p3 ≤ 0
Also, ∆p1 + ∆p2 + ∆p3 = 0. We must show that
UAH (θA , θH ) − UHH (θA , θH ) =

3
X
i=1

18

∆pi u−i < 0.


[form-feed in source extraction]
We consider 2 cases.
Case 1: ∆p2 ≤ 0.
Then, ∆p1 = −(∆p2 + ∆p3 ). This yields
3
X

∆pi u−i = ∆p1 u−1 + ∆p2 u−2 + ∆p3 u−3

i=1

≤ ∆p1 u−1 − ∆p1 min(u−2 , u−3 )
= ∆p1 (λ1 x2 + (1 − λ1 )x3 − min {λ2 x1 + (1 − λ2 )x3 , λ3 x1 + (1 − λ3 )x2 })
≤ ∆p1 (λ1 x2 + (1 − λ1 )x3 − min {λ2 x1 + (1 − λ2 )x3 , x2 })
We can show that this is at most 0 regardless of which term attains the minimum. Because λ2 > λ1 ,
λ1 x2 + (1 − λ1 )x3 − λ2 x1 − (1 − λ2 )x3 = λ1 x2 + x3 − λ1 x3 − λ2 x1 − x3 + λ2 x3
= λ1 x2 − λ1 x3 − λ2 x1 + λ2 x3
= λ1 (x2 − x3 ) + λ2 (x3 − x1 )
< λ1 (x2 − x3 ) + λ1 (x3 − x1 )
= λ1 (x2 − x1 )
<0
For the second term, we have
λ1 x2 + (1 − λ1 )x3 − x2 = (1 − λ1 )(x3 − x2 ) < 0.
Thus,
3
X

∆pi u−i < 0.

i=1

Case 2: ∆p2 > 0. Note that u−1 < x2 < u−3 . Then, using ∆p3 = −(∆p1 + ∆p2 ),
3
X

∆pi u−i = ∆p1 u−1 + ∆p2 u−2 + ∆p3 u−3

i=1

= ∆p1 (u−1 − u−3 ) + ∆p2 (u−2 − u−3 )
≤ ∆p2 (u−1 − u−3 ) + ∆p2 (u−2 − u−3 )

(∆p1 ≥ ∆p2 and u−1 < u−3 )

= ∆p2 (u−1 + u−2 − 2u−3 )
≤ ∆p2 (x2 + x1 − 2(λ3 x1 + (1 − λ3 )x2 ))
[U+0013 control character][U+0013 control character]
[U+0012 control character]
[U+0012 control character]
1
1
x1 + x2
< ∆p2 x2 + x1 − 2
2
2
=0
Thus, UAH (θA , θH ) < UHH (θA , θH ).

[Next verbatim source excerpt]

Next, we show a few basic facts. Let fA (r) be the density function of the joint realization R =
[X1 , . . . , Xn ] = [x1 +ε1 /θA , . . . , xn +εn /θA ] under the algorithmic ranking and fH (r) be the similarly defined
density function under the human-generated ranking. Consider the “contraction” operation r0 = cont(r) such
that ri0 = xi + (ri − xi ) · θθH
. Essentially, the contraction defines a coupling between fA (·) and fH (·), since
A
for r0 = cont(r), fA (r0 ) dr0 = fH (r) dr. Let π(r) be the ranking induced by r. Note that contraction cannot
introduce any new inversions in π(r)—that is, if i is ranked above j in π(r) for i < j, then i is ranked above
j in π(cont(r)). Intuitively, this is because contraction pulls values closer to their means, and can therefore
only correct existing inversions, not introduce new ones. This fact will allow us to prove some useful lemmas.
\end{ReviewVerbatim}


\paragraph{Expanded PaperInterface specification}
\begin{ReviewVerbatim}
∀ (f : ℝ → ℝ) {thetaA thetaH x1 x2 x3 : ℝ},
  Measurable f →
    (∀ {{a b c d : ℝ}}, b < a → d < c → f (a - c) * f (b - d) > f (a - d) * f (b - c)) →
      (∀ (z : ℝ), 0 ≤ f z) →
        ∫⁻ (z : ℝ), ENNReal.ofReal (f z) = 1 →
          0 < thetaH →
            thetaH < thetaA →
              x2 < x1 →
                x3 < x2 →
                  MeasureTheory.IsProbabilityMeasure
                      (MeasureTheory.Measure.pi fun x =>
                        MeasureTheory.volume.withDensity fun z => ENNReal.ofReal (f z)) ∧
                    ∃ rankingLaw,
                      (∀ (theta : ℝ),
                          0 < theta →
                            (rankingLaw theta).toMeasure =
                              MeasureTheory.Measure.map
                                (fun epsilon =>
                                  EconCSLib.SocialChoice.Ranking.rankByScore fun i =>
                                    KR21Monoculture.threeCandidateValueProfile x1 x2 x3 i + epsilon i / theta)
                                (MeasureTheory.Measure.pi fun x =>
                                  MeasureTheory.volume.withDensity fun z => ENNReal.ofReal (f z))) ∧
                        EconCSLib.SocialChoice.Ranking.expectedSecondMoverIndependent (rankingLaw thetaH)
                            (rankingLaw thetaA) (KR21Monoculture.threeCandidateValueProfile x1 x2 x3) <
                          EconCSLib.SocialChoice.Ranking.expectedSecondMoverIndependent (rankingLaw thetaH)
                            (rankingLaw thetaH) (KR21Monoculture.threeCandidateValueProfile x1 x2 x3)
\end{ReviewVerbatim}

\paragraph{Lean proof endpoint (built separately)}\begin{ReviewVerbatim}
KR21Monoculture.PaperInterface.theorem6_source_raw_iid_density_literal_semantic_complete
\end{ReviewVerbatim}

\paragraph{Recorded source-to-Spec screening}
\textbf{Verdict:} \texttt{not recorded}
\\
\textbf{Reason:} not recorded
\reviewmatch{Matches source input}
\noindent\textbf{Reviewer annotation}\par
\reviewerbox
\clearpage
\hypertarget{source-claim-d7ed4659ca9b332e}{}
\section*{Review row 31}
\textbf{Source locator:} {\footnotesize\raggedright cited publication, lines 1293-\allowbreak{}1329\par}
\paragraph{Verbatim source input}
\begin{ReviewVerbatim}
Lemma 1. Both Gaussian and Laplacian distributions are well-ordered.
Proof. The Gaussian noise model is well-ordered:
1
exp(−(a − c)2 − (b − d)2 )
2σ 2 π
1
=
exp(−(a − d)2 − (b − c)2 − 2(ac + bd − ad − bc))
2σ 2 π
= f (a − d)f (b − c) exp(−2((a − b)(c − d)))

f (a − c)f (b − d) =

< f (a − d)f (b − c)
19

(λ3 > 21 )


[form-feed in source extraction]
Laplacian noise is as well:
1
exp(−|a − c| − |b − d|)
4
1
f (a − d)f (b − c) = exp(−|a − d| − |b − c|)
4

f (a − c)f (b − d) =

It suffices to show that for a > b and c > d, |a − c| + |b − d| < |a − d| + |b − c|. To show this, plot
(a, b) and (c, d) in the (x, y) plane. Note that they’re both below the y = x line, and that the `1 distance
between them is |a − c| + |b − d|. Moreover, the `1 distance between any two points must be realized by
some Manhattan path, which is a combination of horizontal and vertical line segments. Consider the point
(b, a), which is above the y = x line. Any Manhattan path from (b, a) to (c, d) must cross the y = x line at
some point (w, w). Since (b, a) and (a, b) are equidistant from (w, w), for any Manhattan path from (b, a)
to (c, d), there exists a Manhattan path from (a, b) to (c, d) passing through (w, w) of the same length,
meaning the `1 distance from (a, b) to (c, d) is smaller than the `1 distance from (b, a) to (c, d). As a result,
|a − c| + |b − d| < |a − d| + |b − c|.
\end{ReviewVerbatim}

\paragraph{Approved corrected review target}
The archival source above is not asserted equivalent to this replacement.
\begin{ReviewVerbatim}
For every kappa > 0 and lambda > 0, the Gaussian kernel G_kappa(r)=exp(-kappa*r^2) satisfies G_kappa(a-c)G_kappa(b-d) > G_kappa(a-d)G_kappa(b-c) for all a>b and c>d. The Laplace kernel L_lambda(r)=exp(-lambda*|r|) satisfies the corresponding >= inequality globally, does not satisfy it strictly for every such quadruple, and satisfies the strict inequality on the explicit overlap domain b<a, d<c, b<c, and d<a.
\end{ReviewVerbatim}

\textbf{Archival source anchor:} {\footnotesize\raggedright cited publication, lines 1293-\allowbreak{}1329\par}
\paragraph{Recorded basis}
User instruction 2026-\allowbreak{}07-\allowbreak{}25:\allowbreak{} assume whatever convention makes the result true, provided the convention is explicit and the archival claim is not represented as proved.\allowbreak{}
\paragraph{Expanded PaperInterface specification}
\begin{ReviewVerbatim}
∀ {kappa lam : ℝ},
  0 < kappa →
    0 < lam →
      KR21Monoculture.StrictlyWellOrderedNoise (KR21Monoculture.gaussianNoiseKernel kappa) ∧
        KR21Monoculture.WeaklyWellOrderedNoise (KR21Monoculture.laplacianNoiseKernel lam) ∧
          ¬KR21Monoculture.StrictlyWellOrderedNoise (KR21Monoculture.laplacianNoiseKernel lam) ∧
            ∀ {a b c d : ℝ},
              b < a →
                d < c →
                  b < c →
                    d < a →
                      KR21Monoculture.laplacianNoiseKernel lam (a - c) *
                          KR21Monoculture.laplacianNoiseKernel lam (b - d) >
                        KR21Monoculture.laplacianNoiseKernel lam (a - d) *
                          KR21Monoculture.laplacianNoiseKernel lam (b - c)
\end{ReviewVerbatim}

\paragraph{Lean proof endpoint (built separately)}\begin{ReviewVerbatim}
KR21Monoculture.PaperInterface.lemma1_corrected_gaussian_laplace_kernel_target
\end{ReviewVerbatim}

\paragraph{Recorded source-to-Spec screening}
\textbf{Verdict:} \texttt{not recorded}
\\
\textbf{Reason:} not recorded
\reviewmatch{Matches approved corrected target}
\noindent\textbf{Reviewer annotation}\par
\reviewerbox
\clearpage
\hypertarget{source-claim-e077005933af7923}{}
\section*{Review row 32}
\textbf{Source locator:} {\footnotesize\raggedright cited publication, lines 1340-\allowbreak{}1364\par}
\paragraph{Verbatim source input}
\begin{ReviewVerbatim}
[No record available.]
\end{ReviewVerbatim}


\paragraph{Expanded PaperInterface specification}
\begin{ReviewVerbatim}
∀ {n : ℕ} {f : ℝ → ℝ},
  Measurable f →
    (∀ (z : ℝ), 0 ≤ f z) →
      ∀ (hnormalized : ∫⁻ (z : ℝ), ENNReal.ofReal (f z) = 1) (value : EconCSLib.SocialChoice.Ranking.Candidate n → ℝ),
        StrictAnti value →
          ∀ {thetaA thetaH : ℝ},
            0 < thetaH →
              thetaH < thetaA →
                (KR21Monoculture.paper_appendixA_scaledNoiseRankingPMF (KR21Monoculture.w11CandidateNoiseLaw f) value
                        thetaA).toMeasure =
                    MeasureTheory.Measure.map
                      (fun epsilon => EconCSLib.SocialChoice.Ranking.rankByScore fun c => value c + epsilon c / thetaA)
                      (MeasureTheory.Measure.pi fun x =>
                        MeasureTheory.volume.withDensity fun z => ENNReal.ofReal (f z)) ∧
                  (KR21Monoculture.paper_appendixA_scaledNoiseRankingPMF (KR21Monoculture.w11CandidateNoiseLaw f) value
                          thetaH).toMeasure =
                      MeasureTheory.Measure.map
                        (fun epsilon =>
                          EconCSLib.SocialChoice.Ranking.rankByScore fun c => value c + epsilon c / thetaH)
                        (MeasureTheory.Measure.pi fun x =>
                          MeasureTheory.volume.withDensity fun z => ENNReal.ofReal (f z)) ∧
                    KR21Monoculture.firstChoiceProb
                        (KR21Monoculture.paper_appendixA_scaledNoiseRankingPMF (KR21Monoculture.w11CandidateNoiseLaw f)
                          value thetaA)
                        (Fin.last (n + 1)) ≤
                      KR21Monoculture.firstChoiceProb
                        (KR21Monoculture.paper_appendixA_scaledNoiseRankingPMF (KR21Monoculture.w11CandidateNoiseLaw f)
                          value thetaH)
                        (Fin.last (n + 1))
\end{ReviewVerbatim}

\paragraph{Lean proof endpoint (built separately)}\begin{ReviewVerbatim}
KR21Monoculture.PaperInterface.lemma2_source_arbitraryFinite_iid_bottom_first_probability_semantic_complete
\end{ReviewVerbatim}

\paragraph{Recorded source-to-Spec screening}
\textbf{Verdict:} \texttt{not recorded}
\\
\textbf{Reason:} not recorded
\reviewmatch{Matches source input}
\noindent\textbf{Reviewer annotation}\par
\reviewerbox
\clearpage
\hypertarget{source-claim-a48d69bf76046da3}{}
\section*{Review row 33}
\textbf{Source locator:} {\footnotesize\raggedright cited publication, lines 1366-\allowbreak{}1464\par}
\paragraph{Verbatim source input}
\begin{ReviewVerbatim}
[No record available.]
\end{ReviewVerbatim}


\paragraph{Expanded PaperInterface specification}
\begin{ReviewVerbatim}
∀ {n : ℕ} {f : ℝ → ℝ},
  KR21Monoculture.StrictlyWellOrderedNoise f →
    Measurable f →
      (∀ (z : ℝ), 0 ≤ f z) →
        ∀ (hnormalized : ∫⁻ (z : ℝ), ENNReal.ofReal (f z) = 1) (value : EconCSLib.SocialChoice.Ranking.Candidate n → ℝ),
          StrictAnti value →
            ∀ {thetaA thetaH : ℝ},
              0 < thetaH →
                thetaH ≤ thetaA →
                  ∀ {i : EconCSLib.SocialChoice.Ranking.Candidate n},
                    i ≠ 0 →
                      (KR21Monoculture.paper_appendixA_scaledNoiseRankingPMF (KR21Monoculture.w11CandidateNoiseLaw f)
                              value thetaA).toMeasure =
                          MeasureTheory.Measure.map
                            (fun epsilon =>
                              EconCSLib.SocialChoice.Ranking.rankByScore fun c => value c + epsilon c / thetaA)
                            (MeasureTheory.Measure.pi fun x =>
                              MeasureTheory.volume.withDensity fun z => ENNReal.ofReal (f z)) ∧
                        (KR21Monoculture.paper_appendixA_scaledNoiseRankingPMF (KR21Monoculture.w11CandidateNoiseLaw f)
                                value thetaH).toMeasure =
                            MeasureTheory.Measure.map
                              (fun epsilon =>
                                EconCSLib.SocialChoice.Ranking.rankByScore fun c => value c + epsilon c / thetaH)
                              (MeasureTheory.Measure.pi fun x =>
                                MeasureTheory.volume.withDensity fun z => ENNReal.ofReal (f z)) ∧
                          KR21Monoculture.firstChoiceProb
                                (KR21Monoculture.paper_appendixA_scaledNoiseRankingPMF
                                  (KR21Monoculture.w11CandidateNoiseLaw f) value thetaA)
                                i -
                              KR21Monoculture.firstChoiceProb
                                (KR21Monoculture.paper_appendixA_scaledNoiseRankingPMF
                                  (KR21Monoculture.w11CandidateNoiseLaw f) value thetaH)
                                i ≤
                            KR21Monoculture.firstChoiceProb
                                (KR21Monoculture.paper_appendixA_scaledNoiseRankingPMF
                                  (KR21Monoculture.w11CandidateNoiseLaw f) value thetaA)
                                0 -
                              KR21Monoculture.firstChoiceProb
                                (KR21Monoculture.paper_appendixA_scaledNoiseRankingPMF
                                  (KR21Monoculture.w11CandidateNoiseLaw f) value thetaH)
                                0
\end{ReviewVerbatim}

\paragraph{Lean proof endpoint (built separately)}\begin{ReviewVerbatim}
KR21Monoculture.PaperInterface.lemma3_source_arbitraryFinite_iid_delta_le_top_delta_semantic_complete
\end{ReviewVerbatim}

\paragraph{Recorded source-to-Spec screening}
\textbf{Verdict:} \texttt{not recorded}
\\
\textbf{Reason:} not recorded
\reviewmatch{Matches source input}
\noindent\textbf{Reviewer annotation}\par
\reviewerbox
\clearpage
\hypertarget{source-claim-aceb5e75cb75d8b8}{}
\section*{Review row 34}
\textbf{Source locator:} {\footnotesize\raggedright cited publication, lines 1467-\allowbreak{}1800\par}
\paragraph{Verbatim source input}
\begin{ReviewVerbatim}
[No record available.]
\end{ReviewVerbatim}


\paragraph{Expanded PaperInterface specification}
\begin{ReviewVerbatim}
∀ {sigma xi xj : ℝ},
  0 < sigma →
    xj < xi →
      have innovation := MeasureTheory.volume.withDensity fun z => ENNReal.ofReal (√2 / 2 * Real.exp (-√2 * |z|));
      have score := fun epsilon c => if c = 0 then xi + sigma * epsilon.1 else xj + sigma * epsilon.2;
      have rank := fun epsilon => EconCSLib.SocialChoice.Ranking.rankByScore (score epsilon);
      have rankLaw := MeasureTheory.Measure.map rank (innovation.prod innovation);
      Measurable rank ∧
        rankLaw Set.univ = 1 ∧
          (∀ (epsilon : ℝ × ℝ),
              xj + sigma * epsilon.2 < xi + sigma * epsilon.1 →
                EconCSLib.SocialChoice.Ranking.firstChoice (rank epsilon) = 0) ∧
            (∀ (a : ℝ),
                0 <
                    ((innovation.prod innovation)
                        {epsilon | xi + sigma * epsilon.1 < a ∧ xj + sigma * epsilon.2 < a}).toReal ∧
                  ∃ d,
                    HasDerivAt
                        (fun u =>
                          ((innovation.prod innovation)
                                {epsilon |
                                  xi + sigma * epsilon.1 < u ∧
                                    xj + sigma * epsilon.2 < u ∧
                                      xj + sigma * epsilon.2 < xi + sigma * epsilon.1}).toReal /
                            ((innovation.prod innovation)
                                {epsilon | xi + sigma * epsilon.1 < u ∧ xj + sigma * epsilon.2 < u}).toReal)
                        d a ∧
                      0 ≤ d) ∧
              ∃ a d,
                HasDerivAt
                    (fun u =>
                      ((innovation.prod innovation)
                            {epsilon |
                              xi + sigma * epsilon.1 < u ∧
                                xj + sigma * epsilon.2 < u ∧ xj + sigma * epsilon.2 < xi + sigma * epsilon.1}).toReal /
                        ((innovation.prod innovation)
                            {epsilon | xi + sigma * epsilon.1 < u ∧ xj + sigma * epsilon.2 < u}).toReal)
                    d a ∧
                  0 < d
\end{ReviewVerbatim}

\paragraph{Lean proof endpoint (built separately)}\begin{ReviewVerbatim}
KR21Monoculture.PaperInterface.theorem7_source_semantic_complete
\end{ReviewVerbatim}

\paragraph{Recorded source-to-Spec screening}
\textbf{Verdict:} \texttt{not recorded}
\\
\textbf{Reason:} not recorded
\reviewmatch{Matches source input}
\noindent\textbf{Reviewer annotation}\par
\reviewerbox
\clearpage
\hypertarget{source-claim-fb611bb534f07b11}{}
\section*{Review row 35}
\textbf{Source locator:} {\footnotesize\raggedright cited publication, lines 1494-\allowbreak{}1764\par}
\paragraph{Verbatim source input}
\begin{ReviewVerbatim}
is piecewise defined, we must consider 3 cases and show that in all 3 cases, (C.2) holds. Note that
Ra
f (x; xi , λ)F (x; xj , λ) dx
(C.3)
Pr[Xi > Xj | Xi < a, Xj < a] = −∞
F (a; xi , λ)F (a; xj , λ)
Case 1: a ≤ xj .
Then, the numerator of (C.3) is
Z a
Z
λ
1
λ a
exp(−λ(xi + xj − 2x)) dx
exp (−λ(xi − x)) · exp(−λ(xj − x)) dx =
2
4 −∞
−∞ 2
Z
λ exp(−λ(xi + xj )) a
exp(2λx) dx
=
4
−∞
λ exp(−λ(xi + xj )) 1
=
exp(2λa)
4
2λ
exp(−λ(xi + xj − 2a))
=
8
The denominator is
1
1
1
exp(−λ(xi − a)) · exp(−λ(xj − a)) = exp(−λ(xi + xj − 2a)).
2
2
4
Thus,
Pr[Xi > Xj | Xi < a, Xj < a] =

1
,
2

so its derivative is trivially nonnegative.
Case 2: xj < a ≤ xi .
Then, the numerator of (C.3) is
[U+0012 control character]
[U+0013 control character]
Z xj
Z a
λ
1
λ
1
exp (−λ(xi − x)) · exp(−λ(xj − x)) dx +
exp (−λ(xi − x)) 1 − exp(−λ(x − xj )) dx
2
2
−∞ 2
xj 2
Z a
Z a
exp(−λ(xi − xj )) λ
λ
+
exp(−λ(xi − x)) dx −
exp(−λ(xi − xj )) dx
=
8
2 xj
4 xj
exp(−λ(xi − xj )) λ 1
λ
+
(exp(−λ(xi − a)) − exp(−λ(xi − xj ))) − (a − xj ) exp(−λ(xi − xj ))
8
2[U+0012 control character]λ
4
[U+0013 control character]
1
3 λ
= exp(−λ(xi − a)) −
+ (a − xj ) exp(−λ(xi − xj ))
2
8
4

=

The denominator is
[U+0012 control character]
[U+0013 control character]
1
1
1
1
1 − exp(−λ(a − xj )) · exp(λ(xj − a)) = exp(−λ(xi − a)) − exp(−λ(xi − xj ))
2
2
2
4
We can factor out 41 exp(−λ(xi − xj )) from both, so
2 exp(λ(a − xj )) − 23 + λ(a − xj )
Pr[Xi > Xj | Xi < a, Xj < a] =
2 exp(λ(a − xj )) − 1
=

[U+0001 control character]

[U+0001 control character]
2 exp(λ(a − xj )) − 1 − 21 + λ(a − xj )
2 exp(λ(a − xj )) − 1

=1−

1
2 + λ(a − xj )

2 exp(λ(a − xj )) − 1

22


[form-feed in source extraction]
Thus,
d
Pr[Xi > Xj | Xi < a, Xj < a] > 0
da
1
d
2 + λ(a − xj )
⇐⇒
<0
da 2 exp(λ(a − xj )) − 1
[U+0012 control character]
[U+0013 control character]
1
⇐⇒(2 exp(λ(a − xj )) − 1)λ <
+ λ(a − xj ) 2λ exp(λ(a − xj ))
2
[U+0012 control character]
[U+0013 control character]
1
⇐⇒2 − exp(−λ(a − xj )) < 2
+ λ(a − xj )
2
⇐⇒1 − exp(−λ(a − xj )) < 2λ(a − xj )
⇐⇒ exp(−λ(a − xj )) > 1 − 2λ(a − xj )
This is true because λ(a − xj ) > 0, and for z > 0,
exp(−z) > 1 − z > 1 − 2z.
Case 3: a > xi .
Then, the numerator of (C.3) is
[U+0012 control character]
[U+0013 control character]
Z xj
Z xi
λ
1
λ
1
exp (−λ(xi − x)) · exp(−λ(xj − x)) dx +
exp (−λ(xi − x)) 1 − exp(−λ(x − xj )) dx
2
2
−∞ 2
xj 2
[U+0012 control character]
[U+0013 control character]
Z a
1
λ
exp(−λ(x − xi )) 1 − exp(−λ(x − xj )) dx
+
2
2
xi
[U+0012 control character]
[U+0013 control character]
Z
1
3 λ
1
λ a
= −
+ (xi − xj ) exp(−λ(xi − xj )) + (1 − exp(−λ(a − xi ))) −
exp(−λ(2x − xi − xj )) dx
2
8
4
2
4 xi
[U+0012 control character]
[U+0013 control character]
3 λ
1
=1−
+ (xi − xj ) exp(−λ(xi − xj )) − exp(−λ(a − xi ))
8
4
2
1
+ exp(λ(xi + xj ))(exp(−2λa) − exp(−2λxi ))
8 [U+0012 control character]
[U+0013 control character]
1
1
1 λ
+ (xi − xj ) exp(−λ(xi − xj )) − exp(−λ(a − xi )) + exp(−λ(2a − xi − xj ))
=1−
2
4
2
8
The denominator is
[U+0012 control character]
[U+0013 control character][U+0012 control character]
[U+0013 control character]
1
1
1 − exp(−λ(t − xi ))
1 − exp(−λ(t − xj ))
2
2
1
1
1
= 1 − exp(−λ(a − xi )) − exp(−λ(a − xj )) + exp(−λ(2a − xi − xj ))
2
2
4
Thus,
Pr[Xi > Xj | Xi < a, Xj < a]
[U+0001 control character]
1 − 21 + λ4 (xi − xj ) exp(−λ(xi − xj )) − 12 exp(−λ(a − xi )) + 81 exp(−λ(2a − xi − xj ))
=
1 − 12 exp(−λ(a − xi )) − 21 exp(−λ(a − xj )) + 14 exp(−λ(2a − xi − xj ))
8 − (4 + 2λ(xi − xj )) exp(−λ(xi − xj )) − 4 exp(−λ(a − xi )) + exp(−λ(2a − xi − xj ))
∝
4 − 2 exp(−λ(a − xi )) − 2 exp(−λ(a − xj )) + exp(−λ(2a − xi − xj ))

23


[form-feed in source extraction]
We’re interested in
d
Pr[Xi > Xj | Xi < a, Xj < a] > 0
da
⇐⇒ (4 − 2 exp(−λ(a − xi )) − 2 exp(−λ(a − xj )) + exp(−λ(2a − xi − xj )
· (4λ exp(−λ(a − xi )) − 2λ exp(−λ(2a − xi − xj )))
> (8 − 4 exp(−λ(a − xi )) + exp(−λ(2a − xi − xj )) − (4 + 2λ(xi − xj )) exp(−λ(xi − xj )))
· (2λ exp(−λ(a − xi )) + 2λ exp(−λ(a − xj )) − 2λ exp(−λ(2a − xi − xj )))
⇐⇒16 exp(−λ(a − xi )) − 8 exp(−λ(2a − xi − xj )) − 8 exp(−2λ(a − xi )) + 4 exp(−λ(3a − 2xi − xj ))
− 8 exp(−λ(2a − xi − xj )) + 4 exp(−λ(3a − xi − 2xj )) + 4 exp(−λ(3a − 2xi − xj ))
− 2 exp(−2λ(2a − xi − xj ))
> 16 exp(−λ(a − xi )) + 16 exp(−λ(a − xj )) − 16 exp(−λ(2a − xi − xj ))
− 8 exp(−2λ(a − xi )) − 8 exp(−λ(2a − xi − xj )) + 8 exp(−λ(3a − 2xi − xj ))
+ 2 exp(−λ(3a − 2xi − xj )) + 2 exp(−λ(3a − xi − 2xj )) − 2 exp(−2λ(2a − xi − xj ))
− 2(4 + 2λ(xi − xj )) exp(−λ(a − xj )) − 2(4 + 2λ(xi − xj )) exp(−λ(a + xi − 2xj ))
+ 2(4 + 2λ(xi − xj )) exp(−2λ(a − xj ))
⇐⇒ exp(−λ(3a − xi − 2xj ))
> 8 exp(−λ(a − xj )) − 4 exp(−λ(2a − xi − xj )) + exp(−λ(3a − 2xi − xj ))
− (4 + 2λ(xi − xj )) exp(−λ(a − xj )) − (4 + 2λ(xi − xj )) exp(−λ(a + xi − 2xj ))
+ (4 + 2λ(xi − xj )) exp(−2λ(a − xj ))
⇐⇒ exp(−λ(2a − xi − xj ))
> 8 − 4 exp(−λ(a − xi )) + exp(−λ(2a − 2xi ))
− (4 + 2λ(xi − xj )) − (4 + 2λ(xi − xj )) exp(−λ(xi − xj )) + (4 + 2λ(xi − xj )) exp(−λ(a − xj ))
⇐⇒ exp(−λ(2a − xi − xj )) − 8 + 4 exp(−λ(a − xi )) − exp(−2λ(a − xi ))
+ (4 + 2λ(xi − xj ))(1 + exp(−λ(xi − xj ))) − (4 + 2λ(xi − xj )) exp(−λ(a − xj ))
>0
\end{ReviewVerbatim}


\paragraph{Expanded PaperInterface specification}
\begin{ReviewVerbatim}
∀ {lam xi xj a : ℝ},
  0 < lam →
    have density := fun location x => lam / 2 * Real.exp (-lam * |x - location|);
    have cdf := fun location x =>
      if x < location then 1 / 2 * Real.exp (-lam * (location - x)) else 1 - 1 / 2 * Real.exp (-lam * (x - location));
    have scoreI := MeasureTheory.volume.withDensity fun x => ENNReal.ofReal (density xi x);
    have scoreJ := MeasureTheory.volume.withDensity fun x => ENNReal.ofReal (density xj x);
    have pairLaw := scoreI.prod scoreJ;
    0 < (pairLaw {p | p.1 < a ∧ p.2 < a}).toReal ∧
      (pairLaw {p | p.1 < a ∧ p.2 < p.1}).toReal / (pairLaw {p | p.1 < a ∧ p.2 < a}).toReal =
        (∫ (x : ℝ) in Set.Iic a, density xi x * cdf xj x) / (cdf xi a * cdf xj a)
\end{ReviewVerbatim}

\paragraph{Lean proof endpoint (built separately)}\begin{ReviewVerbatim}
KR21Monoculture.PaperInterface.equationC3_laplace_strict_conditional_ratio_eq_density_cdf_semantic_complete
\end{ReviewVerbatim}

\paragraph{Recorded source-to-Spec screening}
\textbf{Verdict:} \texttt{not recorded}
\\
\textbf{Reason:} not recorded
\reviewmatch{Matches source input}
\noindent\textbf{Reviewer annotation}\par
\reviewerbox
\clearpage
\hypertarget{source-claim-29d2b7baa5996f3d}{}
\section*{Review row 36}
\textbf{Source locator:} {\footnotesize\raggedright cited publication, lines 1765-\allowbreak{}1800\par}
\paragraph{Verbatim source input}
\begin{ReviewVerbatim}
(C.4)

Note that for any z ≥ 0, we have
(4 + 2z)(1 + e−z ) − 8 ≥ 0 ⇐⇒ (2 + z)(1 + e−z ) ≥ 4
⇐⇒ z + 2e−z + ze−z ≥ 2
For z = 0, this holds with equality, and the left hand side is increasing since
d
z + 2e−z + ze−z ≥ 0 ⇐⇒ 1 − 2e−z + e−z − ze−z ≥ 0
dx
⇐⇒ 1 ≥ e−z + ze−z
1
⇐⇒
≥ e−z
1+z
⇐⇒ 1 + z ≤ ez
Therefore, choosing z = λ(xi − xj ) and plugging back to (C.4), we have
exp(−λ(2a − xi − xj )) − 8 + 4 exp(−λ(a − xi )) − exp(−2λ(a − xi ))
+ (4 + 2λ(xi − xj ))(1 + exp(−λ(xi − xj ))) − (4 + 2λ(xi − xj )) exp(−λ(a − xj )) > 0
⇐= exp(−λ(2a − xi − xj )) + 4 exp(−λ(a − xi )) − exp(−2λ(a − xi )) − (4 + 2λ(xi − xj )) exp(−λ(a − xj )) > 0
⇐⇒ exp(−λ(a − xj )) + 4 − exp(−λ(a − xi )) − (4 + 2λ(xi − xj )) exp(−λ(xi − xj )) > 0
⇐⇒4(1 − exp(−λ(xi − xj ))) + exp(−λ(a − xi ))(exp(−λ(xi − xj )) − 1) − 2λ(xi − xj ) exp(−λ(xi − xj )) > 0
⇐⇒(4 − exp(−λ(a − xi )))(1 − exp(−λ(xi − xj ))) − 2λ(xi − xj ) exp(−λ(xi − xj )) > 0
⇐=3(1 − exp(−λ(xi − xj ))) − 2λ(xi − xj ) exp(−λ(xi − xj )) > 0
24

(exp(−λ(a − xi )) < 1)


[form-feed in source extraction]
Again letting z = λ(xi − xj ), this is true if and only if
3(1 − e−z ) > 2ze−z ⇐⇒ 3(ez − 1) > 2z
⇐⇒ 3ez > 3 + 2z
which is true because ez > 1 + z for z > 0. This completes the proof for Case 3.
As a result, we have that
d
Pr[Xi > Xj | Xi < a, Xj < a] ≥ 0
da
for all a, with strict inequality for some a, which proves the theorem.
\end{ReviewVerbatim}


\paragraph{Expanded PaperInterface specification}
\begin{ReviewVerbatim}
∀ {lam xi xj a : ℝ},
  0 < lam →
    xj < xi →
      xi < a →
        0 <
          Real.exp (-lam * (2 * a - xi - xj)) - 8 + 4 * Real.exp (-lam * (a - xi)) - Real.exp (-2 * lam * (a - xi)) +
              (4 + 2 * lam * (xi - xj)) * (1 + Real.exp (-lam * (xi - xj))) -
            (4 + 2 * lam * (xi - xj)) * Real.exp (-lam * (a - xj))
\end{ReviewVerbatim}

\paragraph{Lean proof endpoint (built separately)}\begin{ReviewVerbatim}
KR21Monoculture.PaperInterface.equationC4_laplace_case3_expression_pos
\end{ReviewVerbatim}

\paragraph{Recorded source-to-Spec screening}
\textbf{Verdict:} \texttt{not recorded}
\\
\textbf{Reason:} not recorded
\reviewmatch{Matches source input}
\noindent\textbf{Reviewer annotation}\par
\reviewerbox
\clearpage
\hypertarget{source-claim-33bdbb546218e99a}{}
\section*{Review row 37}
\textbf{Source locator:} {\footnotesize\raggedright cited publication, lines 1801-\allowbreak{}2160\par}
\paragraph{Verbatim source input}
\begin{ReviewVerbatim}
[No record available.]
\end{ReviewVerbatim}


\paragraph{Expanded PaperInterface specification}
\begin{ReviewVerbatim}
∀ {sigma xi xj : ℝ},
  0 < sigma →
    xj < xi →
      have innovation := ProbabilityTheory.gaussianReal 0 1;
      have score := fun epsilon c => if c = 0 then xi + sigma * epsilon.1 else xj + sigma * epsilon.2;
      have rank := fun epsilon => EconCSLib.SocialChoice.Ranking.rankByScore (score epsilon);
      have rankLaw := MeasureTheory.Measure.map rank (innovation.prod innovation);
      Measurable rank ∧
        rankLaw Set.univ = 1 ∧
          (∀ (epsilon : ℝ × ℝ),
              xj + sigma * epsilon.2 < xi + sigma * epsilon.1 →
                EconCSLib.SocialChoice.Ranking.firstChoice (rank epsilon) = 0) ∧
            ∀ (a : ℝ),
              0 <
                  ((innovation.prod innovation)
                      {epsilon | xi + sigma * epsilon.1 < a ∧ xj + sigma * epsilon.2 < a}).toReal ∧
                ∃ d,
                  HasDerivAt
                      (fun u =>
                        ((innovation.prod innovation)
                              {epsilon |
                                xi + sigma * epsilon.1 < u ∧
                                  xj + sigma * epsilon.2 < u ∧
                                    xj + sigma * epsilon.2 < xi + sigma * epsilon.1}).toReal /
                          ((innovation.prod innovation)
                              {epsilon | xi + sigma * epsilon.1 < u ∧ xj + sigma * epsilon.2 < u}).toReal)
                      d a ∧
                    0 < d
\end{ReviewVerbatim}

\paragraph{Lean proof endpoint (built separately)}\begin{ReviewVerbatim}
KR21Monoculture.PaperInterface.theorem8_source_semantic_complete
\end{ReviewVerbatim}

\paragraph{Recorded source-to-Spec screening}
\textbf{Verdict:} \texttt{not recorded}
\\
\textbf{Reason:} not recorded
\reviewmatch{Matches source input}
\noindent\textbf{Reviewer annotation}\par
\reviewerbox
\clearpage
\hypertarget{source-claim-43364c7dea10aac0}{}
\section*{Review row 38}
\textbf{Source locator:} {\footnotesize\raggedright cited publication, lines 1839-\allowbreak{}1903\par}
\paragraph{Verbatim source input}
\begin{ReviewVerbatim}
(C.5)

Let t = a − xi and δ = xi − xj . Then, using the fact that
Z a

Z a−xi

2

exp(−(x − xi ) )(1 + erf(x − xj )) dx =
−∞

2

Z a−xi

exp(−x ) dx +
−∞

√
=

π
(1 + erf(a − xi )) +
2

exp(−x2 ) erf(x + δ) dx

−∞
Z a−xi

exp(−x2 ) erf(x + δ) dx,

−∞

(C.5) becomes
√
√
Z t
π
(1 + erf(t))(1 + erf(t + δ))2 exp(−t2 )
π
·
>
(1
+
erf(t))
+
exp(−x2 ) erf(x + δ) dx
2 (1 + erf(t)) exp(−(t + δ)2 ) + (1 + erf(t + δ)) exp(−t2 )
2
−∞
Z t
(1 + erf(t))(1 + erf(t + δ))2 exp(−t2 )
2
√
⇐⇒
−
(1
+
erf(t))
−
exp(−x2 ) erf(x + δ) dx > 0
(1 + erf(t)) exp(−(t + δ)2 ) + (1 + erf(t + δ)) exp(−t2 )
π −∞
\end{ReviewVerbatim}


\paragraph{Expanded PaperInterface specification}
\begin{ReviewVerbatim}
∀ {xi xj a : ℝ},
  xj < xi →
    have scoreI := ProbabilityTheory.gaussianReal xi (1 / 2);
    have scoreJ := ProbabilityTheory.gaussianReal xj (1 / 2);
    have denominator := ((scoreI.prod scoreJ) {p | p.1 < a ∧ p.2 < a}).toReal;
    0 < denominator ∧
      ((scoreI.prod scoreJ) {p | p.1 < a ∧ p.2 < a ∧ p.2 < p.1}).toReal / denominator =
          (2 / √Real.pi *
              ∫ (x : ℝ) in Set.Iic a, Real.exp (-(x - xi) ^ 2) * (1 + KR21Monoculture.theorem8Erf (x - xj))) /
            ((1 + KR21Monoculture.theorem8Erf (a - xi)) * (1 + KR21Monoculture.theorem8Erf (a - xj))) ∧
        0 <
          (1 + KR21Monoculture.theorem8Erf (a - xi)) * (1 + KR21Monoculture.theorem8Erf (a - xj)) *
                Real.exp (-(a - xi) ^ 2) *
              (1 + KR21Monoculture.theorem8Erf (a - xj)) -
            (∫ (x : ℝ) in Set.Iic a, Real.exp (-(x - xi) ^ 2) * (1 + KR21Monoculture.theorem8Erf (x - xj))) *
                (2 / √Real.pi) *
              ((1 + KR21Monoculture.theorem8Erf (a - xi)) * Real.exp (-(a - xj) ^ 2) +
                (1 + KR21Monoculture.theorem8Erf (a - xj)) * Real.exp (-(a - xi) ^ 2))
\end{ReviewVerbatim}

\paragraph{Lean proof endpoint (built separately)}\begin{ReviewVerbatim}
KR21Monoculture.PaperInterface.equationC5_gaussian_source_semantic_complete
\end{ReviewVerbatim}

\paragraph{Recorded source-to-Spec screening}
\textbf{Verdict:} \texttt{not recorded}
\\
\textbf{Reason:} not recorded
\reviewmatch{Matches source input}
\noindent\textbf{Reviewer annotation}\par
\reviewerbox
\clearpage
\hypertarget{source-claim-c25f59d1a0e8be35}{}
\section*{Review row 39}
\textbf{Source locator:} {\footnotesize\raggedright cited publication, lines 1874-\allowbreak{}1914\par}
\paragraph{Verbatim source input}
\begin{ReviewVerbatim}
(C.5) becomes
√
√
Z t
π
(1 + erf(t))(1 + erf(t + δ))2 exp(−t2 )
π
·
>
(1
+
erf(t))
+
exp(−x2 ) erf(x + δ) dx
2 (1 + erf(t)) exp(−(t + δ)2 ) + (1 + erf(t + δ)) exp(−t2 )
2
−∞
Z t
(1 + erf(t))(1 + erf(t + δ))2 exp(−t2 )
2
√
⇐⇒
−
(1
+
erf(t))
−
exp(−x2 ) erf(x + δ) dx > 0
(1 + erf(t)) exp(−(t + δ)2 ) + (1 + erf(t + δ)) exp(−t2 )
π −∞
(C.6)
To show that this is true, we will use the fact that f (t) > 0 whenever the following conditions are met:
1. f (t) is continuous and differentiable everywhere
25


[form-feed in source extraction]
2. limt→−∞ f (t) = 0
d
3. dt
f (t) > 0

We’ll show that these conditions hold for the LHS of (C.6).
\end{ReviewVerbatim}


\paragraph{Expanded PaperInterface specification}
\begin{ReviewVerbatim}
∀ {delta t : ℝ},
  0 < delta →
    0 <
      (1 + KR21Monoculture.theorem8Erf t) * (1 + KR21Monoculture.theorem8Erf (t + delta)) ^ 2 * Real.exp (-t ^ 2) /
            ((1 + KR21Monoculture.theorem8Erf t) * Real.exp (-(t + delta) ^ 2) +
              (1 + KR21Monoculture.theorem8Erf (t + delta)) * Real.exp (-t ^ 2)) -
          (1 + KR21Monoculture.theorem8Erf t) -
        2 / √Real.pi * ∫ (x : ℝ) in Set.Iic t, Real.exp (-x ^ 2) * KR21Monoculture.theorem8Erf (x + delta)
\end{ReviewVerbatim}

\paragraph{Lean proof endpoint (built separately)}\begin{ReviewVerbatim}
KR21Monoculture.PaperInterface.equationC6_gaussian_reduced_expression_pos
\end{ReviewVerbatim}

\paragraph{Recorded source-to-Spec screening}
\textbf{Verdict:} \texttt{not recorded}
\\
\textbf{Reason:} not recorded
\reviewmatch{Matches source input}
\noindent\textbf{Reviewer annotation}\par
\reviewerbox
\clearpage
\hypertarget{source-claim-4d39c9a7a45156fa}{}
\section*{Review row 40}
\textbf{Source locator:} {\footnotesize\raggedright cited publication, lines 1914-\allowbreak{}1951\par}
\paragraph{Verbatim source input}
\begin{ReviewVerbatim}
We’ll show that these conditions hold for the LHS of (C.6).
(1 + erf(t))(1 + erf(t + δ))2 exp(−t2 )
2
− (1 + erf(t)) − √
t→−∞ (1 + erf(t)) exp(−(t + δ)2 ) + (1 + erf(t + δ)) exp(−t2 )
π

Z t

lim

(1 + erf(t))(1 + erf(t + δ))2 exp(−t2 )
t→−∞ (1 + erf(t)) exp(−(t + δ)2 ) + (1 + erf(t + δ)) exp(−t2 )

= lim

exp(−x2 ) erf(x + δ) dx

−∞

(C.7)

Observe that both the numerator and denominator of (C.7) are positive, so this limit must be at least 0.
We can upper bound it by
(1 + erf(t))(1 + erf(t + δ))2 exp(−t2 )
(1 + erf(t))(1 + erf(t + δ))2 exp(−t2 )
≤
lim
t→−∞
t→−∞ (1 + erf(t)) exp(−(t + δ)2 ) + (1 + erf(t + δ)) exp(−t2 )
(1 + erf(t + δ)) exp(−t2 )
= lim (1 + erf(t))(1 + erf(t + δ))
lim

t→−∞

=0
Thus, the limit is 0. Now, we must show that the derivative is positive. The derivative is
\end{ReviewVerbatim}


\paragraph{Expanded PaperInterface specification}
\begin{ReviewVerbatim}
∀ (delta : ℝ),
  Filter.Tendsto
    (fun t =>
      (1 + KR21Monoculture.theorem8Erf t) * (1 + KR21Monoculture.theorem8Erf (t + delta)) ^ 2 * Real.exp (-t ^ 2) /
            ((1 + KR21Monoculture.theorem8Erf t) * Real.exp (-(t + delta) ^ 2) +
              (1 + KR21Monoculture.theorem8Erf (t + delta)) * Real.exp (-t ^ 2)) -
          (1 + KR21Monoculture.theorem8Erf t) -
        2 / √Real.pi * ∫ (x : ℝ) in Set.Iic t, Real.exp (-x ^ 2) * KR21Monoculture.theorem8Erf (x + delta))
    Filter.atBot (nhds 0)
\end{ReviewVerbatim}

\paragraph{Lean proof endpoint (built separately)}\begin{ReviewVerbatim}
KR21Monoculture.PaperInterface.equationC7_gaussian_reduced_expression_tendsto_atBot_zero
\end{ReviewVerbatim}

\paragraph{Recorded source-to-Spec screening}
\textbf{Verdict:} \texttt{not recorded}
\\
\textbf{Reason:} not recorded
\reviewmatch{Matches source input}
\noindent\textbf{Reviewer annotation}\par
\reviewerbox
\clearpage
\hypertarget{source-claim-2c9d403cd319b538}{}
\section*{Review row 41}
\textbf{Source locator:} {\footnotesize\raggedright cited publication, lines 1951-\allowbreak{}1985\par}
\paragraph{Verbatim source input}
\begin{ReviewVerbatim}
Thus, the limit is 0. Now, we must show that the derivative is positive. The derivative is
[U+0015 control character]
[U+0014 control character]
Z t
2
(1 + erf(t))(1 + erf(t + δ))2 exp(−t2 )
d
2
− (1 + erf(t)) − √
exp(−x ) erf(x + δ) dx
dt (1 + erf(t)) exp(−(t + δ)2 ) + (1 + erf(t + δ)) exp(−t2 )
π −∞
[U+0014 control character]
[U+0015 control character]
d
(1 + erf(t))(1 + erf(t + δ))2 exp(−t2 )
2
2
=
− √ exp(−t2 ) − √ exp(−t2 ) erf(t + δ)
dt (1 + erf(t)) exp(−(t + δ)2 ) + (1 + erf(t + δ)) exp(−t2 )
π
π
(C.8)
Taking this derivative and factoring out
√

2(1 + erf(t))(1 + erf(t + δ)) exp(4t2 )
2

π (erf (t) + 1) et2 + (erf (δ + t) + 1) e(δ+t)

[U+0001 control character]2 ,

we get that (C.8) is positive if and only if
\end{ReviewVerbatim}


\paragraph{Expanded PaperInterface specification}
\begin{ReviewVerbatim}
∀ (delta t : ℝ),
  HasDerivAt
    (fun u =>
      (1 + KR21Monoculture.theorem8Erf u) * (1 + KR21Monoculture.theorem8Erf (u + delta)) ^ 2 * Real.exp (-u ^ 2) /
            ((1 + KR21Monoculture.theorem8Erf u) * Real.exp (-(u + delta) ^ 2) +
              (1 + KR21Monoculture.theorem8Erf (u + delta)) * Real.exp (-u ^ 2)) -
          (1 + KR21Monoculture.theorem8Erf u) -
        2 / √Real.pi * ∫ (x : ℝ) in Set.Iic u, Real.exp (-x ^ 2) * KR21Monoculture.theorem8Erf (x + delta))
    (deriv
          (fun u =>
            (1 + KR21Monoculture.theorem8Erf u) * (1 + KR21Monoculture.theorem8Erf (u + delta)) ^ 2 *
                Real.exp (-u ^ 2) /
              ((1 + KR21Monoculture.theorem8Erf u) * Real.exp (-(u + delta) ^ 2) +
                (1 + KR21Monoculture.theorem8Erf (u + delta)) * Real.exp (-u ^ 2)))
          t -
        2 / √Real.pi * Real.exp (-t ^ 2) -
      2 / √Real.pi * Real.exp (-t ^ 2) * KR21Monoculture.theorem8Erf (t + delta))
    t
\end{ReviewVerbatim}

\paragraph{Lean proof endpoint (built separately)}\begin{ReviewVerbatim}
KR21Monoculture.PaperInterface.equationC8_gaussian_reduced_expression_hasDerivAt
\end{ReviewVerbatim}

\paragraph{Recorded source-to-Spec screening}
\textbf{Verdict:} \texttt{not recorded}
\\
\textbf{Reason:} not recorded
\reviewmatch{Matches source input}
\noindent\textbf{Reviewer annotation}\par
\reviewerbox
\clearpage
\hypertarget{source-claim-4139a609895b0de6}{}
\section*{Review row 42}
\textbf{Source locator:} {\footnotesize\raggedright cited publication, lines 1801-\allowbreak{}1914\par}
\paragraph{Verbatim source input}
\begin{ReviewVerbatim}
Theorem 8. For any a ∈ R and Xi = xi + σεi where εi ∼ N (0, 1),
d
Pr[Xi > Xj | Xi < a, Xj < a] > 0.
da
√
0
Proof. Assume σ = 1/ 2. This
√ is without loss of generality because for any instance with arbitrary σ , there
is an instance with σ = 1/ 2 that yields the same distribution over rankings (simply by scaling all item
values by σ/σ 0 ). First, we have
Ra
Pr[Xi = x] Pr[Xj < x] dx
Pr[Xi > Xj | Xi < a, Xj < a] = −∞
P r[Xi < a] Pr[Xj < a]
Ra
√
exp(−(x − xi )2 )/ π · (1 + erf(x − xj ))/2 dx
−∞
=
(1 + erf(a − xi ))/2 · (1 + erf(a − xj ))/2
Ra
exp(−(x − xi )2 )(1 + erf(x − xj )) dx
2
= √ −∞
(1 + erf(a − xi )) · (1 + erf(a − xj ))
π
The derivative with respect to a is positive if and only if
(1 + erf(a − xi ))(1 + erf(a − xj )) exp(−(a − xi )2 )(1 + erf(a − xj ))
Z a
>
exp(−(x − xi )2 )(1 + erf(x − xj )) dx
−∞

[U+0001 control character]
2
· √ (1 + erf(a − xi )) exp(−(a − xj )2 ) + (1 + erf(a − xj )) exp(−(a − xi )2 )
π

(C.5)

Let t = a − xi and δ = xi − xj . Then, using the fact that
Z a

Z a−xi

2

exp(−(x − xi ) )(1 + erf(x − xj )) dx =
−∞

2

Z a−xi

exp(−x ) dx +
−∞

√
=

π
(1 + erf(a − xi )) +
2

exp(−x2 ) erf(x + δ) dx

−∞
Z a−xi

exp(−x2 ) erf(x + δ) dx,

−∞

(C.5) becomes
√
√
Z t
π
(1 + erf(t))(1 + erf(t + δ))2 exp(−t2 )
π
·
>
(1
+
erf(t))
+
exp(−x2 ) erf(x + δ) dx
2 (1 + erf(t)) exp(−(t + δ)2 ) + (1 + erf(t + δ)) exp(−t2 )
2
−∞
Z t
(1 + erf(t))(1 + erf(t + δ))2 exp(−t2 )
2
√
⇐⇒
−
(1
+
erf(t))
−
exp(−x2 ) erf(x + δ) dx > 0
(1 + erf(t)) exp(−(t + δ)2 ) + (1 + erf(t + δ)) exp(−t2 )
π −∞
(C.6)
To show that this is true, we will use the fact that f (t) > 0 whenever the following conditions are met:
1. f (t) is continuous and differentiable everywhere
25


[form-feed in source extraction]
2. limt→−∞ f (t) = 0
d
3. dt
f (t) > 0

We’ll show that these conditions hold for the LHS of (C.6).

[Next verbatim source excerpt]

Thus, the limit is 0. Now, we must show that the derivative is positive. The derivative is
[U+0015 control character]
[U+0014 control character]
Z t
2
(1 + erf(t))(1 + erf(t + δ))2 exp(−t2 )
d
2
− (1 + erf(t)) − √
exp(−x ) erf(x + δ) dx
dt (1 + erf(t)) exp(−(t + δ)2 ) + (1 + erf(t + δ)) exp(−t2 )
π −∞
[U+0014 control character]
[U+0015 control character]
d
(1 + erf(t))(1 + erf(t + δ))2 exp(−t2 )
2
2
=
− √ exp(−t2 ) − √ exp(−t2 ) erf(t + δ)
dt (1 + erf(t)) exp(−(t + δ)2 ) + (1 + erf(t + δ)) exp(−t2 )
π
π
(C.8)
Taking this derivative and factoring out
√

2(1 + erf(t))(1 + erf(t + δ)) exp(4t2 )
2

π (erf (t) + 1) et2 + (erf (δ + t) + 1) e(δ+t)

[U+0001 control character]2 ,

we get that (C.8) is positive if and only if
√
δ π exp((t + δ)2 )(1 + erf(t))(1 + erf(t + δ)) − exp(2δt + t2 )(1 + erf(t + δ)) + (1 + erf(t)) > 0
√
1 + erf(t)
⇐⇒δ π exp((t + δ)2 )(1 + erf(t)) +
− exp(2δt + t2 ) > 0
1 + erf(t + δ)
√
1 + erf(t)
⇐⇒δ π exp(t2 )(1 + erf(t)) + exp(−2δt − t2 )
−1>0
1 + erf(t + δ)
[U+0014 control character]
[U+0015 control character]
√
exp(−2δt − δ 2 )
⇐⇒(1 + erf(t)) δ π exp(t2 ) +
−1>0
1 + erf(t + δ)
[U+0015 control character]
[U+0014 control character]
1 + erf(t) √
exp(−(t + δ)2 )
⇐⇒
δ
π
+
−1>0
(C.9)
exp(−t2 )
1 + erf(t + δ)
Define
g(t) ,

1 + erf(t)
.
exp(−t2 )

Then, (C.9) is
[U+0014 control character]
√
g(t) δ π +

[U+0015 control character]
1
−1>0
g(t + δ)
√
1
1
⇐⇒
−
<δ π
g(t) g(t + δ)
26


[form-feed in source extraction]
By the Mean Value Theorem,
1
1
d 1
−
= −δ
g(t) g(t + δ)
dt g(t) t=t∗
for some t ≤ t∗ ≤ t + δ. Thus, it suffices to show that
\end{ReviewVerbatim}

\paragraph{Approved corrected review target}
The archival source above is not asserted equivalent to this replacement.
\begin{ReviewVerbatim}
In the source normalization sigma = 1/sqrt(2), with t = a-x_i and delta = x_i-x_j > 0, let e(u)=erf(u) and H_delta(u) = ((1+e(u))(1+e(u+delta))^2 exp(-u^2))/((1+e(u))exp(-(u+delta)^2)+(1+e(u+delta))exp(-u^2)) - (1+e(u)) - (2/sqrt(pi))*integral_{x<=u} exp(-x^2)e(x+delta) dx. Then H_delta has derivative P_delta(t)*B_delta(t), where P_delta(t) = 2(1+e(t))(1+e(t+delta))^2 exp((t+delta)^2)/(sqrt(pi)*((e(t)+1)exp(t^2)+(e(t+delta)+1)exp((t+delta)^2))^2) and B_delta(t) = ((1+e(t))/exp(-t^2))*(delta*sqrt(pi) + ((1+e(t+delta))/exp(-(t+delta)^2))^(-1)) - 1. Both P_delta(t) and B_delta(t) are strictly positive. This is the corrected factorization and parenthesized C.9 bracket, not the printed prefactor or unparenthesized display.
\end{ReviewVerbatim}

\textbf{Archival source anchor:} {\footnotesize\raggedright cited publication, lines 1951-\allowbreak{}2039\par}
\paragraph{Recorded basis}
User instruction 2026-\allowbreak{}07-\allowbreak{}25:\allowbreak{} assume whatever convention makes the result true, provided the convention is explicit and the archival claim is not represented as proved.\allowbreak{}
\paragraph{Expanded PaperInterface specification}
\begin{ReviewVerbatim}
∀ {delta t : ℝ},
  0 < delta →
    HasDerivAt
        (fun u =>
          (1 + KR21Monoculture.theorem8Erf u) * (1 + KR21Monoculture.theorem8Erf (u + delta)) ^ 2 * Real.exp (-u ^ 2) /
                ((1 + KR21Monoculture.theorem8Erf u) * Real.exp (-(u + delta) ^ 2) +
                  (1 + KR21Monoculture.theorem8Erf (u + delta)) * Real.exp (-u ^ 2)) -
              (1 + KR21Monoculture.theorem8Erf u) -
            2 / √Real.pi * ∫ (x : ℝ) in Set.Iic u, Real.exp (-x ^ 2) * KR21Monoculture.theorem8Erf (x + delta))
        (2 * (1 + KR21Monoculture.theorem8Erf t) * (1 + KR21Monoculture.theorem8Erf (t + delta)) ^ 2 *
              Real.exp ((t + delta) ^ 2) /
            (√Real.pi *
              ((KR21Monoculture.theorem8Erf t + 1) * Real.exp (t ^ 2) +
                  (KR21Monoculture.theorem8Erf (t + delta) + 1) * Real.exp ((t + delta) ^ 2)) ^
                2) *
          ((1 + KR21Monoculture.theorem8Erf t) / Real.exp (-t ^ 2) *
              (delta * √Real.pi + ((1 + KR21Monoculture.theorem8Erf (t + delta)) / Real.exp (-(t + delta) ^ 2))⁻¹) -
            1))
        t ∧
      0 <
          2 * (1 + KR21Monoculture.theorem8Erf t) * (1 + KR21Monoculture.theorem8Erf (t + delta)) ^ 2 *
              Real.exp ((t + delta) ^ 2) /
            (√Real.pi *
              ((KR21Monoculture.theorem8Erf t + 1) * Real.exp (t ^ 2) +
                  (KR21Monoculture.theorem8Erf (t + delta) + 1) * Real.exp ((t + delta) ^ 2)) ^
                2) ∧
        0 <
          (1 + KR21Monoculture.theorem8Erf t) / Real.exp (-t ^ 2) *
              (delta * √Real.pi + ((1 + KR21Monoculture.theorem8Erf (t + delta)) / Real.exp (-(t + delta) ^ 2))⁻¹) -
            1
\end{ReviewVerbatim}

\paragraph{Lean proof endpoint (built separately)}\begin{ReviewVerbatim}
KR21Monoculture.PaperInterface.equationC8_C9_corrected_gaussian_target
\end{ReviewVerbatim}

\paragraph{Recorded source-to-Spec screening}
\textbf{Verdict:} \texttt{not recorded}
\\
\textbf{Reason:} not recorded
\reviewmatch{Matches approved corrected target}
\noindent\textbf{Reviewer annotation}\par
\reviewerbox
\clearpage
\hypertarget{source-claim-7a440699a84fb997}{}
\section*{Review row 43}
\textbf{Source locator:} {\footnotesize\raggedright cited publication, lines 2050-\allowbreak{}2160\par}
\paragraph{Verbatim source input}
\begin{ReviewVerbatim}
for some t ≤ t∗ ≤ t + δ. Thus, it suffices to show that
√
d 1
>− π
dt g(t)

(C.10)

for all t. To do this, consider Mills Ratio [17]
Z ∞

2

R(t) , exp(t /2)

exp(−x2 /2) dx.

t

Note that this is quite similar in functional form to g(t), and with some manipulation, we can relate the two:
Z ∞
2
R(t) = exp(t /2)
exp(−x2 /2) dx
t
Z ∞
√
R( 2t) = exp(t2 ) √ exp(−x2 /2) dx
2t
Z ∞
√
2
= 2 exp(t )
exp(−x2 ) dx
t
Z −t
√
2
= 2 exp(t )
exp(−x2 ) dx
(exp(−x2 ) is symmetric)
−∞

√

R(− 2t) =

√

2 exp(t2 )

Z t

exp(−x2 ) dx

−∞

√

√

π
2 exp(t ) ·
(1 + erf(t))
2
r [U+0012 control character]
[U+0013 control character]
π 1 + erf(t)
=
2 exp(−t2 )
r
√
π
R(− 2t) =
g(t)
2
=

2

d 1
Sampford [21, Eq. (3)] proved that dt
R(t) < 1 for any t. Thus,

d 1
d
1
q
=
=
dt g(t)
dt 2 R(−√2t)
π

r

π d
1
√
>
2 dt R(− 2t)

r

√
√
π
· 1 · − 2 = − π,
2

which proves (C.10) and completes the proof.
\end{ReviewVerbatim}


\paragraph{Expanded PaperInterface specification}
\begin{ReviewVerbatim}
∀ (t : ℝ), ∃ d, HasDerivAt (fun u => ((1 + KR21Monoculture.theorem8Erf u) / Real.exp (-u ^ 2))⁻¹) d t ∧ -√Real.pi < d
\end{ReviewVerbatim}

\paragraph{Lean proof endpoint (built separately)}\begin{ReviewVerbatim}
KR21Monoculture.PaperInterface.equationC10_gaussian_g_inv_derivative_lower_bound
\end{ReviewVerbatim}

\paragraph{Recorded source-to-Spec screening}
\textbf{Verdict:} \texttt{not recorded}
\\
\textbf{Reason:} not recorded
\reviewmatch{Matches source input}
\noindent\textbf{Reviewer annotation}\par
\reviewerbox
\clearpage
\hypertarget{source-claim-d732fbc16f2ecdeb}{}
\section*{Review row 44}
\textbf{Source locator:} {\footnotesize\raggedright cited publication, lines 133-\allowbreak{}138\par}
\paragraph{Verbatim source input}
\begin{ReviewVerbatim}
More formally, we model the n candidates as having intrinsic values x1 , . . . , xn , where any employer would
derive utility xi from hiring candidate i. Throughout the paper, we assume without loss of generality that
x1 > x2 > · · · > xn . These values, however, are unknown to the employer; instead, they must use some
noisy procedure to rank the candidates. We model such a procedure as a randomized mechanism R that
takes in the true candidate values and draws a permutation π over those candidates from some distribution.
Our main results hold for families of distributions over permutations as defined below:

[Next verbatim source excerpt]

Theorem 9. The family of distributions Fθ produced by the Mallows Model with Kendall tau distance with
θ = φ − 1 satisfies the conditions of Definition 1.
Proof. We must show that Fθ satisfies the differentiability, asymptotic optimality, and monotonicity conditions of Definition 1.
Differentiability: Let Π be the set of all permutations on n candidates. The probability of a realizing
a particular permutation π under the Mallows model is
∗

φ−d(π,π )
−d(π 0 ,π ∗ )
π 0 ∈Π φ

Pr[π] = P
θ

27


[form-feed in source extraction]
Both the numerator and denominator are differentiable with respect to θ = φ − 1, so Prθ [π] is differentiable
with respect to θ.
Asymptotic optimality: For the correct ranking π ∗ ,
1
,
Z

Pr[π ∗ ] =
θ

where the normalizing constant Z is
Z=

X

∗

φ−d(π,π )

π∈Π

In the limit,
lim Z = lim Z
φ→∞
X
∗
= lim
φ−d(π,π )

θ→∞

φ→∞

π∈Π

φ→∞

φ−d(π,π )

π6=π ∗ ∈Π
∗

X

=1+

∗

X

= lim 1 +

lim φ−d(π,π )

π6=π ∗ ∈Π

φ→∞

=1
because for any π 6= π ∗ , d(π, π ∗ ) ≥ 1. Therefore,
1
=1
θ→∞ Z

lim Pr[π ∗ ] = lim

θ→∞ θ

(−S)

Monotonicity: We must show that for any S ⊂ x, if π1
denotes the value of the top-ranked candidate
according to π excluding candidates in S,
h
i
h
i
(−S)
(−S)
EFθ0 π1
≥ EFθ π1
.
(−S)

= xj .

(−S)

= xj }

For any i ∈
/ S, let j be the smallest index such that j > i and j ∈
/ S. Consider any π such that π1
(−S)
Then, swapping i and j yields a permutation π̂ such that π̂1
= xi . Moreover,
Pr[π̂] = Pr[π] · φinv(π)−inv(π̂) .
Since i < j, inv(π) − inv(π̂) ≥ 1. Finally, note that swapping i and j is a bijection between {π : π1
(−S)
and {π : π1
= xi }. Thus,
(−S)

= xi ]
=
(−S)
Pr[π1
= xj ]
(−S)
π:π
=x
Pr[π1

1

Note that the terms

Pr[π]

X

j

(−S)
Pr[π1
= xj ]

· φinv(π)−inv(π̂)

Pr[π]
sum to 1, so this is sum is some polynomial in φ with nonnegative weights
(−S)
Pr[π1
=xj ]

and integer powers of φ. As a result, it must have a positive derivative with respect to φ, i.e., for i < j,
(−S)

= xi ]
d Pr[π1
>0
dφ Pr[π (−S) = xj ]
1

0

Let φ > φ. Then,
(−S)

Prφ [π1

= xi ]

(−S)
Prφ [π1
= xj ]

(−S)

<

28

Prφ0 [π1

= xi ]
(−S)
Prφ0 [π1
= xj ]


[form-feed in source extraction]
Rearranging,
(−S)

Prφ [π1

= xi ]

(−S)
Prφ0 [π1
= xi ]

(−S)

Prφ [π1

= xj ]
(−S)
Prφ0 [π1
= xj ]

<

(D.1)

For θ0 − φ0 − 1 and θ = φ − 1,
h
i X h
i
(−S)
−(S)
EFθ π1
=
Pr π1
= xi xi
φ

i∈S
/

h

(−S)

EFθ0 π1

i

=

X
i∈S
/

h
i
−(S)
Pr0 π1
= xi xi
φ

By Lemma 4,
h
i
h
i
(−S)
(−S)
EFθ0 π1
> EFθ π1
,
which completes the proof.
\end{ReviewVerbatim}


\paragraph{Expanded PaperInterface specification}
\begin{ReviewVerbatim}
∀ {n : ℕ} (center : EconCSLib.SocialChoice.Ranking.Ranking n) (value : EconCSLib.SocialChoice.Ranking.Candidate n → ℝ),
  EconCSLib.SocialChoice.Ranking.StrictlyOrderedBy center value →
    (∀ (phi theta : ℝ),
        1 < phi →
          theta = phi - 1 →
            have M := KR21Monoculture.concreteMallowsSpec center theta;
            M.q = phi⁻¹ ∧
              ∀ (pi : EconCSLib.SocialChoice.Ranking.Ranking n),
                (M.law pi).toReal =
                  phi⁻¹ ^ EconCSLib.SocialChoice.Ranking.kendallTau center pi /
                    ∑ tau, phi⁻¹ ^ EconCSLib.SocialChoice.Ranking.kendallTau center tau) ∧
      (∀ (theta : ℝ),
          0 < theta →
            ∀ (pi : EconCSLib.SocialChoice.Ranking.Ranking n),
              ContinuousAt (fun theta' => ((KR21Monoculture.concreteMallowsSpec center theta').law pi).toReal) theta ∧
                DifferentiableAt ℝ (fun theta' => ((KR21Monoculture.concreteMallowsSpec center theta').law pi).toReal)
                  theta) ∧
        Filter.Tendsto (fun theta => ((KR21Monoculture.concreteMallowsSpec center theta).law center).toReal)
            Filter.atTop (nhds 1) ∧
          ∀ (thetaA thetaH : ℝ),
            0 < thetaH →
              thetaH < thetaA →
                (∀ (remaining : Finset (EconCSLib.SocialChoice.Ranking.Candidate n)),
                    remaining.Nonempty →
                      KR21Monoculture.expectedBestInSet (KR21Monoculture.concreteMallowsSpec center thetaH).law value
                          remaining ≤
                        KR21Monoculture.expectedBestInSet (KR21Monoculture.concreteMallowsSpec center thetaA).law value
                          remaining) ∧
                  KR21Monoculture.expectedBestInSet (KR21Monoculture.concreteMallowsSpec center thetaH).law value
                      Finset.univ <
                    KR21Monoculture.expectedBestInSet (KR21Monoculture.concreteMallowsSpec center thetaA).law value
                      Finset.univ
\end{ReviewVerbatim}

\paragraph{Lean proof endpoint (built separately)}\begin{ReviewVerbatim}
KR21Monoculture.PaperInterface.theorem9_source_mallows_definition1_literal_semantic_complete
\end{ReviewVerbatim}

\paragraph{Recorded source-to-Spec screening}
\textbf{Verdict:} \texttt{not recorded}
\\
\textbf{Reason:} not recorded
\reviewmatch{Matches source input}
\noindent\textbf{Reviewer annotation}\par
\reviewerbox
\clearpage
\hypertarget{source-claim-7751231944a77b74}{}
\section*{Review row 45}
\textbf{Source locator:} {\footnotesize\raggedright cited publication, lines 133-\allowbreak{}138\par}
\paragraph{Verbatim source input}
\begin{ReviewVerbatim}
More formally, we model the n candidates as having intrinsic values x1 , . . . , xn , where any employer would
derive utility xi from hiring candidate i. Throughout the paper, we assume without loss of generality that
x1 > x2 > · · · > xn . These values, however, are unknown to the employer; instead, they must use some
noisy procedure to rank the candidates. We model such a procedure as a randomized mechanism R that
takes in the true candidate values and draws a permutation π over those candidates from some distribution.
Our main results hold for families of distributions over permutations as defined below:

[Next verbatim source excerpt]

The Mallows Model also appears frequently in the ranking literature [8, 11], and is much more analytically
tractable than RUMs. Under the Mallows Model, the likelihood of a permutation is related to its distance
from the true ranking π ∗ :
∗
1
(8)
Pr[π] = φ−d(π,π ) ,
Z
where Z is a normalizing constant. In this model, φ > 1 is the accuracy parameter: the larger φ is, the
more likely the ranking procedure is to output a ranking π that is close to the true ranking r. To instantiate
this model, we need a notion of distance d(·, ·) over permutations. For this, we’ll use Kendall tau distance,
another standard notion in the literature, which is simply the number of pairs of elements in π that are
incorrectly ordered [10]. In Appendix D, we verify that the family of distributions Fθ given by the Mallows
Model satisfies Definition 1, defining θ = φ − 1 (for consistency, so θ is well-defined on (0, ∞)).

[Next verbatim source excerpt]

Theorem 9. The family of distributions Fθ produced by the Mallows Model with Kendall tau distance with
θ = φ − 1 satisfies the conditions of Definition 1.
Proof. We must show that Fθ satisfies the differentiability, asymptotic optimality, and monotonicity conditions of Definition 1.
Differentiability: Let Π be the set of all permutations on n candidates. The probability of a realizing
a particular permutation π under the Mallows model is
∗

φ−d(π,π )
−d(π 0 ,π ∗ )
π 0 ∈Π φ

Pr[π] = P
θ

27


[form-feed in source extraction]
Both the numerator and denominator are differentiable with respect to θ = φ − 1, so Prθ [π] is differentiable

[Next verbatim source excerpt]

Monotonicity: We must show that for any S ⊂ x, if π1
denotes the value of the top-ranked candidate
according to π excluding candidates in S,
h
i
h
i
(−S)
(−S)
EFθ0 π1
≥ EFθ π1
.
(−S)

= xj .

(−S)

= xj }

For any i ∈
/ S, let j be the smallest index such that j > i and j ∈
/ S. Consider any π such that π1
(−S)
Then, swapping i and j yields a permutation π̂ such that π̂1
= xi . Moreover,
Pr[π̂] = Pr[π] · φinv(π)−inv(π̂) .
Since i < j, inv(π) − inv(π̂) ≥ 1. Finally, note that swapping i and j is a bijection between {π : π1
(−S)
and {π : π1
= xi }. Thus,
(−S)

= xi ]
=
(−S)
Pr[π1
= xj ]
(−S)
π:π
=x
Pr[π1

1

Note that the terms

Pr[π]

X

j

(−S)
Pr[π1
= xj ]

· φinv(π)−inv(π̂)

Pr[π]
sum to 1, so this is sum is some polynomial in φ with nonnegative weights
(−S)
Pr[π1
=xj ]

and integer powers of φ. As a result, it must have a positive derivative with respect to φ, i.e., for i < j,
(−S)

= xi ]
d Pr[π1
>0
dφ Pr[π (−S) = xj ]
1

0

Let φ > φ. Then,
(−S)

Prφ [π1

= xi ]

(−S)
Prφ [π1
= xj ]

(−S)

<

28

Prφ0 [π1

= xi ]
(−S)
Prφ0 [π1
= xj ]


[form-feed in source extraction]
Rearranging,
(−S)

Prφ [π1

= xi ]

(−S)
Prφ0 [π1
= xi ]

(−S)

Prφ [π1

= xj ]
(−S)
Prφ0 [π1
= xj ]

<

(D.1)
\end{ReviewVerbatim}

\paragraph{Approved corrected review target}
The archival source above is not asserted equivalent to this replacement.
\begin{ReviewVerbatim}
For a common center on exactly three candidates, phi_accurate > phi_noisy > 1, theta_r = phi_r-1, and any nonempty remaining set containing center-ordered candidates better before worse, the two source Mallows laws have q_r = phi_r^(-1) and normalized mass phi_r^(-KendallTau(center,pi)) / sum_tau phi_r^(-KendallTau(center,tau)). Writing P_r(c) for the probability that c is best in that remaining set under law r, P_accurate(better)P_noisy(worse) > P_accurate(worse)P_noisy(better). No arbitrary-universe claim or printed reverse D.1 inequality is asserted.
\end{ReviewVerbatim}

\textbf{Archival source anchor:} {\footnotesize\raggedright cited publication, lines 2317-\allowbreak{}2376\par}
\paragraph{Recorded basis}
User instruction 2026-\allowbreak{}07-\allowbreak{}25:\allowbreak{} assume whatever convention makes the result true, provided the convention is explicit and the archival claim is not represented as proved.\allowbreak{}
\paragraph{Expanded PaperInterface specification}
\begin{ReviewVerbatim}
∀ (center : EconCSLib.SocialChoice.Ranking.Ranking 1) (phiNoisy thetaNoisy phiAccurate thetaAccurate : ℝ),
  1 < phiNoisy →
    phiNoisy < phiAccurate →
      thetaNoisy = phiNoisy - 1 →
        thetaAccurate = phiAccurate - 1 →
          ∀ {remaining : Finset (EconCSLib.SocialChoice.Ranking.Candidate 1)},
            remaining.Nonempty →
              ∀ {better worse : EconCSLib.SocialChoice.Ranking.Candidate 1},
                better ∈ remaining →
                  worse ∈ remaining →
                    EconCSLib.SocialChoice.Ranking.rankOf center better <
                        EconCSLib.SocialChoice.Ranking.rankOf center worse →
                      have MAccurate := KR21Monoculture.concreteMallowsSpec center thetaAccurate;
                      have MNoisy := KR21Monoculture.concreteMallowsSpec center thetaNoisy;
                      MAccurate.q = phiAccurate⁻¹ ∧
                        MNoisy.q = phiNoisy⁻¹ ∧
                          (∀ (pi : EconCSLib.SocialChoice.Ranking.Ranking 1),
                              (MAccurate.law pi).toReal =
                                phiAccurate⁻¹ ^ EconCSLib.SocialChoice.Ranking.kendallTau center pi /
                                  ∑ tau, phiAccurate⁻¹ ^ EconCSLib.SocialChoice.Ranking.kendallTau center tau) ∧
                            (∀ (pi : EconCSLib.SocialChoice.Ranking.Ranking 1),
                                (MNoisy.law pi).toReal =
                                  phiNoisy⁻¹ ^ EconCSLib.SocialChoice.Ranking.kendallTau center pi /
                                    ∑ tau, phiNoisy⁻¹ ^ EconCSLib.SocialChoice.Ranking.kendallTau center tau) ∧
                              0 <
                                ((EconCSLib.pmfProb MAccurate.law fun pi =>
                                      better = KR21Monoculture.bestInSet pi remaining) *
                                    EconCSLib.pmfProb MNoisy.law fun pi =>
                                      worse = KR21Monoculture.bestInSet pi remaining) -
                                  (EconCSLib.pmfProb MAccurate.law fun pi =>
                                      worse = KR21Monoculture.bestInSet pi remaining) *
                                    EconCSLib.pmfProb MNoisy.law fun pi =>
                                      better = KR21Monoculture.bestInSet pi remaining
\end{ReviewVerbatim}

\paragraph{Lean proof endpoint (built separately)}\begin{ReviewVerbatim}
KR21Monoculture.PaperInterface.equationD1_source_phi_likelihood_ratio_complete
\end{ReviewVerbatim}

\paragraph{Recorded source-to-Spec screening}
\textbf{Verdict:} \texttt{not recorded}
\\
\textbf{Reason:} not recorded
\reviewmatch{Matches approved corrected target}
\noindent\textbf{Reviewer annotation}\par
\reviewerbox
\clearpage
\hypertarget{source-claim-c7560a9682e8906d}{}
\section*{Review row 46}
\textbf{Source locator:} {\footnotesize\raggedright cited publication, lines 133-\allowbreak{}138\par}
\paragraph{Verbatim source input}
\begin{ReviewVerbatim}
More formally, we model the n candidates as having intrinsic values x1 , . . . , xn , where any employer would
derive utility xi from hiring candidate i. Throughout the paper, we assume without loss of generality that
x1 > x2 > · · · > xn . These values, however, are unknown to the employer; instead, they must use some
noisy procedure to rank the candidates. We model such a procedure as a randomized mechanism R that
takes in the true candidate values and draws a permutation π over those candidates from some distribution.
Our main results hold for families of distributions over permutations as defined below:

[Next verbatim source excerpt]

The Mallows Model also appears frequently in the ranking literature [8, 11], and is much more analytically
tractable than RUMs. Under the Mallows Model, the likelihood of a permutation is related to its distance
from the true ranking π ∗ :
∗
1
(8)
Pr[π] = φ−d(π,π ) ,
Z
where Z is a normalizing constant. In this model, φ > 1 is the accuracy parameter: the larger φ is, the
more likely the ranking procedure is to output a ranking π that is close to the true ranking r. To instantiate
this model, we need a notion of distance d(·, ·) over permutations. For this, we’ll use Kendall tau distance,
another standard notion in the literature, which is simply the number of pairs of elements in π that are
incorrectly ordered [10]. In Appendix D, we verify that the family of distributions Fθ given by the Mallows
Model satisfies Definition 1, defining θ = φ − 1 (for consistency, so θ is well-defined on (0, ∞)).

[Next verbatim source excerpt]

E.1

Verifying Definition 2

We must show that when π, τ ∼ Fθ ,
E [π1 − π2 | π1 6= τ1 ] > 0.

(E.1)
\end{ReviewVerbatim}

\paragraph{Approved corrected review target}
The archival source above is not asserted equivalent to this replacement.
\begin{ReviewVerbatim}
For Candidate n = Fin(n+2) with n > 0, a center ranking, phi > 1, theta = phi-1, and a value profile strictly ordered by that center, the source Mallows law has q = phi^(-1) and normalized mass phi^(-KendallTau(center,pi)) / sum_tau phi^(-KendallTau(center,tau)). For two iid draws pi,tau, the top-disagreement event has positive probability and E[(value(first(pi))-value(second(pi)))*1_{first(pi) != first(tau)}] divided by that probability is strictly positive.
\end{ReviewVerbatim}

\textbf{Archival source anchor:} {\footnotesize\raggedright cited publication, lines 2444-\allowbreak{}2447\par}
\paragraph{Recorded basis}
User instruction 2026-\allowbreak{}07-\allowbreak{}25:\allowbreak{} assume whatever convention makes the result true, provided the convention is explicit and the archival claim is not represented as proved.\allowbreak{}
\paragraph{Expanded PaperInterface specification}
\begin{ReviewVerbatim}
∀ {n : ℕ} (center : EconCSLib.SocialChoice.Ranking.Ranking n) (phi theta : ℝ),
  1 < phi →
    theta = phi - 1 →
      0 < n →
        ∀ {value : EconCSLib.SocialChoice.Ranking.Candidate n → ℝ},
          EconCSLib.SocialChoice.Ranking.StrictlyOrderedBy center value →
            have M := KR21Monoculture.concreteMallowsSpec center theta;
            M.q = phi⁻¹ ∧
              (∀ (pi : EconCSLib.SocialChoice.Ranking.Ranking n),
                  (M.law pi).toReal =
                    phi⁻¹ ^ EconCSLib.SocialChoice.Ranking.kendallTau center pi /
                      ∑ tau, phi⁻¹ ^ EconCSLib.SocialChoice.Ranking.kendallTau center tau) ∧
                0 < EconCSLib.pmfPairProb M.law M.law KR21Monoculture.disagreementEvent ∧
                  0 <
                    (EconCSLib.pmfPairIndicatorExp M.law M.law KR21Monoculture.disagreementEvent fun pair =>
                        value (EconCSLib.SocialChoice.Ranking.firstChoice pair.1) -
                          value (EconCSLib.SocialChoice.Ranking.secondChoice pair.1)) /
                      EconCSLib.pmfPairProb M.law M.law KR21Monoculture.disagreementEvent
\end{ReviewVerbatim}

\paragraph{Lean proof endpoint (built separately)}\begin{ReviewVerbatim}
KR21Monoculture.PaperInterface.equationE1_source_phi_literal_conditional_semantic_complete
\end{ReviewVerbatim}

\paragraph{Recorded source-to-Spec screening}
\textbf{Verdict:} \texttt{not recorded}
\\
\textbf{Reason:} not recorded
\reviewmatch{Matches approved corrected target}
\noindent\textbf{Reviewer annotation}\par
\reviewerbox
\clearpage
\hypertarget{source-claim-73dd5d1917ec1a7c}{}
\section*{Review row 47}
\textbf{Source locator:} {\footnotesize\raggedright cited publication, lines 133-\allowbreak{}138\par}
\paragraph{Verbatim source input}
\begin{ReviewVerbatim}
More formally, we model the n candidates as having intrinsic values x1 , . . . , xn , where any employer would
derive utility xi from hiring candidate i. Throughout the paper, we assume without loss of generality that
x1 > x2 > · · · > xn . These values, however, are unknown to the employer; instead, they must use some
noisy procedure to rank the candidates. We model such a procedure as a randomized mechanism R that
takes in the true candidate values and draws a permutation π over those candidates from some distribution.
Our main results hold for families of distributions over permutations as defined below:

[Next verbatim source excerpt]

The Mallows Model also appears frequently in the ranking literature [8, 11], and is much more analytically
tractable than RUMs. Under the Mallows Model, the likelihood of a permutation is related to its distance
from the true ranking π ∗ :
∗
1
(8)
Pr[π] = φ−d(π,π ) ,
Z
where Z is a normalizing constant. In this model, φ > 1 is the accuracy parameter: the larger φ is, the
more likely the ranking procedure is to output a ranking π that is close to the true ranking r. To instantiate
this model, we need a notion of distance d(·, ·) over permutations. For this, we’ll use Kendall tau distance,
another standard notion in the literature, which is simply the number of pairs of elements in π that are
incorrectly ordered [10]. In Appendix D, we verify that the family of distributions Fθ given by the Mallows
Model satisfies Definition 1, defining θ = φ − 1 (for consistency, so θ is well-defined on (0, ∞)).

[Next verbatim source excerpt]

We begin by expanding:
E [π1 − π2 | π1 6= τ1 ] =

n
n X
X

(xi − xj ) Pr[π1 = xi ∩ π2 = xj | π1 6= τ1 ]

i=1 j=1

=

n−1
XX

(xi − xj ) (Pr[π1 = xi ∩ π2 = xj | π1 6= τ1 ] − Pr[π1 = xj ∩ π2 = xi | π1 6= τ1 ])

i=1 j>i

Since xi > xj for i < j, it suffices to show that for all i < j,
Pr[π1 = xi ∩ π2 = xj | π1 6= τ1 ] ≥ Pr[π1 = xj ∩ π2 = xi | π1 6= τ1 ],

(E.2)

and that this holds strictly for some i < j. We simplify (E.2) as follows:
\end{ReviewVerbatim}

\paragraph{Approved corrected review target}
The archival source above is not asserted equivalent to this replacement.
\begin{ReviewVerbatim}
For Candidate n = Fin(n+2) with n > 0, a center ranking, phi > 1, and theta = phi-1, the source Mallows law has q = phi^(-1), normalized mass phi^(-KendallTau(center,pi)) / sum_tau phi^(-KendallTau(center,tau)), and positive iid top-disagreement probability. Conditional on top disagreement, for every center-ordered pair c before d, the probability of first-two order d,c is at most that of c,d; for at least one center-ordered pair the inequality is strict.
\end{ReviewVerbatim}

\textbf{Archival source anchor:} {\footnotesize\raggedright cited publication, lines 2469-\allowbreak{}2474\par}
\paragraph{Recorded basis}
User instruction 2026-\allowbreak{}07-\allowbreak{}25:\allowbreak{} assume whatever convention makes the result true, provided the convention is explicit and the archival claim is not represented as proved.\allowbreak{}
\paragraph{Expanded PaperInterface specification}
\begin{ReviewVerbatim}
∀ {n : ℕ} (center : EconCSLib.SocialChoice.Ranking.Ranking n) (phi theta : ℝ),
  1 < phi →
    theta = phi - 1 →
      0 < n →
        have M := KR21Monoculture.concreteMallowsSpec center theta;
        M.q = phi⁻¹ ∧
          (∀ (pi : EconCSLib.SocialChoice.Ranking.Ranking n),
              (M.law pi).toReal =
                phi⁻¹ ^ EconCSLib.SocialChoice.Ranking.kendallTau center pi /
                  ∑ tau, phi⁻¹ ^ EconCSLib.SocialChoice.Ranking.kendallTau center tau) ∧
            0 < EconCSLib.pmfPairProb M.law M.law KR21Monoculture.disagreementEvent ∧
              (∀ (c d : EconCSLib.SocialChoice.Ranking.Candidate n),
                  EconCSLib.SocialChoice.Ranking.rankOf M.center c < EconCSLib.SocialChoice.Ranking.rankOf M.center d →
                    (EconCSLib.pmfPairIndicatorExp M.law M.law KR21Monoculture.disagreementEvent fun pair =>
                          if
                              d = EconCSLib.SocialChoice.Ranking.firstChoice pair.1 ∧
                                c = EconCSLib.SocialChoice.Ranking.secondChoice pair.1 then
                            1
                          else 0) /
                        EconCSLib.pmfPairProb M.law M.law KR21Monoculture.disagreementEvent ≤
                      (EconCSLib.pmfPairIndicatorExp M.law M.law KR21Monoculture.disagreementEvent fun pair =>
                          if
                              c = EconCSLib.SocialChoice.Ranking.firstChoice pair.1 ∧
                                d = EconCSLib.SocialChoice.Ranking.secondChoice pair.1 then
                            1
                          else 0) /
                        EconCSLib.pmfPairProb M.law M.law KR21Monoculture.disagreementEvent) ∧
                ∃ c d,
                  EconCSLib.SocialChoice.Ranking.rankOf M.center c < EconCSLib.SocialChoice.Ranking.rankOf M.center d ∧
                    (EconCSLib.pmfPairIndicatorExp M.law M.law KR21Monoculture.disagreementEvent fun pair =>
                          if
                              d = EconCSLib.SocialChoice.Ranking.firstChoice pair.1 ∧
                                c = EconCSLib.SocialChoice.Ranking.secondChoice pair.1 then
                            1
                          else 0) /
                        EconCSLib.pmfPairProb M.law M.law KR21Monoculture.disagreementEvent <
                      (EconCSLib.pmfPairIndicatorExp M.law M.law KR21Monoculture.disagreementEvent fun pair =>
                          if
                              c = EconCSLib.SocialChoice.Ranking.firstChoice pair.1 ∧
                                d = EconCSLib.SocialChoice.Ranking.secondChoice pair.1 then
                            1
                          else 0) /
                        EconCSLib.pmfPairProb M.law M.law KR21Monoculture.disagreementEvent
\end{ReviewVerbatim}

\paragraph{Lean proof endpoint (built separately)}\begin{ReviewVerbatim}
KR21Monoculture.PaperInterface.equationE2_source_phi_literal_conditional_semantic_complete
\end{ReviewVerbatim}

\paragraph{Recorded source-to-Spec screening}
\textbf{Verdict:} \texttt{not recorded}
\\
\textbf{Reason:} not recorded
\reviewmatch{Matches approved corrected target}
\noindent\textbf{Reviewer annotation}\par
\reviewerbox
\clearpage
\hypertarget{source-claim-027748e89e133285}{}
\section*{Review row 48}
\textbf{Source locator:} {\footnotesize\raggedright cited publication, lines 133-\allowbreak{}138\par}
\paragraph{Verbatim source input}
\begin{ReviewVerbatim}
More formally, we model the n candidates as having intrinsic values x1 , . . . , xn , where any employer would
derive utility xi from hiring candidate i. Throughout the paper, we assume without loss of generality that
x1 > x2 > · · · > xn . These values, however, are unknown to the employer; instead, they must use some
noisy procedure to rank the candidates. We model such a procedure as a randomized mechanism R that
takes in the true candidate values and draws a permutation π over those candidates from some distribution.
Our main results hold for families of distributions over permutations as defined below:

[Next verbatim source excerpt]

The Mallows Model also appears frequently in the ranking literature [8, 11], and is much more analytically
tractable than RUMs. Under the Mallows Model, the likelihood of a permutation is related to its distance
from the true ranking π ∗ :
∗
1
(8)
Pr[π] = φ−d(π,π ) ,
Z
where Z is a normalizing constant. In this model, φ > 1 is the accuracy parameter: the larger φ is, the
more likely the ranking procedure is to output a ranking π that is close to the true ranking r. To instantiate
this model, we need a notion of distance d(·, ·) over permutations. For this, we’ll use Kendall tau distance,
another standard notion in the literature, which is simply the number of pairs of elements in π that are
incorrectly ordered [10]. In Appendix D, we verify that the family of distributions Fθ given by the Mallows
Model satisfies Definition 1, defining θ = φ − 1 (for consistency, so θ is well-defined on (0, ∞)).

[Next verbatim source excerpt]

and that this holds strictly for some i < j. We simplify (E.2) as follows:
Pr[π1 = xi ∩ π2 = xj | π1 6= τ1 ] > Pr[π1 = xj ∩ π2 = xi | π1 6= τ1 ]
Pr[π1 = xi ∩ π2 = xj ∩ π1 6= τ1 ]
Pr[π1 = xj ∩ π2 = xi ∩ π1 6= τ1 ]
>
Pr[π1 6= τ1 ]
Pr[π1 6= τ1 ]
⇐⇒ Pr[π1 = xi ∩ π2 = xj ∩ π1 6= τ1 ] > Pr[π1 = xj ∩ π2 = xi ∩ π1 6= τ1 ]
⇐⇒

⇐⇒ Pr[π1 = xi ∩ π2 = xj ∩ τ1 6= xi ] > Pr[π1 = xj ∩ π2 = xi ∩ τ1 6= xj ]
⇐⇒ Pr[π1 = xi ∩ π2 = xj ] Pr[τ1 6= xi ] > Pr[π1 = xj ∩ π2 = xi ] Pr[τ1 6= xj ]

(E.3)

We can simplify (E.3) using Lemmas 5 and 6. Let |i − j| denote the difference in rank between xi and xj .
Pr[π1 = xi ∩ π2 = xj ] Pr[τ1 6= xi ] − Pr[π1 = xj ∩ π2 = xi ] Pr[τ1 6= xj ]
= Pr[π1 = xi ∩ π2 = xj ](1 − Pr[τ1 = xi ]) − φ−1 Pr[π1 = xi ∩ π2 = xj ](1 − Pr[τ1 = xj ])
= Pr[π1 = xi ∩ π2 = xj ](1 − Pr[τ1 = xi ]) − φ−1 Pr[π1 = xi ∩ π2 = xj ](1 − φ−|i−j| Pr[τ1 = xi ])
= Pr[π1 = xi ∩ π2 = xj ](1 − Pr[τ1 = xi ] − φ−1 − φ−|i−j|−1 Pr[τ1 = xi ]))
29


[form-feed in source extraction]
This is positive if and only if
1 − Pr[τ1 = xi ] − φ−1 − φ−|i−j|−1 Pr[τ1 = xi ] > 0
[U+0011 control character]
[U+0010 control character]
⇐⇒ Pr[τ1 = xi ] 1 − φ−|i−j|−1 < 1 − φ−1
1 − φ−1
1 − φ−|i−j|−1
−1
1−φ
1 − φ−1
⇐⇒ i−1
<
−n
φ (1 − φ )
1 − φ−|i−j|−1
⇐⇒ Pr[τ1 = xi ] <

⇐⇒ φi−1 (1 − φ−n ) > 1 − φ−|i−j|−1
This is weakly true for any i < j because φi−1 ≥ 1 and |i − j| + 1 ≤ n, and it is strictly true for any i, j
other than 1 and n. Thus, E [π1 − π2 | π1 6= τ1 ] > 0.
\end{ReviewVerbatim}

\paragraph{Approved corrected review target}
The archival source above is not asserted equivalent to this replacement.
\begin{ReviewVerbatim}
For Candidate n = Fin(n+2) with n > 0, a center ranking, phi > 1, and theta = phi-1, the source Mallows law has q = phi^(-1) and normalized mass phi^(-KendallTau(center,pi)) / sum_tau phi^(-KendallTau(center,tau)). For every center-ordered pair c before d, P(first=c,second=d)P(first!=c) >= P(first=d,second=c)P(first!=d), and this inequality is strict for at least one center-ordered pair. The target uses the corrected positive final product term; it does not certify the printed negative-sign intermediate expansion.
\end{ReviewVerbatim}

\textbf{Archival source anchor:} {\footnotesize\raggedright cited publication, lines 2474-\allowbreak{}2515\par}
\paragraph{Recorded basis}
User instruction 2026-\allowbreak{}07-\allowbreak{}25:\allowbreak{} assume whatever convention makes the result true, provided the convention is explicit and the archival claim is not represented as proved.\allowbreak{}
\paragraph{Expanded PaperInterface specification}
\begin{ReviewVerbatim}
∀ {n : ℕ} (center : EconCSLib.SocialChoice.Ranking.Ranking n) (phi theta : ℝ),
  1 < phi →
    theta = phi - 1 →
      0 < n →
        have M := KR21Monoculture.concreteMallowsSpec center theta;
        M.q = phi⁻¹ ∧
          (∀ (pi : EconCSLib.SocialChoice.Ranking.Ranking n),
              (M.law pi).toReal =
                phi⁻¹ ^ EconCSLib.SocialChoice.Ranking.kendallTau center pi /
                  ∑ tau, phi⁻¹ ^ EconCSLib.SocialChoice.Ranking.kendallTau center tau) ∧
            (∀ (c d : EconCSLib.SocialChoice.Ranking.Candidate n),
                EconCSLib.SocialChoice.Ranking.rankOf M.center c < EconCSLib.SocialChoice.Ranking.rankOf M.center d →
                  0 ≤
                    M.firstSecondChoiceProb c d * KR21Monoculture.firstChoiceMissProb M.law c -
                      M.firstSecondChoiceProb d c * KR21Monoculture.firstChoiceMissProb M.law d) ∧
              ∃ c d,
                EconCSLib.SocialChoice.Ranking.rankOf M.center c < EconCSLib.SocialChoice.Ranking.rankOf M.center d ∧
                  0 <
                    M.firstSecondChoiceProb c d * KR21Monoculture.firstChoiceMissProb M.law c -
                      M.firstSecondChoiceProb d c * KR21Monoculture.firstChoiceMissProb M.law d
\end{ReviewVerbatim}

\paragraph{Lean proof endpoint (built separately)}\begin{ReviewVerbatim}
KR21Monoculture.PaperInterface.equationE3_source_phi_complete
\end{ReviewVerbatim}

\paragraph{Recorded source-to-Spec screening}
\textbf{Verdict:} \texttt{not recorded}
\\
\textbf{Reason:} not recorded
\reviewmatch{Matches approved corrected target}
\noindent\textbf{Reviewer annotation}\par
\reviewerbox
\clearpage
\hypertarget{source-claim-f5654e8ee63c4d5d}{}
\section*{Review row 49}
\textbf{Source locator:} {\footnotesize\raggedright cited publication, lines 2703-\allowbreak{}2850\par}
\paragraph{Verbatim source input}
\begin{ReviewVerbatim}
Lemma 4. Let {yi }ni=1 , {pi }ni=1 , and {qi }ni=1 be sequences such that
• yi is strictly increasing.
Pn
Pn
•
i=1 pi =
i=1 qi = 1.
• pqii is decreasing.
Pn
Pn
Then, i=1 pi yi < i=1 qi yi .
Proof. First, note that there exists j such that pi > qi for i < j and pi ≤ qi for i ≥ j. To see this, let j be
the smallest index such that pj ≤ qj . Such a j must exist because pi and qi both sum to 1, so it cannot be
the case that pi > qi for all i. This implies pi /qi ≤ 1, and since pi /qi is decreasing, pi ≤ qi for i ≥ j.
Next, note that
0=

n
X

(pi − qi )

i=1

=

j−1
X

(pi − qi ) +

i=1

n
X
(pi − qi ),
i=j

meaning
j−1
X

(pi − qi ) =

i=1

n
X
i=j

31

(qi − pi ).


[form-feed in source extraction]
Using this choice of j, we can write
n
X
i=1

p i yi −

n
X

q i yi =

i=1

n
X

(pi − qi )yi

i=1

=

j−1
X

(pi − qi )yi −

i=1

≤

j−1
X

n
X

(pi − qi )yj−1 −

i=1

=

j−1
X

j−1
X

(pi − qi )yj−1 −

j−1
X

(qi − pi )yj

n
X

(qi − pi )yj

i=j

(pi − qi )yj−1 −

i=1

=

n
X
i=j

i=1

=

(qi − pi )yi

i=j

j−1
X

(pi − qi )yj

i=1

(pi − qi )(yj−1 − yj )

i=1

<0
\end{ReviewVerbatim}

\paragraph{Approved corrected review target}
The archival source above is not asserted equivalent to this replacement.
\begin{ReviewVerbatim}
For Candidate n = Fin(n+2), real sequences p, q, and a strictly increasing real sequence y, assume sum_i p_i = sum_i q_i = 1, p_i*q_j-p_j*q_i >= 0 for every i<j, and there exist i<j with p_i*q_j-p_j*q_i > 0. Then sum_i p_i*y_i < sum_i q_i*y_i. The explicit strict cross-ratio witness is required; strictness is not inferred from the non-strict pairwise condition alone.
\end{ReviewVerbatim}

\textbf{Archival source anchor:} {\footnotesize\raggedright cited publication, lines 2703-\allowbreak{}2713\par}
\paragraph{Recorded basis}
User instruction 2026-\allowbreak{}07-\allowbreak{}25:\allowbreak{} assume whatever convention makes the result true, provided the convention is explicit and the archival claim is not represented as proved.\allowbreak{}
\paragraph{Expanded PaperInterface specification}
\begin{ReviewVerbatim}
∀ (n : ℕ) {p q y : EconCSLib.SocialChoice.Ranking.Candidate n → ℝ},
  ∑ i, p i = 1 →
    ∑ i, q i = 1 →
      (∀ (i j : EconCSLib.SocialChoice.Ranking.Candidate n), i < j → 0 ≤ p i * q j - p j * q i) →
        (∃ i j, i < j ∧ 0 < p i * q j - p j * q i) → StrictMono y → ∑ i, p i * y i < ∑ i, q i * y i
\end{ReviewVerbatim}

\paragraph{Lean proof endpoint (built separately)}\begin{ReviewVerbatim}
KR21Monoculture.PaperInterface.lemma4_weighted_average_lt_of_cross_ratio
\end{ReviewVerbatim}

\paragraph{Recorded source-to-Spec screening}
\textbf{Verdict:} \texttt{not recorded}
\\
\textbf{Reason:} not recorded
\reviewmatch{Matches approved corrected target}
\noindent\textbf{Reviewer annotation}\par
\reviewerbox
\clearpage
\hypertarget{source-claim-7e7ac58d3e77e0cc}{}
\section*{Review row 50}
\textbf{Source locator:} {\footnotesize\raggedright cited publication, lines 133-\allowbreak{}138\par}
\paragraph{Verbatim source input}
\begin{ReviewVerbatim}
More formally, we model the n candidates as having intrinsic values x1 , . . . , xn , where any employer would
derive utility xi from hiring candidate i. Throughout the paper, we assume without loss of generality that
x1 > x2 > · · · > xn . These values, however, are unknown to the employer; instead, they must use some
noisy procedure to rank the candidates. We model such a procedure as a randomized mechanism R that
takes in the true candidate values and draws a permutation π over those candidates from some distribution.
Our main results hold for families of distributions over permutations as defined below:

[Next verbatim source excerpt]

The Mallows Model also appears frequently in the ranking literature [8, 11], and is much more analytically
tractable than RUMs. Under the Mallows Model, the likelihood of a permutation is related to its distance
from the true ranking π ∗ :
∗
1
(8)
Pr[π] = φ−d(π,π ) ,
Z
where Z is a normalizing constant. In this model, φ > 1 is the accuracy parameter: the larger φ is, the
more likely the ranking procedure is to output a ranking π that is close to the true ranking r. To instantiate
this model, we need a notion of distance d(·, ·) over permutations. For this, we’ll use Kendall tau distance,
another standard notion in the literature, which is simply the number of pairs of elements in π that are
incorrectly ordered [10]. In Appendix D, we verify that the family of distributions Fθ given by the Mallows
Model satisfies Definition 1, defining θ = φ − 1 (for consistency, so θ is well-defined on (0, ∞)).

[Next verbatim source excerpt]

Lemma 5. For xi > xj ,
Pr[π1 = xi ∩ π2 = xj ] = φ Pr[π1 = xj ∩ π2 = xi ].

(F.1)
\end{ReviewVerbatim}


\paragraph{Expanded PaperInterface specification}
\begin{ReviewVerbatim}
∀ {n : ℕ} (center : EconCSLib.SocialChoice.Ranking.Ranking n) (phi theta : ℝ),
  1 < phi →
    theta = phi - 1 →
      ∀ {c d : EconCSLib.SocialChoice.Ranking.Candidate n},
        EconCSLib.SocialChoice.Ranking.rankOf center c < EconCSLib.SocialChoice.Ranking.rankOf center d →
          have M := KR21Monoculture.concreteMallowsSpec center theta;
          M.q = phi⁻¹ ∧
            (∀ (pi : EconCSLib.SocialChoice.Ranking.Ranking n),
                (M.law pi).toReal =
                  phi⁻¹ ^ EconCSLib.SocialChoice.Ranking.kendallTau center pi /
                    ∑ tau, phi⁻¹ ^ EconCSLib.SocialChoice.Ranking.kendallTau center tau) ∧
              M.firstSecondChoiceProb c d = phi * M.firstSecondChoiceProb d c
\end{ReviewVerbatim}

\paragraph{Lean proof endpoint (built separately)}\begin{ReviewVerbatim}
KR21Monoculture.PaperInterface.equationF1_source_phi_complete
\end{ReviewVerbatim}

\paragraph{Recorded source-to-Spec screening}
\textbf{Verdict:} \texttt{not recorded}
\\
\textbf{Reason:} not recorded
\reviewmatch{Matches source input}
\noindent\textbf{Reviewer annotation}\par
\reviewerbox
\clearpage
\hypertarget{source-claim-7ef5e5254a6c73ab}{}
\section*{Review row 51}
\textbf{Source locator:} {\footnotesize\raggedright cited publication, lines 133-\allowbreak{}138\par}
\paragraph{Verbatim source input}
\begin{ReviewVerbatim}
More formally, we model the n candidates as having intrinsic values x1 , . . . , xn , where any employer would
derive utility xi from hiring candidate i. Throughout the paper, we assume without loss of generality that
x1 > x2 > · · · > xn . These values, however, are unknown to the employer; instead, they must use some
noisy procedure to rank the candidates. We model such a procedure as a randomized mechanism R that
takes in the true candidate values and draws a permutation π over those candidates from some distribution.
Our main results hold for families of distributions over permutations as defined below:

[Next verbatim source excerpt]

The Mallows Model also appears frequently in the ranking literature [8, 11], and is much more analytically
tractable than RUMs. Under the Mallows Model, the likelihood of a permutation is related to its distance
from the true ranking π ∗ :
∗
1
(8)
Pr[π] = φ−d(π,π ) ,
Z
where Z is a normalizing constant. In this model, φ > 1 is the accuracy parameter: the larger φ is, the
more likely the ranking procedure is to output a ranking π that is close to the true ranking r. To instantiate
this model, we need a notion of distance d(·, ·) over permutations. For this, we’ll use Kendall tau distance,
another standard notion in the literature, which is simply the number of pairs of elements in π that are
incorrectly ordered [10]. In Appendix D, we verify that the family of distributions Fθ given by the Mallows
Model satisfies Definition 1, defining θ = φ − 1 (for consistency, so θ is well-defined on (0, ∞)).

[Next verbatim source excerpt]

Lemma 6. For 1 ≤ i ≤ n,
Pr[π1 = xi ] =

1 − φ−1
.
φi−1 (1 − φ−n )

(F.2)

Proof. Let π−i be a permutation over all items except i. Then,
Pr[π1 = xi ] =

X

Pr[π1 = xi | π−i ] Pr[π−i ]

π−i

=

X

φ−(i−1) Pr[π−i ]

π−i

= φ−(i−1)

X

Pr[π−i ]

π−i

Note that
PnPr[π−i ] doesn’t depend on which n − 1 items are being ranked, so this term appears for any i.
Moreover, i=1 Pr[π1 = xi ] = 1. Therefore, we have
32


[form-feed in source extraction]
Pr[π1 = xi ] ∝ φ−(i−1) .
Normalizing, we get
φ−(i−1)
Pr[π1 = xi ] = Pn
−(j−1)
j=1 φ
=
=

φ−(i−1)
1−φ−n
1−φ−1

1 − φ−1
φi−1 (1 − φ−n )

Intuitively, any permutation over n−1 items is equally likely regardless of what those items are, and inserting
any item at the front of this permutation yields a likelihood proportional to the number of additional
inversions this causes, which is equal to the item’s position on the list.1
\end{ReviewVerbatim}


\paragraph{Expanded PaperInterface specification}
\begin{ReviewVerbatim}
∀ {n : ℕ} (center : EconCSLib.SocialChoice.Ranking.Ranking n) (phi theta : ℝ),
  1 < phi →
    theta = phi - 1 →
      ∀ (c : EconCSLib.SocialChoice.Ranking.Candidate n),
        have M := KR21Monoculture.concreteMallowsSpec center theta;
        M.q = phi⁻¹ ∧
          (∀ (pi : EconCSLib.SocialChoice.Ranking.Ranking n),
              (M.law pi).toReal =
                phi⁻¹ ^ EconCSLib.SocialChoice.Ranking.kendallTau center pi /
                  ∑ tau, phi⁻¹ ^ EconCSLib.SocialChoice.Ranking.kendallTau center tau) ∧
            KR21Monoculture.firstChoiceProb M.law c =
              phi⁻¹ ^ ↑(EconCSLib.SocialChoice.Ranking.rankOf M.center c) /
                KR21Monoculture.candidateRankPowerSum n phi⁻¹
\end{ReviewVerbatim}

\paragraph{Lean proof endpoint (built separately)}\begin{ReviewVerbatim}
KR21Monoculture.PaperInterface.lemma6_source_phi_rank_power_complete
\end{ReviewVerbatim}

\paragraph{Recorded source-to-Spec screening}
\textbf{Verdict:} \texttt{not recorded}
\\
\textbf{Reason:} not recorded
\reviewmatch{Matches source input}
\noindent\textbf{Reviewer annotation}\par
\reviewerbox
\clearpage
\hypertarget{source-claim-4e3a949fd1ab09f3}{}
\section*{Review row 52}
\textbf{Source locator:} {\footnotesize\raggedright cited publication, lines 133-\allowbreak{}138\par}
\paragraph{Verbatim source input}
\begin{ReviewVerbatim}
More formally, we model the n candidates as having intrinsic values x1 , . . . , xn , where any employer would
derive utility xi from hiring candidate i. Throughout the paper, we assume without loss of generality that
x1 > x2 > · · · > xn . These values, however, are unknown to the employer; instead, they must use some
noisy procedure to rank the candidates. We model such a procedure as a randomized mechanism R that
takes in the true candidate values and draws a permutation π over those candidates from some distribution.
Our main results hold for families of distributions over permutations as defined below:

[Next verbatim source excerpt]

The Mallows Model also appears frequently in the ranking literature [8, 11], and is much more analytically
tractable than RUMs. Under the Mallows Model, the likelihood of a permutation is related to its distance
from the true ranking π ∗ :
∗
1
(8)
Pr[π] = φ−d(π,π ) ,
Z
where Z is a normalizing constant. In this model, φ > 1 is the accuracy parameter: the larger φ is, the
more likely the ranking procedure is to output a ranking π that is close to the true ranking r. To instantiate
this model, we need a notion of distance d(·, ·) over permutations. For this, we’ll use Kendall tau distance,
another standard notion in the literature, which is simply the number of pairs of elements in π that are
incorrectly ordered [10]. In Appendix D, we verify that the family of distributions Fθ given by the Mallows
Model satisfies Definition 1, defining θ = φ − 1 (for consistency, so θ is well-defined on (0, ∞)).

[Next verbatim source excerpt]

Lemma 6. For 1 ≤ i ≤ n,
Pr[π1 = xi ] =

1 − φ−1
.
φi−1 (1 − φ−n )

(F.2)
\end{ReviewVerbatim}


\paragraph{Expanded PaperInterface specification}
\begin{ReviewVerbatim}
∀ {n : ℕ} (center : EconCSLib.SocialChoice.Ranking.Ranking n) (phi theta : ℝ),
  1 < phi →
    theta = phi - 1 →
      ∀ (c : EconCSLib.SocialChoice.Ranking.Candidate n),
        have M := KR21Monoculture.concreteMallowsSpec center theta;
        M.q = phi⁻¹ ∧
          (∀ (pi : EconCSLib.SocialChoice.Ranking.Ranking n),
              (M.law pi).toReal =
                phi⁻¹ ^ EconCSLib.SocialChoice.Ranking.kendallTau center pi /
                  ∑ tau, phi⁻¹ ^ EconCSLib.SocialChoice.Ranking.kendallTau center tau) ∧
            KR21Monoculture.firstChoiceProb M.law c =
              (1 - phi⁻¹) / (phi ^ ↑(EconCSLib.SocialChoice.Ranking.rankOf M.center c) * (1 - phi⁻¹ ^ (n + 2)))
\end{ReviewVerbatim}

\paragraph{Lean proof endpoint (built separately)}\begin{ReviewVerbatim}
KR21Monoculture.PaperInterface.equationF2_source_phi_closed_form_complete
\end{ReviewVerbatim}

\paragraph{Recorded source-to-Spec screening}
\textbf{Verdict:} \texttt{not recorded}
\\
\textbf{Reason:} not recorded
\reviewmatch{Matches source input}
\noindent\textbf{Reviewer annotation}\par
\reviewerbox
\clearpage
\hypertarget{source-claim-1e4808a386b72297}{}
\section*{Review row 53}
\textbf{Source locator:} {\footnotesize\raggedright cited publication, lines 133-\allowbreak{}138\par}
\paragraph{Verbatim source input}
\begin{ReviewVerbatim}
More formally, we model the n candidates as having intrinsic values x1 , . . . , xn , where any employer would
derive utility xi from hiring candidate i. Throughout the paper, we assume without loss of generality that
x1 > x2 > · · · > xn . These values, however, are unknown to the employer; instead, they must use some
noisy procedure to rank the candidates. We model such a procedure as a randomized mechanism R that
takes in the true candidate values and draws a permutation π over those candidates from some distribution.
Our main results hold for families of distributions over permutations as defined below:

[Next verbatim source excerpt]

The Mallows Model also appears frequently in the ranking literature [8, 11], and is much more analytically
tractable than RUMs. Under the Mallows Model, the likelihood of a permutation is related to its distance
from the true ranking π ∗ :
∗
1
(8)
Pr[π] = φ−d(π,π ) ,
Z
where Z is a normalizing constant. In this model, φ > 1 is the accuracy parameter: the larger φ is, the
more likely the ranking procedure is to output a ranking π that is close to the true ranking r. To instantiate
this model, we need a notion of distance d(·, ·) over permutations. For this, we’ll use Kendall tau distance,
another standard notion in the literature, which is simply the number of pairs of elements in π that are
incorrectly ordered [10]. In Appendix D, we verify that the family of distributions Fθ given by the Mallows
Model satisfies Definition 1, defining θ = φ − 1 (for consistency, so θ is well-defined on (0, ∞)).

[Next verbatim source excerpt]

Lemma 7. For the Mallows Model, UH (θA , θH ) > UHH (θA , θH ).
Proof. Intuitively, this is because selecting first is better than selecting second. To prove this, let π and τ
be ranking generated by independent human evaluators under the Mallows Model, i.e., π, τ ∼ FθH .
UH (θA , θH ) − UHH (θA , θH ) = E [π1 ] − E [τ1 · 1π1 6=τ1 + τ2 · 1π1 =τ1 ]
= E [(π1 − τ2 ) · 1π1 =τ1 ]

= E [(π1 − π2 ) · 1π1 =τ1 ]
For any i < j, conditioned on π1 = τ1 , they are more likely to be correctly ordered than not:
X
E [(π1 − π2 ) · 1π1 =τ1 ] =
(Pr[π1 = xi ∩ τ1 = xi ∩ π2 = xj ] − Pr[π1 = xj ∩ τ1 = xj ∩ π2 = xi ]) (xi − xj )
i<j

=

X

(Pr[π1 = xi ∩ π2 = xj ] Pr[τ1 = xi ] − Pr[π1 = xj ∩ π2 = xi ] Pr[τ1 = xj ]) (xi − xj )

i<j

>

X

(Pr[π1 = xi ∩ π2 = xj ] Pr[τ1 = xj ] − Pr[π1 = xj ∩ π2 = xi ] Pr[τ1 = xj ]) (xi − xj )

i<j

=

X

(Pr[π1 = xi ∩ π2 = xj ] − Pr[π1 = xj ∩ π2 = xi ]) (xi − xj )

i<j

≥

X

(φH Pr[π1 = xj ∩ π2 = xi ] − Pr[π1 = xj ∩ π2 = xi ]) (xi − xj )

i<j

>0
\end{ReviewVerbatim}


\paragraph{Expanded PaperInterface specification}
\begin{ReviewVerbatim}
∀ {n : ℕ} (center : EconCSLib.SocialChoice.Ranking.Ranking n) (phi theta : ℝ),
  1 < phi →
    theta = phi - 1 →
      ∀ {value : EconCSLib.SocialChoice.Ranking.Candidate n → ℝ},
        EconCSLib.SocialChoice.Ranking.StrictlyOrderedBy center value →
          have M := KR21Monoculture.concreteMallowsSpec center theta;
          M.q = phi⁻¹ ∧
            (∀ (pi : EconCSLib.SocialChoice.Ranking.Ranking n),
                (M.law pi).toReal =
                  phi⁻¹ ^ EconCSLib.SocialChoice.Ranking.kendallTau center pi /
                    ∑ tau, phi⁻¹ ^ EconCSLib.SocialChoice.Ranking.kendallTau center tau) ∧
              EconCSLib.SocialChoice.Ranking.expectedSecondMoverIndependent M.law M.law value <
                EconCSLib.SocialChoice.Ranking.expectedFirstMoverUtility M.law value
\end{ReviewVerbatim}

\paragraph{Lean proof endpoint (built separately)}\begin{ReviewVerbatim}
KR21Monoculture.PaperInterface.lemma7_source_phi_complete
\end{ReviewVerbatim}

\paragraph{Recorded source-to-Spec screening}
\textbf{Verdict:} \texttt{not recorded}
\\
\textbf{Reason:} not recorded
\reviewmatch{Matches source input}
\noindent\textbf{Reviewer annotation}\par
\reviewerbox
\clearpage
\hypertarget{source-claim-961880e0b1d1c772}{}
\section*{Review row 54}
\textbf{Source locator:} {\footnotesize\raggedright cited publication, lines 133-\allowbreak{}138\par}
\paragraph{Verbatim source input}
\begin{ReviewVerbatim}
More formally, we model the n candidates as having intrinsic values x1 , . . . , xn , where any employer would
derive utility xi from hiring candidate i. Throughout the paper, we assume without loss of generality that
x1 > x2 > · · · > xn . These values, however, are unknown to the employer; instead, they must use some
noisy procedure to rank the candidates. We model such a procedure as a randomized mechanism R that
takes in the true candidate values and draws a permutation π over those candidates from some distribution.
Our main results hold for families of distributions over permutations as defined below:

[Next verbatim source excerpt]

Lemma 8. Under the Mallows model, the probability that any two items i < j are correctly ranked increases
monotonically with the accuracy parameter φ.
1 Alternatively, we could prove this by showing that for any permutation with i in front, the permutation in which i and i − 1
are swapped is φ times more likely, and thus, i − 1 is φ times more likely to be in front than i.

33


[form-feed in source extraction]
Proof. Let inv(π) be the number of inversions in a permutation π. Under the Mallows model, the probability
of observing π is proportional to φ− inv(π) . Let Si[U+001F control character]j (resp. Sj[U+001F control character]i ) be the set of permutations where i is ranked
before j (resp. j is ranked before i). Then, the probability i is ranked before j is
P
− inv(π)
π∈Si[U+001F control character]j φ
P
.
Pr[i [U+001F control character] j] = P
− inv(π) +
− inv(π)
π∈Si[U+001F control character]j φ
π∈Sj[U+001F control character]i φ
d
d Pr[i[U+001F control character]j]
We will show that dφ
Pr[i [U+001F control character] j] > 0. Note that this is equivalent to showing dφ
Pr[j[U+001F control character]i] > 0. Note that
P
− inv(π)
Pr[i [U+001F control character] j]
π∈Si[U+001F control character]j φ
.
=P
− inv(π)
Pr[j [U+001F control character] i]
π∈Sj[U+001F control character]i φ

Let πi:j be the subsequence of π containing elements i through j. Then, we have
P
− inv(π)
Pr[i [U+001F control character] j]
π∈Si[U+001F control character]j φ
=P
− inv(π)
Pr[j [U+001F control character] i]
π∈Sj[U+001F control character]i φ
P
P
0
− inv(πi:j )
φinv(πi:j )−inv(π )
0 =π
πi:j :π∈Si[U+001F control character]j φ
π 0 :πi:j
i:j
P
=P
0
− inv(πi:j )
φinv(πi:j )−inv(π )
0 =π
πi:j :π∈Sj[U+001F control character]i φ
π 0 :πi:j
i:j
P
− inv(πi:j )
πi:j :π∈Si[U+001F control character]j φ
=P
− inv(πi:j )
πi:j :π∈Sj[U+001F control character]i φ
P
0
Intuitively, the term
φinv(πi:j )−inv(π ) does not depend on πi:j because for any πi:j , if we fix
0 =π
π 0 :πi:j
i:j
the order and positions of the remaining elements, the number of inversions involving at least one element
outside of i : j (i.e., inv(π 0 )−inv(πi:j )) is a constant. For fixed πi:j , there is a bijection between permutations
0
π 0 : πi:j
= πi:j and a fixed order and position of the remaining elements (excluding i : j), meaning this sum
does not depend on πi:j . Thus, for the remainder of this proof, we can assume without loss of generality
that i = 1 and j = n. The quantity of interest becomes
P
− inv(π1:n )
Pr[1 [U+001F control character] n]
π :π∈S1[U+001F control character]n φ
= P 1:n
− inv(π1:n )
Pr[n [U+001F control character] 1]
πi:j :π∈Sn[U+001F control character]1 φ
P
− inv(π)
π∈S1[U+001F control character]n φ
=P
− inv(π)
π∈Sn[U+001F control character]1 φ
Next, we observe that we can similarly ignore inversions between two elements that are neither 1 nor n.
To see this, let inv1,n (π) be the number of inversions involving at least one of 1 and n. Then, if we fix the
order and positions of 1 and n, all possible permutations of the remaining elements 2 through n − 1 produce
the same number of inversions inv1,n (π). More formally, let π(1) and π(n) be the respective positions of
elements 1 and n. Then, this we have
X
X
X
φ− inv(π) =
φ− inv(π)
π∈S1[U+001F control character]n

k<` π:π(1) =k,π(n) =`

=

X

=

X

X

φ− inv1,n (φ) · φinv1,n (φ)−inv(π)

k<` π:π(1) =k,π(n) =`

φ−(k−1)−(n−`)

k<`

X

φinv1,n (φ)−inv(π)

π:π(1) =k,π(n) =`

inv1,n (φ)−inv(π)
As noted above,
does not depend on k or `, since every permutation of
π:π(1) =k,π(n) =` φ
the remaining elements yields the same number of inversions among them regardless of k and `. A similar
argument yields
X
X
X
φ− inv(π) =
φ−(k−1)−(n−`)+1
φinv1,n (φ)−inv(π)

P

π∈Sn[U+001F control character]1

k>`

π:π(1) =k,π(n) =`

34


[form-feed in source extraction]
Putting these together, we have
P
P
−(k−1)−(n−`)
inv1,n (φ)−inv(π)
Pr[1 [U+001F control character] n]
k<` φ
π:π(1) =k,π(n) =` φ
P
=P
inv1,n (φ)−inv(π)
−(k−1)−(n−`)+1
Pr[n [U+001F control character] 1]
π:π(1) =k,π(n) =` φ
k>` φ
P
φ−(k−1)−(n−`)
= P k<` −(k−1)−(n−`)+1
k>` φ
P
φ−(k−1)−(n−`)
φn−1
= P k<` −(k−1)−(n−`)+1 · n−1
φ
k>` φ
P
φ`−k
= P k<` `−k+1
k>` φ
Note that each term in the numerator is strictly increasing in φ, while each term in the denominator is
d
d Pr[1[U+001F control character]n]
weakly decreasing in φ. As a result, dφ
Pr[n[U+001F control character]1] > 0, meaning for any i < j, dφ Pr[i [U+001F control character] j] > 0.
\end{ReviewVerbatim}


\paragraph{Expanded PaperInterface specification}
\begin{ReviewVerbatim}
∀ {n : ℕ} (center : EconCSLib.SocialChoice.Ranking.Ranking n) {phiMore phiLess : ℝ},
  1 < phiLess →
    phiLess < phiMore →
      ∀ {c d : EconCSLib.SocialChoice.Ranking.Candidate n},
        EconCSLib.SocialChoice.Ranking.rankOf center c < EconCSLib.SocialChoice.Ranking.rankOf center d →
          have Mless := KR21Monoculture.concreteMallowsSpec center (phiLess - 1);
          have Mmore := KR21Monoculture.concreteMallowsSpec center (phiMore - 1);
          Mless.q = phiLess⁻¹ ∧
            Mmore.q = phiMore⁻¹ ∧
              (∀ (pi : EconCSLib.SocialChoice.Ranking.Ranking n),
                  (Mless.law pi).toReal =
                    phiLess⁻¹ ^ EconCSLib.SocialChoice.Ranking.kendallTau center pi /
                      ∑ tau, phiLess⁻¹ ^ EconCSLib.SocialChoice.Ranking.kendallTau center tau) ∧
                (∀ (pi : EconCSLib.SocialChoice.Ranking.Ranking n),
                    (Mmore.law pi).toReal =
                      phiMore⁻¹ ^ EconCSLib.SocialChoice.Ranking.kendallTau center pi /
                        ∑ tau, phiMore⁻¹ ^ EconCSLib.SocialChoice.Ranking.kendallTau center tau) ∧
                  (EconCSLib.pmfProb Mless.law fun pi =>
                      EconCSLib.SocialChoice.Ranking.rankOf pi c < EconCSLib.SocialChoice.Ranking.rankOf pi d) <
                    EconCSLib.pmfProb Mmore.law fun pi =>
                      EconCSLib.SocialChoice.Ranking.rankOf pi c < EconCSLib.SocialChoice.Ranking.rankOf pi d
\end{ReviewVerbatim}

\paragraph{Lean proof endpoint (built separately)}\begin{ReviewVerbatim}
KR21Monoculture.PaperInterface.lemma8_source_mallows_phi_pairwise_correct_probability_lt
\end{ReviewVerbatim}

\paragraph{Recorded source-to-Spec screening}
\textbf{Verdict:} \texttt{not recorded}
\\
\textbf{Reason:} not recorded
\reviewmatch{Matches source input}
\noindent\textbf{Reviewer annotation}\par
\reviewerbox
\clearpage
\hypertarget{source-claim-2d4c1840ef47dce5}{}
\section*{Review row 55}
\textbf{Source locator:} {\footnotesize\raggedright cited publication, lines 497-\allowbreak{}500\par}
\paragraph{Verbatim source input}
\begin{ReviewVerbatim}
First, we must define a family of RUMs that satisfies the conditions of Definition 1. Assume without loss of
generality that the noise distribution E has unit variance. Then, consider the family of RUMs parameterized
by θ in which candidates are ranked according to xi + εθi . By this definition, the standard deviation of the
noise for a particular value of θ is simply 1/θ. Intuitively, larger values of θ reduce the effect of the noise,

[Next verbatim source excerpt]

Finally, there exist noise distributions that violate Definition 2 for any candidate distribution D. In
particular, the RUM family defined by the Gumbel distribution is well-known to be equivalent to the PlackettLuce model of ranking, which is generated by sequentially selecting candidate i with probability
exp(θxi )
,
j∈S exp(θxj )

P

(7)

where S is the set of remaining candidates [12, 4]. Under the Plackett-Luce model, for any θ, UAH (θ, θ) =
\end{ReviewVerbatim}

\paragraph{Approved corrected review target}
The archival source above is not asserted equivalent to this replacement.
\begin{ReviewVerbatim}
For any finite candidate set, common location, innovation scale s > 0, and theta > 0, let arrival_i be iid unit-rate exponential and set epsilon_i = location - s*log(arrival_i). The arrival law is the iid exponential product law, arrivals are almost surely positive, the innovation law is its coordinatewise pushforward, and ranking value_i + epsilon_i/theta has exactly the Plackett-Luce ranking law at inverse temperature theta/s. If the visible analytic premise Var(-log T)=pi^2/6 holds for unit-rate exponential T, then the specialization s=sqrt(6)/pi has coordinate variance one and Plackett-Luce inverse temperature theta/(sqrt(6)/pi). The archival unit-variance/same-theta formula is not asserted.
\end{ReviewVerbatim}

\textbf{Archival source anchor:} {\footnotesize\raggedright cited publication, lines 497-\allowbreak{}548\par}
\paragraph{Recorded basis}
User instruction 2026-\allowbreak{}07-\allowbreak{}25:\allowbreak{} assume whatever convention makes the result true, provided the convention is explicit and the archival claim is not represented as proved.\allowbreak{}
\paragraph{Expanded PaperInterface specification}
\begin{ReviewVerbatim}
∀ {n : ℕ} (location : ℝ) {s theta : ℝ},
  0 < s →
    0 < theta →
      ∀ (value : EconCSLib.SocialChoice.Ranking.Candidate n → ℝ),
        (KR21Monoculture.gumbelArrivalLaw n = MeasureTheory.Measure.pi fun x => ProbabilityTheory.expMeasure 1) ∧
          (∀ᵐ (arrival : EconCSLib.SocialChoice.Ranking.Candidate n → ℝ) ∂KR21Monoculture.gumbelArrivalLaw n,
              ∀ (i : EconCSLib.SocialChoice.Ranking.Candidate n), 0 < arrival i) ∧
            (∀ (arrival : EconCSLib.SocialChoice.Ranking.Candidate n → ℝ)
                (i : EconCSLib.SocialChoice.Ranking.Candidate n),
                KR21Monoculture.commonLocationPositiveScaleGumbelNoise location s arrival i =
                  location + -s * Real.log (arrival i)) ∧
              (∀ (arrival : EconCSLib.SocialChoice.Ranking.Candidate n → ℝ)
                  (i : EconCSLib.SocialChoice.Ranking.Candidate n),
                  KR21Monoculture.commonLocationPositiveScaleGumbelScores location s theta value arrival i =
                    value i + (location + -s * Real.log (arrival i)) / theta) ∧
                (MeasureTheory.Measure.map (KR21Monoculture.commonLocationPositiveScaleGumbelNoise location s)
                      (KR21Monoculture.gumbelArrivalLaw n) =
                    MeasureTheory.Measure.pi fun x =>
                      MeasureTheory.Measure.map (fun arrival => location + -s * Real.log arrival)
                        (ProbabilityTheory.expMeasure 1)) ∧
                  KR21Monoculture.commonLocationPositiveScaleGumbelRUMRankingPMF location s theta value =
                      KR21Monoculture.plackettLuceRankingPMF (theta / s) value ∧
                    (ProbabilityTheory.variance id KR21Monoculture.scaleOneGumbelMeasure = Real.pi ^ 2 / 6 →
                      (KR21Monoculture.sourceUnitVarianceGumbelNoiseLaw location =
                          MeasureTheory.Measure.pi fun x =>
                            MeasureTheory.Measure.map (fun arrival => location + -(√6 / Real.pi) * Real.log arrival)
                              (ProbabilityTheory.expMeasure 1)) ∧
                        (∀ (epsilon : EconCSLib.SocialChoice.Ranking.Candidate n → ℝ)
                            (i : EconCSLib.SocialChoice.Ranking.Candidate n),
                            KR21Monoculture.sourceUnitVarianceGumbelScores theta value epsilon i =
                              value i + epsilon i / theta) ∧
                          (∀ (i : EconCSLib.SocialChoice.Ranking.Candidate n),
                              ProbabilityTheory.variance (fun epsilon => epsilon i)
                                  (KR21Monoculture.sourceUnitVarianceGumbelNoiseLaw location) =
                                1) ∧
                            KR21Monoculture.sourceUnitVarianceGumbelRUMRankingPMF location theta value =
                              KR21Monoculture.plackettLuceRankingPMF (theta / (√6 / Real.pi)) value)
\end{ReviewVerbatim}

\paragraph{Lean proof endpoint (built separately)}\begin{ReviewVerbatim}
KR21Monoculture.PaperInterface.section31_corrected_gumbel_plackett_luce_target
\end{ReviewVerbatim}

\paragraph{Recorded source-to-Spec screening}
\textbf{Verdict:} \texttt{not recorded}
\\
\textbf{Reason:} not recorded
\reviewmatch{Matches approved corrected target}
\noindent\textbf{Reviewer annotation}\par
\reviewerbox
\clearpage
\hypertarget{source-claim-e12ec3045f35c0cc}{}
\section*{Review row 56}
\textbf{Source locator:} {\footnotesize\raggedright cited publication, lines 176-\allowbreak{}182\par}
\paragraph{Verbatim source input}
\begin{ReviewVerbatim}
Each firm in our model has access to the same underlying pool of n candidates, which they rank using
a randomized mechanism R to get a permutation π as described above. Then, in a random order, each
firm hires the highest-ranked remaining candidate according to their ranking. Thus, if two firms both rank
candidate i first, only one of them can hire i; the other must hire the next available candidate according to
their ranking. In our model, candidates automatically accept the offer they get from a firm. For the sake
of simplicity, throughout this paper, we restrict ourselves to the case where there are two firms hiring one
candidate each, although our model readily generalizes to more complex cases.

[Next verbatim source excerpt]

where S is the set of remaining candidates [12, 4]. Under the Plackett-Luce model, for any θ, UAH (θ, θ) =
UAA (θ, θ). To see this, suppose the firm that hires first selects candidate i∗ . Then, the firm that hires second
8


[form-feed in source extraction]
gets each candidate i with probability given by (7) with S = {1, . . . , n}\i∗ . As a result, by (3), if π, σ ∼ Fθ ,
E [π1 − π2 | π1 6= σ1 ] = 0
for any candidate distribution D, meaning the Plackett-Luce model never meets Definition 2. Thus, under
the Plackett-Luce model, monoculture has no effect—the optimal strategy is always to use the best available
ranking, regardless of competitors’ strategies.
\end{ReviewVerbatim}


\paragraph{Expanded PaperInterface specification}
\begin{ReviewVerbatim}
∀ {n : ℕ} (D : MeasureTheory.Measure (KR21Monoculture.ValueProfile n)) [MeasureTheory.IsProbabilityMeasure D]
  (center : EconCSLib.SocialChoice.Ranking.Ranking n) {thetaA thetaH : ℝ},
  0 < thetaA →
    0 < thetaH →
      (∀ (c : EconCSLib.SocialChoice.Ranking.Candidate n), MeasureTheory.Integrable (fun value => value c) D) →
        (∀ᵐ (value : EconCSLib.SocialChoice.Ranking.Candidate n → ℝ) ∂D,
            EconCSLib.SocialChoice.Ranking.StrictlyOrderedBy center value) →
          have M := fun value =>
            { algorithmRanking := KR21Monoculture.plackettLuceRankingPMF thetaA value,
              humanRanking := KR21Monoculture.plackettLuceRankingPMF thetaH value, value := value };
          (∀ (self other : KR21Monoculture.Strategy),
              MeasureTheory.Integrable (fun value => (M value).payoffAgainst self other) D) ∧
            (thetaH ≤ thetaA →
                ∫ (value : KR21Monoculture.ValueProfile n),
                      (M value).payoffAgainst KR21Monoculture.Strategy.algorithm KR21Monoculture.Strategy.algorithm ∂D ≥
                    ∫ (value : KR21Monoculture.ValueProfile n),
                      (M value).payoffAgainst KR21Monoculture.Strategy.human KR21Monoculture.Strategy.algorithm ∂D ∧
                  ∫ (value : KR21Monoculture.ValueProfile n),
                      (M value).payoffAgainst KR21Monoculture.Strategy.algorithm KR21Monoculture.Strategy.human ∂D ≥
                    ∫ (value : KR21Monoculture.ValueProfile n),
                      (M value).payoffAgainst KR21Monoculture.Strategy.human KR21Monoculture.Strategy.human ∂D) ∧
              (thetaA ≤ thetaH →
                ∫ (value : KR21Monoculture.ValueProfile n),
                      (M value).payoffAgainst KR21Monoculture.Strategy.human KR21Monoculture.Strategy.algorithm ∂D ≥
                    ∫ (value : KR21Monoculture.ValueProfile n),
                      (M value).payoffAgainst KR21Monoculture.Strategy.algorithm KR21Monoculture.Strategy.algorithm ∂D ∧
                  ∫ (value : KR21Monoculture.ValueProfile n),
                      (M value).payoffAgainst KR21Monoculture.Strategy.human KR21Monoculture.Strategy.human ∂D ≥
                    ∫ (value : KR21Monoculture.ValueProfile n),
                      (M value).payoffAgainst KR21Monoculture.Strategy.algorithm KR21Monoculture.Strategy.human ∂D)
\end{ReviewVerbatim}

\paragraph{Lean proof endpoint (built separately)}\begin{ReviewVerbatim}
KR21Monoculture.PaperInterface.section31_plackettLuce_outer_best_available_weakly_dominates
\end{ReviewVerbatim}

\paragraph{Recorded source-to-Spec screening}
\textbf{Verdict:} \texttt{not recorded}
\\
\textbf{Reason:} not recorded
\reviewmatch{Matches source input}
\noindent\textbf{Reviewer annotation}\par
\reviewerbox
\clearpage
\hypertarget{source-claim-4cfc28215a7bcccf}{}
\section*{Review row 57}
\textbf{Source locator:} {\footnotesize\raggedright cited publication, lines 139-\allowbreak{}165\par}
\paragraph{Verbatim source input}
\begin{ReviewVerbatim}
Definition 1 (Noisy permutation family). A noisy permutation family Fθ is a family of distributions over
permutations that satisfies the following conditions for any θ > 0 and set of candidates x:
1. (Differentiability) For any permutation π, PrFθ [π] is continuous and differentiable in θ.
2. (Asymptotic optimality) For the true ranking π ∗ , limθ→∞ PrFθ [π ∗ ] = 1.
3. (Monotonicity) For any (possibly empty) S ⊂ x, let π (−S) be the partial ranking produced by removing
(−S)
the items in S from π. Let π1
denote the value of the top-ranked candidate according to π (−S) . For
0
any θ > θ,
h
i
h
i
(−S)
(−S)
EFθ0 π1
≥ EFθ π1
.
(1)
Moreover, for S = ∅, (1) holds with strict inequality.
θ serves as an “accuracy parameter”: for large θ, the noisy ranking converges to the true ranking over
candidates. The monotonicity condition states that a higher value of θ leads to a better first choice, even
if some of the candidates are removed after ranking. Removal after ranking (as opposed to before) is
important because some of the ranking models we will consider later do not satisfy Independence of Irrelevant
Alternatives. Examples of noisy permutation families include Random Utility Models [23] and the Mallows
Model [14], both of which we will discuss in detail later.

[Next verbatim source excerpt]

B

3-candidate RUM Counterexamples

B.1

Violating Definition 2

Here, we provide a noise mode E, accuracy parameter θ, and candidate distribution D such that UAH < UAA .
Choose the noise distribution E and accuracy parameter θ such that
[U+F8F1 delimiter glyph]
[U+F8F4 delimiter glyph]1
w.p. 2δ
ε [U+F8F2 delimiter glyph]
= 0
w.p.1 − δ
θ [U+F8F4 delimiter glyph]
[U+F8F3 delimiter glyph]
−1 w.p. 2δ
Note that this distribution does not satisfy Definition 1 because it is neither differentiable nor supported on
(−∞, ∞); however, we can provide a “smooth” approximation to this distribution by expressing it as the
sum of arbitrarily tightly concentrated Gaussians with the same results.

15


[form-feed in source extraction]
We choose the candidate distribution D such that x1 − 1 > x2 > x3 > x1 − 2. For example,
7
4
1
x2 =
2
x3 = 0
x1 =

Under this condition, assuming x3 = 0 without loss of generality,
UAH (θ, θ) − UAA (θ, θ) =

[U+0001 control character]
δ2 3
δ x1 − 4δ 2 x1 + 4δx1 + 2δ 3 x2 − 14δ 2 x2 + 20δx2 − 8x2
32
2

Notice that the lowest-power δ term is − δ 4x2 . Therefore, for sufficiently small δ, this is negative. For
example, plugging in the values given above with δ = .1, UAH (θ, θ) − UAA (θ, θ) ≈ −0.00076.

B.2

[Next verbatim source excerpt]

Violating Definition 3

Next, we’ll give a 3-candidate RUM for which UAH < UHH does not hold in general. Consider the following
3-candidate example.
x1 = 3
x2 = 2
x3 = 0
Choose E and θ such that

[U+F8F1 delimiter glyph]
1
[U+F8F4 delimiter glyph]
[U+F8F4 delimiter glyph]
[U+F8F4 delimiter glyph]
[U+F8F2 delimiter glyph]
−1
ε
=
θ [U+F8F4 delimiter glyph]
10
[U+F8F4 delimiter glyph]
[U+F8F4 delimiter glyph]
[U+F8F3 delimiter glyph]
−10

w.p. 1−δ
2
w.p. 1−δ
2
w.p. 2δ
w.p. 2δ

Again, while this noise model doesn’t satisfy Definition 1, we can approximate it arbitrarily closely with the
sum of tightly concentrated Gaussians. Let the θA = 1.1θ and θH = 0.9θ.
We will show that for these parameters, UAH (θA , θH ) > UHH (θA , θH ), i.e., it is somehow better to choose
after a better opponent than after a worse opponent. At a high level, the reasoning for this is as follows:
1. When choosing first, the only difference between the algorithm and the human evaluator is that the
algorithm is more likely to choose x2 than x3 . Both strategies have identical probabilities of selecting
x1 .
2. When choosing second, the human evaluator’s utility is higher when x2 has already been chosen than
when x3 has already been chosen. This is because when x2 is unavailable, the human evaluator
is almost guaranteed to get x1 ; when x3 is unavailable, the human evaluator will choose x2 with
probability ≈ 1/4.
Let τ and π be rankings generated by the algorithm and human evaluator respectively. First, we will
show that
Pr[τ1 = x1 ] = Pr[π1 = x1 ]

(B.1)

Pr[τ1 = x2 ] > Pr[π1 = x2 ]

(B.2)

To do so, consider the realizations of ε1 , ε2 , ε3 that result in different rankings under θA and θH . In fact,
the only set of realizations that result in different rankings are when ε2 /θ = −1 and ε3 /θ = 1. Thus,
the algorithm and human evaluator always rank x1 in the same position, conditioned on a realization,
which proves (B.1); the only difference is that the algorithm sometimes ranks x2 above x3 when the human
16


[form-feed in source extraction]
evaluator does not. Moreover, whenever ε1 /θ = −10, x2 is more strictly more likely to be ranked first under
the algorithm than the human evaluator, which proves (B.2).
Next, we must show that when choosing second, the human evaluator is better off when x2 is unavailable
than when x3 is unavailable. This is clearly true because for the human evaluator,
[U+0014 control character]
[U+0015 control character]
ε1 θH
ε3 θH
Pr x1 +
> x3 +
≈ 1 − O(δ)
θ
θ
[U+0014 control character]
[U+0015 control character]
ε1 θH
ε3 θH
3
Pr x1 +
> x2 +
≈
θ
θ
4
Thus, conditioned on x2 being unavailable, the human evaluator gets utility ≈ 3, whereas when x3 is
unavailable, the human evaluator gets utility ≈ 2.75. Let u−i be the expected utility for the human evaluator
when xi is unavailable. Putting this together, we get
UAH (θA , θH ) − UHH (θA , θH )
=

3
X

(Pr[τ1 = xi ] − Pr[π1 = xi ])u−i

i=1

= (Pr[τ1 = x1 ] − Pr[π1 = x1 ])u−1 + (Pr[τ1 = x2 ] − Pr[π1 = x2 ])u−2 + (Pr[τ1 = x3 ] − Pr[π1 = x3 ])u−3
= (Pr[τ1 = x2 ] − Pr[π1 = x2 ])u−2 + (Pr[τ1 = x3 ] − Pr[π1 = x3 ])u−3
(Pr[τ1 = x1 = Pr[π1 = x1 ])
P3
P3
= (Pr[τ1 = x2 ] − Pr[π1 = x2 ])(u−2 − u−3 )
( i=1 Pr[τ1 = xi ] = i=1 Pr[π1 = xi ])
>0
The last step follows from (B.2) and because u−2 > u−3 .
\end{ReviewVerbatim}


\paragraph{Expanded PaperInterface specification}
\begin{ReviewVerbatim}
(KR21Monoculture.appendixB1Value 0 = 7 / 4 ∧
    KR21Monoculture.appendixB1Value 1 = 1 / 2 ∧ KR21Monoculture.appendixB1Value 2 = 0) ∧
  (KR21Monoculture.appendixB1NoiseValue KR21Monoculture.AppendixB1NoiseAtom.plusOne = 1 ∧
      KR21Monoculture.appendixB1NoiseValue KR21Monoculture.AppendixB1NoiseAtom.zero = 0 ∧
        KR21Monoculture.appendixB1NoiseValue KR21Monoculture.AppendixB1NoiseAtom.minusOne = -1) ∧
    ((KR21Monoculture.appendixB1NoisePMF KR21Monoculture.AppendixB1NoiseAtom.plusOne).toReal = 1 / 20 ∧
        (KR21Monoculture.appendixB1NoisePMF KR21Monoculture.AppendixB1NoiseAtom.zero).toReal = 9 / 10 ∧
          (KR21Monoculture.appendixB1NoisePMF KR21Monoculture.AppendixB1NoiseAtom.minusOne).toReal = 1 / 20) ∧
      KR21Monoculture.appendixB1GaussianLatentMeasure =
          ((KR21Monoculture.appendixB1NoisePMF.toMeasure.prod KR21Monoculture.appendixB1NoisePMF.toMeasure).prod
                KR21Monoculture.appendixB1NoisePMF.toMeasure).prod
            ((ProbabilityTheory.gaussianReal 0 1).prod
              ((ProbabilityTheory.gaussianReal 0 1).prod (ProbabilityTheory.gaussianReal 0 1))) ∧
        (∀ (s theta : ℝ),
            0 < s →
              0 < theta →
                (KR21Monoculture.appendixB1SourceGaussianMixtureFamily s).dist theta =
                  KR21Monoculture.rumRankingPMFOfMeasure KR21Monoculture.appendixB1GaussianLatentMeasure
                    (fun omega =>
                      EconCSLib.SocialChoice.Ranking.rankByScore fun c =>
                        KR21Monoculture.appendixB1Value c +
                          (KR21Monoculture.appendixB1NoiseValue
                                (KR21Monoculture.appendixB1NoiseTripleFunction omega.1 c) +
                              s * KR21Monoculture.appendixBGaussianTripleFunction omega.2 c) /
                            theta)
                    ⋯) ∧
          (KR21Monoculture.appendixB2Value 0 = 3 ∧
              KR21Monoculture.appendixB2Value 1 = 2 ∧ KR21Monoculture.appendixB2Value 2 = 0) ∧
            (KR21Monoculture.appendixB2NoiseValue KR21Monoculture.AppendixB2NoiseAtom.plusOne = 1 ∧
                KR21Monoculture.appendixB2NoiseValue KR21Monoculture.AppendixB2NoiseAtom.minusOne = -1 ∧
                  KR21Monoculture.appendixB2NoiseValue KR21Monoculture.AppendixB2NoiseAtom.plusTen = 10 ∧
                    KR21Monoculture.appendixB2NoiseValue KR21Monoculture.AppendixB2NoiseAtom.minusTen = -10) ∧
              ((KR21Monoculture.appendixB2NoisePMF KR21Monoculture.AppendixB2NoiseAtom.plusOne).toReal = 9 / 20 ∧
                  (KR21Monoculture.appendixB2NoisePMF KR21Monoculture.AppendixB2NoiseAtom.minusOne).toReal = 9 / 20 ∧
                    (KR21Monoculture.appendixB2NoisePMF KR21Monoculture.AppendixB2NoiseAtom.plusTen).toReal = 1 / 20 ∧
                      (KR21Monoculture.appendixB2NoisePMF KR21Monoculture.AppendixB2NoiseAtom.minusTen).toReal =
                        1 / 20) ∧
                KR21Monoculture.appendixB2GaussianLatentMeasure =
                    ((KR21Monoculture.appendixB2NoisePMF.toMeasure.prod
                              KR21Monoculture.appendixB2NoisePMF.toMeasure).prod
                          KR21Monoculture.appendixB2NoisePMF.toMeasure).prod
                      ((ProbabilityTheory.gaussianReal 0 1).prod
                        ((ProbabilityTheory.gaussianReal 0 1).prod (ProbabilityTheory.gaussianReal 0 1))) ∧
                  (∀ (s theta : ℝ),
                      0 < s →
                        0 < theta →
                          (KR21Monoculture.appendixB2SourceGaussianMixtureFamily s).dist theta =
                            KR21Monoculture.rumRankingPMFOfMeasure KR21Monoculture.appendixB2GaussianLatentMeasure
                              (fun omega =>
                                EconCSLib.SocialChoice.Ranking.rankByScore fun c =>
                                  KR21Monoculture.appendixB2Value c +
                                    (KR21Monoculture.appendixB2NoiseValue
                                          (KR21Monoculture.appendixB2NoiseTripleFunction omega.1 c) +
                                        s * KR21Monoculture.appendixBGaussianTripleFunction omega.2 c) /
                                      theta)
                              ⋯) ∧
                    (∃ delta,
                        0 < delta ∧
                          ∀ (s : ℝ),
                            0 < s →
                              s < delta →
                                EconCSLib.SocialChoice.Ranking.expectedSecondMoverIndependent
                                      ((KR21Monoculture.appendixB1SourceGaussianMixtureFamily s).dist 1)
                                      ((KR21Monoculture.appendixB1SourceGaussianMixtureFamily s).dist 1)
                                      KR21Monoculture.appendixB1Value -
                                    EconCSLib.SocialChoice.Ranking.expectedSecondMoverShared
                                      ((KR21Monoculture.appendixB1SourceGaussianMixtureFamily s).dist 1)
                                      KR21Monoculture.appendixB1Value <
                                  0) ∧
                      ∃ delta,
                        0 < delta ∧
                          ∀ (s : ℝ),
                            0 < s →
                              s < delta →
                                0 <
                                  EconCSLib.SocialChoice.Ranking.expectedSecondMoverIndependent
                                      ((KR21Monoculture.appendixB2SourceGaussianMixtureFamily s).dist (9 / 10))
                                      ((KR21Monoculture.appendixB2SourceGaussianMixtureFamily s).dist (11 / 10))
                                      KR21Monoculture.appendixB2Value -
                                    EconCSLib.SocialChoice.Ranking.expectedSecondMoverIndependent
                                      ((KR21Monoculture.appendixB2SourceGaussianMixtureFamily s).dist (9 / 10))
                                      ((KR21Monoculture.appendixB2SourceGaussianMixtureFamily s).dist (9 / 10))
                                      KR21Monoculture.appendixB2Value
\end{ReviewVerbatim}

\paragraph{Lean proof endpoint (built separately)}\begin{ReviewVerbatim}
KR21Monoculture.PaperInterface.appendixB_smoothing_source_core
\end{ReviewVerbatim}

\paragraph{Recorded source-to-Spec screening}
\textbf{Verdict:} \texttt{not recorded}
\\
\textbf{Reason:} not recorded
\reviewmatch{Matches source input}
\noindent\textbf{Reviewer annotation}\par
\reviewerbox

\end{document}
